\documentclass{l3in2edoc}

\usepackage{xcolor,csquotes}

\ExplSyntaxOn \makeatletter
\cs_set_protected:Npn \__codedoc_names_parse_one_aux:nnNn #1#2#3#4
  {
    \bool_if:NTF #3
      {
        \tl_if_head_eq_charcode:nNTF {#2} :
          { \__codedoc_names_parse_aux:nnn {#4} {#4} { \scan_stop: } }
          {
            \exp_args:Ne \__codedoc_names_parse_aux:nnn
              { \__codedoc_base_form_aux:nnN {#1} {#2} #3 }
              {#1} {#2}
          }
      }
      {
%        \bool_if:NT \l__codedoc_macro_TF_bool
%          { \msg_error:nne { l3doc } { no-signature-TF } {#4} }
        \__codedoc_names_parse_aux:nnn {#4} {#4} { \scan_stop: }
      }
  }


% Next lines make "int" acceptable in the optional arg to function but
% it doesn't do much yet. How should internal functions be marked?
% For now I simply marked the commands in red as in (don't use or rely
% on).

\keys_define:nn { l3doc/function }
  {
    int .value_forbidden:n = true ,
    int .code:n =
      {
        \bool_set_true:N \l__codedoc_macro_internal_bool
      } ,
    user .value_forbidden:n = true ,
    user .code:n =
      {
        \bool_set_true:N \l__codedoc_macro_user_bool
      } ,
  }

\bool_new:N \l__codedoc_macro_user_bool

\cs_set_protected:Npn \__codedoc_typeset_functions:
  {
    \small\ttfamily
    \__codedoc_target:
    \Hy@MakeCurrentHref { HD. \int_use:N \c@HD@hypercount }
    \__codedoc_date_compare:VNVF \l__codedoc_date_added_tl > \l__codedoc_base_date_tl
      { \tl_clear:N \l__codedoc_date_added_tl }
    \__codedoc_date_compare:VNVF \l__codedoc_date_updated_tl > \l__codedoc_base_date_tl
      { \tl_clear:N \l__codedoc_date_updated_tl }
    \group_begin:
      \bool_if:NT \l__codedoc_macro_internal_bool { \color{red}  }
      \bool_if:NT \l__codedoc_macro_user_bool     { \color{blue} }
      \begin{tabular} [t] { @{} l @{} >{\hspace{\tabcolsep}} r @{} }
        \toprule
        \__codedoc_function_extra_labels:
        \__codedoc_names_typeset:
        \__codedoc_typeset_dates:
        \bottomrule
      \end{tabular}
    \group_end:
    \normalfont\normalsize
  }

\ExplSyntaxOff \makeatother

\begin{document}

\title{\LaTeXe{} Interfaces and non-interfaces}
\author{%
 The \LaTeX{} Project\thanks
   {%
     E-mail:
       \href{mailto:latex-team@latex-project.org}
         {latex-team@latex-project.org}%
   }%
}

\date{Released 2026-09-16}

\pagenumbering{roman}

\maketitle

\begin{abstract}
  I spent a bit this afternoon on making a list of all \LaTeXe{}
  commands in the kernel as input for something like the
  \texttt{cmdguide} or as a reference manual named
  \texttt{interface2}.

  So far this only covers commands from \texttt{latex.ltx}. Those from
  related files would need adding too.

  In its present form it is not yet useful, but once more commands are
  sorted and documented \ldots\ Frank
\end{abstract}


\tableofcontents

\newpage


\section{Conventions}

\begin{function}[]
  {
    \foo, \foovariant
  }
  \begin{syntax}
    \cs{foo} \oarg{optional arg} \Arg{mandatory arg}
  \end{syntax}
  General form of documentation for commands. Commands that share a
  common behavior/interface may be documented together like this.

  The syntax box gives one example with arguments. If the command is
  not yet documented it contains only the command name amd no
  explanation afterwards.
  
\end{function}

\begin{function}[TF]
  {
    \foo
  }
  \begin{syntax}
    \cs{fooF} \Arg{text arg} \Arg{false branch}
  \end{syntax}
  Conditionals using the \texttt{TF} convention are documented like
  this.  In \LaTeXe{} commands normally do not indicate the arguments
  they take (\texttt{TF} is one of the exeptions. But any other
  arguments aren't indicated so one has to \enquote{know}.

  In contrast,
  commands from the L3 programming layer always denote their argument
  signature in the name (via the \texttt{:} convention).
  
\end{function}

\begin{function}[int]
  {
    \foo
  }
  \begin{syntax}
    \cs{foo} 
  \end{syntax}
  Internal commands that should not be relied on are marked in red.
  
\end{function}

\begin{function}[user]
  {
    \foo
  }
  \begin{syntax}
    \cs{foo} 
  \end{syntax}
  %
  Document-level (user) commands are shown in blue. In addition
  commands that are mainly for class and package writer but can also
  be useful in the document preamble or inside the document are also
  marked this way. This means that basically all core commands
  documented in \emph{The LaTeX{} Companion}\cite{tlc3} are shown in
  blue.
  
\end{function}

\begin{function}[EXP]
  {
    \foo
  }
  \begin{syntax}
    \cs{foo} 
  \end{syntax}
  Fully expandable functions, i.e., those that can be used
  inside \cs{edef} or \cs{expanded} are indicated with a star.

  Most user commands can be safely put into \cs{edef} but then do
  nothing at all (because they are internally protected either
  by \cs{protect} or by being defined with \cs{protected}). Such
  commands are not marked with a star!
  
\end{function}



\newpage


\section{Command reference}

\pagenumbering{arabic}

Most of what follows is just a list of all commands defined by
the \LaTeX{} kernel, but at this point not documented or sorted.

%===============================================

\subsection{Low-level \TeX{} registers and streams}


%-----------------------------------------------
\subsubsection{\TeX{} registers}

\begin{function}[]
  {
    \newcount,
    \newdimen,
    \newskip,
    \newmuskip,
    \newtoks
  }
  \begin{syntax}
    \cs{newcount} \meta{cmd}
  \end{syntax}
  A new integer (count), dimen, skip, muskip, toks register is allocated and
  the csname \meta{cmd} refers to it.

\end{function}


\begin{function}[]
  {
    \newbox
  }
  \begin{syntax}
    \cs{newbox}
  \end{syntax}

\end{function}


\begin{function}[]
  {
    \newlanguage
  }
  \begin{syntax}
    \cs{newlanguage}
  \end{syntax}

\end{function}

\begin{function}[]
  {
    \newinsert
  }
  \begin{syntax}
    \cs{newinsert}
  \end{syntax}

\end{function}

\begin{function}[]
  {
    \newmarks
  }
  \begin{syntax}
    \cs{newmarks}
  \end{syntax}

\end{function}



\begin{function}[]
  {
    \newattribute
  }
  \begin{syntax}
    \cs{newattribute}
  \end{syntax}

\end{function}




\begin{function}[]
  {
    \advance,
    \divide
  }
  \begin{syntax}
    \cs{advance} \meta{cmd} \texttt{by} \meta{value}
  \end{syntax}

  Change the register value by adding or dividing. The \texttt{by} is
  optional.
  
\end{function}



%-----------------------------------------------
\subsubsection{\TeX{} input and output streams}

\begin{function}[]
  {
    \newread,
    \newwrite
  }
  \begin{syntax}
    \cs{newwrite} cs{\meta{cmd}}
  \end{syntax}

  A new read or write stream is allocated and the csname \cs{\meta{cmd}}
  refers to it.

\end{function}


\begin{function}[]
  {
    \openin,
    \openout
  }
  \begin{syntax}
    \cs{openin}
  \end{syntax}

\end{function}

\begin{function}[]
  {
    \read
  }
  \begin{syntax}
    \cs{read}
  \end{syntax}

\end{function}


\begin{function}[]
  {
    \wlog,
    \protected@wlog
  }
  \begin{syntax}
    \cs{wlog} \Arg{message}
  \end{syntax}
  Writes \meta{message} to the transcript file on a line by
  itself. Commands inside \meta{message} are expanded, which may
  result in errors depending on the content. \cs{protected@wlog} is
  like \cs{wlog} but sets \cs{protect} in a suitable way so that
  robust commands behave as in \cs{typeout}.
\end{function}


\begin{function}[]
  {
    \message
  }
  \begin{syntax}
    \cs{message}
  \end{syntax}
  A \TeX{} primitive that writes \meta{message} to the terminal and
  transcript file. It doesn't produce new lines, so if one is wanted
  one can use \verb=^^J=. The \meta{message} is expanded.
\end{function}




%===============================================

\subsection{\LaTeX{} counters and lengths commands}

%-----------------------------------------------
\subsubsection{\LaTeX{} counters}

Each \LaTeX{} counter \meta{name} has a corresponding \TeX{} counter
register (defined internally with \cs{newcount}) with the name
\cs{c@\meta{name}}.

\begin{function}[user]
  {
    \newcounter
  }
  \begin{syntax}
    \cs{newcounter} \Arg{name} \oarg{within}
  \end{syntax}

\end{function}

\begin{function}[user]
  {
    \setcounter,
    \addtocounter,
    \stepcounter
  }
  \begin{syntax}
    \cs{setcounter} \Arg{name} \oarg{integer}
  \end{syntax}
  Set a \LaTeX{} counter \meta{name} to a \meta{value} or add a value or
  step the counter by one.
\end{function}


\begin{function}[user]
  {
    \refstepcounter
  }
  \begin{syntax}
    \cs{refstepcounter} \Arg{name}
  \end{syntax}
  Like \cs{stepcounter} but also sets \cs{@currentlabel} and \cs{@currentlabelname}.
\end{function}


\begin{variable}[]
  {
    \@currentlabel
  }

\end{variable}



\begin{variable}[]
  {
    \@currentlabelname
  }

\end{variable}


\begin{function}[]
  {
    \@currentcounter
  }
  \begin{syntax}
    \cs{@currentcounter}
  \end{syntax}

\end{function}





\begin{function}[]
  {
    \@nocounterr
  }
  \begin{syntax}
    \cs{@nocounterr} \Arg{name}
  \end{syntax}
  Error message that \meta{name} is not a \LaTeX{} counter.
\end{function}



%

%-----------------------------------------------
\subsubsection{\LaTeX{} lengths}

\LaTeX{} doesn't distinguishes between skips (variabl lengths) and
dimens (fixed lengths), the high-level user commands only provide
skips (called \enquote{lengths} in the \LaTeX{} manual.

\begin{function}[user]
  {
    \newlength
  }
  \begin{syntax}
    \cs{newlength} \meta{cmd}
  \end{syntax}

\end{function}

\begin{function}[user]
  {
    \addtolength,
    \setlength
  }
  \begin{syntax}
    \cs{addtolength} \meta{cmd} \Arg{length}
  \end{syntax}

\end{function}


\begin{function}[user]
  {
    \settowidth,
    \settoheight
  }
  \begin{syntax}
    \cs{settowidth} \meta{cmd} \Arg{material}
  \end{syntax}

\end{function}


%===============================================

\subsection{Text commands}


\begin{function}[user]
  {
    \ ,
    \space,
  }
  \begin{syntax}
    \cs{\space} \cs{\space} \cs{\space}
  \end{syntax}
  %
  The command \verb*=\ = typesets a space character, same as just
  using a \verb*= =, but it can be used in places where a normal space
  gets ignored, e.g., after a command such as \verb*=\LaTeX\ =. Several
  such commands in a row produce multiple spaces in contrast to using
  several normal spaces in a row, e.g., the example above would produce 3
  spaces.

  The \verb*=\ = command is a \TeX{} primitive for typsetting a space
  and for some operations, such as assigning a \cs{catcode} to the
  space character or displaying a space in \cs{typeout} one needs to
  use \cs{\space} to refer to this character.

  The command \cs{space} is a command that contains a space character
  as its replacement material. For typesetting this is identical to
  using \verb*=\ = but in other cases there are differences:
  \verb=\csname X\space\endcsname= produces a macro name with a space
  character in the name. In that position \cs{\space} (as a
  typesetting instruction) would not be allowed. For the same reason
  you need to use \cs{space} and not \cs{\space} inside \cs{typeout}.
  
\end{function}

\begin{function}[]
  {
    \@spaces

  }
  \begin{syntax}
    \cs{@spaces}
  \end{syntax}
  %
  The command \cs{@spaces} is defined to as four \cs{space} commands,
  so it can be used in \cs{typeout} to indent lines which is what
  \LaTeX{} does in many of its error and warning messages.

\end{function}





%-----------------------------------------------

\subsubsection{Text accents}

\begin{function}[user,EXP]
  {
    \",
    \',
    \.,
    \~,
    \b,    
    \c,
  }
  \begin{syntax}
    \cs{"}\meta{base letter}
  \end{syntax}
  %
  these commands typeset various diacritcals. These days \LaTeX{} is
  fully Unicode-aware (even in 8-bit engines like \pdfTeX{}) so that
  one normally simply writes \texttt{ä} instead of \verb=\"a= (or even
  \verb=\"{a}= as it was recommended in Lamport's manual).

  However, this remains a valid
  input syntax and it makes it possible to even produce glyphs not
  part of Unicode, e.g., \"\$ (produced with \verb=\"\$=).

\end{function}




%-----------------------------------------------


\subsubsection{Text characters and symbols}

These days characters and symbols can entered directly as UTF8 encoded
characters (and typeset if the used font supports them). Many of them
are also available via command names, largely for historical reasons,
when \TeX{} was still 8-bit based. In certain situation these are
still useful, for example, if your keyboard doesn't allow you easily
to produce a certain Unicode character.

There are also a few that need command names because they serve as
syntax characters in \TeX{}.

\begin{function}[user,EXP]
  {
    \$,
    \#,
    \&,
    \textasciitilde
  }
%
  They produce \$, \#, \& using them directly as characters is not
  easily possible as the have syntax functions. For
  \textasciitilde\ one needs to use the long command name, because
  \cs{~} is used for a title accent.

  There is also \cs{\%} for a percent sign (without the backslash it
  would start a line comment. Due to some limitation of the macros we
  use here it can't be displayed in the syntax box.

\end{function}

%\begin{function}[]
%  {
%    \%
%  }
%  \begin{syntax}
%    \cs{\%}
%  \end{syntax}
%
%\end{function}


\typeout{The command \string\% fails in function and in \string\cs }

\begin{function}[]
  {
    \@percentchar
  }
  \begin{syntax}
    \cs{@percentchar}
  \end{syntax}

\end{function}



%===============================================

\subsection{Math commands}

------------------------------------------------
\subsubsection{Math environments}


\begin{function}[user]
  {
    \(,
    \)
  }
  \begin{syntax}
    \cs{(} \meta{inline math} \cs{)}
  \end{syntax}
  %
  \LaTeX's alternative to \verb=$...$= (which I personally prefer as
  it is easier to spot in the source, even though it means you can't
  programtically check then that start and stop of math mode are
  properly matchd.
\end{function}


\begin{function}[user]
  {
    \[,
    \]
  }
  \begin{syntax}
    \cs{[} \meta{dislay math} \cs{]}
  \end{syntax}
  %
  \LaTeX's short commands for display formulas. It is not equivalent
  to the plain \TeX{} syntax \verb=$$...$$=. The latter should not be
  used in a \LaTeX{} document!
\end{function}

\begin{function}[user]
  {
    \math,
    \endmath
  }
  \begin{syntax}
    \cs{begin}\Arg{math}  \meta{inline math} \cs{end}\Arg{math}
  \end{syntax}
  Long environment form of \verb=\(...\)=.
\end{function}



\begin{function}[user]
  {
    \displaymath,
    \enddisplaymath
  }
  \begin{syntax}
    \cs{begin}\Arg{displaymath}  \meta{display math} \cs{end}\Arg{displaymath}
  \end{syntax}
  Long environment form of \verb=\[...\]=.
\end{function}


\begin{function}[int]
  {
    \@kernel@eqno,
    \@kernel@leqno
  }
  %
  Saved version of the \TeX{} primitives \cs{eqno} and \cs{leqno}
  which are then redefined by \LaTeX{}. They are not supposed to be
  used directly.
\end{function}


\begin{function}[]
  {
    \eqno,
    \leqno
  }
  \begin{syntax}
    \cs{eqno}
  \end{syntax}
  %
  These are not the \TeX{} primitives to add an equation number to a
  displayed equation, but are redefined by \LaTeX{} to do more work
  and only then internally call the \TeX{} primitives (which have been
  renamed to \cs{@kernel@eqno} and \cs{@kernel@leqno}).
\end{function}


\begin{function}[]
  {
    \equation,
    \endequation

  }
  \begin{syntax}
    \cs{equation} \meta{display formula}  \cs{endequation}
  \end{syntax}

\end{function}



------------------------------------------------
\subsubsection{Math operators}

\begin{function}[user]
  {
    \cos,
    \sin,
    \tan,
    \cosh,
    \sinh,
    \tanh,
    \cot,
    \coth,
    \arccos,
    \arcsin,
    \arctan,
    \log,
    \lg,
    \ln,
    \lim,
    \limsup,
    \liminf,
    \sec,
    \csc,
    \max,
    \min,
  }
  \begin{syntax}
    \cs{lim}\_\Arg{x \cs{to} \cs{infty}}
  \end{syntax}

\end{function}






\begin{function}[user]
  {
    \limits
  }
  \begin{syntax}
    \cs{limits}
  \end{syntax}

\end{function}



------------------------------------------------
\subsubsection{Other math commands}

\begin{function}[]
  {
    \atopwithdelims
  }
  \begin{syntax}
    \cs{atopwithdelims}
  \end{syntax}

\end{function}



\begin{function}[user]
  {
    \sqrt
  }
  \begin{syntax}
    \cs{sqrt}
  \end{syntax}

\end{function}




\begin{function}[]
  {
    \pmatrix
  }
  \begin{syntax}
    \cs{begin}\Arg{pmatrix} ...
  \end{syntax}

\end{function}



\begin{function}[user]
  {
    \pmod
  }
  \begin{syntax}
    \cs{pmod}
  \end{syntax}

\end{function}





\begin{function}[user]
  {
    \big,
    \bigg
  }
  \begin{syntax}
    \cs{big}
  \end{syntax}

\end{function}


\begin{function}[user]
  {
    \bigl,
    \bigm,
    \bigr
  }
  \begin{syntax}
    \cs{bigl}
  \end{syntax}

\end{function}



\begin{function}[user]
  {
    \biggl,
    \biggm,
    \biggr
  }
  \begin{syntax}
    \cs{biggl}
  \end{syntax}

\end{function}




------------------------------------------------
\subsubsection{Math environments}


------------------------------------------------
\subsubsection{Math symbols}


\begin{function}[]
  {
    \sqrtsign
  }
  \begin{syntax}
    \cs{sqrtsign}
  \end{syntax}

\end{function}



%===============================================

\subsection{Plain \TeX{} leftovers}

\LaTeX\,2.09 started out as a format on top of plain \TeX{} and kept a
number of plain \TeX{} commands. Some of them have been useful in
programming \LaTeXe{} macros, though the L3 programming layer provide
much better alternatives.  A few ended up in regular use user
documents, but with a few exceptions this is not recommended!


\begin{function}[user]
  {
    \centerline,
    \leftline,
    \rightline
  }
  \begin{syntax}
    \cs{centerline} \Arg{material}
  \end{syntax}
  %
  Produces a \cs{hbox} of size \cs{hsize} and typesets the
  \meta{material} inside, either centered, flush left, or flush
  right. As a simple \cs{hbox} it does not start a paragraph (as
  \cs{mbox} or \cs{makebox} do)
  
\end{function}

\begin{function}[int, deprecated]
  {
    \@@line
  }
  \begin{syntax}
    \cs{@@line} \Arg{material}
  \end{syntax}
  Command to implement \cs{centerline}, \cs{leftline}, and
  \cs{rightline}.  Deprecated, do not use directly.
\end{function}


\begin{function}[]
  {
    \clap,
    \llap,
    \rlap
  }
  \begin{syntax}
    \cs{rlap} \Arg{material}
  \end{syntax}
  %
  Puts the \meta{material} into an \cs{hbox} of zero width together
  with \cs{hss} on one or both sides. As a result, the material sticks
  out of the box without taking up any size in \LaTeX{}'s eyes. Thus
  one can overprint material or have something sticking out into the
  margin, etc. However, as this is using an \cs{hbox} no paragraph is
  started so if this command is used just in front of a paragraph it
  will appear on a line before the paragraph---not where most people
  would expect it.

  The safer \LaTeX{} equivalent would be
  \cs{makebox}\oarg{position}\verb={0pt}=\Arg{material} as this would
  always start a paragraph, if necessary.
\end{function}





%===============================================

\section{The unsorted rest}

The headings are simply taken from the \texttt{latex.ltx} file to help
me with initial sorting. Eventually though would all vanish and get
replaced by headings based on functionality, because the two do
not necessarily match.


%-----------------------------------------------
\subsection{From File: ltdirchk.dtx}

\leavevmode


%-----------------------------------------------
\subsection{From File: ltplain.dtx}

\begin{function}[]
  {
    \catcode
  }
  \begin{syntax}
    \cs{catcode} `^^I=10  % ascii tab is a blank space
  \end{syntax}
  \TeX{} primitive to set the catcode of a character.
\end{function}

\begin{function}[]
  {
    \chardef
  }
  \begin{syntax}
    \cs{chardef}
  \end{syntax}

\end{function}


\begin{function}[]
  {
    \countdef
  }
  \begin{syntax}
    \cs{countdef}
  \end{syntax}

\end{function}


\begin{function}[]
  {
    \dimendef
  }
  \begin{syntax}
    \cs{dimendef}
  \end{syntax}

\end{function}



\begin{function}[]
  {
    \active
  }
  \begin{syntax}
    \cs{catcode} `\cs{~}=\cs{active}
  \end{syntax}
  A plain \TeX{} macro representing the number \texttt{13} which is
  the catcode for \emph{active} characters.
\end{function}


\begin{function}[]
  {
    \dospecials
  }
  \begin{syntax}
    \def\do#1{...}    \cs{dospecials}
  \end{syntax}
  \cs{dospecials} calls \cs{do} for a list of typical \LaTeX syntax characters so that with an appropriate definition of \cs{do} something can be to for each of them. The list is:
  \verb=\  \\ \{ \} \$ \& \# \^ \_ \% \~ \^^I=
\end{function}


%-----------------------------------------------
\subsubsection{Handling \TeX{}-level register allocations}

\begin{function}[int]
  {
    \insc@unt,
    \allocationnumber
  }
  \begin{syntax}
    \cs{insc@unt}
  \end{syntax}
  Internal counters to mange \TeX{} registers.
\end{function}

\begin{function}[int]
  {
    \alloc@
  }
  \begin{syntax}
    \cs{alloc@}
  \end{syntax}

\end{function}

\begin{function}[int]
  {
    \e@alloc
  }
  \begin{syntax}
    \cs{e@alloc}
  \end{syntax}

\end{function}

\begin{function}[int]
  {
    \e@alloc@top
  }
  \begin{syntax}
    \cs{e@alloc@top}
  \end{syntax}

\end{function}



\begin{function}[int]
  {
    \e@ch@ck
  }
  \begin{syntax}
    \cs{e@ch@ck}
  \end{syntax}

\end{function}



\begin{function}[int]
  {
    \e@insert@top
  }
  \begin{syntax}
    \cs{e@insert@top}
  \end{syntax}

\end{function}







%-----------------------------------------------
\subsection{From File: ltvers.dtx}

\begin{function}[]
  {
    \fmtname
  }
  \begin{syntax}
    \cs{fmtname}
  \end{syntax}

\end{function}

\begin{function}[]
  {
    \development@branch@name
  }
  \begin{syntax}
    \cs{development@branch@name}
  \end{syntax}

\end{function}


\begin{function}[]
  {
    \LaTeXReleaseInfo
  }
  \begin{syntax}
    \cs{LaTeXReleaseInfo}
  \end{syntax}

\end{function}


\begin{function}[]
  {
    \DeclareCurrentRelease
  }
  \begin{syntax}
    \cs{DeclareCurrentRelease}
  \end{syntax}

\end{function}


\begin{function}[]
  {
    \DeclareRelease
  }
  \begin{syntax}
    \cs{DeclareRelease}
  \end{syntax}

\end{function}

\begin{function}[]
  {
    \EndIncludeInRelease
  }
  \begin{syntax}
    \cs{EndIncludeInRelease}
  \end{syntax}

\end{function}



\begin{function}[]
  {
    \NewModuleRelease
  }
  \begin{syntax}
    \cs{NewModuleRelease}
  \end{syntax}

\end{function}




\begin{function}[]
  {
    \EndModuleRelease
  }
  \begin{syntax}
    \cs{EndModuleRelease}
  \end{syntax}

\end{function}




\begin{function}[int]
  {
    \if@skipping@module,
    \@skipping@modulefalse,
    \@skipping@moduletrue
  }
  \begin{syntax}
    \cs{if@skipping@module} ... \cs{else} ... \cs{fi}
  \end{syntax}

\end{function}


\begin{function}[int]
  {
    \finish@module@release
  }
  \begin{syntax}
    \cs{finish@module@release}
  \end{syntax}

\end{function}


\begin{function}[int]
  {
    \gobble@finish@module@release
  }
  \begin{syntax}
    \cs{gobble@finish@module@release}
  \end{syntax}

\end{function}

\begin{function}[int]
  {
    \new@module@skip
  }
  \begin{syntax}
    \cs{new@module@skip}
  \end{syntax}

\end{function}



\begin{function}[int]
  {
    \new@moduledate
  }
  \begin{syntax}
    \cs{new@moduledate}
  \end{syntax}

\end{function}



\begin{function}[int]
  {
    \new@modulename
  }
  \begin{syntax}
    \cs{new@modulename}
  \end{syntax}

\end{function}



\begin{function}[int]
  {
    \@IncludeInRelease
  }
  \begin{syntax}
    \cs{@IncludeInRelease}
  \end{syntax}

\end{function}

\begin{function}[int]
  {
    \@check@IncludeInRelease
  }
  \begin{syntax}
    \cs{@check@IncludeInRelease}
  \end{syntax}

\end{function}



\begin{function}[int]
  {
    \@end@check@IncludeInRelease
  }
  \begin{syntax}
    \cs{@end@check@IncludeInRelease}
  \end{syntax}

\end{function}

\begin{function}[int]
  {
    \@gobble@IncludeInRelease
  }
  \begin{syntax}
    \cs{@gobble@IncludeInRelease}
  \end{syntax}

\end{function}




\begin{function}[user]
  {
    \filecontents,
    \endfilecontents
  }
  \begin{syntax}
    \cs{begin}\Arg{filecontents}\oarg{options}\Arg{filename}
    \ \  ...
    \cs{end}\Arg{filecontents}
  \end{syntax}

\end{function}


%-----------------------------------------------
\subsubsection{Requesting a specific release}

\begin{function}[]
  {
    \requestedLaTeXdate
  }
  \begin{syntax}
    \cs{requestedLaTeXdate}
  \end{syntax}

\end{function}



\begin{function}[]
  {
    \IfTargetDateBefore
  }
  \begin{syntax}
    \cs{IfTargetDateBefore} \Arg{target date} \Arg{true branch} \Arg{false branch}
  \end{syntax}

\end{function}


\begin{function}[int]
  {
    \pkgcls@arg
  }
  \begin{syntax}
    \cs{pkgcls@arg}
  \end{syntax}

\end{function}



\begin{function}[int]
  {
    \pkgcls@candidate
  }
  \begin{syntax}
    \cs{pkgcls@candidate}
  \end{syntax}

\end{function}



\begin{function}[int]
  {
    \pkgcls@debug
  }
  \begin{syntax}
    \cs{pkgcls@debug}
  \end{syntax}

\end{function}



\begin{function}[int]
  {
    \pkgcls@ext
  }
  \begin{syntax}
    \cs{pkgcls@ext}
  \end{syntax}

\end{function}



\begin{function}[int]
  {
    \pkgcls@innerdate
  }
  \begin{syntax}
    \cs{pkgcls@innerdate}
  \end{syntax}

\end{function}



\begin{function}[int]
  {
    \pkgcls@mindate
  }
  \begin{syntax}
    \cs{pkgcls@mindate}
  \end{syntax}

\end{function}



\begin{function}[int]
  {
    \pkgcls@name
  }
  \begin{syntax}
    \cs{pkgcls@name}
  \end{syntax}

\end{function}



\begin{function}[int]
  {
    \pkgcls@parse@date@arg
  }
  \begin{syntax}
    \cs{pkgcls@parse@date@arg}
  \end{syntax}

\end{function}



\begin{function}[int]
  {
    \pkgcls@parse@date@arg@
  }
  \begin{syntax}
    \cs{pkgcls@parse@date@arg@}
  \end{syntax}

\end{function}



\begin{function}[int]
  {
    \pkgcls@parse@date@arg@version
  }
  \begin{syntax}
    \cs{pkgcls@parse@date@arg@version}
  \end{syntax}

\end{function}



\begin{function}[int]
  {
    \pkgcls@releasedate
  }
  \begin{syntax}
    \cs{pkgcls@releasedate}
  \end{syntax}

\end{function}



\begin{function}[int]
  {
    \pkgcls@rollbackdate@error
  }
  \begin{syntax}
    \cs{pkgcls@rollbackdate@error}
  \end{syntax}

\end{function}



\begin{function}[int]
  {
    \pkgcls@show@selection
  }
  \begin{syntax}
    \cs{pkgcls@show@selection}
  \end{syntax}

\end{function}



\begin{variable}[int]
  {
    \pkgcls@targetdate
  }
  %
  Internal \TeX{} counter register
%  holding zero (no target release
%  date requested) or the target date in ISO format but with all
%  \texttt{-} stripped, e.g., millenium's New Year would be
%  \texttt{20000101}.

\end{variable}



\begin{function}[int]
  {
    \pkgcls@targetlabel
  }
  \begin{syntax}
    \cs{pkgcls@targetlabel}
  \end{syntax}

\end{function}



\begin{function}[int]
  {
    \pkgcls@use@this@release
  }
  \begin{syntax}
    \cs{pkgcls@use@this@release}
  \end{syntax}

\end{function}



%-----------------------------------------------
\subsection{From File: ltluatex.dtx}

\begin{variable}[]
  {
    \luatexversion
  }

\end{variable}



\begin{function}[]
  {
    \@gobble,
    \@gobbletwo,
    \@gobblethree,
    \@gobblefour
  }
  \begin{syntax}
    \cs{@gobble} \Arg{arg}
    \cs{@gobblefour}\Arg{arg\textsubscript{1}}\Arg{arg\textsubscript{2}}\Arg{arg\textsubscript{3}}\Arg{arg\textsubscript{4}}
  \end{syntax}

\end{function}



\begin{function}[]
  {
    \setattribute,
    \unsetattribute
  }
  \begin{syntax}
    \cs{setattribute}
  \end{syntax}

\end{function}




\begin{function}[]
  {
    \newcatcodetable
  }
  \begin{syntax}
    \cs{newcatcodetable}
  \end{syntax}

\end{function}




\begin{variable}[]
  {
    \catcodetable@initex,
    \catcodetable@string,
    \catcodetable@latex,
    \catcodetable@atletter
  }

\end{variable}



\begin{function}[]
  {
    \setrangecatcode
  }
  \begin{syntax}
    \cs{setrangecatcode}
  \end{syntax}

\end{function}



\begin{function}[]
  {
    \newluafunction
  }
  \begin{syntax}
    \cs{newluafunction}
  \end{syntax}

\end{function}



\begin{function}[]
  {
    \newluacmd
  }
  \begin{syntax}
    \cs{newluacmd}
  \end{syntax}

\end{function}



\begin{function}[]
  {
    \newprotectedluacmd
  }
  \begin{syntax}
    \cs{newprotectedluacmd}
  \end{syntax}

\end{function}



\begin{function}[]
  {
    \newwhatsit
  }
  \begin{syntax}
    \cs{newwhatsit}
  \end{syntax}

\end{function}


\begin{function}[]
  {
    \newluabytecode
  }
  \begin{syntax}
    \cs{newluabytecode}
  \end{syntax}

\end{function}



\begin{function}[]
  {
    \newluachunkname
  }
  \begin{syntax}
    \cs{newluachunkname}
  \end{syntax}

\end{function}



\begin{function}[]
  {
    \now@and@everyjob
  }
  \begin{syntax}
    \cs{now@and@everyjob}
  \end{syntax}

\end{function}





\begin{variable}[int]
  {
    \e@alloc@attribute@count,
    \e@alloc@whatsit@count
  }

\end{variable}




%-----------------------------------------------
\subsubsection{Parsing Unicode data}


\begin{function}[int]
  {
    \parseunicodedataI,
    \parseunicodedataII,
    \parseunicodedataIII,
    \parseunicodedataIV,
    \parseunicodedataV
  }

\end{function}



\begin{function}[int]
  {
    \storedpar
  }
  \begin{syntax}
    \cs{storedpar}
  \end{syntax}

\end{function}



\begin{function}[int]
  {
    \unicoderead
  }
  Read stream for parsing Unicode data.
\end{function}



%-----------------------------------------------
\subsection{From File: ltexpl.dtx}



\begin{function}[int]
  {
    \@kernel@after@enddocument
  }
  \begin{syntax}
    \cs{@kernel@after@enddocument}
  \end{syntax}

\end{function}


\begin{function}[int]
  {
    \@kernel@after@enddocument@afterlastpage
  }
  \begin{syntax}
    \cs{@kernel@after@enddocument@afterlastpage}
  \end{syntax}

\end{function}


\begin{function}[int]
  {
    \@kernel@before@begindocument
  }
  \begin{syntax}
    \cs{@kernel@before@begindocument}
  \end{syntax}

\end{function}



\begin{function}[int]
  {
    \@kernel@after@begindocument
  }
  \begin{syntax}
    \cs{@kernel@after@begindocument}
  \end{syntax}

\end{function}



\begin{function}[int]
  {
    \@kernel@after@begindocument@before
  }
  \begin{syntax}
    \cs{@kernel@after@begindocument@before}
  \end{syntax}

\end{function}

\begin{function}[int]
  {
    \@kernel@before@enddocument@afterlastpage
  }
  \begin{syntax}
    \cs{@kernel@before@enddocument@afterlastpage}
  \end{syntax}

\end{function}


\begin{function}[int]
  {
    \@expl@sys@load@backend@@
  }
  \begin{syntax}
    \cs{@expl@sys@load@backend@@}
  \end{syntax}

\end{function}


\begin{function}[int]
  {
    \@expl@push@filename@@
  }
  \begin{syntax}
    \cs{@expl@push@filename@@}
  \end{syntax}

\end{function}


\begin{function}[int]
  {
    \@expl@push@filename@aux@@
  }
  \begin{syntax}
    \cs{@expl@push@filename@aux@@}
  \end{syntax}

\end{function}


\begin{function}[int]
  {
    \@expl@pop@filename@@
  }
  \begin{syntax}
    \cs{@expl@pop@filename@@}
  \end{syntax}

\end{function}



\begin{function}[user]
  {
    \IfFileExists
  }
  \begin{syntax}
    \cs{IfFileExists}
  \end{syntax}

\end{function}



\begin{function}[int]
  {
    \IfFileExists@,
    \IfFileExists@@
  }

\end{function}






%-----------------------------------------------
\subsection{From File: ltdefns.dtx}

\leavevmode


%-----------------------------------------------
\subsection{From File: ltcmd.dtx}

\leavevmode


%-----------------------------------------------
\subsection{From File: lthooks.dtx}

\leavevmode


%-----------------------------------------------
\subsection{From File: ltcmdhooks.dtx}

\leavevmode


%-----------------------------------------------
\subsection{From File: lthooks.dtx}

\leavevmode


%-----------------------------------------------
\subsection{From File: ltsockets.dtx}

\leavevmode


%-----------------------------------------------
\subsection{From File: lttemplates.dtx}

\leavevmode


%-----------------------------------------------
\subsection{From File: ltalloc.dtx}

\leavevmode


%-----------------------------------------------
\subsection{From File: ltcntrl.dtx}

\leavevmode


%-----------------------------------------------
\subsection{From File: lterror.dtx}

\leavevmode


%-----------------------------------------------
\subsection{From File: ltpar.dtx}

\leavevmode


%-----------------------------------------------
\subsection{From File: ltpara.dtx}

\leavevmode


%-----------------------------------------------
\subsection{From File: ltmeta.dtx}

\leavevmode


%-----------------------------------------------
\subsection{From File: ltspace.dtx}



\begin{function}[user]
  {
    \bigskip,
    \medskip,
    \smallskip
  }
  \begin{syntax}
    \cs{bigskip}
  \end{syntax}

\end{function}



\begin{variable}[user]
  {
    \bigskipamount,
    \medskipamount,
    \smallskipamount
  }

\end{variable}



\begin{function}[user]
  {
    \nopagebreak
  }
  \begin{syntax}
    \cs{nopagebreak}
  \end{syntax}

\end{function}


\begin{function}[user]
  {
    \pagebreak
  }
  \begin{syntax}
    \cs{pagebreak}
  \end{syntax}

\end{function}

\begin{function}[user]
  {
    \linebreak
  }
  \begin{syntax}
    \cs{linebreak}
  \end{syntax}

\end{function}



\begin{function}[user]
  {
    \nolinebreak
  }
  \begin{syntax}
    \cs{nolinebreak}
  \end{syntax}

\end{function}


\begin{function}[user]
  {
    \samepage
  }
  \begin{syntax}
    \cs{samepage}
  \end{syntax}

\end{function}


\begin{function}[int]
  {
    \@no@pgbk
  }
  \begin{syntax}
    \cs{@no@pgbk}
  \end{syntax}

\end{function}

\begin{function}[int]
  {
    \@no@lnbk
  }
  \begin{syntax}
    \cs{@no@lnbk}
  \end{syntax}

\end{function}

\begin{function}[]
  {
    \@normalcr
  }
  \begin{syntax}
    \cs{let} \cs{\textbackslash} \cs{@normalcr}
  \end{syntax}
  Default definition for \verb=\\=.
\end{function}


\begin{function}[int]
  {
    \@vspace@calcify
  }
  \begin{syntax}
    \cs{@vspace@calcify}
  \end{syntax}
  Internal macro to support \pkg{calc} syntax in arguments of
  \cs{vspace} and similar places.
\end{function}



\begin{function}[int]
  {
    \@newline,
    \@gnewline,
    \@xnewline
  }
  Internal macros implementing \cs{newline}. Might vanish!
\end{function}




\begin{function}[]
  {
    \@getpen
  }
  \begin{syntax}
    \cs{@getpen} \Arg{number}
  \end{syntax}
  \LaTeXe's way of turning the number 0--4 into zero, low, medium, or high penalty.
\end{function}





\begin{function}[]
  {
    \if@skipping@module,
    \@skipping@modulefalse,
    \@skipping@moduletrue
  }
  \begin{syntax}
    \cs{if@skipping@module} 
  \end{syntax}

\end{function}



\begin{function}[]
  {
    \if@nobreak,
    \@nobreakfalse,
    \@nobreaktrue
  }
  \begin{syntax}
    \cs{if@nobreak} ... \cs{else} ... \cs{fi}
  \end{syntax}
  Despite its name this is a global \TeX{} switch. True when when
  \LaTeX{} signals that an upcoming breakpoint should be suppressed,
  e.g., directly after a heading etc. Set and checked in various places.
\end{function}


\begin{function}[]
  {
    \@bsphack,
    \@esphack
  }
  \begin{syntax}
    \cs{@bsphack} \meta{material} {@esphack}
  \end{syntax}

\end{function}



\begin{function}[]
  {
    \@Esphack
  }
  \begin{syntax}
    \cs{@bsphack} \meta{material} \cs{@Esphack}
  \end{syntax}
  Variant of \cs{@esphack} used in environments that want to remain
  invisible.

  This needs some checking if it is still needed (and/or  correct)!
\end{function}



\begin{function}[int]
  {
    \@savsk,
    \@savsf
  }
%
  Internal dimen and count registers used in the implementation of
  \cs{@bsphack} and \cs{@esphack}. Not for direct use.
\end{function}



\begin{function}[user]
  {
    \addvspace
  }
  \begin{syntax}
    \cs{addvspace}
  \end{syntax}

\end{function}


\begin{function}[user]
  {
    \addpenalty
  }
  \begin{syntax}
    \cs{addpenalty}
  \end{syntax}

\end{function}


\begin{function}[user]
  {
    \vspace
  }
  \begin{syntax}
    \cs{vspace}
  \end{syntax}

\end{function}


\begin{function}[user]
  {
    \nobreakdashes
  }
  \begin{syntax}
    \cs{nobreakdashes}
  \end{syntax}

\end{function}


\begin{function}[user]
  {
    \nobreakspace
  }
  \begin{syntax}
    \cs{nobreakspace}
  \end{syntax}

\end{function}


\begin{function}[user]
  {
    \@
  }
  \begin{syntax}
    \cs{@}
  \end{syntax}

\end{function}

\begin{function}[user]
  {
    \hspace
  }
  \begin{syntax}
    \cs{hspace}
  \end{syntax}

\end{function}





\begin{function}[int]
  {
    \@xobeysp
  }
  \begin{syntax}
    \cs{@xobeysp}
  \end{syntax}

\end{function}

\begin{function}[int]
  {
    \@xobeytab
  }
  \begin{syntax}
    \cs{@xobeytab}
  \end{syntax}

\end{function}



\begin{function}[int]
  {
    \@xaddvskip
  }
  \begin{syntax}
    \cs{@xaddvskip}
  \end{syntax}

\end{function}


\begin{function}[int]
  {
    \@vspace,
    \@vspacer
  }
  Internal commands to implement \cs{vspace}. Might vanish.
\end{function}


\begin{function}[int]
  {
    \@hspace,
    \@hspacer
  }
  Internal commands to implement \cs{hspace}. Might vanish.
\end{function}







%-----------------------------------------------
\subsection{From File: ltlogos.dtx}


\begin{function}[user]
  {
    \LaTeX,
    \LaTeXe,
    \TeX
  }
  \begin{syntax}
    \cs{LaTeX}
  \end{syntax}

\end{function}




%-----------------------------------------------
\subsection{From File: ltfiles.dtx}

\leavevmode


%-----------------------------------------------
\subsection{From File: ltoutenc.dtx}

\leavevmode


%-----------------------------------------------
\subsection{From File: ltcounts.dtx}

\leavevmode


%-----------------------------------------------
\subsection{From File: ltlength.dtx}

\leavevmode


%-----------------------------------------------
\subsection{From File: ltfssbas.dtx}

\leavevmode


%-----------------------------------------------
\subsection{From File: ltfssaxes.dtx}

\leavevmode


%-----------------------------------------------
\subsection{From File: ltfsstrc.dtx}

\leavevmode


%-----------------------------------------------
\subsection{From File: ltfssdcl.dtx}

\leavevmode


%-----------------------------------------------
\subsection{From File: ltfssini.dtx}

\leavevmode


%-----------------------------------------------
\subsection{From File: ltfntcmd.dtx}

\leavevmode


%-----------------------------------------------
\subsection{From File: lttextcomp.dtx}

\leavevmode


%-----------------------------------------------
\subsection{From File: ltpageno.dtx}

\leavevmode


%-----------------------------------------------
\subsection{From File: ltxref.dtx}

\leavevmode


%-----------------------------------------------
\subsection{From File: ltproperties.dtx}

\leavevmode


%-----------------------------------------------
\subsection{From File: ltmiscen.dtx}


\begin{function}[user]
  {
    \begin,
    \end
  }
  \begin{syntax}
    \cs{begin}\Arg{name}\meta{argument(s)} ... \cs{end}\Arg{name}
  \end{syntax}

\end{function}


\begin{function}[]
  {
    \if@ignore,
    \@ignorefalse,
    \@ignoretrue
  }
  \begin{syntax}
    \cs{if@ignore} ... \cs{else} ... \cs{fi}
  \end{syntax}
  %
  If this \TeX{}-level boolean is set to \texttt{true} then spaces
  after the next \cs{end} are ignored, i.e., \cs{ignorespaces} is
  applied there.
  
  This \TeX{}-level boolean is globally set, i.e., \cs{@ignorefalse}
  and \cs{@ignoretrue} change it globally.
\end{function}


\begin{function}[user]
  {
    \ignorespacesafterend
  }
  \begin{syntax}
    \cs{ignorespacesafterend}
  \end{syntax}
  %
  For use in the code executed at the end of an environment, a public
  name for \cs{@ignoretrue}.
\end{function}




\begin{variable}[]
  {
    \@currenvir
  }
%
  Name of the current environment, i.e., the \meta{name} from the
  inner most open \cs{begin}\Arg{name}.
\end{variable}



\begin{function}[user]
  {
    \document,
    \enddocument
  }
  \begin{syntax}
    \cs{begin}\Arg{document}  \meta{document body} \cs{end}\Arg{document}
  \end{syntax}

\end{function}




%-----------------------------------------------
\subsection{From File: ltmath.dtx}

\begin{function}[]
  {
    \medspace
  }
  \begin{syntax}
    \cs{medspace}
  \end{syntax}

\end{function}





%-----------------------------------------------
\subsection{From File: ltlists.dtx}


\begin{function}[]
  {
    \list,
    \endlist
  }
  \begin{syntax}
    \cs{begin}\Arg{list}\Arg{default-label}\Arg{declarations}  \cs{item} ... \cs{end}\Arg{list}
  \end{syntax}

\end{function}


\begin{variable}[]
  {
    \topsep,
    \partopsep,
    \itemsep,
    \parsep
  }

\end{variable}


\begin{variable}[]
  {
    \leftmargin,
    \rightmargin,
    \listparindent,
  }
  \begin{syntax}
    \cs{leftmargin}
  \end{syntax}

\end{variable}




\begin{variable}[]
  {
    \leftmargini,
    \leftmarginii,
    \leftmarginiii,
    \leftmarginiv,
    \leftmarginv,
    \leftmarginvi
  }

\end{variable}



\begin{variable}[]
  {
    \@beginparpenalty,
    \@itempenalty,
    \@endparpenalty,
  }

\end{variable}


\begin{variable}[]
  {
    \@listdepth
  }

\end{variable}



\begin{variable}[int]
  {
    \@topsep
  }

\end{variable}

\begin{variable}[int]
  {
    \@topsepadd
  }

\end{variable}



\begin{function}[int,deprecated]
  {
    \if@noparitem,
    \@noparitemfalse,
    \@noparitemtrue
  }
  \begin{syntax}
    \cs{if@noparitem} ...
  \end{syntax}
  %
  No longer used in the new template-based list implementations.
\end{function}



%-----------------------------------------------
\subsection{From File: ltboxes.dtx}

\leavevmode


%-----------------------------------------------
\subsection{From File: lttab.dtx}

\leavevmode


%-----------------------------------------------
\subsection{From File: ltpictur.dtx}

\leavevmode


%-----------------------------------------------
\subsection{From File: ltthm.dtx}

\leavevmode


%-----------------------------------------------
\subsection{From File: ltsect.dtx}

\leavevmode


%-----------------------------------------------
\subsection{From File: ltfloat.dtx}


\begin{function}[]
  {
    \extrafloats
  }
  \begin{syntax}
    \cs{extrafloats}
  \end{syntax}

\end{function}



\begin{variable}[]
  {
    \c@bottomnumber
  }
  
\end{variable}



%-----------------------------------------------
\subsection{From File: ltidxglo.dtx}

\leavevmode


%-----------------------------------------------
\subsection{From File: ltbibl.dtx}

\leavevmode


%-----------------------------------------------
\subsection{From File: ltmarks.dtx}

\leavevmode


%-----------------------------------------------
\subsection{From File: ltpage.dtx}

\leavevmode


%-----------------------------------------------
\subsection{From File: ltclass.dtx}

\leavevmode


%-----------------------------------------------
\subsection{From File: ltkeys.dtx}

\leavevmode


%-----------------------------------------------
\subsection{From File: ltfilehook.dtx}

\leavevmode


%-----------------------------------------------
\subsection{From File: ltshipout.dtx}

\leavevmode


%-----------------------------------------------
\subsection{From File: ltoutput.dtx}

\leavevmode


%-----------------------------------------------
\subsection{From File: lttagging.dtx}


\begin{function}[]
  {
    \SuspendTagging,
    \ResumeTagging
  }
  \begin{syntax}
    \cs{SuspendTagging} \meta{material} \cs{ResumeTagging}
  \end{syntax}

\end{function}


\begin{function}[]
  {
    \NewTaggingSocket
  }
  \begin{syntax}
    \cs{NewTaggingSocket}
  \end{syntax}

\end{function}



\begin{function}[]
  {
    \NewTaggingSocketPlug
  }
  \begin{syntax}
    \cs{NewTaggingSocketPlug}
  \end{syntax}

\end{function}


\begin{function}[]
  {
    \AssignTaggingSocketPlug
  }
  \begin{syntax}
    \cs{AssignTaggingSocketPlug}
  \end{syntax}

\end{function}


\begin{function}[]
  {
    \UseTaggingSocket
  }
  \begin{syntax}
    \cs{UseTaggingSocket}
  \end{syntax}

\end{function}


\begin{function}[]
  {
    \UseExpandableTaggingSocket
  }
  \begin{syntax}
    \cs{UseExpandableTaggingSocket}
  \end{syntax}

\end{function}


\begin{function}[]
  {
    \MathCollectFalse,
    \MathCollectTrue
  }
  \begin{syntax}
    \cs{MathCollectFalse}
  \end{syntax}

\end{function}





\begin{function}[user]
  {
    \MathMLintent,
    \MathMLarg
  }
  \begin{syntax}
    \cs{MathMLintent}
  \end{syntax}

\end{function}



\begin{function}[]
  {
    \NewStructureName
  }
  \begin{syntax}
    \cs{NewStructureName}
  \end{syntax}

\end{function}


\begin{function}[]
  {
    \UseStructureName
  }
  \begin{syntax}
    \cs{UseStructureName}
  \end{syntax}

\end{function}



\begin{function}[user]
  {
    \DebugTablesOn,
    \DebugTablesOff
  }
  \begin{syntax}
    \cs{DebugTablesOn}
  \end{syntax}

\end{function}


\begin{function}[int]
  {
    \hyper@nopatch@longtable
  }

\end{function}



%-----------------------------------------------
\subsection{From File: ltfinal.dtx}


\begin{variable}[]
  {
    \@lowpenalty,
    \@medpenalty,
    \@highpenalty
  }
  \begin{syntax}
    \cs{@lowpenalty}
  \end{syntax}

\end{variable}


\begin{function}[]
  {
    \newXeTeXintercharclass
  }
  \begin{syntax}
    \cs{newXeTeXintercharclass}
  \end{syntax}

\end{function}


\begin{function}[]
  {
    \XeTeXuseglyphmetrics
  }
  \begin{syntax}
    \cs{XeTeXuseglyphmetrics}
  \end{syntax}

\end{function}

\begin{function}[]
  {
    \XeTeXdashbreakstate
  }
  \begin{syntax}
    \cs{XeTeXdashbreakstate}
  \end{syntax}

\end{function}


\begin{function}[]
  {
    \IfPDFManagementActiveTF,
    \IfPDFManagementActiveT,
    \IfPDFManagementActiveF
  }
  \begin{syntax}
    \cs{IfPDFManagementActiveTF}
  \end{syntax}

\end{function}


\begin{function}[]
  {
    \dump
  }
  \begin{syntax}
    \cs{dump}
  \end{syntax}

\end{function}




\begin{variable}[int]
  {
    \xe@alloc@intercharclass
  }

\end{variable}

\begin{function}[int]
  {
    \@expl@finalise@setup@@
  }
  \begin{syntax}
    \cs{@expl@finalise@setup@@}
  \end{syntax}

\end{function}



\begin{function}[int]
  {
    \@addtofilelist,
    \@addtofilelist@aux
  }

\end{function}





%-----------------------------------------------


%-----------------------------------------------
\subsection{All remaining commands in alphabetical order}


\begin{function}[user]
  {
    \!
  }
  \begin{syntax}
    \cs{!}
  \end{syntax}

\end{function}















\begin{function}[]
  {
    \*
  }
  \begin{syntax}
    \cs{*}
  \end{syntax}

\end{function}



\begin{function}[]
  {
    \+
  }
  \begin{syntax}
    \cs{+}
  \end{syntax}

\end{function}



\begin{function}[]
  {
    \,
  }
  \begin{syntax}
    \cs{,}
  \end{syntax}

\end{function}



\begin{function}[]
  {
    \-
  }
  \begin{syntax}
    \cs{-}
  \end{syntax}

\end{function}



\begin{function}[]
  {
    \/
  }
  \begin{syntax}
    \cs{/}
  \end{syntax}

\end{function}



\begin{function}[]
  {
    \:
  }
  \begin{syntax}
    \cs{:}
  \end{syntax}

\end{function}



\begin{function}[]
  {
    \;
  }
  \begin{syntax}
    \cs{;}
  \end{syntax}

\end{function}



\begin{function}[]
  {
    \<
  }
  \begin{syntax}
    \cs{<}
  \end{syntax}

\end{function}



\begin{function}[]
  {
    \=
  }
  \begin{syntax}
    \cs{=}
  \end{syntax}

\end{function}



\begin{function}[]
  {
    \>
  }
  \begin{syntax}
    \cs{>}
  \end{syntax}

\end{function}



\begin{function}[]
  {
    \?
  }
  \begin{syntax}
    \cs{?}
  \end{syntax}

\end{function}



\begin{function}[]
  {
    \@
  }
  \begin{syntax}
    \cs{@}
  \end{syntax}

\end{function}



\begin{function}[]
  {
    \@@
  }
  \begin{syntax}
    \cs{@@}
  \end{syntax}

\end{function}



\begin{function}[]
  {
    \@@data@glyphtounicode
  }
  \begin{syntax}
    \cs{@@data@glyphtounicode}
  \end{syntax}

\end{function}



\begin{function}[]
  {
    \@@data@glyphtounicode@cmex
  }
  \begin{syntax}
    \cs{@@data@glyphtounicode@cmex}
  \end{syntax}

\end{function}



\begin{function}[]
  {
    \@@enc@update
  }
  \begin{syntax}
    \cs{@@enc@update}
  \end{syntax}

\end{function}



\begin{function}[]
  {
    \@@end
  }
  \begin{syntax}
    \cs{@@end}
  \end{syntax}

\end{function}



\begin{function}[]
  {
    \@@endpbox
  }
  \begin{syntax}
    \cs{@@endpbox}
  \end{syntax}

\end{function}



\begin{function}[]
  {
    \@@eqncr
  }
  \begin{syntax}
    \cs{@@eqncr}
  \end{syntax}

\end{function}



\begin{function}[]
  {
    \@@fileswith@pti@ns
  }
  \begin{syntax}
    \cs{@@fileswith@pti@ns}
  \end{syntax}

\end{function}



\begin{function}[]
  {
    \@@hyph
  }
  \begin{syntax}
    \cs{@@hyph}
  \end{syntax}

\end{function}



\begin{function}[]
  {
    \@@if@newlist
  }
  \begin{syntax}
    \cs{@@if@newlist}
  \end{syntax}

\end{function}



\begin{function}[]
  {
    \@@ifdefinable
  }
  \begin{syntax}
    \cs{@@ifdefinable}
  \end{syntax}

\end{function}



\begin{function}[]
  {
    \@@input
  }
  \begin{syntax}
    \cs{@@input}
  \end{syntax}

\end{function}



\begin{function}[]
  {
    \@@italiccorr
  }
  \begin{syntax}
    \cs{@@italiccorr}
  \end{syntax}

\end{function}




\begin{function}[]
  {
    \@@math@bgroup
  }
  \begin{syntax}
    \cs{@@math@bgroup}
  \end{syntax}

\end{function}



\begin{function}[]
  {
    \@@par
  }
  \begin{syntax}
    \cs{@@par}
  \end{syntax}

\end{function}



\begin{function}[]
  {
    \@@protect
  }
  \begin{syntax}
    \cs{@@protect}
  \end{syntax}

\end{function}



\begin{function}[]
  {
    \@@startpbox
  }
  \begin{syntax}
    \cs{@@startpbox}
  \end{syntax}

\end{function}



\begin{function}[]
  {
    \@@sverb
  }
  \begin{syntax}
    \cs{@@sverb}
  \end{syntax}

\end{function}



\begin{function}[]
  {
    \@@text@case@aux
  }
  \begin{syntax}
    \cs{@@text@case@aux}
  \end{syntax}

\end{function}



\begin{function}[int]
  {
    \@@underline
  }
  \begin{syntax}
    \cs{@@underline}
  \end{syntax}

  Saved version of the \TeX{} primitve \tn{underline}. \LaTeX{}
  replaces the primitive with a robust version that also works in text
  mode.

\end{function}



\begin{function}[]
  {
    \@@unprocessedoptions
  }
  \begin{syntax}
    \cs{@@unprocessedoptions}
  \end{syntax}

\end{function}



\begin{function}[int,deprecated]
  {
    \@@warning
  }
  \begin{syntax}
    \cs{@@warning} \Arg{message}
  \end{syntax}

  Another name for \cs{@latex@warning@no@line}. Deprecated.
\end{function}



\begin{function}[]
  {
    \@@write
  }
  \begin{syntax}
    \cs{@@write}
  \end{syntax}

\end{function}



\begin{function}[]
  {
    \@Alph
  }
  \begin{syntax}
    \cs{@Alph}
  \end{syntax}

\end{function}



\begin{function}[]
  {
    \@DeclareEncodingSubset
  }
  \begin{syntax}
    \cs{@DeclareEncodingSubset}
  \end{syntax}

\end{function}



\begin{function}[]
  {
    \@DeclareMathDelimiter
  }
  \begin{syntax}
    \cs{@DeclareMathDelimiter}
  \end{syntax}

\end{function}



\begin{function}[]
  {
    \@DeclareMathSizes
  }
  \begin{syntax}
    \cs{@DeclareMathSizes}
  \end{syntax}

\end{function}





\begin{function}[]
  {
    \@IncludeInRele@se
  }
  \begin{syntax}
    \cs{@IncludeInRele@se}
  \end{syntax}

\end{function}





\begin{function}[]
  {
    \@LRmoderr
  }
  \begin{syntax}
    \cs{@LRmoderr}
  \end{syntax}

\end{function}



\begin{function}[]
  {
    \@M
  }
  \begin{syntax}
    \cs{@M}
  \end{syntax}

\end{function}



\begin{function}[]
  {
    \@MM
  }
  \begin{syntax}
    \cs{@MM}
  \end{syntax}

\end{function}



\begin{function}[]
  {
    \@Mi
  }
  \begin{syntax}
    \cs{@Mi}
  \end{syntax}

\end{function}



\begin{function}[]
  {
    \@Mii
  }
  \begin{syntax}
    \cs{@Mii}
  \end{syntax}

\end{function}



\begin{function}[]
  {
    \@Miii
  }
  \begin{syntax}
    \cs{@Miii}
  \end{syntax}

\end{function}



\begin{function}[]
  {
    \@Miv
  }
  \begin{syntax}
    \cs{@Miv}
  \end{syntax}

\end{function}



\begin{function}[]
  {
    \@Roman
  }
  \begin{syntax}
    \cs{@Roman}
  \end{syntax}

\end{function}



\begin{function}[]
  {
    \@TeXversion
  }
  \begin{syntax}
    \cs{@TeXversion}
  \end{syntax}

\end{function}



\begin{function}[]
  {
    \@abspage@last
  }
  \begin{syntax}
    \cs{@abspage@last}
  \end{syntax}

\end{function}



\begin{function}[]
  {
    \@acci
  }
  \begin{syntax}
    \cs{@acci}
  \end{syntax}

\end{function}



\begin{function}[]
  {
    \@accii
  }
  \begin{syntax}
    \cs{@accii}
  \end{syntax}

\end{function}



\begin{function}[]
  {
    \@acciii
  }
  \begin{syntax}
    \cs{@acciii}
  \end{syntax}

\end{function}



\begin{function}[]
  {
    \@acol
  }
  \begin{syntax}
    \cs{@acol}
  \end{syntax}

\end{function}



\begin{function}[]
  {
    \@acolampacol
  }
  \begin{syntax}
    \cs{@acolampacol}
  \end{syntax}

\end{function}



\begin{function}[]
  {
    \@activechar@info
  }
  \begin{syntax}
    \cs{@activechar@info}
  \end{syntax}

\end{function}



\begin{function}[]
  {
    \@addamp
  }
  \begin{syntax}
    \cs{@addamp}
  \end{syntax}

\end{function}



\begin{function}[]
  {
    \@addfield
  }
  \begin{syntax}
    \cs{@addfield}
  \end{syntax}

\end{function}



\begin{function}[]
  {
    \@addmarginpar
  }
  \begin{syntax}
    \cs{@addmarginpar}
  \end{syntax}

\end{function}



\begin{function}[]
  {
    \@addtobot
  }
  \begin{syntax}
    \cs{@addtobot}
  \end{syntax}

\end{function}



\begin{function}[]
  {
    \@addtocurcol
  }
  \begin{syntax}
    \cs{@addtocurcol}
  \end{syntax}

\end{function}



\begin{function}[]
  {
    \@addtodblcol
  }
  \begin{syntax}
    \cs{@addtodblcol}
  \end{syntax}

\end{function}



\begin{function}[]
  {
    \@addtofilelist
  }
  \begin{syntax}
    \cs{@addtofilelist}
  \end{syntax}

\end{function}



\begin{function}[]
  {
    \@addtofilelist@aux
  }
  \begin{syntax}
    \cs{@addtofilelist@aux}
  \end{syntax}

\end{function}



\begin{function}[]
  {
    \@addtonextcol
  }
  \begin{syntax}
    \cs{@addtonextcol}
  \end{syntax}

\end{function}



\begin{function}[]
  {
    \@addtopreamble
  }
  \begin{syntax}
    \cs{@addtopreamble}
  \end{syntax}

\end{function}



\begin{function}[]
  {
    \@addtoreset
  }
  \begin{syntax}
    \cs{@addtoreset}
  \end{syntax}

\end{function}



\begin{function}[]
  {
    \@addtotoporbot
  }
  \begin{syntax}
    \cs{@addtotoporbot}
  \end{syntax}

\end{function}



\begin{function}[]
  {
    \@afterheading
  }
  \begin{syntax}
    \cs{@afterheading}
  \end{syntax}

\end{function}



\begin{function}[]
  {
    \@afterindentfalse
  }
  \begin{syntax}
    \cs{@afterindentfalse}
  \end{syntax}

\end{function}



\begin{function}[]
  {
    \@afterindenttrue
  }
  \begin{syntax}
    \cs{@afterindenttrue}
  \end{syntax}

\end{function}



\begin{function}[]
  {
    \@alph
  }
  \begin{syntax}
    \cs{@alph}
  \end{syntax}

\end{function}



\begin{function}[]
  {
    \@ampacol
  }
  \begin{syntax}
    \cs{@ampacol}
  \end{syntax}

\end{function}



\begin{function}[]
  {
    \@arabic
  }
  \begin{syntax}
    \cs{@arabic}
  \end{syntax}

\end{function}



\begin{function}[]
  {
    \@argarraycr
  }
  \begin{syntax}
    \cs{@argarraycr}
  \end{syntax}

\end{function}



\begin{function}[]
  {
    \@argdef
  }
  \begin{syntax}
    \cs{@argdef}
  \end{syntax}

\end{function}



\begin{function}[]
  {
    \@argtabularcr
  }
  \begin{syntax}
    \cs{@argtabularcr}
  \end{syntax}

\end{function}



\begin{function}[]
  {
    \@array
  }
  \begin{syntax}
    \cs{@array}
  \end{syntax}

\end{function}



\begin{function}[]
  {
    \@arrayacol
  }
  \begin{syntax}
    \cs{@arrayacol}
  \end{syntax}

\end{function}



\begin{function}[]
  {
    \@arrayclassiv
  }
  \begin{syntax}
    \cs{@arrayclassiv}
  \end{syntax}

\end{function}



\begin{function}[]
  {
    \@arrayclassz
  }
  \begin{syntax}
    \cs{@arrayclassz}
  \end{syntax}

\end{function}



\begin{function}[]
  {
    \@arraycr
  }
  \begin{syntax}
    \cs{@arraycr}
  \end{syntax}

\end{function}



\begin{function}[]
  {
    \@arrayparboxrestore
  }
  \begin{syntax}
    \cs{@arrayparboxrestore}
  \end{syntax}

\end{function}



\begin{function}[]
  {
    \@arrayrule
  }
  \begin{syntax}
    \cs{@arrayrule}
  \end{syntax}

\end{function}



\begin{function}[]
  {
    \@arstrut
  }
  \begin{syntax}
    \cs{@arstrut}
  \end{syntax}

\end{function}



\begin{function}[]
  {
    \@arstrutbox
  }
  \begin{syntax}
    \cs{@arstrutbox}
  \end{syntax}

\end{function}



\begin{function}[]
  {
    \@author
  }
  \begin{syntax}
    \cs{@author}
  \end{syntax}

\end{function}



\begin{function}[]
  {
    \@auxout
  }
  \begin{syntax}
    \cs{@auxout}
  \end{syntax}

\end{function}



\begin{function}[]
  {
    \@backslashchar
  }
  \begin{syntax}
    \cs{@backslashchar}
  \end{syntax}

\end{function}



\begin{function}[]
  {
    \@backup@outputbox@depth
  }
  \begin{syntax}
    \cs{@backup@outputbox@depth}
  \end{syntax}

\end{function}



\begin{function}[]
  {
    \@badend
  }
  \begin{syntax}
    \cs{@badend}
  \end{syntax}

\end{function}



\begin{function}[]
  {
    \@badlinearg
  }
  \begin{syntax}
    \cs{@badlinearg}
  \end{syntax}

\end{function}



\begin{function}[]
  {
    \@badmath
  }
  \begin{syntax}
    \cs{@badmath}
  \end{syntax}

\end{function}



\begin{function}[]
  {
    \@badpoptabs
  }
  \begin{syntax}
    \cs{@badpoptabs}
  \end{syntax}

\end{function}



\begin{function}[]
  {
    \@badrequireerror
  }
  \begin{syntax}
    \cs{@badrequireerror}
  \end{syntax}

\end{function}



\begin{function}[]
  {
    \@badtab
  }
  \begin{syntax}
    \cs{@badtab}
  \end{syntax}

\end{function}



\begin{function}[]
  {
    \@begin@tempboxa
  }
  \begin{syntax}
    \cs{@begin@tempboxa}
  \end{syntax}

\end{function}



\begin{function}[]
  {
    \@begindocumenthook
  }
  \begin{syntax}
    \cs{@begindocumenthook}
  \end{syntax}

\end{function}



\begin{function}[]
  {
    \@begindvi
  }
  \begin{syntax}
    \cs{@begindvi}
  \end{syntax}

\end{function}



\begin{function}[]
  {
    \@begindvibox
  }
  \begin{syntax}
    \cs{@begindvibox}
  \end{syntax}

\end{function}



\begin{function}[]
  {
    \@begintheorem
  }
  \begin{syntax}
    \cs{@begintheorem}
  \end{syntax}

\end{function}



\begin{function}[]
  {
    \@bezier
  }
  \begin{syntax}
    \cs{@bezier}
  \end{syntax}

\end{function}



\begin{function}[]
  {
    \@bibitem
  }
  \begin{syntax}
    \cs{@bibitem}
  \end{syntax}

\end{function}



\begin{function}[]
  {
    \@biblabel
  }
  \begin{syntax}
    \cs{@biblabel}
  \end{syntax}

\end{function}



\begin{function}[]
  {
    \@bitor
  }
  \begin{syntax}
    \cs{@bitor}
  \end{syntax}

\end{function}



\begin{function}[]
  {
    \@botlist
  }
  \begin{syntax}
    \cs{@botlist}
  \end{syntax}

\end{function}



\begin{function}[]
  {
    \@botnum
  }
  \begin{syntax}
    \cs{@botnum}
  \end{syntax}

\end{function}



\begin{function}[]
  {
    \@botroom
  }
  \begin{syntax}
    \cs{@botroom}
  \end{syntax}

\end{function}



\begin{function}[]
  {
    \@boxfpsbit
  }
  \begin{syntax}
    \cs{@boxfpsbit}
  \end{syntax}

\end{function}



\begin{function}[]
  {
    \@break@tfor
  }
  \begin{syntax}
    \cs{@break@tfor}
  \end{syntax}

\end{function}





\begin{function}[]
  {
    \@caption
  }
  \begin{syntax}
    \cs{@caption}
  \end{syntax}

\end{function}



\begin{function}[]
  {
    \@captype
  }
  \begin{syntax}
    \cs{@captype}
  \end{syntax}

\end{function}



\begin{function}[]
  {
    \@car
  }
  \begin{syntax}
    \cs{@car}
  \end{syntax}

\end{function}



\begin{function}[]
  {
    \@carcube
  }
  \begin{syntax}
    \cs{@carcube}
  \end{syntax}

\end{function}



\begin{function}[]
  {
    \@cclv
  }
  \begin{syntax}
    \cs{@cclv}
  \end{syntax}

\end{function}



\begin{function}[]
  {
    \@cclvi
  }
  \begin{syntax}
    \cs{@cclvi}
  \end{syntax}

\end{function}



\begin{function}[]
  {
    \@cdr
  }
  \begin{syntax}
    \cs{@cdr}
  \end{syntax}

\end{function}



\begin{function}[]
  {
    \@centercr
  }
  \begin{syntax}
    \cs{@centercr}
  \end{syntax}

\end{function}



\begin{function}[]
  {
    \@centering
  }
  \begin{syntax}
    \cs{@centering}
  \end{syntax}

\end{function}



\begin{function}[]
  {
    \@cflb
  }
  \begin{syntax}
    \cs{@cflb}
  \end{syntax}

\end{function}



\begin{function}[]
  {
    \@cflt
  }
  \begin{syntax}
    \cs{@cflt}
  \end{syntax}

\end{function}



\begin{function}[]
  {
    \@changed@cmd
  }
  \begin{syntax}
    \cs{@changed@cmd}
  \end{syntax}

\end{function}



\begin{function}[]
  {
    \@charlb
  }
  \begin{syntax}
    \cs{@charlb}
  \end{syntax}

\end{function}



\begin{function}[]
  {
    \@charrb
  }
  \begin{syntax}
    \cs{@charrb}
  \end{syntax}

\end{function}



\begin{function}[]
  {
    \@chclass
  }
  \begin{syntax}
    \cs{@chclass}
  \end{syntax}

\end{function}





\begin{function}[]
  {
    \@check@c
  }
  \begin{syntax}
    \cs{@check@c}
  \end{syntax}

\end{function}



\begin{function}[]
  {
    \@check@eq
  }
  \begin{syntax}
    \cs{@check@eq}
  \end{syntax}

\end{function}



\begin{function}[]
  {
    \@checkend
  }
  \begin{syntax}
    \cs{@checkend}
  \end{syntax}

\end{function}



\begin{function}[]
  {
    \@chnum
  }
  \begin{syntax}
    \cs{@chnum}
  \end{syntax}

\end{function}



\begin{function}[]
  {
    \@circ
  }
  \begin{syntax}
    \cs{@circ}
  \end{syntax}

\end{function}



\begin{function}[]
  {
    \@circle
  }
  \begin{syntax}
    \cs{@circle}
  \end{syntax}

\end{function}



\begin{function}[]
  {
    \@circlefnt
  }
  \begin{syntax}
    \cs{@circlefnt}
  \end{syntax}

\end{function}



\begin{function}[]
  {
    \@cite
  }
  \begin{syntax}
    \cs{@cite}
  \end{syntax}

\end{function}



\begin{function}[]
  {
    \@cite@ofmt
  }
  \begin{syntax}
    \cs{@cite@ofmt}
  \end{syntax}

\end{function}



\begin{function}[]
  {
    \@citea
  }
  \begin{syntax}
    \cs{@citea}
  \end{syntax}

\end{function}



\begin{function}[]
  {
    \@citeb
  }
  \begin{syntax}
    \cs{@citeb}
  \end{syntax}

\end{function}





\begin{function}[]
  {
    \@citex
  }
  \begin{syntax}
    \cs{@citex}
  \end{syntax}

\end{function}



\begin{function}[]
  {
    \@citex@checkblank
  }
  \begin{syntax}
    \cs{@citex@checkblank}
  \end{syntax}

\end{function}



\begin{function}[]
  {
    \@classi
  }
  \begin{syntax}
    \cs{@classi}
  \end{syntax}

\end{function}



\begin{function}[]
  {
    \@classii
  }
  \begin{syntax}
    \cs{@classii}
  \end{syntax}

\end{function}



\begin{function}[]
  {
    \@classiii
  }
  \begin{syntax}
    \cs{@classiii}
  \end{syntax}

\end{function}



\begin{function}[]
  {
    \@classiv
  }
  \begin{syntax}
    \cs{@classiv}
  \end{syntax}

\end{function}



\begin{function}[]
  {
    \@classoptionslist
  }
  \begin{syntax}
    \cs{@classoptionslist}
  \end{syntax}

\end{function}



\begin{function}[]
  {
    \@classv
  }
  \begin{syntax}
    \cs{@classv}
  \end{syntax}

\end{function}



\begin{function}[]
  {
    \@classz
  }
  \begin{syntax}
    \cs{@classz}
  \end{syntax}

\end{function}



\begin{function}[]
  {
    \@cline
  }
  \begin{syntax}
    \cs{@cline}
  \end{syntax}

\end{function}



\begin{function}[]
  {
    \@clnht
  }
  \begin{syntax}
    \cs{@clnht}
  \end{syntax}

\end{function}



\begin{function}[]
  {
    \@clnwd
  }
  \begin{syntax}
    \cs{@clnwd}
  \end{syntax}

\end{function}



\begin{function}[]
  {
    \@cls@pkg
  }
  \begin{syntax}
    \cs{@cls@pkg}
  \end{syntax}

\end{function}



\begin{function}[]
  {
    \@clsextension
  }
  \begin{syntax}
    \cs{@clsextension}
  \end{syntax}

\end{function}



\begin{function}[]
  {
    \@clubpenalty
  }
  \begin{syntax}
    \cs{@clubpenalty}
  \end{syntax}

\end{function}



\begin{function}[]
  {
    \@colht
  }
  \begin{syntax}
    \cs{@colht}
  \end{syntax}

\end{function}



\begin{function}[]
  {
    \@colnum
  }
  \begin{syntax}
    \cs{@colnum}
  \end{syntax}

\end{function}



\begin{function}[]
  {
    \@colroom
  }
  \begin{syntax}
    \cs{@colroom}
  \end{syntax}

\end{function}



\begin{function}[]
  {
    \@combinedblfloats
  }
  \begin{syntax}
    \cs{@combinedblfloats}
  \end{syntax}

\end{function}



\begin{function}[]
  {
    \@combinefloats
  }
  \begin{syntax}
    \cs{@combinefloats}
  \end{syntax}

\end{function}



\begin{function}[]
  {
    \@comdblflelt
  }
  \begin{syntax}
    \cs{@comdblflelt}
  \end{syntax}

\end{function}



\begin{function}[]
  {
    \@comflelt
  }
  \begin{syntax}
    \cs{@comflelt}
  \end{syntax}

\end{function}



\begin{function}[]
  {
    \@cons
  }
  \begin{syntax}
    \cs{@cons}
  \end{syntax}

\end{function}



\begin{function}[]
  {
    \@contentsline@destination
  }
  \begin{syntax}
    \cs{@contentsline@destination}
  \end{syntax}

\end{function}



\begin{function}[]
  {
    \@contfield
  }
  \begin{syntax}
    \cs{@contfield}
  \end{syntax}

\end{function}



\begin{function}[]
  {
    \@copy@DeclareRobustCommand
  }
  \begin{syntax}
    \cs{@copy@DeclareRobustCommand}
  \end{syntax}

\end{function}



\begin{function}[]
  {
    \@copy@newcommand
  }
  \begin{syntax}
    \cs{@copy@newcommand}
  \end{syntax}

\end{function}



\begin{function}[]
  {
    \@ctrerr
  }
  \begin{syntax}
    \cs{@ctrerr}
  \end{syntax}

\end{function}



\begin{function}[]
  {
    \@curfield
  }
  \begin{syntax}
    \cs{@curfield}
  \end{syntax}

\end{function}



\begin{function}[]
  {
    \@curline
  }
  \begin{syntax}
    \cs{@curline}
  \end{syntax}

\end{function}



\begin{function}[]
  {
    \@curr@enc
  }
  \begin{syntax}
    \cs{@curr@enc}
  \end{syntax}

\end{function}



\begin{function}[]
  {
    \@curr@file
  }
  \begin{syntax}
    \cs{@curr@file}
  \end{syntax}

\end{function}



\begin{function}[]
  {
    \@curr@file@reqd
  }
  \begin{syntax}
    \cs{@curr@file@reqd}
  \end{syntax}

\end{function}



\begin{function}[]
  {
    \@currbox
  }
  \begin{syntax}
    \cs{@currbox}
  \end{syntax}

\end{function}



\begin{function}[]
  {
    \@currbox1sp
  }
  \begin{syntax}
    \cs{@currbox1sp}
  \end{syntax}

\end{function}



\begin{function}[]
  {
    \@currdir
  }
  \begin{syntax}
    \cs{@currdir}
  \end{syntax}

\end{function}



\begin{function}[]
  {
    \@current@cmd
  }
  \begin{syntax}
    \cs{@current@cmd}
  \end{syntax}

\end{function}



\begin{function}[]
  {
    \@currentHpage
  }
  \begin{syntax}
    \cs{@currentHpage}
  \end{syntax}

\end{function}



\begin{function}[]
  {
    \@currentHref
  }
  \begin{syntax}
    \cs{@currentHref}
  \end{syntax}

\end{function}






\begin{function}[]
  {
    \@currentmetafamily
  }
  \begin{syntax}
    \cs{@currentmetafamily}
  \end{syntax}

\end{function}





\begin{function}[]
  {
    \@currenvline
  }
  \begin{syntax}
    \cs{@currenvline}
  \end{syntax}

\end{function}



\begin{function}[]
  {
    \@currext
  }
  \begin{syntax}
    \cs{@currext}
  \end{syntax}

\end{function}



\begin{function}[]
  {
    \@currlist
  }
  \begin{syntax}
    \cs{@currlist}
  \end{syntax}

\end{function}



\begin{function}[]
  {
    \@currname
  }
  \begin{syntax}
    \cs{@currname}
  \end{syntax}

\end{function}



\begin{function}[]
  {
    \@currnamestack
  }
  \begin{syntax}
    \cs{@currnamestack}
  \end{syntax}

\end{function}



\begin{function}[]
  {
    \@curroptions
  }
  \begin{syntax}
    \cs{@curroptions}
  \end{syntax}

\end{function}



\begin{function}[]
  {
    \@currpath
  }
  \begin{syntax}
    \cs{@currpath}
  \end{syntax}

\end{function}



\begin{function}[]
  {
    \@currpkg@reqd
  }
  \begin{syntax}
    \cs{@currpkg@reqd}
  \end{syntax}

\end{function}



\begin{function}[]
  {
    \@currsize
  }
  \begin{syntax}
    \cs{@currsize}
  \end{syntax}

\end{function}



\begin{function}[]
  {
    \@currtype
  }
  \begin{syntax}
    \cs{@currtype}
  \end{syntax}

\end{function}



\begin{function}[]
  {
    \@curtab
  }
  \begin{syntax}
    \cs{@curtab}
  \end{syntax}

\end{function}



\begin{function}[]
  {
    \@curtabmar
  }
  \begin{syntax}
    \cs{@curtabmar}
  \end{syntax}

\end{function}



\begin{function}[]
  {
    \@dashbox
  }
  \begin{syntax}
    \cs{@dashbox}
  \end{syntax}

\end{function}



\begin{function}[]
  {
    \@dashcnt
  }
  \begin{syntax}
    \cs{@dashcnt}
  \end{syntax}

\end{function}



\begin{function}[]
  {
    \@dashdim
  }
  \begin{syntax}
    \cs{@dashdim}
  \end{syntax}

\end{function}



\begin{function}[]
  {
    \@date
  }
  \begin{syntax}
    \cs{@date}
  \end{syntax}

\end{function}



\begin{function}[]
  {
    \@dbflt
  }
  \begin{syntax}
    \cs{@dbflt}
  \end{syntax}

\end{function}



\begin{function}[]
  {
    \@dblarg
  }
  \begin{syntax}
    \cs{@dblarg}
  \end{syntax}

\end{function}



\begin{function}[]
  {
    \@dbldeferlist
  }
  \begin{syntax}
    \cs{@dbldeferlist}
  \end{syntax}

\end{function}



\begin{function}[]
  {
    \@dblfloat
  }
  \begin{syntax}
    \cs{@dblfloat}
  \end{syntax}

\end{function}



\begin{function}[]
  {
    \@dblfloatplacement
  }
  \begin{syntax}
    \cs{@dblfloatplacement}
  \end{syntax}

\end{function}



\begin{function}[]
  {
    \@dblfpbot
  }
  \begin{syntax}
    \cs{@dblfpbot}
  \end{syntax}

\end{function}



\begin{function}[]
  {
    \@dblfpsep
  }
  \begin{syntax}
    \cs{@dblfpsep}
  \end{syntax}

\end{function}



\begin{function}[]
  {
    \@dblfptop
  }
  \begin{syntax}
    \cs{@dblfptop}
  \end{syntax}

\end{function}



\begin{function}[]
  {
    \@dbltoplist
  }
  \begin{syntax}
    \cs{@dbltoplist}
  \end{syntax}

\end{function}



\begin{function}[]
  {
    \@dbltopnum
  }
  \begin{syntax}
    \cs{@dbltopnum}
  \end{syntax}

\end{function}



\begin{function}[]
  {
    \@dbltoproom
  }
  \begin{syntax}
    \cs{@dbltoproom}
  \end{syntax}

\end{function}



\begin{function}[]
  {
    \@dec@text@cmd
  }
  \begin{syntax}
    \cs{@dec@text@cmd}
  \end{syntax}

\end{function}



\begin{function}[]
  {
    \@declarecommandcopylisthook
  }
  \begin{syntax}
    \cs{@declarecommandcopylisthook}
  \end{syntax}

\end{function}



\begin{function}[]
  {
    \@declaredoptions
  }
  \begin{syntax}
    \cs{@declaredoptions}
  \end{syntax}

\end{function}




\begin{function}[]
  {
    \@declareoption
  }
  \begin{syntax}
    \cs{@declareoption}
  \end{syntax}

\end{function}



\begin{function}[]
  {
    \@defaultfamilyhook
  }
  \begin{syntax}
    \cs{@defaultfamilyhook}
  \end{syntax}

\end{function}



\begin{function}[]
  {
    \@defaultsubs
  }
  \begin{syntax}
    \cs{@defaultsubs}
  \end{syntax}

\end{function}



\begin{function}[]
  {
    \@defaultunits
  }
  \begin{syntax}
    \cs{@defaultunits}
  \end{syntax}

\end{function}



\begin{function}[]
  {
    \@defaultunitsset
  }
  \begin{syntax}
    \cs{@defaultunitsset}
  \end{syntax}

\end{function}



\begin{function}[]
  {
    \@defdefault@ds
  }
  \begin{syntax}
    \cs{@defdefault@ds}
  \end{syntax}

\end{function}



\begin{function}[]
  {
    \@deferlist
  }
  \begin{syntax}
    \cs{@deferlist}
  \end{syntax}

\end{function}



\begin{function}[]
  {
    \@definecounter
  }
  \begin{syntax}
    \cs{@definecounter}
  \end{syntax}

\end{function}



\begin{function}[]
  {
    \@depth
  }
  \begin{syntax}
    \cs{@depth}
  \end{syntax}

\end{function}



\begin{function}[]
  {
    \@disable@packageload@do
  }
  \begin{syntax}
    \cs{@disable@packageload@do}
  \end{syntax}

\end{function}



\begin{function}[]
  {
    \@dischyph
  }
  \begin{syntax}
    \cs{@dischyph}
  \end{syntax}

\end{function}



\begin{function}[]
  {
    \@doclearpage
  }
  \begin{syntax}
    \cs{@doclearpage}
  \end{syntax}

\end{function}



\begin{function}[]
  {
    \@documentclasshook
  }
  \begin{syntax}
    \cs{@documentclasshook}
  \end{syntax}

\end{function}



\begin{function}[]
  {
    \@doendpe
  }
  \begin{syntax}
    \cs{@doendpe}
  \end{syntax}

\end{function}



\begin{function}[]
  {
    \@dofilelist
  }
  \begin{syntax}
    \cs{@dofilelist}
  \end{syntax}

\end{function}



\begin{function}[]
  {
    \@dofilelist@hash
  }
  \begin{syntax}
    \cs{@dofilelist@hash}
  \end{syntax}

\end{function}



\begin{function}[]
  {
    \@dofilelist@size
  }
  \begin{syntax}
    \cs{@dofilelist@size}
  \end{syntax}

\end{function}



\begin{function}[]
  {
    \@donoparitem
  }
  \begin{syntax}
    \cs{@donoparitem}
  \end{syntax}

\end{function}



\begin{function}[]
  {
    \@dot
  }
  \begin{syntax}
    \cs{@dot}
  \end{syntax}

\end{function}



\begin{function}[]
  {
    \@dotsep
  }
  \begin{syntax}
    \cs{@dotsep}
  \end{syntax}

\end{function}



\begin{function}[]
  {
    \@dottedtocline
  }
  \begin{syntax}
    \cs{@dottedtocline}
  \end{syntax}

\end{function}



\begin{function}[]
  {
    \@downline
  }
  \begin{syntax}
    \cs{@downline}
  \end{syntax}

\end{function}



\begin{function}[]
  {
    \@downvector
  }
  \begin{syntax}
    \cs{@downvector}
  \end{syntax}

\end{function}



\begin{function}[]
  {
    \@eha
  }
  \begin{syntax}
    \cs{@eha}
  \end{syntax}

\end{function}



\begin{function}[]
  {
    \@ehb
  }
  \begin{syntax}
    \cs{@ehb}
  \end{syntax}

\end{function}



\begin{function}[]
  {
    \@ehc
  }
  \begin{syntax}
    \cs{@ehc}
  \end{syntax}

\end{function}



\begin{function}[]
  {
    \@ehd
  }
  \begin{syntax}
    \cs{@ehd}
  \end{syntax}

\end{function}



\begin{function}[]
  {
    \@elt
  }
  \begin{syntax}
    \cs{@elt}
  \end{syntax}

\end{function}



\begin{function}[]
  {
    \@empty
  }
  \begin{syntax}
    \cs{@empty}
  \end{syntax}

\end{function}



\begin{function}[]
  {
    \@emptycol
  }
  \begin{syntax}
    \cs{@emptycol}
  \end{syntax}

\end{function}



\begin{function}[]
  {
    \@enc@info
  }
  \begin{syntax}
    \cs{@enc@info}
  \end{syntax}

\end{function}




\begin{function}[]
  {
    \@end@tempboxa
  }
  \begin{syntax}
    \cs{@end@tempboxa}
  \end{syntax}

\end{function}



\begin{function}[]
  {
    \@enddocument@kernel@warnings
  }
  \begin{syntax}
    \cs{@enddocument@kernel@warnings}
  \end{syntax}

\end{function}



\begin{function}[]
  {
    \@enddocumenthook
  }
  \begin{syntax}
    \cs{@enddocumenthook}
  \end{syntax}

\end{function}



\begin{function}[]
  {
    \@endfloatbox
  }
  \begin{syntax}
    \cs{@endfloatbox}
  \end{syntax}

\end{function}



\begin{function}[]
  {
    \@endparenv
  }
  \begin{syntax}
    \cs{@endparenv}
  \end{syntax}

\end{function}





\begin{function}[]
  {
    \@endpbox
  }
  \begin{syntax}
    \cs{@endpbox}
  \end{syntax}

\end{function}



\begin{function}[]
  {
    \@endpefalse
  }
  \begin{syntax}
    \cs{@endpefalse}
  \end{syntax}

\end{function}



\begin{function}[]
  {
    \@endpetrue
  }
  \begin{syntax}
    \cs{@endpetrue}
  \end{syntax}

\end{function}



\begin{function}[]
  {
    \@endtheorem
  }
  \begin{syntax}
    \cs{@endtheorem}
  \end{syntax}

\end{function}



\begin{function}[]
  {
    \@enlargepage
  }
  \begin{syntax}
    \cs{@enlargepage}
  \end{syntax}

\end{function}



\begin{function}[]
  {
    \@ensuredmath
  }
  \begin{syntax}
    \cs{@ensuredmath}
  \end{syntax}

\end{function}



\begin{function}[]
  {
    \@enumctr
  }
  \begin{syntax}
    \cs{@enumctr}
  \end{syntax}

\end{function}



\begin{function}[]
  {
    \@enumdepth
  }
  \begin{syntax}
    \cs{@enumdepth}
  \end{syntax}

\end{function}



\begin{function}[]
  {
    \@eqcnt
  }
  \begin{syntax}
    \cs{@eqcnt}
  \end{syntax}

\end{function}



\begin{function}[]
  {
    \@eqncr
  }
  \begin{syntax}
    \cs{@eqncr}
  \end{syntax}

\end{function}



\begin{function}[]
  {
    \@eqnnum
  }
  \begin{syntax}
    \cs{@eqnnum}
  \end{syntax}

\end{function}



\begin{function}[]
  {
    \@eqnsel
  }
  \begin{syntax}
    \cs{@eqnsel}
  \end{syntax}

\end{function}



\begin{function}[]
  {
    \@eqnswfalse
  }
  \begin{syntax}
    \cs{@eqnswfalse}
  \end{syntax}

\end{function}



\begin{function}[]
  {
    \@eqnswtrue
  }
  \begin{syntax}
    \cs{@eqnswtrue}
  \end{syntax}

\end{function}



\begin{function}[]
  {
    \@eqpen
  }
  \begin{syntax}
    \cs{@eqpen}
  \end{syntax}

\end{function}



\begin{function}[]
  {
    \@err@
  }
  \begin{syntax}
    \cs{@err@}
  \end{syntax}

\end{function}





\begin{function}[]
  {
    \@evenfoot
  }
  \begin{syntax}
    \cs{@evenfoot}
  \end{syntax}

\end{function}



\begin{function}[]
  {
    \@evenhead
  }
  \begin{syntax}
    \cs{@evenhead}
  \end{syntax}

\end{function}



\begin{function}[]
  {
    \@execute@begin@hook
  }
  \begin{syntax}
    \cs{@execute@begin@hook}
  \end{syntax}

\end{function}



\begin{function}[]
  {
    \@expandtwoargs
  }
  \begin{syntax}
    \cs{@expandtwoargs}
  \end{syntax}

\end{function}



\begin{function}[]
  {
    \@expast
  }
  \begin{syntax}
    \cs{@expast}
  \end{syntax}

\end{function}



\begin{function}[]
  {
    \@expl@@@filehook@clear@replacement@flag@@
  }
  \begin{syntax}
    \cs{@expl@@@filehook@clear@replacement@flag@@}
  \end{syntax}

\end{function}



\begin{function}[]
  {
    \@expl@@@filehook@drop@extension@@N
  }
  \begin{syntax}
    \cs{@expl@@@filehook@drop@extension@@N}
  \end{syntax}

\end{function}



\begin{function}[]
  {
    \@expl@@@filehook@file@pop@@
  }
  \begin{syntax}
    \cs{@expl@@@filehook@file@pop@@}
  \end{syntax}

\end{function}



\begin{function}[]
  {
    \@expl@@@filehook@file@pop@assign@@nnnn
  }
  \begin{syntax}
    \cs{@expl@@@filehook@file@pop@assign@@nnnn}
  \end{syntax}

\end{function}



\begin{function}[]
  {
    \@expl@@@filehook@file@push@@
  }
  \begin{syntax}
    \cs{@expl@@@filehook@file@push@@}
  \end{syntax}

\end{function}



\begin{function}[]
  {
    \@expl@@@filehook@if@file@replaced@@TF
  }
  \begin{syntax}
    \cs{@expl@@@filehook@if@file@replaced@@TF}
  \end{syntax}

\end{function}



\begin{function}[]
  {
    \@expl@@@filehook@if@no@extension@@nTF
  }
  \begin{syntax}
    \cs{@expl@@@filehook@if@no@extension@@nTF}
  \end{syntax}

\end{function}



\begin{function}[]
  {
    \@expl@@@filehook@normalize@file@name@@w
  }
  \begin{syntax}
    \cs{@expl@@@filehook@normalize@file@name@@w}
  \end{syntax}

\end{function}



\begin{function}[]
  {
    \@expl@@@filehook@resolve@file@subst@@w
  }
  \begin{syntax}
    \cs{@expl@@@filehook@resolve@file@subst@@w}
  \end{syntax}

\end{function}



\begin{function}[]
  {
    \@expl@@@filehook@set@curr@file@@nNN
  }
  \begin{syntax}
    \cs{@expl@@@filehook@set@curr@file@@nNN}
  \end{syntax}

\end{function}



\begin{function}[]
  {
    \@expl@@@hook@curr@name@pop@@
  }
  \begin{syntax}
    \cs{@expl@@@hook@curr@name@pop@@}
  \end{syntax}

\end{function}



\begin{function}[]
  {
    \@expl@@@initialize@all@@
  }
  \begin{syntax}
    \cs{@expl@@@initialize@all@@}
  \end{syntax}

\end{function}



\begin{function}[]
  {
    \@expl@@@mark@update@dblcol@structures@@
  }
  \begin{syntax}
    \cs{@expl@@@mark@update@dblcol@structures@@}
  \end{syntax}

\end{function}



\begin{function}[]
  {
    \@expl@@@mark@update@singlecol@structures@@
  }
  \begin{syntax}
    \cs{@expl@@@mark@update@singlecol@structures@@}
  \end{syntax}

\end{function}



\begin{function}[]
  {
    \@expl@@@shipout@add@background@box@@n
  }
  \begin{syntax}
    \cs{@expl@@@shipout@add@background@box@@n}
  \end{syntax}

\end{function}



\begin{function}[]
  {
    \@expl@@@shipout@add@background@picture@@n
  }
  \begin{syntax}
    \cs{@expl@@@shipout@add@background@picture@@n}
  \end{syntax}

\end{function}



\begin{function}[]
  {
    \@expl@@@shipout@add@firstpage@material@@Nn
  }
  \begin{syntax}
    \cs{@expl@@@shipout@add@firstpage@material@@Nn}
  \end{syntax}

\end{function}



\begin{function}[]
  {
    \@expl@@@shipout@add@foreground@box@@n
  }
  \begin{syntax}
    \cs{@expl@@@shipout@add@foreground@box@@n}
  \end{syntax}

\end{function}



\begin{function}[]
  {
    \@expl@@@shipout@add@foreground@picture@@n
  }
  \begin{syntax}
    \cs{@expl@@@shipout@add@foreground@picture@@n}
  \end{syntax}

\end{function}



\begin{function}[]
  {
    \@expl@char@generate@@nn
  }
  \begin{syntax}
    \cs{@expl@char@generate@@nn}
  \end{syntax}

\end{function}



\begin{function}[]
  {
    \@expl@cs@parameter@spec@@N
  }
  \begin{syntax}
    \cs{@expl@cs@parameter@spec@@N}
  \end{syntax}

\end{function}



\begin{function}[]
  {
    \@expl@cs@prefix@spec@@N
  }
  \begin{syntax}
    \cs{@expl@cs@prefix@spec@@N}
  \end{syntax}

\end{function}



\begin{function}[]
  {
    \@expl@cs@replacement@spec@@N
  }
  \begin{syntax}
    \cs{@expl@cs@replacement@spec@@N}
  \end{syntax}

\end{function}



\begin{function}[]
  {
    \@expl@cs@to@str@@N
  }
  \begin{syntax}
    \cs{@expl@cs@to@str@@N}
  \end{syntax}

\end{function}



\begin{function}[]
  {
    \@expl@finalise@setup@@
  }
  \begin{syntax}
    \cs{@expl@finalise@setup@@}
  \end{syntax}

\end{function}



\begin{function}[]
  {
    \@expl@pop@filename@@
  }
  \begin{syntax}
    \cs{@expl@pop@filename@@}
  \end{syntax}

\end{function}



\begin{function}[]
  {
    \@expl@push@filename@@
  }
  \begin{syntax}
    \cs{@expl@push@filename@@}
  \end{syntax}

\end{function}



\begin{function}[]
  {
    \@expl@push@filename@aux@@
  }
  \begin{syntax}
    \cs{@expl@push@filename@aux@@}
  \end{syntax}

\end{function}



\begin{function}[]
  {
    \@expl@str@if@eq@@nnTF
  }
  \begin{syntax}
    \cs{@expl@str@if@eq@@nnTF}
  \end{syntax}

\end{function}



\begin{function}[]
  {
    \@expl@str@map@function@@NN
  }
  \begin{syntax}
    \cs{@expl@str@map@function@@NN}
  \end{syntax}

\end{function}



\begin{function}[]
  {
    \@expl@sys@load@backend@@
  }
  \begin{syntax}
    \cs{@expl@sys@load@backend@@}
  \end{syntax}

\end{function}



\begin{function}[]
  {
    \@extra@page@added
  }
  \begin{syntax}
    \cs{@extra@page@added}
  \end{syntax}

\end{function}



\begin{function}[]
  {
    \@failedlist
  }
  \begin{syntax}
    \cs{@failedlist}
  \end{syntax}

\end{function}



\begin{function}[]
  {
    \@fcolmadefalse
  }
  \begin{syntax}
    \cs{@fcolmadefalse}
  \end{syntax}

\end{function}



\begin{function}[]
  {
    \@fcolmadetrue
  }
  \begin{syntax}
    \cs{@fcolmadetrue}
  \end{syntax}

\end{function}



\begin{function}[]
  {
    \@filef@und
  }
  \begin{syntax}
    \cs{@filef@und}
  \end{syntax}

\end{function}



\begin{function}[]
  {
    \@filef@und"
  }
  \begin{syntax}
    \cs{@filef@und"}
  \end{syntax}

\end{function}



\begin{function}[]
  {
    \@filehook@set@CurrentFile
  }
  \begin{syntax}
    \cs{@filehook@set@CurrentFile}
  \end{syntax}

\end{function}



\begin{function}[]
  {
    \@filelist
  }
  \begin{syntax}
    \cs{@filelist}
  \end{syntax}

\end{function}




\begin{function}[]
  {
    \@fileswfalse
  }
  \begin{syntax}
    \cs{@fileswfalse}
  \end{syntax}

\end{function}



\begin{function}[]
  {
    \@fileswith@pti@ns
  }
  \begin{syntax}
    \cs{@fileswith@pti@ns}
  \end{syntax}

\end{function}



\begin{function}[]
  {
    \@fileswith@ptions
  }
  \begin{syntax}
    \cs{@fileswith@ptions}
  \end{syntax}

\end{function}



\begin{function}[]
  {
    \@fileswithoptions
  }
  \begin{syntax}
    \cs{@fileswithoptions}
  \end{syntax}

\end{function}



\begin{function}[]
  {
    \@fileswtrue
  }
  \begin{syntax}
    \cs{@fileswtrue}
  \end{syntax}

\end{function}



\begin{function}[]
  {
    \@finalstrut
  }
  \begin{syntax}
    \cs{@finalstrut}
  \end{syntax}

\end{function}



\begin{function}[]
  {
    \@firstampfalse
  }
  \begin{syntax}
    \cs{@firstampfalse}
  \end{syntax}

\end{function}



\begin{function}[]
  {
    \@firstamptrue
  }
  \begin{syntax}
    \cs{@firstamptrue}
  \end{syntax}

\end{function}



\begin{function}[]
  {
    \@firstcolumnfalse
  }
  \begin{syntax}
    \cs{@firstcolumnfalse}
  \end{syntax}

\end{function}



\begin{function}[]
  {
    \@firstcolumntrue
  }
  \begin{syntax}
    \cs{@firstcolumntrue}
  \end{syntax}

\end{function}



\begin{function}[]
  {
    \@firstoffive
  }
  \begin{syntax}
    \cs{@firstoffive}
  \end{syntax}

\end{function}



\begin{function}[]
  {
    \@firstofone
  }
  \begin{syntax}
    \cs{@firstofone}
  \end{syntax}

\end{function}



\begin{function}[]
  {
    \@firstoftwo
  }
  \begin{syntax}
    \cs{@firstoftwo}
  \end{syntax}

\end{function}



\begin{function}[]
  {
    \@firsttab
  }
  \begin{syntax}
    \cs{@firsttab}
  \end{syntax}

\end{function}



\begin{function}[]
  {
    \@flcheckspace
  }
  \begin{syntax}
    \cs{@flcheckspace}
  \end{syntax}

\end{function}



\begin{function}[]
  {
    \@flfail
  }
  \begin{syntax}
    \cs{@flfail}
  \end{syntax}

\end{function}



\begin{function}[]
  {
    \@float
  }
  \begin{syntax}
    \cs{@float}
  \end{syntax}

\end{function}



\begin{function}[]
  {
    \@floatboxreset
  }
  \begin{syntax}
    \cs{@floatboxreset}
  \end{syntax}

\end{function}



\begin{function}[]
  {
    \@floatpenalty
  }
  \begin{syntax}
    \cs{@floatpenalty}
  \end{syntax}

\end{function}



\begin{function}[]
  {
    \@floatplacement
  }
  \begin{syntax}
    \cs{@floatplacement}
  \end{syntax}

\end{function}



\begin{function}[]
  {
    \@flsetnum
  }
  \begin{syntax}
    \cs{@flsetnum}
  \end{syntax}

\end{function}



\begin{function}[]
  {
    \@flsettextmin
  }
  \begin{syntax}
    \cs{@flsettextmin}
  \end{syntax}

\end{function}



\begin{function}[]
  {
    \@flstop
  }
  \begin{syntax}
    \cs{@flstop}
  \end{syntax}

\end{function}



\begin{function}[]
  {
    \@flsucceed
  }
  \begin{syntax}
    \cs{@flsucceed}
  \end{syntax}

\end{function}



\begin{function}[]
  {
    \@fltovf
  }
  \begin{syntax}
    \cs{@fltovf}
  \end{syntax}

\end{function}



\begin{function}[]
  {
    \@flupdates
  }
  \begin{syntax}
    \cs{@flupdates}
  \end{syntax}

\end{function}



\begin{function}[]
  {
    \@flushglue
  }
  \begin{syntax}
    \cs{@flushglue}
  \end{syntax}

\end{function}



\begin{function}[]
  {
    \@fnsymbol
  }
  \begin{syntax}
    \cs{@fnsymbol}
  \end{syntax}

\end{function}



\begin{function}[]
  {
    \@font@aliasinfo
  }
  \begin{syntax}
    \cs{@font@aliasinfo}
  \end{syntax}

\end{function}



\begin{function}[]
  {
    \@font@info
  }
  \begin{syntax}
    \cs{@font@info}
  \end{syntax}

\end{function}



\begin{function}[]
  {
    \@font@series@contextfalse
  }
  \begin{syntax}
    \cs{@font@series@contextfalse}
  \end{syntax}

\end{function}



\begin{function}[]
  {
    \@font@series@contexttrue
  }
  \begin{syntax}
    \cs{@font@series@contexttrue}
  \end{syntax}

\end{function}



\begin{function}[]
  {
    \@font@shape@subst@warning
  }
  \begin{syntax}
    \cs{@font@shape@subst@warning}
  \end{syntax}

\end{function}



\begin{function}[]
  {
    \@font@warning
  }
  \begin{syntax}
    \cs{@font@warning}
  \end{syntax}

\end{function}



\begin{function}[]
  {
    \@fontenc@load@list
  }
  \begin{syntax}
    \cs{@fontenc@load@list}
  \end{syntax}

\end{function}



\begin{function}[]
  {
    \@fontswitch
  }
  \begin{syntax}
    \cs{@fontswitch}
  \end{syntax}

\end{function}



\begin{function}[]
  {
    \@footnotemark
  }
  \begin{syntax}
    \cs{@footnotemark}
  \end{syntax}

\end{function}



\begin{function}[]
  {
    \@footnotetext
  }
  \begin{syntax}
    \cs{@footnotetext}
  \end{syntax}

\end{function}



\begin{function}[]
  {
    \@for
  }
  \begin{syntax}
    \cs{@for}
  \end{syntax}

\end{function}



\begin{function}[]
  {
    \@forced@seriesfalse
  }
  \begin{syntax}
    \cs{@forced@seriesfalse}
  \end{syntax}

\end{function}



\begin{function}[]
  {
    \@forced@seriestrue
  }
  \begin{syntax}
    \cs{@forced@seriestrue}
  \end{syntax}

\end{function}



\begin{function}[]
  {
    \@forloop
  }
  \begin{syntax}
    \cs{@forloop}
  \end{syntax}

\end{function}



\begin{function}[]
  {
    \@fornoop
  }
  \begin{syntax}
    \cs{@fornoop}
  \end{syntax}

\end{function}



\begin{function}[]
  {
    \@fortmp
  }
  \begin{syntax}
    \cs{@fortmp}
  \end{syntax}

\end{function}



\begin{function}[]
  {
    \@fpbot
  }
  \begin{syntax}
    \cs{@fpbot}
  \end{syntax}

\end{function}



\begin{function}[]
  {
    \@fpmin
  }
  \begin{syntax}
    \cs{@fpmin}
  \end{syntax}

\end{function}



\begin{function}[]
  {
    \@fps
  }
  \begin{syntax}
    \cs{@fps}
  \end{syntax}

\end{function}



\begin{function}[]
  {
    \@fpsadddefault
  }
  \begin{syntax}
    \cs{@fpsadddefault}
  \end{syntax}

\end{function}



\begin{function}[]
  {
    \@fpsep
  }
  \begin{syntax}
    \cs{@fpsep}
  \end{syntax}

\end{function}



\begin{function}[]
  {
    \@fpstype
  }
  \begin{syntax}
    \cs{@fpstype}
  \end{syntax}

\end{function}



\begin{function}[]
  {
    \@fptop
  }
  \begin{syntax}
    \cs{@fptop}
  \end{syntax}

\end{function}



\begin{function}[]
  {
    \@frameb@x
  }
  \begin{syntax}
    \cs{@frameb@x}
  \end{syntax}

\end{function}



\begin{function}[]
  {
    \@framebox
  }
  \begin{syntax}
    \cs{@framebox}
  \end{syntax}

\end{function}



\begin{function}[]
  {
    \@framepicbox
  }
  \begin{syntax}
    \cs{@framepicbox}
  \end{syntax}

\end{function}



\begin{function}[]
  {
    \@freelist
  }
  \begin{syntax}
    \cs{@freelist}
  \end{syntax}

\end{function}



\begin{function}[]
  {
    \@getcirc
  }
  \begin{syntax}
    \cs{@getcirc}
  \end{syntax}

\end{function}



\begin{function}[]
  {
    \@getfpsbit
  }
  \begin{syntax}
    \cs{@getfpsbit}
  \end{syntax}

\end{function}



\begin{function}[]
  {
    \@getlarrow
  }
  \begin{syntax}
    \cs{@getlarrow}
  \end{syntax}

\end{function}



\begin{function}[]
  {
    \@getlinechar
  }
  \begin{syntax}
    \cs{@getlinechar}
  \end{syntax}

\end{function}






\begin{function}[]
  {
    \@getrarrow
  }
  \begin{syntax}
    \cs{@getrarrow}
  \end{syntax}

\end{function}



\begin{function}[]
  {
    \@glossaryfile
  }
  \begin{syntax}
    \cs{@glossaryfile}
  \end{syntax}

\end{function}





\begin{function}[]
  {
    \@gobble
  }
  \begin{syntax}
    \cs{@gobble}
  \end{syntax}

\end{function}





\begin{function}[]
  {
    \@gobble@om
  }
  \begin{syntax}
    \cs{@gobble@om}
  \end{syntax}

\end{function}



\begin{function}[]
  {
    \@gobble@som
  }
  \begin{syntax}
    \cs{@gobble@som}
  \end{syntax}

\end{function}



\begin{function}[]
  {
    \@gobble@with@sphack@om
  }
  \begin{syntax}
    \cs{@gobble@with@sphack@om}
  \end{syntax}

\end{function}



\begin{function}[]
  {
    \@gobble@with@sphack@som
  }
  \begin{syntax}
    \cs{@gobble@with@sphack@som}
  \end{syntax}

\end{function}



\begin{function}[]
  {
    \@gobblecr
  }
  \begin{syntax}
    \cs{@gobblecr}
  \end{syntax}

\end{function}



\begin{function}[]
  {
    \@gobblefour
  }
  \begin{syntax}
    \cs{@gobblefour}
  \end{syntax}

\end{function}



\begin{function}[]
  {
    \@gobblethree
  }
  \begin{syntax}
    \cs{@gobblethree}
  \end{syntax}

\end{function}



\begin{function}[]
  {
    \@gobbletwo
  }
  \begin{syntax}
    \cs{@gobbletwo}
  \end{syntax}

\end{function}



\begin{function}[]
  {
    \@gtempa
  }
  \begin{syntax}
    \cs{@gtempa}
  \end{syntax}

\end{function}



\begin{function}[]
  {
    \@halfwidth
  }
  \begin{syntax}
    \cs{@halfwidth}
  \end{syntax}

\end{function}



\begin{function}[]
  {
    \@halignto
  }
  \begin{syntax}
    \cs{@halignto}
  \end{syntax}

\end{function}



\begin{function}[]
  {
    \@hangfrom
  }
  \begin{syntax}
    \cs{@hangfrom}
  \end{syntax}

\end{function}



\begin{function}[]
  {
    \@height
  }
  \begin{syntax}
    \cs{@height}
  \end{syntax}

\end{function}



\begin{function}[]
  {
    \@height.3ex
  }
  \begin{syntax}
    \cs{@height.3ex}
  \end{syntax}

\end{function}



\begin{function}[]
  {
    \@height.7
  }
  \begin{syntax}
    \cs{@height.7}
  \end{syntax}

\end{function}



\begin{function}[]
  {
    \@highpenalty
  }
  \begin{syntax}
    \cs{@highpenalty}
  \end{syntax}

\end{function}



\begin{function}[]
  {
    \@hightab
  }
  \begin{syntax}
    \cs{@hightab}
  \end{syntax}

\end{function}



\begin{function}[]
  {
    \@hline
  }
  \begin{syntax}
    \cs{@hline}
  \end{syntax}

\end{function}



\begin{function}[]
  {
    \@holdpg
  }
  \begin{syntax}
    \cs{@holdpg}
  \end{syntax}

\end{function}



\begin{function}[]
  {
    \@hspace
  }
  \begin{syntax}
    \cs{@hspace}
  \end{syntax}

\end{function}



\begin{function}[]
  {
    \@hspacer
  }
  \begin{syntax}
    \cs{@hspacer}
  \end{syntax}

\end{function}



\begin{function}[]
  {
    \@hvector
  }
  \begin{syntax}
    \cs{@hvector}
  \end{syntax}

\end{function}



\begin{function}[]
  {
    \@icentercr
  }
  \begin{syntax}
    \cs{@icentercr}
  \end{syntax}

\end{function}



\begin{function}[]
  {
    \@iden
  }
  \begin{syntax}
    \cs{@iden}
  \end{syntax}

\end{function}



\begin{function}[]
  {
    \@if
  }
  \begin{syntax}
    \cs{@if}
  \end{syntax}

\end{function}



\begin{function}[]
  {
    \@if@DeclareRobustCommand
  }
  \begin{syntax}
    \cs{@if@DeclareRobustCommand}
  \end{syntax}

\end{function}



\begin{function}[]
  {
    \@if@bottomfloats@TF
  }
  \begin{syntax}
    \cs{@if@bottomfloats@TF}
  \end{syntax}

\end{function}



\begin{function}[]
  {
    \@if@flushbottom@TF
  }
  \begin{syntax}
    \cs{@if@flushbottom@TF}
  \end{syntax}

\end{function}



\begin{function}[]
  {
    \@if@footnotes@TF
  }
  \begin{syntax}
    \cs{@if@footnotes@TF}
  \end{syntax}

\end{function}



\begin{function}[]
  {
    \@if@newcommand
  }
  \begin{syntax}
    \cs{@if@newcommand}
  \end{syntax}

\end{function}



\begin{function}[]
  {
    \@if@pti@ns
  }
  \begin{syntax}
    \cs{@if@pti@ns}
  \end{syntax}

\end{function}



\begin{function}[]
  {
    \@if@ptions
  }
  \begin{syntax}
    \cs{@if@ptions}
  \end{syntax}

\end{function}



\begin{function}[]
  {
    \@ifatmargin
  }
  \begin{syntax}
    \cs{@ifatmargin}
  \end{syntax}

\end{function}



\begin{function}[]
  {
    \@ifbothcounters
  }
  \begin{syntax}
    \cs{@ifbothcounters}
  \end{syntax}

\end{function}



\begin{function}[]
  {
    \@ifclasslater
  }
  \begin{syntax}
    \cs{@ifclasslater}
  \end{syntax}

\end{function}



\begin{function}[]
  {
    \@ifclassloaded
  }
  \begin{syntax}
    \cs{@ifclassloaded}
  \end{syntax}

\end{function}



\begin{function}[]
  {
    \@ifclasswith
  }
  \begin{syntax}
    \cs{@ifclasswith}
  \end{syntax}

\end{function}



\begin{function}[]
  {
    \@ifdefinable
  }
  \begin{syntax}
    \cs{@ifdefinable}
  \end{syntax}

\end{function}



\begin{function}[]
  {
    \@iffileonpath
  }
  \begin{syntax}
    \cs{@iffileonpath}
  \end{syntax}

\end{function}



\begin{function}[]
  {
    \@ifl@aded
  }
  \begin{syntax}
    \cs{@ifl@aded}
  \end{syntax}

\end{function}



\begin{function}[]
  {
    \@ifl@t@r
  }
  \begin{syntax}
    \cs{@ifl@t@r}
  \end{syntax}

\end{function}



\begin{function}[]
  {
    \@ifl@ter
  }
  \begin{syntax}
    \cs{@ifl@ter}
  \end{syntax}

\end{function}



\begin{function}[]
  {
    \@ifnch
  }
  \begin{syntax}
    \cs{@ifnch}
  \end{syntax}

\end{function}



\begin{function}[]
  {
    \@ifnextchar
  }
  \begin{syntax}
    \cs{@ifnextchar}
  \end{syntax}

\end{function}



\begin{function}[]
  {
    \@iforloop
  }
  \begin{syntax}
    \cs{@iforloop}
  \end{syntax}

\end{function}



\begin{function}[]
  {
    \@ifpackagelater
  }
  \begin{syntax}
    \cs{@ifpackagelater}
  \end{syntax}

\end{function}



\begin{function}[]
  {
    \@ifpackageloaded
  }
  \begin{syntax}
    \cs{@ifpackageloaded}
  \end{syntax}

\end{function}



\begin{function}[]
  {
    \@ifpackagewith
  }
  \begin{syntax}
    \cs{@ifpackagewith}
  \end{syntax}

\end{function}



\begin{function}[]
  {
    \@iframebox
  }
  \begin{syntax}
    \cs{@iframebox}
  \end{syntax}

\end{function}



\begin{function}[]
  {
    \@iframepicbox
  }
  \begin{syntax}
    \cs{@iframepicbox}
  \end{syntax}

\end{function}



\begin{function}[]
  {
    \@ifstar
  }
  \begin{syntax}
    \cs{@ifstar}
  \end{syntax}

\end{function}



\begin{function}[]
  {
    \@ifundefin@d@i
  }
  \begin{syntax}
    \cs{@ifundefin@d@i}
  \end{syntax}

\end{function}



\begin{function}[]
  {
    \@ifundefin@d@ii
  }
  \begin{syntax}
    \cs{@ifundefin@d@ii}
  \end{syntax}

\end{function}



\begin{function}[]
  {
    \@ifundefined
  }
  \begin{syntax}
    \cs{@ifundefined}
  \end{syntax}

\end{function}






\begin{function}[]
  {
    \@iiiminipage
  }
  \begin{syntax}
    \cs{@iiiminipage}
  \end{syntax}

\end{function}



\begin{function}[]
  {
    \@iiiparbox
  }
  \begin{syntax}
    \cs{@iiiparbox}
  \end{syntax}

\end{function}



\begin{function}[]
  {
    \@iiminipage
  }
  \begin{syntax}
    \cs{@iiminipage}
  \end{syntax}

\end{function}



\begin{function}[]
  {
    \@iinput
  }
  \begin{syntax}
    \cs{@iinput}
  \end{syntax}

\end{function}



\begin{function}[]
  {
    \@iiparbox
  }
  \begin{syntax}
    \cs{@iiparbox}
  \end{syntax}

\end{function}



\begin{function}[]
  {
    \@iirsbox
  }
  \begin{syntax}
    \cs{@iirsbox}
  \end{syntax}

\end{function}



\begin{function}[]
  {
    \@imakebox
  }
  \begin{syntax}
    \cs{@imakebox}
  \end{syntax}

\end{function}



\begin{function}[]
  {
    \@imakepicbox
  }
  \begin{syntax}
    \cs{@imakepicbox}
  \end{syntax}

\end{function}



\begin{function}[]
  {
    \@iminipage
  }
  \begin{syntax}
    \cs{@iminipage}
  \end{syntax}

\end{function}



\begin{function}[]
  {
    \@in@minipage@envtrue
  }
  \begin{syntax}
    \cs{@in@minipage@envtrue}
  \end{syntax}

\end{function}



\begin{function}[]
  {
    \@include
  }
  \begin{syntax}
    \cs{@include}
  \end{syntax}

\end{function}



\begin{function}[]
  {
    \@includeinreleasefalse
  }
  \begin{syntax}
    \cs{@includeinreleasefalse}
  \end{syntax}

\end{function}



\begin{function}[]
  {
    \@includeinreleasetrue
  }
  \begin{syntax}
    \cs{@includeinreleasetrue}
  \end{syntax}

\end{function}



\begin{function}[]
  {
    \@index
  }
  \begin{syntax}
    \cs{@index}
  \end{syntax}

\end{function}



\begin{function}[]
  {
    \@indexfile
  }
  \begin{syntax}
    \cs{@indexfile}
  \end{syntax}

\end{function}



\begin{function}[]
  {
    \@inlabelfalse
  }
  \begin{syntax}
    \cs{@inlabelfalse}
  \end{syntax}

\end{function}



\begin{function}[]
  {
    \@inlabeltrue
  }
  \begin{syntax}
    \cs{@inlabeltrue}
  \end{syntax}

\end{function}



\begin{function}[]
  {
    \@inmatherr
  }
  \begin{syntax}
    \cs{@inmatherr}
  \end{syntax}

\end{function}



\begin{function}[]
  {
    \@inmathwarn
  }
  \begin{syntax}
    \cs{@inmathwarn}
  \end{syntax}

\end{function}



\begin{function}[]
  {
    \@inpenc@test
  }
  \begin{syntax}
    \cs{@inpenc@test}
  \end{syntax}

\end{function}



\begin{function}[]
  {
    \@input
  }
  \begin{syntax}
    \cs{@input}
  \end{syntax}

\end{function}



\begin{function}[]
  {
    \@input@
  }
  \begin{syntax}
    \cs{@input@}
  \end{syntax}

\end{function}



\begin{function}[]
  {
    \@input@file@exists@with@hooks
  }
  \begin{syntax}
    \cs{@input@file@exists@with@hooks}
  \end{syntax}

\end{function}



\begin{function}[]
  {
    \@inputcheck
  }
  \begin{syntax}
    \cs{@inputcheck}
  \end{syntax}

\end{function}



\begin{function}[]
  {
    \@inputcheck"
  }
  \begin{syntax}
    \cs{@inputcheck"}
  \end{syntax}

\end{function}



\begin{function}[]
  {
    \@inputcheck0
  }
  \begin{syntax}
    \cs{@inputcheck0}
  \end{syntax}

\end{function}



\begin{function}[]
  {
    \@insertfalse
  }
  \begin{syntax}
    \cs{@insertfalse}
  \end{syntax}

\end{function}



\begin{function}[]
  {
    \@inserttrue
  }
  \begin{syntax}
    \cs{@inserttrue}
  \end{syntax}

\end{function}



\begin{function}[]
  {
    \@iparbox
  }
  \begin{syntax}
    \cs{@iparbox}
  \end{syntax}

\end{function}



\begin{function}[]
  {
    \@irsbox
  }
  \begin{syntax}
    \cs{@irsbox}
  \end{syntax}

\end{function}



\begin{function}[]
  {
    \@isavebox
  }
  \begin{syntax}
    \cs{@isavebox}
  \end{syntax}

\end{function}



\begin{function}[]
  {
    \@isavepicbox
  }
  \begin{syntax}
    \cs{@isavepicbox}
  \end{syntax}

\end{function}



\begin{function}[]
  {
    \@ishortstack
  }
  \begin{syntax}
    \cs{@ishortstack}
  \end{syntax}

\end{function}



\begin{function}[]
  {
    \@istackcr
  }
  \begin{syntax}
    \cs{@istackcr}
  \end{syntax}

\end{function}



\begin{function}[]
  {
    \@itabcr
  }
  \begin{syntax}
    \cs{@itabcr}
  \end{syntax}

\end{function}



\begin{function}[]
  {
    \@item
  }
  \begin{syntax}
    \cs{@item}
  \end{syntax}

\end{function}



\begin{function}[]
  {
    \@itemdepth
  }
  \begin{syntax}
    \cs{@itemdepth}
  \end{syntax}

\end{function}



\begin{function}[]
  {
    \@itemfudge
  }
  \begin{syntax}
    \cs{@itemfudge}
  \end{syntax}

\end{function}



\begin{function}[]
  {
    \@itemitem
  }
  \begin{syntax}
    \cs{@itemitem}
  \end{syntax}

\end{function}



\begin{function}[]
  {
    \@itemlabel
  }
  \begin{syntax}
    \cs{@itemlabel}
  \end{syntax}

\end{function}



\begin{function}[]
  {
    \@iwhiledim
  }
  \begin{syntax}
    \cs{@iwhiledim}
  \end{syntax}

\end{function}



\begin{function}[]
  {
    \@iwhilenum
  }
  \begin{syntax}
    \cs{@iwhilenum}
  \end{syntax}

\end{function}



\begin{function}[]
  {
    \@iwhilesw
  }
  \begin{syntax}
    \cs{@iwhilesw}
  \end{syntax}

\end{function}



\begin{function}[]
  {
    \@ixpt
  }
  \begin{syntax}
    \cs{@ixpt}
  \end{syntax}

\end{function}



\begin{function}[]
  {
    \@ixstackcr
  }
  \begin{syntax}
    \cs{@ixstackcr}
  \end{syntax}

\end{function}



\begin{function}[int]
  {
    \@kernel@Ref
  }
  \begin{syntax}
    \cs{@kernel@Ref}
  \end{syntax}

\end{function}



\begin{function}[int]
  {
    \@kernel@after@begindocument
  }
  \begin{syntax}
    \cs{@kernel@after@begindocument}
  \end{syntax}

\end{function}



\begin{function}[int]
  {
    \@kernel@after@begindocument@before
  }
  \begin{syntax}
    \cs{@kernel@after@begindocument@before}
  \end{syntax}

\end{function}



\begin{function}[int]
  {
    \@kernel@after@enddocument
  }
  \begin{syntax}
    \cs{@kernel@after@enddocument}
  \end{syntax}

\end{function}



\begin{function}[int]
  {
    \@kernel@after@enddocument@afterlastpage
  }
  \begin{syntax}
    \cs{@kernel@after@enddocument@afterlastpage}
  \end{syntax}

\end{function}



\begin{function}[int]
  {
    \@kernel@after@para@after
  }
  \begin{syntax}
    \cs{@kernel@after@para@after}
  \end{syntax}

\end{function}



\begin{function}[int]
  {
    \@kernel@after@para@end
  }
  \begin{syntax}
    \cs{@kernel@after@para@end}
  \end{syntax}

\end{function}



\begin{function}[int]
  {
    \@kernel@after@shipout@background
  }
  \begin{syntax}
    \cs{@kernel@after@shipout@background}
  \end{syntax}

\end{function}



\begin{function}[int]
  {
    \@kernel@after@shipout@lastpage
  }
  \begin{syntax}
    \cs{@kernel@after@shipout@lastpage}
  \end{syntax}

\end{function}



\begin{function}[int]
  {
    \@kernel@before@begindocument
  }
  \begin{syntax}
    \cs{@kernel@before@begindocument}
  \end{syntax}

\end{function}



\begin{function}[int]
  {
    \@kernel@before@enddocument
  }
  \begin{syntax}
    \cs{@kernel@before@enddocument}
  \end{syntax}

\end{function}



\begin{function}[int]
  {
    \@kernel@before@enddocument@afterlastpage
  }
  \begin{syntax}
    \cs{@kernel@before@enddocument@afterlastpage}
  \end{syntax}

\end{function}



\begin{function}[int]
  {
    \@kernel@before@insertmark
  }
  \begin{syntax}
    \cs{@kernel@before@insertmark}
  \end{syntax}

\end{function}



\begin{function}[int]
  {
    \@kernel@before@para@before
  }
  \begin{syntax}
    \cs{@kernel@before@para@before}
  \end{syntax}

\end{function}



\begin{function}[int]
  {
    \@kernel@before@para@begin
  }
  \begin{syntax}
    \cs{@kernel@before@para@begin}
  \end{syntax}

\end{function}



\begin{function}[int]
  {
    \@kernel@before@shipout@background
  }
  \begin{syntax}
    \cs{@kernel@before@shipout@background}
  \end{syntax}

\end{function}



\begin{function}[int]
  {
    \@kernel@currpathstack
  }
  \begin{syntax}
    \cs{@kernel@currpathstack}
  \end{syntax}

\end{function}




\begin{function}[]
  {
    \@kernel@get@ctr@root
  }
  \begin{syntax}
    \cs{@kernel@get@ctr@root}
  \end{syntax}

\end{function}





\begin{function}[]
  {
    \@kernel@make@file@csname
  }
  \begin{syntax}
    \cs{@kernel@make@file@csname}
  \end{syntax}

\end{function}



\begin{function}[]
  {
    \@kernel@new@label@record@testdef
  }
  \begin{syntax}
    \cs{@kernel@new@label@record@testdef}
  \end{syntax}

\end{function}



\begin{function}[]
  {
    \@kernel@pageref
  }
  \begin{syntax}
    \cs{@kernel@pageref}
  \end{syntax}

\end{function}



\begin{function}[]
  {
    \@kernel@pageref@exp
  }
  \begin{syntax}
    \cs{@kernel@pageref@exp}
  \end{syntax}

\end{function}



\begin{function}[]
  {
    \@kernel@ref
  }
  \begin{syntax}
    \cs{@kernel@ref}
  \end{syntax}

\end{function}



\begin{function}[]
  {
    \@kernel@ref@exp
  }
  \begin{syntax}
    \cs{@kernel@ref@exp}
  \end{syntax}

\end{function}



\begin{function}[]
  {
    \@kernel@refstepcounter
  }
  \begin{syntax}
    \cs{@kernel@refstepcounter}
  \end{syntax}

\end{function}



\begin{function}[]
  {
    \@kernel@rename@newcommand
  }
  \begin{syntax}
    \cs{@kernel@rename@newcommand}
  \end{syntax}

\end{function}



\begin{function}[]
  {
    \@kernel@reserved@label@data
  }
  \begin{syntax}
    \cs{@kernel@reserved@label@data}
  \end{syntax}

\end{function}



\begin{function}[]
  {
    \@kernel@sRef
  }
  \begin{syntax}
    \cs{@kernel@sRef}
  \end{syntax}

\end{function}



\begin{function}[]
  {
    \@kernel@spageref
  }
  \begin{syntax}
    \cs{@kernel@spageref}
  \end{syntax}

\end{function}



\begin{function}[]
  {
    \@kernel@sref
  }
  \begin{syntax}
    \cs{@kernel@sref}
  \end{syntax}

\end{function}



\begin{function}[]
  {
    \@kernel@whatsit
  }
  \begin{syntax}
    \cs{@kernel@whatsit}
  \end{syntax}

\end{function}



\begin{function}[]
  {
    \@killglue
  }
  \begin{syntax}
    \cs{@killglue}
  \end{syntax}

\end{function}



\begin{function}[]
  {
    \@kludgeins
  }
  \begin{syntax}
    \cs{@kludgeins}
  \end{syntax}

\end{function}



\begin{function}[]
  {
    \@labels
  }
  \begin{syntax}
    \cs{@labels}
  \end{syntax}

\end{function}



\begin{function}[]
  {
    \@largefloatcheck
  }
  \begin{syntax}
    \cs{@largefloatcheck}
  \end{syntax}

\end{function}



\begin{function}[]
  {
    \@lastchclass
  }
  \begin{syntax}
    \cs{@lastchclass}
  \end{syntax}

\end{function}



\begin{function}[]
  {
    \@lastchclass6
  }
  \begin{syntax}
    \cs{@lastchclass6}
  \end{syntax}

\end{function}



\begin{function}[]
  {
    \@latex@error
  }
  \begin{syntax}
    \cs{@latex@error}
  \end{syntax}

\end{function}



\begin{function}[]
  {
    \@latex@info
  }
  \begin{syntax}
    \cs{@latex@info}
  \end{syntax}

\end{function}



\begin{function}[]
  {
    \@latex@info@no@line
  }
  \begin{syntax}
    \cs{@latex@info@no@line}
  \end{syntax}

\end{function}



\begin{function}[]
  {
    \@latex@note
  }
  \begin{syntax}
    \cs{@latex@note}
  \end{syntax}

\end{function}



\begin{function}[]
  {
    \@latex@note@no@line
  }
  \begin{syntax}
    \cs{@latex@note@no@line}
  \end{syntax}

\end{function}



\begin{function}[]
  {
    \@latex@warning
  }
  \begin{syntax}
    \cs{@latex@warning}
  \end{syntax}

\end{function}



\begin{function}[]
  {
    \@latex@warning@no@line
  }
  \begin{syntax}
    \cs{@latex@warning@no@line}
  \end{syntax}

\end{function}



\begin{function}[]
  {
    \@latexbug
  }
  \begin{syntax}
    \cs{@latexbug}
  \end{syntax}

\end{function}



\begin{function}[int,deprecated]
  {
    \@latexerr
  }
  \begin{syntax}
    \cs{@latexerr} \Arg{message} \Arg{help text}
  \end{syntax}
  Another name for \cs{@latex@error}. Deprecated.

\end{function}



\begin{function}[]
  {
    \@lbibitem
  }
  \begin{syntax}
    \cs{@lbibitem}
  \end{syntax}

\end{function}



\begin{function}[]
  {
    \@leftcolumn
  }
  \begin{syntax}
    \cs{@leftcolumn}
  \end{syntax}

\end{function}



\begin{function}[]
  {
    \@leftmark
  }
  \begin{syntax}
    \cs{@leftmark}
  \end{syntax}

\end{function}



\begin{function}[]
  {
    \@let@token
  }
  \begin{syntax}
    \cs{@let@token}
  \end{syntax}

\end{function}



\begin{function}[]
  {
    \@lign
  }
  \begin{syntax}
    \cs{@lign}
  \end{syntax}

\end{function}



\begin{function}[]
  {
    \@linechar
  }
  \begin{syntax}
    \cs{@linechar}
  \end{syntax}

\end{function}



\begin{function}[]
  {
    \@linefnt
  }
  \begin{syntax}
    \cs{@linefnt}
  \end{syntax}

\end{function}



\begin{function}[]
  {
    \@linelen
  }
  \begin{syntax}
    \cs{@linelen}
  \end{syntax}

\end{function}



\begin{function}[]
  {
    \@listctr
  }
  \begin{syntax}
    \cs{@listctr}
  \end{syntax}

\end{function}





\begin{function}[]
  {
    \@listfiles
  }
  \begin{syntax}
    \cs{@listfiles}
  \end{syntax}

\end{function}



\begin{function}[]
  {
    \@loadwithoptions
  }
  \begin{syntax}
    \cs{@loadwithoptions}
  \end{syntax}

\end{function}



\begin{function}[]
  {
    \@lowpenalty
  }
  \begin{syntax}
    \cs{@lowpenalty}
  \end{syntax}

\end{function}



\begin{function}[]
  {
    \@ltab
  }
  \begin{syntax}
    \cs{@ltab}
  \end{syntax}

\end{function}



\begin{function}[]
  {
    \@m
  }
  \begin{syntax}
    \cs{@m}
  \end{syntax}

\end{function}



\begin{function}[]
  {
    \@mainaux
  }
  \begin{syntax}
    \cs{@mainaux}
  \end{syntax}

\end{function}



\begin{function}[]
  {
    \@make@normalcolbox
  }
  \begin{syntax}
    \cs{@make@normalcolbox}
  \end{syntax}

\end{function}



\begin{function}[]
  {
    \@make@specialcolbox
  }
  \begin{syntax}
    \cs{@make@specialcolbox}
  \end{syntax}

\end{function}



\begin{function}[]
  {
    \@makebox
  }
  \begin{syntax}
    \cs{@makebox}
  \end{syntax}

\end{function}



\begin{function}[]
  {
    \@makecaption
  }
  \begin{syntax}
    \cs{@makecaption}
  \end{syntax}

\end{function}



\begin{function}[]
  {
    \@makecol
  }
  \begin{syntax}
    \cs{@makecol}
  \end{syntax}

\end{function}



\begin{function}[]
  {
    \@makecol@handlesplitfootnotes
  }
  \begin{syntax}
    \cs{@makecol@handlesplitfootnotes}
  \end{syntax}

\end{function}



\begin{function}[]
  {
    \@makefcolumn
  }
  \begin{syntax}
    \cs{@makefcolumn}
  \end{syntax}

\end{function}



\begin{function}[]
  {
    \@makefnmark
  }
  \begin{syntax}
    \cs{@makefnmark}
  \end{syntax}

\end{function}



\begin{function}[]
  {
    \@makefntext
  }
  \begin{syntax}
    \cs{@makefntext}
  \end{syntax}

\end{function}



\begin{function}[]
  {
    \@makeother
  }
  \begin{syntax}
    \cs{@makeother}
  \end{syntax}

\end{function}



\begin{function}[]
  {
    \@makepicbox
  }
  \begin{syntax}
    \cs{@makepicbox}
  \end{syntax}

\end{function}



\begin{function}[]
  {
    \@marbox
  }
  \begin{syntax}
    \cs{@marbox}
  \end{syntax}

\end{function}



\begin{function}[]
  {
    \@marginparreset
  }
  \begin{syntax}
    \cs{@marginparreset}
  \end{syntax}

\end{function}



\begin{function}[]
  {
    \@math@level
  }
  \begin{syntax}
    \cs{@math@level}
  \end{syntax}

\end{function}



\begin{function}[]
  {
    \@maxdepth
  }
  \begin{syntax}
    \cs{@maxdepth}
  \end{syntax}

\end{function}



\begin{function}[]
  {
    \@maxtab
  }
  \begin{syntax}
    \cs{@maxtab}
  \end{syntax}

\end{function}



\begin{function}[]
  {
    \@medpenalty
  }
  \begin{syntax}
    \cs{@medpenalty}
  \end{syntax}

\end{function}



\begin{function}[]
  {
    \@meta@family@list
  }
  \begin{syntax}
    \cs{@meta@family@list}
  \end{syntax}

\end{function}



\begin{function}[]
  {
    \@midlist
  }
  \begin{syntax}
    \cs{@midlist}
  \end{syntax}

\end{function}



\begin{function}[]
  {
    \@minipagefalse
  }
  \begin{syntax}
    \cs{@minipagefalse}
  \end{syntax}

\end{function}



\begin{function}[]
  {
    \@minipagerestore
  }
  \begin{syntax}
    \cs{@minipagerestore}
  \end{syntax}

\end{function}



\begin{function}[]
  {
    \@minipagetrue
  }
  \begin{syntax}
    \cs{@minipagetrue}
  \end{syntax}

\end{function}



\begin{function}[]
  {
    \@minus
  }
  \begin{syntax}
    \cs{@minus}
  \end{syntax}

\end{function}



\begin{function}[]
  {
    \@missing@onefilewithoptions
  }
  \begin{syntax}
    \cs{@missing@onefilewithoptions}
  \end{syntax}

\end{function}



\begin{function}[]
  {
    \@missingfile@area
  }
  \begin{syntax}
    \cs{@missingfile@area}
  \end{syntax}

\end{function}



\begin{function}[]
  {
    \@missingfile@base
  }
  \begin{syntax}
    \cs{@missingfile@base}
  \end{syntax}

\end{function}



\begin{function}[]
  {
    \@missingfile@ext
  }
  \begin{syntax}
    \cs{@missingfile@ext}
  \end{syntax}

\end{function}



\begin{function}[]
  {
    \@missingfileerror
  }
  \begin{syntax}
    \cs{@missingfileerror}
  \end{syntax}

\end{function}



\begin{function}[]
  {
    \@mkboth
  }
  \begin{syntax}
    \cs{@mkboth}
  \end{syntax}

\end{function}



\begin{function}[]
  {
    \@mklab
  }
  \begin{syntax}
    \cs{@mklab}
  \end{syntax}

\end{function}



\begin{function}[]
  {
    \@mkpream
  }
  \begin{syntax}
    \cs{@mkpream}
  \end{syntax}

\end{function}



\begin{function}[]
  {
    \@mparbottom
  }
  \begin{syntax}
    \cs{@mparbottom}
  \end{syntax}

\end{function}



\begin{function}[]
  {
    \@mpargs
  }
  \begin{syntax}
    \cs{@mpargs}
  \end{syntax}

\end{function}



\begin{function}[]
  {
    \@mparswitchfalse
  }
  \begin{syntax}
    \cs{@mparswitchfalse}
  \end{syntax}

\end{function}



\begin{function}[]
  {
    \@mpfn
  }
  \begin{syntax}
    \cs{@mpfn}
  \end{syntax}

\end{function}



\begin{function}[]
  {
    \@mpfootins
  }
  \begin{syntax}
    \cs{@mpfootins}
  \end{syntax}

\end{function}



\begin{function}[]
  {
    \@mpfootnotetext
  }
  \begin{syntax}
    \cs{@mpfootnotetext}
  \end{syntax}

\end{function}



\begin{function}[]
  {
    \@mplistdepth
  }
  \begin{syntax}
    \cs{@mplistdepth}
  \end{syntax}

\end{function}



\begin{function}[]
  {
    \@multicnt
  }
  \begin{syntax}
    \cs{@multicnt}
  \end{syntax}

\end{function}



\begin{function}[]
  {
    \@multiplelabels
  }
  \begin{syntax}
    \cs{@multiplelabels}
  \end{syntax}

\end{function}



\begin{function}[]
  {
    \@multiput
  }
  \begin{syntax}
    \cs{@multiput}
  \end{syntax}

\end{function}



\begin{function}[]
  {
    \@multispan
  }
  \begin{syntax}
    \cs{@multispan}
  \end{syntax}

\end{function}



\begin{function}[]
  {
    \@namedef
  }
  \begin{syntax}
    \cs{@namedef}
  \end{syntax}

\end{function}



\begin{function}[]
  {
    \@nameuse
  }
  \begin{syntax}
    \cs{@nameuse}
  \end{syntax}

\end{function}



\begin{function}[]
  {
    \@nbitem
  }
  \begin{syntax}
    \cs{@nbitem}
  \end{syntax}

\end{function}



\begin{function}[]
  {
    \@ne
  }
  \begin{syntax}
    \cs{@ne}
  \end{syntax}

\end{function}



\begin{function}[]
  {
    \@needsf@rmat
  }
  \begin{syntax}
    \cs{@needsf@rmat}
  \end{syntax}

\end{function}



\begin{function}[]
  {
    \@needsformat
  }
  \begin{syntax}
    \cs{@needsformat}
  \end{syntax}

\end{function}



\begin{function}[]
  {
    \@negargfalse
  }
  \begin{syntax}
    \cs{@negargfalse}
  \end{syntax}

\end{function}



\begin{function}[]
  {
    \@negargtrue
  }
  \begin{syntax}
    \cs{@negargtrue}
  \end{syntax}

\end{function}



\begin{function}[]
  {
    \@newcommand
  }
  \begin{syntax}
    \cs{@newcommand}
  \end{syntax}

\end{function}



\begin{function}[]
  {
    \@newctr
  }
  \begin{syntax}
    \cs{@newctr}
  \end{syntax}

\end{function}



\begin{function}[]
  {
    \@newenv
  }
  \begin{syntax}
    \cs{@newenv}
  \end{syntax}

\end{function}



\begin{function}[]
  {
    \@newenva
  }
  \begin{syntax}
    \cs{@newenva}
  \end{syntax}

\end{function}



\begin{function}[]
  {
    \@newenvb
  }
  \begin{syntax}
    \cs{@newenvb}
  \end{syntax}

\end{function}



\begin{function}[]
  {
    \@newl@bel
  }
  \begin{syntax}
    \cs{@newl@bel}
  \end{syntax}

\end{function}





\begin{function}[]
  {
    \@newlistfalse
  }
  \begin{syntax}
    \cs{@newlistfalse}
  \end{syntax}

\end{function}



\begin{function}[]
  {
    \@newlisttrue
  }
  \begin{syntax}
    \cs{@newlisttrue}
  \end{syntax}

\end{function}



\begin{function}[]
  {
    \@next
  }
  \begin{syntax}
    \cs{@next}
  \end{syntax}

\end{function}



\begin{function}[]
  {
    \@nextchar
  }
  \begin{syntax}
    \cs{@nextchar}
  \end{syntax}

\end{function}



\begin{function}[]
  {
    \@nil
  }
  \begin{syntax}
    \cs{@nil}
  \end{syntax}

\end{function}



\begin{function}[]
  {
    \@nmbrlistfalse
  }
  \begin{syntax}
    \cs{@nmbrlistfalse}
  \end{syntax}

\end{function}



\begin{function}[]
  {
    \@nmbrlisttrue
  }
  \begin{syntax}
    \cs{@nmbrlisttrue}
  \end{syntax}

\end{function}



\begin{function}[]
  {
    \@nnil
  }
  \begin{syntax}
    \cs{@nnil}
  \end{syntax}

\end{function}








\begin{function}[]
  {
    \@nocnterr
  }
  \begin{syntax}
    \cs{@nocnterr}
  \end{syntax}

\end{function}





\begin{function}[]
  {
    \@nodocument
  }
  \begin{syntax}
    \cs{@nodocument}
  \end{syntax}

\end{function}



\begin{function}[]
  {
    \@noitemargfalse
  }
  \begin{syntax}
    \cs{@noitemargfalse}
  \end{syntax}

\end{function}



\begin{function}[]
  {
    \@noitemargtrue
  }
  \begin{syntax}
    \cs{@noitemargtrue}
  \end{syntax}

\end{function}



\begin{function}[]
  {
    \@noitemerr
  }
  \begin{syntax}
    \cs{@noitemerr}
  \end{syntax}

\end{function}



\begin{function}[]
  {
    \@noligs
  }
  \begin{syntax}
    \cs{@noligs}
  \end{syntax}

\end{function}



\begin{function}[]
  {
    \@nolnerr
  }
  \begin{syntax}
    \cs{@nolnerr}
  \end{syntax}

\end{function}



\begin{function}[]
  {
    \@nomath
  }
  \begin{syntax}
    \cs{@nomath}
  \end{syntax}

\end{function}







\begin{function}[]
  {
    \@noparlistfalse
  }
  \begin{syntax}
    \cs{@noparlistfalse}
  \end{syntax}

\end{function}



\begin{function}[]
  {
    \@noparlisttrue
  }
  \begin{syntax}
    \cs{@noparlisttrue}
  \end{syntax}

\end{function}




\begin{function}[]
  {
    \@normalsize
  }
  \begin{syntax}
    \cs{@normalsize}
  \end{syntax}

\end{function}



\begin{function}[]
  {
    \@noskipsecfalse
  }
  \begin{syntax}
    \cs{@noskipsecfalse}
  \end{syntax}

\end{function}



\begin{function}[]
  {
    \@noskipsectrue
  }
  \begin{syntax}
    \cs{@noskipsectrue}
  \end{syntax}

\end{function}



\begin{function}[]
  {
    \@notdefinable
  }
  \begin{syntax}
    \cs{@notdefinable}
  \end{syntax}

\end{function}



\begin{function}[]
  {
    \@notprerr
  }
  \begin{syntax}
    \cs{@notprerr}
  \end{syntax}

\end{function}



\begin{function}[]
  {
    \@nthm
  }
  \begin{syntax}
    \cs{@nthm}
  \end{syntax}

\end{function}



\begin{function}[]
  {
    \@nxttabmar
  }
  \begin{syntax}
    \cs{@nxttabmar}
  \end{syntax}

\end{function}



\begin{function}[]
  {
    \@obsoletefile
  }
  \begin{syntax}
    \cs{@obsoletefile}
  \end{syntax}

\end{function}



\begin{function}[]
  {
    \@oddfoot
  }
  \begin{syntax}
    \cs{@oddfoot}
  \end{syntax}

\end{function}



\begin{function}[]
  {
    \@oddhead
  }
  \begin{syntax}
    \cs{@oddhead}
  \end{syntax}

\end{function}



\begin{function}[]
  {
    \@onefilewithoptions
  }
  \begin{syntax}
    \cs{@onefilewithoptions}
  \end{syntax}

\end{function}



\begin{function}[]
  {
    \@onefilewithoptions@clashchk
  }
  \begin{syntax}
    \cs{@onefilewithoptions@clashchk}
  \end{syntax}

\end{function}



\begin{function}[]
  {
    \@onelevel@sanitize
  }
  \begin{syntax}
    \cs{@onelevel@sanitize}
  \end{syntax}

\end{function}



\begin{function}[]
  {
    \@onlypreamble
  }
  \begin{syntax}
    \cs{@onlypreamble}
  \end{syntax}

\end{function}



\begin{function}[]
  {
    \@opargbegintheorem
  }
  \begin{syntax}
    \cs{@opargbegintheorem}
  \end{syntax}

\end{function}



\begin{function}[]
  {
    \@opcol
  }
  \begin{syntax}
    \cs{@opcol}
  \end{syntax}

\end{function}



\begin{function}[]
  {
    \@options
  }
  \begin{syntax}
    \cs{@options}
  \end{syntax}

\end{function}



\begin{function}[]
  {
    \@othm
  }
  \begin{syntax}
    \cs{@othm}
  \end{syntax}

\end{function}



\begin{function}[]
  {
    \@outerparskip
  }
  \begin{syntax}
    \cs{@outerparskip}
  \end{syntax}

\end{function}



\begin{function}[]
  {
    \@outputbox
  }
  \begin{syntax}
    \cs{@outputbox}
  \end{syntax}

\end{function}



\begin{function}[]
  {
    \@outputbox@append
  }
  \begin{syntax}
    \cs{@outputbox@append}
  \end{syntax}

\end{function}



\begin{function}[]
  {
    \@outputbox@appendfil
  }
  \begin{syntax}
    \cs{@outputbox@appendfil}
  \end{syntax}

\end{function}



\begin{function}[]
  {
    \@outputbox@appendfootnotes
  }
  \begin{syntax}
    \cs{@outputbox@appendfootnotes}
  \end{syntax}

\end{function}



\begin{function}[]
  {
    \@outputbox@attachbottomfloats
  }
  \begin{syntax}
    \cs{@outputbox@attachbottomfloats}
  \end{syntax}

\end{function}



\begin{function}[]
  {
    \@outputbox@attachfloats
  }
  \begin{syntax}
    \cs{@outputbox@attachfloats}
  \end{syntax}

\end{function}



\begin{function}[]
  {
    \@outputbox@attachtopfloats
  }
  \begin{syntax}
    \cs{@outputbox@attachtopfloats}
  \end{syntax}

\end{function}



\begin{function}[]
  {
    \@outputbox@depth
  }
  \begin{syntax}
    \cs{@outputbox@depth}
  \end{syntax}

\end{function}



\begin{function}[]
  {
    \@outputbox@reinsertbskip
  }
  \begin{syntax}
    \cs{@outputbox@reinsertbskip}
  \end{syntax}

\end{function}



\begin{function}[]
  {
    \@outputbox@removebskip
  }
  \begin{syntax}
    \cs{@outputbox@removebskip}
  \end{syntax}

\end{function}



\begin{function}[]
  {
    \@outputdblcol
  }
  \begin{syntax}
    \cs{@outputdblcol}
  \end{syntax}

\end{function}



\begin{function}[]
  {
    \@outputpage
  }
  \begin{syntax}
    \cs{@outputpage}
  \end{syntax}

\end{function}



\begin{function}[]
  {
    \@oval
  }
  \begin{syntax}
    \cs{@oval}
  \end{syntax}

\end{function}



\begin{function}[]
  {
    \@ovbtrue
  }
  \begin{syntax}
    \cs{@ovbtrue}
  \end{syntax}

\end{function}



\begin{function}[]
  {
    \@ovdx
  }
  \begin{syntax}
    \cs{@ovdx}
  \end{syntax}

\end{function}



\begin{function}[]
  {
    \@ovdy
  }
  \begin{syntax}
    \cs{@ovdy}
  \end{syntax}

\end{function}



\begin{function}[]
  {
    \@ovhlinefalse
  }
  \begin{syntax}
    \cs{@ovhlinefalse}
  \end{syntax}

\end{function}



\begin{function}[]
  {
    \@ovhlinetrue
  }
  \begin{syntax}
    \cs{@ovhlinetrue}
  \end{syntax}

\end{function}



\begin{function}[]
  {
    \@ovhorz
  }
  \begin{syntax}
    \cs{@ovhorz}
  \end{syntax}

\end{function}



\begin{function}[]
  {
    \@ovltrue
  }
  \begin{syntax}
    \cs{@ovltrue}
  \end{syntax}

\end{function}



\begin{function}[]
  {
    \@ovri
  }
  \begin{syntax}
    \cs{@ovri}
  \end{syntax}

\end{function}



\begin{function}[]
  {
    \@ovro
  }
  \begin{syntax}
    \cs{@ovro}
  \end{syntax}

\end{function}



\begin{function}[]
  {
    \@ovrtrue
  }
  \begin{syntax}
    \cs{@ovrtrue}
  \end{syntax}

\end{function}



\begin{function}[]
  {
    \@ovttrue
  }
  \begin{syntax}
    \cs{@ovttrue}
  \end{syntax}

\end{function}



\begin{function}[]
  {
    \@ovvert
  }
  \begin{syntax}
    \cs{@ovvert}
  \end{syntax}

\end{function}



\begin{function}[]
  {
    \@ovvert01
  }
  \begin{syntax}
    \cs{@ovvert01}
  \end{syntax}

\end{function}



\begin{function}[]
  {
    \@ovvert32
  }
  \begin{syntax}
    \cs{@ovvert32}
  \end{syntax}

\end{function}



\begin{function}[]
  {
    \@ovvlinefalse
  }
  \begin{syntax}
    \cs{@ovvlinefalse}
  \end{syntax}

\end{function}



\begin{function}[]
  {
    \@ovvlinetrue
  }
  \begin{syntax}
    \cs{@ovvlinetrue}
  \end{syntax}

\end{function}



\begin{function}[]
  {
    \@ovxx
  }
  \begin{syntax}
    \cs{@ovxx}
  \end{syntax}

\end{function}



\begin{function}[]
  {
    \@ovyy
  }
  \begin{syntax}
    \cs{@ovyy}
  \end{syntax}

\end{function}



\begin{function}[]
  {
    \@p@pfilename
  }
  \begin{syntax}
    \cs{@p@pfilename}
  \end{syntax}

\end{function}



\begin{function}[]
  {
    \@p@pfilepath
  }
  \begin{syntax}
    \cs{@p@pfilepath}
  \end{syntax}

\end{function}



\begin{function}[]
  {
    \@p@pfilepath@aux
  }
  \begin{syntax}
    \cs{@p@pfilepath@aux}
  \end{syntax}

\end{function}



\begin{function}[]
  {
    \@pagedp
  }
  \begin{syntax}
    \cs{@pagedp}
  \end{syntax}

\end{function}



\begin{function}[]
  {
    \@pageht
  }
  \begin{syntax}
    \cs{@pageht}
  \end{syntax}

\end{function}



\begin{function}[]
  {
    \@par
  }
  \begin{syntax}
    \cs{@par}
  \end{syntax}

\end{function}



\begin{function}[]
  {
    \@parboxrestore
  }
  \begin{syntax}
    \cs{@parboxrestore}
  \end{syntax}

\end{function}



\begin{function}[]
  {
    \@parboxto
  }
  \begin{syntax}
    \cs{@parboxto}
  \end{syntax}

\end{function}



\begin{function}[]
  {
    \@parmoderr
  }
  \begin{syntax}
    \cs{@parmoderr}
  \end{syntax}

\end{function}



\begin{function}[]
  {
    \@parse@version
  }
  \begin{syntax}
    \cs{@parse@version}
  \end{syntax}

\end{function}



\begin{function}[]
  {
    \@parse@version0
  }
  \begin{syntax}
    \cs{@parse@version0}
  \end{syntax}

\end{function}



\begin{function}[]
  {
    \@parse@version@
  }
  \begin{syntax}
    \cs{@parse@version@}
  \end{syntax}

\end{function}



\begin{function}[]
  {
    \@parse@version@dash
  }
  \begin{syntax}
    \cs{@parse@version@dash}
  \end{syntax}

\end{function}



\begin{function}[]
  {
    \@partaux
  }
  \begin{syntax}
    \cs{@partaux}
  \end{syntax}

\end{function}



\begin{function}[]
  {
    \@partlist
  }
  \begin{syntax}
    \cs{@partlist}
  \end{syntax}

\end{function}



\begin{function}[]
  {
    \@partswfalse
  }
  \begin{syntax}
    \cs{@partswfalse}
  \end{syntax}

\end{function}



\begin{function}[]
  {
    \@partswtrue
  }
  \begin{syntax}
    \cs{@partswtrue}
  \end{syntax}

\end{function}



\begin{function}[]
  {
    \@pass@ptions
  }
  \begin{syntax}
    \cs{@pass@ptions}
  \end{syntax}

\end{function}



\begin{function}[]
  {
    \@pboxswfalse
  }
  \begin{syntax}
    \cs{@pboxswfalse}
  \end{syntax}

\end{function}



\begin{function}[]
  {
    \@penup
  }
  \begin{syntax}
    \cs{@penup}
  \end{syntax}

\end{function}





\begin{function}[]
  {
    \@picbox
  }
  \begin{syntax}
    \cs{@picbox}
  \end{syntax}

\end{function}



\begin{function}[]
  {
    \@picht
  }
  \begin{syntax}
    \cs{@picht}
  \end{syntax}

\end{function}



\begin{function}[]
  {
    \@picture
  }
  \begin{syntax}
    \cs{@picture}
  \end{syntax}

\end{function}



\begin{function}[]
  {
    \@picture@warn
  }
  \begin{syntax}
    \cs{@picture@warn}
  \end{syntax}

\end{function}



\begin{function}[]
  {
    \@pkgextension
  }
  \begin{syntax}
    \cs{@pkgextension}
  \end{syntax}

\end{function}



\begin{function}[]
  {
    \@plus
  }
  \begin{syntax}
    \cs{@plus}
  \end{syntax}

\end{function}



\begin{function}[]
  {
    \@pnumwidth
  }
  \begin{syntax}
    \cs{@pnumwidth}
  \end{syntax}

\end{function}



\begin{function}[]
  {
    \@popfilename
  }
  \begin{syntax}
    \cs{@popfilename}
  \end{syntax}

\end{function}



\begin{function}[]
  {
    \@pr@videpackage
  }
  \begin{syntax}
    \cs{@pr@videpackage}
  \end{syntax}

\end{function}



\begin{function}[]
  {
    \@preamble
  }
  \begin{syntax}
    \cs{@preamble}
  \end{syntax}

\end{function}



\begin{function}[]
  {
    \@preamblecmds
  }
  \begin{syntax}
    \cs{@preamblecmds}
  \end{syntax}

\end{function}



\begin{function}[]
  {
    \@preamerr
  }
  \begin{syntax}
    \cs{@preamerr}
  \end{syntax}

\end{function}



\begin{function}[]
  {
    \@process@pti@ns
  }
  \begin{syntax}
    \cs{@process@pti@ns}
  \end{syntax}

\end{function}



\begin{function}[]
  {
    \@process@ptions
  }
  \begin{syntax}
    \cs{@process@ptions}
  \end{syntax}

\end{function}



\begin{function}[]
  {
    \@protected@testopt
  }
  \begin{syntax}
    \cs{@protected@testopt}
  \end{syntax}

\end{function}



\begin{function}[]
  {
    \@providesfile
  }
  \begin{syntax}
    \cs{@providesfile}
  \end{syntax}

\end{function}



\begin{function}[]
  {
    \@ptionlist
  }
  \begin{syntax}
    \cs{@ptionlist}
  \end{syntax}

\end{function}



\begin{function}[]
  {
    \@pushfilename
  }
  \begin{syntax}
    \cs{@pushfilename}
  \end{syntax}

\end{function}



\begin{function}[]
  {
    \@put
  }
  \begin{syntax}
    \cs{@put}
  \end{syntax}

\end{function}



\begin{function}[]
  {
    \@qend
  }
  \begin{syntax}
    \cs{@qend}
  \end{syntax}

\end{function}



\begin{function}[]
  {
    \@qrelax
  }
  \begin{syntax}
    \cs{@qrelax}
  \end{syntax}

\end{function}



\begin{function}[]
  {
    \@raw@classoptionslist
  }
  \begin{syntax}
    \cs{@raw@classoptionslist}
  \end{syntax}

\end{function}



\begin{function}[]
  {
    \@rc@ifdefinable
  }
  \begin{syntax}
    \cs{@rc@ifdefinable}
  \end{syntax}

\end{function}



\begin{function}[]
  {
    \@reargdef
  }
  \begin{syntax}
    \cs{@reargdef}
  \end{syntax}

\end{function}



\begin{function}[]
  {
    \@refundefined
  }
  \begin{syntax}
    \cs{@refundefined}
  \end{syntax}

\end{function}



\begin{function}[]
  {
    \@reinserts
  }
  \begin{syntax}
    \cs{@reinserts}
  \end{syntax}

\end{function}



\begin{function}[]
  {
    \@remove@eq@value
  }
  \begin{syntax}
    \cs{@remove@eq@value}
  \end{syntax}

\end{function}



\begin{function}[]
  {
    \@removeelement
  }
  \begin{syntax}
    \cs{@removeelement}
  \end{syntax}

\end{function}



\begin{function}[]
  {
    \@removefromreset
  }
  \begin{syntax}
    \cs{@removefromreset}
  \end{syntax}

\end{function}



\begin{function}[]
  {
    \@reqcolroom
  }
  \begin{syntax}
    \cs{@reqcolroom}
  \end{syntax}

\end{function}



\begin{function}[]
  {
    \@reset@ptions
  }
  \begin{syntax}
    \cs{@reset@ptions}
  \end{syntax}

\end{function}



\begin{function}[]
  {
    \@resetactivechars
  }
  \begin{syntax}
    \cs{@resetactivechars}
  \end{syntax}

\end{function}



\begin{function}[]
  {
    \@resethfps
  }
  \begin{syntax}
    \cs{@resethfps}
  \end{syntax}

\end{function}



\begin{function}[]
  {
    \@restoremetafamily
  }
  \begin{syntax}
    \cs{@restoremetafamily}
  \end{syntax}

\end{function}



\begin{function}[]
  {
    \@restorepar
  }
  \begin{syntax}
    \cs{@restorepar}
  \end{syntax}

\end{function}



\begin{function}[]
  {
    \@reversemarginfalse
  }
  \begin{syntax}
    \cs{@reversemarginfalse}
  \end{syntax}

\end{function}



\begin{function}[]
  {
    \@reversemargintrue
  }
  \begin{syntax}
    \cs{@reversemargintrue}
  \end{syntax}

\end{function}



\begin{function}[]
  {
    \@rightmark
  }
  \begin{syntax}
    \cs{@rightmark}
  \end{syntax}

\end{function}



\begin{function}[]
  {
    \@rightskip
  }
  \begin{syntax}
    \cs{@rightskip}
  \end{syntax}

\end{function}



\begin{function}[]
  {
    \@rjfieldfalse
  }
  \begin{syntax}
    \cs{@rjfieldfalse}
  \end{syntax}

\end{function}



\begin{function}[]
  {
    \@rjfieldtrue
  }
  \begin{syntax}
    \cs{@rjfieldtrue}
  \end{syntax}

\end{function}



\begin{function}[]
  {
    \@rmfamilyhook
  }
  \begin{syntax}
    \cs{@rmfamilyhook}
  \end{syntax}

\end{function}



\begin{function}[]
  {
    \@roman
  }
  \begin{syntax}
    \cs{@roman}
  \end{syntax}

\end{function}



\begin{function}[]
  {
    \@rsbox
  }
  \begin{syntax}
    \cs{@rsbox}
  \end{syntax}

\end{function}



\begin{function}[]
  {
    \@rtab
  }
  \begin{syntax}
    \cs{@rtab}
  \end{syntax}

\end{function}



\begin{function}[]
  {
    \@rule
  }
  \begin{syntax}
    \cs{@rule}
  \end{syntax}

\end{function}



\begin{function}[]
  {
    \@sanitize
  }
  \begin{syntax}
    \cs{@sanitize}
  \end{syntax}

\end{function}



\begin{function}[]
  {
    \@savebox
  }
  \begin{syntax}
    \cs{@savebox}
  \end{syntax}

\end{function}



\begin{function}[]
  {
    \@savemarbox
  }
  \begin{syntax}
    \cs{@savemarbox}
  \end{syntax}

\end{function}



\begin{function}[]
  {
    \@savepicbox
  }
  \begin{syntax}
    \cs{@savepicbox}
  \end{syntax}

\end{function}








\begin{function}[]
  {
    \@scolelt
  }
  \begin{syntax}
    \cs{@scolelt}
  \end{syntax}

\end{function}



\begin{function}[]
  {
    \@sdblcolelt
  }
  \begin{syntax}
    \cs{@sdblcolelt}
  \end{syntax}

\end{function}



\begin{function}[]
  {
    \@seccntformat
  }
  \begin{syntax}
    \cs{@seccntformat}
  \end{syntax}

\end{function}



\begin{function}[]
  {
    \@secondoffive
  }
  \begin{syntax}
    \cs{@secondoffive}
  \end{syntax}

\end{function}



\begin{function}[]
  {
    \@secondoftwo
  }
  \begin{syntax}
    \cs{@secondoftwo}
  \end{syntax}

\end{function}



\begin{function}[]
  {
    \@secpenalty
  }
  \begin{syntax}
    \cs{@secpenalty}
  \end{syntax}

\end{function}



\begin{function}[]
  {
    \@sect
  }
  \begin{syntax}
    \cs{@sect}
  \end{syntax}

\end{function}



\begin{function}[]
  {
    \@seqncr
  }
  \begin{syntax}
    \cs{@seqncr}
  \end{syntax}

\end{function}



\begin{function}[]
  {
    \@set@curr@file@aux
  }
  \begin{syntax}
    \cs{@set@curr@file@aux}
  \end{syntax}

\end{function}



\begin{function}[]
  {
    \@setckpt
  }
  \begin{syntax}
    \cs{@setckpt}
  \end{syntax}

\end{function}



\begin{function}[]
  {
    \@setfloattypecounts
  }
  \begin{syntax}
    \cs{@setfloattypecounts}
  \end{syntax}

\end{function}



\begin{function}[]
  {
    \@setfontsize
  }
  \begin{syntax}
    \cs{@setfontsize}
  \end{syntax}

\end{function}



\begin{function}[]
  {
    \@setfpsbit
  }
  \begin{syntax}
    \cs{@setfpsbit}
  \end{syntax}

\end{function}



\begin{function}[]
  {
    \@setminipage
  }
  \begin{syntax}
    \cs{@setminipage}
  \end{syntax}

\end{function}



\begin{function}[]
  {
    \@setnobreak
  }
  \begin{syntax}
    \cs{@setnobreak}
  \end{syntax}

\end{function}



\begin{function}[]
  {
    \@setpar
  }
  \begin{syntax}
    \cs{@setpar}
  \end{syntax}

\end{function}



\begin{function}[]
  {
    \@setref
  }
  \begin{syntax}
    \cs{@setref}
  \end{syntax}

\end{function}



\begin{function}[]
  {
    \@setsize
  }
  \begin{syntax}
    \cs{@setsize}
  \end{syntax}

\end{function}



\begin{function}[]
  {
    \@settab
  }
  \begin{syntax}
    \cs{@settab}
  \end{syntax}

\end{function}



\begin{function}[]
  {
    \@settodim
  }
  \begin{syntax}
    \cs{@settodim}
  \end{syntax}

\end{function}



\begin{function}[]
  {
    \@settopoint
  }
  \begin{syntax}
    \cs{@settopoint}
  \end{syntax}

\end{function}



\begin{function}[]
  {
    \@setupverbvisiblespace
  }
  \begin{syntax}
    \cs{@setupverbvisiblespace}
  \end{syntax}

\end{function}



\begin{function}[]
  {
    \@setupverbvisibletab
  }
  \begin{syntax}
    \cs{@setupverbvisibletab}
  \end{syntax}

\end{function}



\begin{function}[]
  {
    \@sffamilyhook
  }
  \begin{syntax}
    \cs{@sffamilyhook}
  \end{syntax}

\end{function}



\begin{function}[]
  {
    \@sharp
  }
  \begin{syntax}
    \cs{@sharp}
  \end{syntax}

\end{function}



\begin{function}[]
  {
    \@shortstack
  }
  \begin{syntax}
    \cs{@shortstack}
  \end{syntax}

\end{function}



\begin{function}[]
  {
    \@show@DeclareRobustCommand
  }
  \begin{syntax}
    \cs{@show@DeclareRobustCommand}
  \end{syntax}

\end{function}



\begin{function}[]
  {
    \@show@DeclareRobustCommand@env
  }
  \begin{syntax}
    \cs{@show@DeclareRobustCommand@env}
  \end{syntax}

\end{function}



\begin{function}[]
  {
    \@show@environment
  }
  \begin{syntax}
    \cs{@show@environment}
  \end{syntax}

\end{function}



\begin{function}[]
  {
    \@show@environment@begin
  }
  \begin{syntax}
    \cs{@show@environment@begin}
  \end{syntax}

\end{function}



\begin{function}[]
  {
    \@show@environment@end
  }
  \begin{syntax}
    \cs{@show@environment@end}
  \end{syntax}

\end{function}



\begin{function}[]
  {
    \@show@environment@end@aux
  }
  \begin{syntax}
    \cs{@show@environment@end@aux}
  \end{syntax}

\end{function}



\begin{function}[]
  {
    \@show@macro
  }
  \begin{syntax}
    \cs{@show@macro}
  \end{syntax}

\end{function}



\begin{function}[]
  {
    \@show@newcommand
  }
  \begin{syntax}
    \cs{@show@newcommand}
  \end{syntax}

\end{function}



\begin{function}[]
  {
    \@show@newcommand@aux
  }
  \begin{syntax}
    \cs{@show@newcommand@aux}
  \end{syntax}

\end{function}



\begin{function}[]
  {
    \@show@newcommand@env
  }
  \begin{syntax}
    \cs{@show@newcommand@env}
  \end{syntax}

\end{function}



\begin{function}[]
  {
    \@show@nonstop
  }
  \begin{syntax}
    \cs{@show@nonstop}
  \end{syntax}

\end{function}



\begin{function}[]
  {
    \@show@normalenv
  }
  \begin{syntax}
    \cs{@show@normalenv}
  \end{syntax}

\end{function}



\begin{function}[]
  {
    \@show@tokens
  }
  \begin{syntax}
    \cs{@show@tokens}
  \end{syntax}

\end{function}



\begin{function}[]
  {
    \@show@typeout
  }
  \begin{syntax}
    \cs{@show@typeout}
  \end{syntax}

\end{function}



\begin{function}[]
  {
    \@showcommandlisthook
  }
  \begin{syntax}
    \cs{@showcommandlisthook}
  \end{syntax}

\end{function}



\begin{function}[]
  {
    \@showenvironmentlisthook
  }
  \begin{syntax}
    \cs{@showenvironmentlisthook}
  \end{syntax}

\end{function}




\begin{function}[]
  {
    \@sline
  }
  \begin{syntax}
    \cs{@sline}
  \end{syntax}

\end{function}



\begin{function}[]
  {
    \@slowromancap
  }
  \begin{syntax}
    \cs{@slowromancap}
  \end{syntax}

\end{function}





\begin{function}[]
  {
    \@specialoutput
  }
  \begin{syntax}
    \cs{@specialoutput}
  \end{syntax}

\end{function}



\begin{function}[]
  {
    \@specialpagefalse
  }
  \begin{syntax}
    \cs{@specialpagefalse}
  \end{syntax}

\end{function}



\begin{function}[]
  {
    \@specialpagetrue
  }
  \begin{syntax}
    \cs{@specialpagetrue}
  \end{syntax}

\end{function}



\begin{function}[]
  {
    \@specialstyle
  }
  \begin{syntax}
    \cs{@specialstyle}
  \end{syntax}

\end{function}



\begin{function}[]
  {
    \@sptoken
  }
  \begin{syntax}
    \cs{@sptoken}
  \end{syntax}

\end{function}



\begin{function}[]
  {
    \@sqrt
  }
  \begin{syntax}
    \cs{@sqrt}
  \end{syntax}

\end{function}



\begin{function}[]
  {
    \@ssect
  }
  \begin{syntax}
    \cs{@ssect}
  \end{syntax}

\end{function}



\begin{function}[]
  {
    \@stackcr
  }
  \begin{syntax}
    \cs{@stackcr}
  \end{syntax}

\end{function}



\begin{function}[]
  {
    \@star@or@long
  }
  \begin{syntax}
    \cs{@star@or@long}
  \end{syntax}

\end{function}



\begin{function}[]
  {
    \@startcolumn
  }
  \begin{syntax}
    \cs{@startcolumn}
  \end{syntax}

\end{function}



\begin{function}[]
  {
    \@startdblcolumn
  }
  \begin{syntax}
    \cs{@startdblcolumn}
  \end{syntax}

\end{function}



\begin{function}[]
  {
    \@startfield
  }
  \begin{syntax}
    \cs{@startfield}
  \end{syntax}

\end{function}



\begin{function}[]
  {
    \@startline
  }
  \begin{syntax}
    \cs{@startline}
  \end{syntax}

\end{function}



\begin{function}[]
  {
    \@startpbox
  }
  \begin{syntax}
    \cs{@startpbox}
  \end{syntax}

\end{function}



\begin{function}[]
  {
    \@startsection
  }
  \begin{syntax}
    \cs{@startsection}
  \end{syntax}

\end{function}



\begin{function}[]
  {
    \@starttoc
  }
  \begin{syntax}
    \cs{@starttoc}
  \end{syntax}

\end{function}



\begin{function}[]
  {
    \@stopfield
  }
  \begin{syntax}
    \cs{@stopfield}
  \end{syntax}

\end{function}



\begin{function}[]
  {
    \@stopline
  }
  \begin{syntax}
    \cs{@stopline}
  \end{syntax}

\end{function}



\begin{function}[]
  {
    \@stpelt
  }
  \begin{syntax}
    \cs{@stpelt}
  \end{syntax}

\end{function}



\begin{function}[]
  {
    \@string@makeletter
  }
  \begin{syntax}
    \cs{@string@makeletter}
  \end{syntax}

\end{function}



\begin{function}[]
  {
    \@strip@tex@ext
  }
  \begin{syntax}
    \cs{@strip@tex@ext}
  \end{syntax}

\end{function}



\begin{function}[]
  {
    \@strip@tex@ext@aux
  }
  \begin{syntax}
    \cs{@strip@tex@ext@aux}
  \end{syntax}

\end{function}



\begin{function}[]
  {
    \@svector
  }
  \begin{syntax}
    \cs{@svector}
  \end{syntax}

\end{function}



\begin{function}[]
  {
    \@sverb
  }
  \begin{syntax}
    \cs{@sverb}
  \end{syntax}

\end{function}



\begin{function}[]
  {
    \@svsec
  }
  \begin{syntax}
    \cs{@svsec}
  \end{syntax}

\end{function}



\begin{function}[]
  {
    \@svsechd
  }
  \begin{syntax}
    \cs{@svsechd}
  \end{syntax}

\end{function}



\begin{function}[]
  {
    \@swaptwoargs
  }
  \begin{syntax}
    \cs{@swaptwoargs}
  \end{syntax}

\end{function}



\begin{function}[]
  {
    \@sxverbatim
  }
  \begin{syntax}
    \cs{@sxverbatim}
  \end{syntax}

\end{function}



\begin{function}[]
  {
    \@tabacckludge
  }
  \begin{syntax}
    \cs{@tabacckludge}
  \end{syntax}

\end{function}



\begin{function}[]
  {
    \@tabacol
  }
  \begin{syntax}
    \cs{@tabacol}
  \end{syntax}

\end{function}



\begin{function}[]
  {
    \@tabarray
  }
  \begin{syntax}
    \cs{@tabarray}
  \end{syntax}

\end{function}



\begin{function}[]
  {
    \@tabclassiv
  }
  \begin{syntax}
    \cs{@tabclassiv}
  \end{syntax}

\end{function}



\begin{function}[]
  {
    \@tabclassz
  }
  \begin{syntax}
    \cs{@tabclassz}
  \end{syntax}

\end{function}



\begin{function}[]
  {
    \@tabcr
  }
  \begin{syntax}
    \cs{@tabcr}
  \end{syntax}

\end{function}



\begin{function}[]
  {
    \@tabfbox
  }
  \begin{syntax}
    \cs{@tabfbox}
  \end{syntax}

\end{function}



\begin{function}[]
  {
    \@tablab
  }
  \begin{syntax}
    \cs{@tablab}
  \end{syntax}

\end{function}



\begin{function}[]
  {
    \@tabminus
  }
  \begin{syntax}
    \cs{@tabminus}
  \end{syntax}

\end{function}



\begin{function}[]
  {
    \@tabplus
  }
  \begin{syntax}
    \cs{@tabplus}
  \end{syntax}

\end{function}



\begin{function}[]
  {
    \@tabpush
  }
  \begin{syntax}
    \cs{@tabpush}
  \end{syntax}

\end{function}



\begin{function}[]
  {
    \@tabrj
  }
  \begin{syntax}
    \cs{@tabrj}
  \end{syntax}

\end{function}



\begin{function}[]
  {
    \@tabular
  }
  \begin{syntax}
    \cs{@tabular}
  \end{syntax}

\end{function}



\begin{function}[]
  {
    \@tabularcr
  }
  \begin{syntax}
    \cs{@tabularcr}
  \end{syntax}

\end{function}



\begin{function}[]
  {
    \@tempboxa
  }
  \begin{syntax}
    \cs{@tempboxa}
  \end{syntax}

\end{function}



\begin{function}[]
  {
    \@tempcnta
  }
  \begin{syntax}
    \cs{@tempcnta}
  \end{syntax}

\end{function}



\begin{function}[]
  {
    \@tempcntb
  }
  \begin{syntax}
    \cs{@tempcntb}
  \end{syntax}

\end{function}



\begin{function}[]
  {
    \@tempdima
  }
  \begin{syntax}
    \cs{@tempdima}
  \end{syntax}

\end{function}





\begin{function}[]
  {
    \@tempdimb
  }
  \begin{syntax}
    \cs{@tempdimb}
  \end{syntax}

\end{function}



\begin{function}[]
  {
    \@tempdimc
  }
  \begin{syntax}
    \cs{@tempdimc}
  \end{syntax}

\end{function}



\begin{function}[]
  {
    \@tempskipa
  }
  \begin{syntax}
    \cs{@tempskipa}
  \end{syntax}

\end{function}



\begin{function}[]
  {
    \@tempskipb
  }
  \begin{syntax}
    \cs{@tempskipb}
  \end{syntax}

\end{function}



\begin{function}[]
  {
    \@tempswafalse
  }
  \begin{syntax}
    \cs{@tempswafalse}
  \end{syntax}

\end{function}



\begin{function}[]
  {
    \@tempswatrue
  }
  \begin{syntax}
    \cs{@tempswatrue}
  \end{syntax}

\end{function}



\begin{function}[]
  {
    \@temptokena
  }
  \begin{syntax}
    \cs{@temptokena}
  \end{syntax}

\end{function}



\begin{function}[]
  {
    \@testdef
  }
  \begin{syntax}
    \cs{@testdef}
  \end{syntax}

\end{function}



\begin{function}[]
  {
    \@testfalse
  }
  \begin{syntax}
    \cs{@testfalse}
  \end{syntax}

\end{function}



\begin{function}[]
  {
    \@testfp
  }
  \begin{syntax}
    \cs{@testfp}
  \end{syntax}

\end{function}



\begin{function}[]
  {
    \@testopt
  }
  \begin{syntax}
    \cs{@testopt}
  \end{syntax}

\end{function}



\begin{function}[]
  {
    \@testpach
  }
  \begin{syntax}
    \cs{@testpach}
  \end{syntax}

\end{function}



\begin{function}[]
  {
    \@testtrue
  }
  \begin{syntax}
    \cs{@testtrue}
  \end{syntax}

\end{function}



\begin{function}[]
  {
    \@testwrongwidth
  }
  \begin{syntax}
    \cs{@testwrongwidth}
  \end{syntax}

\end{function}



\begin{function}[]
  {
    \@text@composite
  }
  \begin{syntax}
    \cs{@text@composite}
  \end{syntax}

\end{function}



\begin{function}[]
  {
    \@text@composite@x
  }
  \begin{syntax}
    \cs{@text@composite@x}
  \end{syntax}

\end{function}



\begin{function}[]
  {
    \@textbottom
  }
  \begin{syntax}
    \cs{@textbottom}
  \end{syntax}

\end{function}



\begin{function}[]
  {
    \@textfloatsheight
  }
  \begin{syntax}
    \cs{@textfloatsheight}
  \end{syntax}

\end{function}



\begin{function}[]
  {
    \@textmin
  }
  \begin{syntax}
    \cs{@textmin}
  \end{syntax}

\end{function}



\begin{function}[]
  {
    \@textsubscript
  }
  \begin{syntax}
    \cs{@textsubscript}
  \end{syntax}

\end{function}



\begin{function}[]
  {
    \@textsuperscript
  }
  \begin{syntax}
    \cs{@textsuperscript}
  \end{syntax}

\end{function}



\begin{function}[]
  {
    \@texttop
  }
  \begin{syntax}
    \cs{@texttop}
  \end{syntax}

\end{function}



\begin{function}[]
  {
    \@tf@r
  }
  \begin{syntax}
    \cs{@tf@r}
  \end{syntax}

\end{function}



\begin{function}[]
  {
    \@tfor
  }
  \begin{syntax}
    \cs{@tfor}
  \end{syntax}

\end{function}



\begin{function}[]
  {
    \@tforloop
  }
  \begin{syntax}
    \cs{@tforloop}
  \end{syntax}

\end{function}



\begin{function}[]
  {
    \@thanks
  }
  \begin{syntax}
    \cs{@thanks}
  \end{syntax}

\end{function}



\begin{function}[]
  {
    \@thefnmark
  }
  \begin{syntax}
    \cs{@thefnmark}
  \end{syntax}

\end{function}



\begin{function}[]
  {
    \@thefoot
  }
  \begin{syntax}
    \cs{@thefoot}
  \end{syntax}

\end{function}



\begin{function}[]
  {
    \@thehead
  }
  \begin{syntax}
    \cs{@thehead}
  \end{syntax}

\end{function}



\begin{function}[]
  {
    \@themargin
  }
  \begin{syntax}
    \cs{@themargin}
  \end{syntax}

\end{function}



\begin{function}[]
  {
    \@thirdofthree
  }
  \begin{syntax}
    \cs{@thirdofthree}
  \end{syntax}

\end{function}



\begin{function}[]
  {
    \@thm
  }
  \begin{syntax}
    \cs{@thm}
  \end{syntax}

\end{function}



\begin{function}[]
  {
    \@thmcounter
  }
  \begin{syntax}
    \cs{@thmcounter}
  \end{syntax}

\end{function}



\begin{function}[]
  {
    \@thmcountersep
  }
  \begin{syntax}
    \cs{@thmcountersep}
  \end{syntax}

\end{function}



\begin{function}[]
  {
    \@title
  }
  \begin{syntax}
    \cs{@title}
  \end{syntax}

\end{function}



\begin{function}[]
  {
    \@tocrmarg
  }
  \begin{syntax}
    \cs{@tocrmarg}
  \end{syntax}

\end{function}



\begin{function}[]
  {
    \@toodeep
  }
  \begin{syntax}
    \cs{@toodeep}
  \end{syntax}

\end{function}



\begin{function}[]
  {
    \@toplist
  }
  \begin{syntax}
    \cs{@toplist}
  \end{syntax}

\end{function}



\begin{function}[]
  {
    \@topnewpage
  }
  \begin{syntax}
    \cs{@topnewpage}
  \end{syntax}

\end{function}



\begin{function}[]
  {
    \@topnum
  }
  \begin{syntax}
    \cs{@topnum}
  \end{syntax}

\end{function}



\begin{function}[]
  {
    \@toproom
  }
  \begin{syntax}
    \cs{@toproom}
  \end{syntax}

\end{function}







\begin{function}[]
  {
    \@totalleftmargin
  }
  \begin{syntax}
    \cs{@totalleftmargin}
  \end{syntax}

\end{function}



\begin{function}[]
  {
    \@trivlist
  }
  \begin{syntax}
    \cs{@trivlist}
  \end{syntax}

\end{function}



\begin{function}[]
  {
    \@tryfcolumn
  }
  \begin{syntax}
    \cs{@tryfcolumn}
  \end{syntax}

\end{function}



\begin{function}[]
  {
    \@trylist
  }
  \begin{syntax}
    \cs{@trylist}
  \end{syntax}

\end{function}



\begin{function}[]
  {
    \@ttfamilyhook
  }
  \begin{syntax}
    \cs{@ttfamilyhook}
  \end{syntax}

\end{function}



\begin{function}[]
  {
    \@twoclasseserror
  }
  \begin{syntax}
    \cs{@twoclasseserror}
  \end{syntax}

\end{function}



\begin{function}[]
  {
    \@twocolumnfalse
  }
  \begin{syntax}
    \cs{@twocolumnfalse}
  \end{syntax}

\end{function}



\begin{function}[]
  {
    \@twocolumntrue
  }
  \begin{syntax}
    \cs{@twocolumntrue}
  \end{syntax}

\end{function}



\begin{function}[]
  {
    \@twoloadclasserror
  }
  \begin{syntax}
    \cs{@twoloadclasserror}
  \end{syntax}

\end{function}



\begin{function}[]
  {
    \@twosidefalse
  }
  \begin{syntax}
    \cs{@twosidefalse}
  \end{syntax}

\end{function}



\begin{function}[]
  {
    \@typein
  }
  \begin{syntax}
    \cs{@typein}
  \end{syntax}

\end{function}



\begin{function}[]
  {
    \@typeset@protect
  }
  \begin{syntax}
    \cs{@typeset@protect}
  \end{syntax}

\end{function}



\begin{function}[]
  {
    \@uclclist
  }
  \begin{syntax}
    \cs{@uclclist}
  \end{syntax}

\end{function}



\begin{function}[]
  {
    \@undefined
  }
  \begin{syntax}
    \cs{@undefined}
  \end{syntax}

\end{function}



\begin{function}[]
  {
    \@unexpandable@protect
  }
  \begin{syntax}
    \cs{@unexpandable@protect}
  \end{syntax}

\end{function}



\begin{function}[]
  {
    \@unknownoptionerror
  }
  \begin{syntax}
    \cs{@unknownoptionerror}
  \end{syntax}

\end{function}



\begin{function}[]
  {
    \@unprocessedoptions
  }
  \begin{syntax}
    \cs{@unprocessedoptions}
  \end{syntax}

\end{function}



\begin{function}[]
  {
    \@unused
  }
  \begin{syntax}
    \cs{@unused}
  \end{syntax}

\end{function}



\begin{function}[]
  {
    \@unusedoptionlist
  }
  \begin{syntax}
    \cs{@unusedoptionlist}
  \end{syntax}

\end{function}



\begin{function}[]
  {
    \@upline
  }
  \begin{syntax}
    \cs{@upline}
  \end{syntax}

\end{function}



\begin{function}[]
  {
    \@upordown
  }
  \begin{syntax}
    \cs{@upordown}
  \end{syntax}

\end{function}



\begin{function}[]
  {
    \@upvector
  }
  \begin{syntax}
    \cs{@upvector}
  \end{syntax}

\end{function}



\begin{function}[]
  {
    \@use@ption
  }
  \begin{syntax}
    \cs{@use@ption}
  \end{syntax}

\end{function}



\begin{function}[]
  {
    \@use@text@encoding
  }
  \begin{syntax}
    \cs{@use@text@encoding}
  \end{syntax}

\end{function}



\begin{function}[]
  {
    \@verb
  }
  \begin{syntax}
    \cs{@verb}
  \end{syntax}

\end{function}



\begin{function}[]
  {
    \@verbatim
  }
  \begin{syntax}
    \cs{@verbatim}
  \end{syntax}

\end{function}



\begin{function}[]
  {
    \@verbvisiblespacebox
  }
  \begin{syntax}
    \cs{@verbvisiblespacebox}
  \end{syntax}

\end{function}



\begin{function}[]
  {
    \@viiipt
  }
  \begin{syntax}
    \cs{@viiipt}
  \end{syntax}

\end{function}



\begin{function}[]
  {
    \@viipt
  }
  \begin{syntax}
    \cs{@viipt}
  \end{syntax}

\end{function}



\begin{function}[]
  {
    \@vipt
  }
  \begin{syntax}
    \cs{@vipt}
  \end{syntax}

\end{function}



\begin{function}[]
  {
    \@vline
  }
  \begin{syntax}
    \cs{@vline}
  \end{syntax}

\end{function}



\begin{function}[]
  {
    \@vobeyspaces
  }
  \begin{syntax}
    \cs{@vobeyspaces}
  \end{syntax}

\end{function}



\begin{function}[]
  {
    \@vobeytabs
  }
  \begin{syntax}
    \cs{@vobeytabs}
  \end{syntax}

\end{function}



\begin{function}[]
  {
    \@vpt
  }
  \begin{syntax}
    \cs{@vpt}
  \end{syntax}

\end{function}



\begin{function}[]
  {
    \@vspace
  }
  \begin{syntax}
    \cs{@vspace}
  \end{syntax}

\end{function}




\begin{function}[]
  {
    \@vspacer
  }
  \begin{syntax}
    \cs{@vspacer}
  \end{syntax}

\end{function}



\begin{function}[]
  {
    \@vtryfc
  }
  \begin{syntax}
    \cs{@vtryfc}
  \end{syntax}

\end{function}



\begin{function}[]
  {
    \@vvector
  }
  \begin{syntax}
    \cs{@vvector}
  \end{syntax}

\end{function}



\begin{function}[int,deprecated]
  {
    \@warning
  }
  \begin{syntax}
    \cs{@warning} \Arg{message}
  \end{syntax}
  Another name for \cs{@latex@warning}. Deprecated.

\end{function}



\begin{function}[]
  {
    \@wckptelt
  }
  \begin{syntax}
    \cs{@wckptelt}
  \end{syntax}

\end{function}



\begin{function}[]
  {
    \@whiledim
  }
  \begin{syntax}
    \cs{@whiledim}
  \end{syntax}

\end{function}



\begin{function}[]
  {
    \@whilenum
  }
  \begin{syntax}
    \cs{@whilenum}
  \end{syntax}

\end{function}



\begin{function}[]
  {
    \@whilesw
  }
  \begin{syntax}
    \cs{@whilesw}
  \end{syntax}

\end{function}



\begin{function}[]
  {
    \@wholewidth
  }
  \begin{syntax}
    \cs{@wholewidth}
  \end{syntax}

\end{function}



\begin{function}[]
  {
    \@width
  }
  \begin{syntax}
    \cs{@width}
  \end{syntax}

\end{function}



\begin{function}[]
  {
    \@wrglossary
  }
  \begin{syntax}
    \cs{@wrglossary}
  \end{syntax}

\end{function}



\begin{function}[]
  {
    \@wrindex
  }
  \begin{syntax}
    \cs{@wrindex}
  \end{syntax}

\end{function}



\begin{function}[]
  {
    \@writeckpt
  }
  \begin{syntax}
    \cs{@writeckpt}
  \end{syntax}

\end{function}



\begin{function}[]
  {
    \@writefile
  }
  \begin{syntax}
    \cs{@writefile}
  \end{syntax}

\end{function}



\begin{function}[]
  {
    \@wrong@font@char
  }
  \begin{syntax}
    \cs{@wrong@font@char}
  \end{syntax}

\end{function}



\begin{function}[]
  {
    \@wtryfc
  }
  \begin{syntax}
    \cs{@wtryfc}
  \end{syntax}

\end{function}



\begin{function}[]
  {
    \@x@protect
  }
  \begin{syntax}
    \cs{@x@protect}
  \end{syntax}

\end{function}



\begin{function}[]
  {
    \@x@sf
  }
  \begin{syntax}
    \cs{@x@sf}
  \end{syntax}

\end{function}



\begin{function}[]
  {
    \@xDeclareMathDelimiter
  }
  \begin{syntax}
    \cs{@xDeclareMathDelimiter}
  \end{syntax}

\end{function}



\begin{function}[]
  {
    \@xaddvskip
  }
  \begin{syntax}
    \cs{@xaddvskip}
  \end{syntax}

\end{function}



\begin{function}[]
  {
    \@xarg
  }
  \begin{syntax}
    \cs{@xarg}
  \end{syntax}

\end{function}



\begin{function}[]
  {
    \@xargarraycr
  }
  \begin{syntax}
    \cs{@xargarraycr}
  \end{syntax}

\end{function}



\begin{function}[]
  {
    \@xargdef
  }
  \begin{syntax}
    \cs{@xargdef}
  \end{syntax}

\end{function}



\begin{function}[]
  {
    \@xarraycr
  }
  \begin{syntax}
    \cs{@xarraycr}
  \end{syntax}

\end{function}



\begin{function}[]
  {
    \@xbitor
  }
  \begin{syntax}
    \cs{@xbitor}
  \end{syntax}

\end{function}



\begin{function}[]
  {
    \@xcentercr
  }
  \begin{syntax}
    \cs{@xcentercr}
  \end{syntax}

\end{function}



\begin{function}[]
  {
    \@xdblarg
  }
  \begin{syntax}
    \cs{@xdblarg}
  \end{syntax}

\end{function}



\begin{function}[]
  {
    \@xdblfloat
  }
  \begin{syntax}
    \cs{@xdblfloat}
  \end{syntax}

\end{function}



\begin{function}[]
  {
    \@xdim
  }
  \begin{syntax}
    \cs{@xdim}
  \end{syntax}

\end{function}



\begin{function}[]
  {
    \@xeqncr
  }
  \begin{syntax}
    \cs{@xeqncr}
  \end{syntax}

\end{function}



\begin{function}[]
  {
    \@xexnoop
  }
  \begin{syntax}
    \cs{@xexnoop}
  \end{syntax}

\end{function}



\begin{function}[]
  {
    \@xexpast
  }
  \begin{syntax}
    \cs{@xexpast}
  \end{syntax}

\end{function}



\begin{function}[]
  {
    \@xfloat
  }
  \begin{syntax}
    \cs{@xfloat}
  \end{syntax}

\end{function}



\begin{function}[]
  {
    \@xfootnote
  }
  \begin{syntax}
    \cs{@xfootnote}
  \end{syntax}

\end{function}



\begin{function}[]
  {
    \@xfootnotemark
  }
  \begin{syntax}
    \cs{@xfootnotemark}
  \end{syntax}

\end{function}



\begin{function}[]
  {
    \@xfootnotenext
  }
  \begin{syntax}
    \cs{@xfootnotenext}
  \end{syntax}

\end{function}



\begin{function}[]
  {
    \@xhline
  }
  \begin{syntax}
    \cs{@xhline}
  \end{syntax}

\end{function}



\begin{function}[]
  {
    \@xifnch
  }
  \begin{syntax}
    \cs{@xifnch}
  \end{syntax}

\end{function}



\begin{function}[]
  {
    \@xiipt
  }
  \begin{syntax}
    \cs{@xiipt}
  \end{syntax}

\end{function}



\begin{function}[]
  {
    \@xipt
  }
  \begin{syntax}
    \cs{@xipt}
  \end{syntax}

\end{function}



\begin{function}[]
  {
    \@xivpt
  }
  \begin{syntax}
    \cs{@xivpt}
  \end{syntax}

\end{function}



\begin{function}[]
  {
    \@xmpar
  }
  \begin{syntax}
    \cs{@xmpar}
  \end{syntax}

\end{function}





\begin{function}[]
  {
    \@xnext
  }
  \begin{syntax}
    \cs{@xnext}
  \end{syntax}

\end{function}



\begin{function}[]
  {
    \@xnthm
  }
  \begin{syntax}
    \cs{@xnthm}
  \end{syntax}

\end{function}



\begin{function}[]
  {
    \@xobeysp
  }
  \begin{syntax}
    \cs{@xobeysp}
  \end{syntax}

\end{function}



\begin{function}[]
  {
    \@xobeytab
  }
  \begin{syntax}
    \cs{@xobeytab}
  \end{syntax}

\end{function}



\begin{function}[]
  {
    \@xprocess@ptions
  }
  \begin{syntax}
    \cs{@xprocess@ptions}
  \end{syntax}

\end{function}



\begin{function}[]
  {
    \@xpt
  }
  \begin{syntax}
    \cs{@xpt}
  \end{syntax}

\end{function}



\begin{function}[]
  {
    \@xsect
  }
  \begin{syntax}
    \cs{@xsect}
  \end{syntax}

\end{function}



\begin{function}[]
  {
    \@xtabcr
  }
  \begin{syntax}
    \cs{@xtabcr}
  \end{syntax}

\end{function}



\begin{function}[]
  {
    \@xtabularcr
  }
  \begin{syntax}
    \cs{@xtabularcr}
  \end{syntax}

\end{function}



\begin{function}[]
  {
    \@xthm
  }
  \begin{syntax}
    \cs{@xthm}
  \end{syntax}

\end{function}



\begin{function}[]
  {
    \@xtryfc
  }
  \begin{syntax}
    \cs{@xtryfc}
  \end{syntax}

\end{function}



\begin{function}[]
  {
    \@xtypein
  }
  \begin{syntax}
    \cs{@xtypein}
  \end{syntax}

\end{function}



\begin{function}[]
  {
    \@xverbatim
  }
  \begin{syntax}
    \cs{@xverbatim}
  \end{syntax}

\end{function}



\begin{function}[]
  {
    \@xviipt
  }
  \begin{syntax}
    \cs{@xviipt}
  \end{syntax}

\end{function}



\begin{function}[]
  {
    \@xxDeclareMathDelimiter
  }
  \begin{syntax}
    \cs{@xxDeclareMathDelimiter}
  \end{syntax}

\end{function}



\begin{function}[]
  {
    \@xxpt
  }
  \begin{syntax}
    \cs{@xxpt}
  \end{syntax}

\end{function}



\begin{function}[]
  {
    \@xxvpt
  }
  \begin{syntax}
    \cs{@xxvpt}
  \end{syntax}

\end{function}



\begin{function}[]
  {
    \@xxxii
  }
  \begin{syntax}
    \cs{@xxxii}
  \end{syntax}

\end{function}



\begin{function}[]
  {
    \@xympar
  }
  \begin{syntax}
    \cs{@xympar}
  \end{syntax}

\end{function}



\begin{function}[]
  {
    \@yarg
  }
  \begin{syntax}
    \cs{@yarg}
  \end{syntax}

\end{function}



\begin{function}[]
  {
    \@yargarraycr
  }
  \begin{syntax}
    \cs{@yargarraycr}
  \end{syntax}

\end{function}



\begin{function}[]
  {
    \@yargd@f
  }
  \begin{syntax}
    \cs{@yargd@f}
  \end{syntax}

\end{function}



\begin{function}[]
  {
    \@yargdef
  }
  \begin{syntax}
    \cs{@yargdef}
  \end{syntax}

\end{function}



\begin{function}[]
  {
    \@ydim
  }
  \begin{syntax}
    \cs{@ydim}
  \end{syntax}

\end{function}



\begin{function}[]
  {
    \@yeqncr
  }
  \begin{syntax}
    \cs{@yeqncr}
  \end{syntax}

\end{function}



\begin{function}[]
  {
    \@ympar
  }
  \begin{syntax}
    \cs{@ympar}
  \end{syntax}

\end{function}



\begin{function}[]
  {
    \@ynthm
  }
  \begin{syntax}
    \cs{@ynthm}
  \end{syntax}

\end{function}



\begin{function}[]
  {
    \@ythm
  }
  \begin{syntax}
    \cs{@ythm}
  \end{syntax}

\end{function}



\begin{function}[]
  {
    \@ytryfc
  }
  \begin{syntax}
    \cs{@ytryfc}
  \end{syntax}

\end{function}



\begin{function}[]
  {
    \@yyarg
  }
  \begin{syntax}
    \cs{@yyarg}
  \end{syntax}

\end{function}



\begin{function}[]
  {
    \@ztryfc
  }
  \begin{syntax}
    \cs{@ztryfc}
  \end{syntax}

\end{function}



\begin{function}[]
  {
    \A
  }
  \begin{syntax}
    \cs{A}
  \end{syntax}

\end{function}



\begin{function}[]
  {
    \AA
  }
  \begin{syntax}
    \cs{AA}
  \end{syntax}

\end{function}



\begin{function}[]
  {
    \AE
  }
  \begin{syntax}
    \cs{AE}
  \end{syntax}

\end{function}



\begin{function}[]
  {
    \ActivateGenericHook
  }
  \begin{syntax}
    \cs{ActivateGenericHook}
  \end{syntax}

\end{function}



\begin{function}[]
  {
    \AddToDocumentProperties
  }
  \begin{syntax}
    \cs{AddToDocumentProperties}
  \end{syntax}

\end{function}



\begin{function}[]
  {
    \AddToHook
  }
  \begin{syntax}
    \cs{AddToHook}
  \end{syntax}

\end{function}



\begin{function}[]
  {
    \AddToHookNext
  }
  \begin{syntax}
    \cs{AddToHookNext}
  \end{syntax}

\end{function}



\begin{function}[]
  {
    \AddToHookNextWithArguments
  }
  \begin{syntax}
    \cs{AddToHookNextWithArguments}
  \end{syntax}

\end{function}



\begin{function}[]
  {
    \AddToHookWithArguments
  }
  \begin{syntax}
    \cs{AddToHookWithArguments}
  \end{syntax}

\end{function}



\begin{function}[]
  {
    \AddToNoCaseChangeList
  }
  \begin{syntax}
    \cs{AddToNoCaseChangeList}
  \end{syntax}

\end{function}



\begin{function}[]
  {
    \AfterEndEnvironment
  }
  \begin{syntax}
    \cs{AfterEndEnvironment}
  \end{syntax}

\end{function}



\begin{function}[]
  {
    \Alph
  }
  \begin{syntax}
    \cs{Alph}
  \end{syntax}

\end{function}



\begin{function}[]
  {
    \AssignSocketPlug
  }
  \begin{syntax}
    \cs{AssignSocketPlug}
  \end{syntax}

\end{function}



\begin{function}[]
  {
    \AssignStructureRole
  }
  \begin{syntax}
    \cs{AssignStructureRole}
  \end{syntax}

\end{function}



\begin{function}[]
  {
    \AssignTaggingSocketPlug
  }
  \begin{syntax}
    \cs{AssignTaggingSocketPlug}
  \end{syntax}

\end{function}



\begin{function}[]
  {
    \AssignTemplateKeys
  }
  \begin{syntax}
    \cs{AssignTemplateKeys}
  \end{syntax}

\end{function}



\begin{function}[]
  {
    \AtBeginDocument
  }
  \begin{syntax}
    \cs{AtBeginDocument}
  \end{syntax}

\end{function}



\begin{function}[]
  {
    \AtBeginDvi
  }
  \begin{syntax}
    \cs{AtBeginDvi}
  \end{syntax}

\end{function}



\begin{function}[]
  {
    \AtBeginEnvironment
  }
  \begin{syntax}
    \cs{AtBeginEnvironment}
  \end{syntax}

\end{function}



\begin{function}[]
  {
    \AtEndDocument
  }
  \begin{syntax}
    \cs{AtEndDocument}
  \end{syntax}

\end{function}



\begin{function}[]
  {
    \AtEndDvi
  }
  \begin{syntax}
    \cs{AtEndDvi}
  \end{syntax}

\end{function}



\begin{function}[]
  {
    \AtEndEnvironment
  }
  \begin{syntax}
    \cs{AtEndEnvironment}
  \end{syntax}

\end{function}



\begin{function}[]
  {
    \AtEndOfClass
  }
  \begin{syntax}
    \cs{AtEndOfClass}
  \end{syntax}

\end{function}



\begin{function}[]
  {
    \AtEndOfPackage
  }
  \begin{syntax}
    \cs{AtEndOfPackage}
  \end{syntax}

\end{function}



\begin{function}[]
  {
    \BCPdata
  }
  \begin{syntax}
    \cs{BCPdata}
  \end{syntax}

\end{function}



\begin{function}[]
  {
    \BeforeBeginEnvironment
  }
  \begin{syntax}
    \cs{BeforeBeginEnvironment}
  \end{syntax}

\end{function}



\begin{function}[]
  {
    \Big
  }
  \begin{syntax}
    \cs{Big}
  \end{syntax}

\end{function}



\begin{function}[]
  {
    \Bigg
  }
  \begin{syntax}
    \cs{Bigg}
  \end{syntax}

\end{function}



\begin{function}[]
  {
    \Biggl
  }
  \begin{syntax}
    \cs{Biggl}
  \end{syntax}

\end{function}



\begin{function}[]
  {
    \Biggm
  }
  \begin{syntax}
    \cs{Biggm}
  \end{syntax}

\end{function}



\begin{function}[]
  {
    \Biggr
  }
  \begin{syntax}
    \cs{Biggr}
  \end{syntax}

\end{function}



\begin{function}[]
  {
    \Bigl
  }
  \begin{syntax}
    \cs{Bigl}
  \end{syntax}

\end{function}



\begin{function}[]
  {
    \Bigm
  }
  \begin{syntax}
    \cs{Bigm}
  \end{syntax}

\end{function}



\begin{function}[]
  {
    \Bigr
  }
  \begin{syntax}
    \cs{Bigr}
  \end{syntax}

\end{function}



\begin{function}[]
  {
    \BooleanFalse
  }
  \begin{syntax}
    \cs{BooleanFalse}
  \end{syntax}

\end{function}



\begin{function}[]
  {
    \BooleanTrue
  }
  \begin{syntax}
    \cs{BooleanTrue}
  \end{syntax}

\end{function}



\begin{function}[]
  {
    \Box
  }
  \begin{syntax}
    \cs{Box}
  \end{syntax}

\end{function}



\begin{function}[]
  {
    \CaseSwitch
  }
  \begin{syntax}
    \cs{CaseSwitch}
  \end{syntax}

\end{function}



\begin{function}[]
  {
    \CheckCommand
  }
  \begin{syntax}
    \cs{CheckCommand}
  \end{syntax}

\end{function}



\begin{function}[]
  {
    \CheckEncodingSubset
  }
  \begin{syntax}
    \cs{CheckEncodingSubset}
  \end{syntax}

\end{function}



\begin{function}[]
  {
    \ClassError
  }
  \begin{syntax}
    \cs{ClassError}
  \end{syntax}

\end{function}



\begin{function}[]
  {
    \ClassInfo
  }
  \begin{syntax}
    \cs{ClassInfo}
  \end{syntax}

\end{function}



\begin{function}[]
  {
    \ClassNote
  }
  \begin{syntax}
    \cs{ClassNote}
  \end{syntax}

\end{function}



\begin{function}[]
  {
    \ClassNoteNoLine
  }
  \begin{syntax}
    \cs{ClassNoteNoLine}
  \end{syntax}

\end{function}



\begin{function}[]
  {
    \ClassWarning
  }
  \begin{syntax}
    \cs{ClassWarning}
  \end{syntax}

\end{function}



\begin{function}[]
  {
    \ClassWarningNoLine
  }
  \begin{syntax}
    \cs{ClassWarningNoLine}
  \end{syntax}

\end{function}



\begin{function}[]
  {
    \ClearHookNext
  }
  \begin{syntax}
    \cs{ClearHookNext}
  \end{syntax}

\end{function}



\begin{function}[]
  {
    \ClearHookRule
  }
  \begin{syntax}
    \cs{ClearHookRule}
  \end{syntax}

\end{function}



\begin{function}[]
  {
    \CountZero
  }
  \begin{syntax}
    \cs{CountZero}
  \end{syntax}

\end{function}



\begin{function}[]
  {
    \CurrentFile
  }
  \begin{syntax}
    \cs{CurrentFile}
  \end{syntax}

\end{function}



\begin{function}[]
  {
    \CurrentFilePath
  }
  \begin{syntax}
    \cs{CurrentFilePath}
  \end{syntax}

\end{function}



\begin{function}[]
  {
    \CurrentFilePathUsed
  }
  \begin{syntax}
    \cs{CurrentFilePathUsed}
  \end{syntax}

\end{function}



\begin{function}[]
  {
    \CurrentFileUsed
  }
  \begin{syntax}
    \cs{CurrentFileUsed}
  \end{syntax}

\end{function}



\begin{function}[]
  {
    \CurrentOption
  }
  \begin{syntax}
    \cs{CurrentOption}
  \end{syntax}

\end{function}



\begin{function}[]
  {
    \DH
  }
  \begin{syntax}
    \cs{DH}
  \end{syntax}

\end{function}



\begin{function}[]
  {
    \DJ
  }
  \begin{syntax}
    \cs{DJ}
  \end{syntax}

\end{function}



\begin{function}[]
  {
    \DebugHooksOff
  }
  \begin{syntax}
    \cs{DebugHooksOff}
  \end{syntax}

\end{function}



\begin{function}[]
  {
    \DebugHooksOn
  }
  \begin{syntax}
    \cs{DebugHooksOn}
  \end{syntax}

\end{function}



\begin{function}[]
  {
    \DebugMarksOff
  }
  \begin{syntax}
    \cs{DebugMarksOff}
  \end{syntax}

\end{function}



\begin{function}[]
  {
    \DebugMarksOn
  }
  \begin{syntax}
    \cs{DebugMarksOn}
  \end{syntax}

\end{function}



\begin{function}[]
  {
    \DebugShipoutsOff
  }
  \begin{syntax}
    \cs{DebugShipoutsOff}
  \end{syntax}

\end{function}



\begin{function}[]
  {
    \DebugShipoutsOn
  }
  \begin{syntax}
    \cs{DebugShipoutsOn}
  \end{syntax}

\end{function}



\begin{function}[]
  {
    \DebugSocketsOff
  }
  \begin{syntax}
    \cs{DebugSocketsOff}
  \end{syntax}

\end{function}



\begin{function}[]
  {
    \DebugSocketsOn
  }
  \begin{syntax}
    \cs{DebugSocketsOn}
  \end{syntax}

\end{function}



\begin{function}[]
  {
    \DebugTablesOff
  }
  \begin{syntax}
    \cs{DebugTablesOff}
  \end{syntax}

\end{function}



\begin{function}[]
  {
    \DebugTablesOn
  }
  \begin{syntax}
    \cs{DebugTablesOn}
  \end{syntax}

\end{function}



\begin{function}[]
  {
    \DebugTemplatesOff
  }
  \begin{syntax}
    \cs{DebugTemplatesOff}
  \end{syntax}

\end{function}



\begin{function}[]
  {
    \DebugTemplatesOn
  }
  \begin{syntax}
    \cs{DebugTemplatesOn}
  \end{syntax}

\end{function}



\begin{function}[]
  {
    \DeclareCaseChangeEquivalent
  }
  \begin{syntax}
    \cs{DeclareCaseChangeEquivalent}
  \end{syntax}

\end{function}



\begin{function}[]
  {
    \DeclareCommandCopy
  }
  \begin{syntax}
    \cs{DeclareCommandCopy}
  \end{syntax}

\end{function}




\begin{function}[]
  {
    \DeclareDefaultHookRule
  }
  \begin{syntax}
    \cs{DeclareDefaultHookRule}
  \end{syntax}

\end{function}



\begin{function}[]
  {
    \DeclareDocumentCommand
  }
  \begin{syntax}
    \cs{DeclareDocumentCommand}
  \end{syntax}

\end{function}



\begin{function}[]
  {
    \DeclareDocumentEnvironment
  }
  \begin{syntax}
    \cs{DeclareDocumentEnvironment}
  \end{syntax}

\end{function}



\begin{function}[]
  {
    \DeclareEmphSequence
  }
  \begin{syntax}
    \cs{DeclareEmphSequence}
  \end{syntax}

\end{function}



\begin{function}[]
  {
    \DeclareEncodingSubset
  }
  \begin{syntax}
    \cs{DeclareEncodingSubset}
  \end{syntax}

\end{function}



\begin{function}[]
  {
    \DeclareEncodingSubset@aux
  }
  \begin{syntax}
    \cs{DeclareEncodingSubset@aux}
  \end{syntax}

\end{function}



\begin{function}[]
  {
    \DeclareEnvironmentCopy
  }
  \begin{syntax}
    \cs{DeclareEnvironmentCopy}
  \end{syntax}

\end{function}



\begin{function}[]
  {
    \DeclareErrorFont
  }
  \begin{syntax}
    \cs{DeclareErrorFont}
  \end{syntax}

\end{function}



\begin{function}[]
  {
    \DeclareExpandableDocumentCommand
  }
  \begin{syntax}
    \cs{DeclareExpandableDocumentCommand}
  \end{syntax}

\end{function}



\begin{function}[]
  {
    \DeclareFixedFont
  }
  \begin{syntax}
    \cs{DeclareFixedFont}
  \end{syntax}

\end{function}



\begin{function}[]
  {
    \DeclareFontEncoding
  }
  \begin{syntax}
    \cs{DeclareFontEncoding}
  \end{syntax}

\end{function}



\begin{function}[]
  {
    \DeclareFontEncoding@
  }
  \begin{syntax}
    \cs{DeclareFontEncoding@}
  \end{syntax}

\end{function}



\begin{function}[]
  {
    \DeclareFontEncoding@saved
  }
  \begin{syntax}
    \cs{DeclareFontEncoding@saved}
  \end{syntax}

\end{function}



\begin{function}[]
  {
    \DeclareFontEncodingDefaults
  }
  \begin{syntax}
    \cs{DeclareFontEncodingDefaults}
  \end{syntax}

\end{function}



\begin{function}[]
  {
    \DeclareFontFamily
  }
  \begin{syntax}
    \cs{DeclareFontFamily}
  \end{syntax}

\end{function}



\begin{function}[]
  {
    \DeclareFontFamilySubstitution
  }
  \begin{syntax}
    \cs{DeclareFontFamilySubstitution}
  \end{syntax}

\end{function}



\begin{function}[]
  {
    \DeclareFontSeriesChangeRule
  }
  \begin{syntax}
    \cs{DeclareFontSeriesChangeRule}
  \end{syntax}

\end{function}



\begin{function}[]
  {
    \DeclareFontSeriesDefault
  }
  \begin{syntax}
    \cs{DeclareFontSeriesDefault}
  \end{syntax}

\end{function}



\begin{function}[]
  {
    \DeclareFontShape
  }
  \begin{syntax}
    \cs{DeclareFontShape}
  \end{syntax}

\end{function}



\begin{function}[]
  {
    \DeclareFontShape@
  }
  \begin{syntax}
    \cs{DeclareFontShape@}
  \end{syntax}

\end{function}



\begin{function}[]
  {
    \DeclareFontShapeChangeRule
  }
  \begin{syntax}
    \cs{DeclareFontShapeChangeRule}
  \end{syntax}

\end{function}



\begin{function}[]
  {
    \DeclareFontSubstitution
  }
  \begin{syntax}
    \cs{DeclareFontSubstitution}
  \end{syntax}

\end{function}



\begin{function}[]
  {
    \DeclareHookRule
  }
  \begin{syntax}
    \cs{DeclareHookRule}
  \end{syntax}

\end{function}



\begin{function}[]
  {
    \DeclareInstance
  }
  \begin{syntax}
    \cs{DeclareInstance}
  \end{syntax}

\end{function}



\begin{function}[]
  {
    \DeclareInstanceCopy
  }
  \begin{syntax}
    \cs{DeclareInstanceCopy}
  \end{syntax}

\end{function}



\begin{function}[]
  {
    \DeclareKeys
  }
  \begin{syntax}
    \cs{DeclareKeys}
  \end{syntax}

\end{function}



\begin{function}[]
  {
    \DeclareLowercaseExclusions
  }
  \begin{syntax}
    \cs{DeclareLowercaseExclusions}
  \end{syntax}

\end{function}



\begin{function}[]
  {
    \DeclareLowercaseMapping
  }
  \begin{syntax}
    \cs{DeclareLowercaseMapping}
  \end{syntax}

\end{function}



\begin{function}[]
  {
    \DeclareMathAccent
  }
  \begin{syntax}
    \cs{DeclareMathAccent}
  \end{syntax}

\end{function}



\begin{function}[]
  {
    \DeclareMathAlphabet
  }
  \begin{syntax}
    \cs{DeclareMathAlphabet}
  \end{syntax}

\end{function}



\begin{function}[]
  {
    \DeclareMathDelimiter
  }
  \begin{syntax}
    \cs{DeclareMathDelimiter}
  \end{syntax}

\end{function}



\begin{function}[]
  {
    \DeclareMathRadical
  }
  \begin{syntax}
    \cs{DeclareMathRadical}
  \end{syntax}

\end{function}



\begin{function}[]
  {
    \DeclareMathScriptfontMapping
  }
  \begin{syntax}
    \cs{DeclareMathScriptfontMapping}
  \end{syntax}

\end{function}



\begin{function}[]
  {
    \DeclareMathSizes
  }
  \begin{syntax}
    \cs{DeclareMathSizes}
  \end{syntax}

\end{function}



\begin{function}[]
  {
    \DeclareMathSymbol
  }
  \begin{syntax}
    \cs{DeclareMathSymbol}
  \end{syntax}

\end{function}



\begin{function}[]
  {
    \DeclareMathVersion
  }
  \begin{syntax}
    \cs{DeclareMathVersion}
  \end{syntax}

\end{function}



\begin{function}[]
  {
    \DeclareOldFontCommand
  }
  \begin{syntax}
    \cs{DeclareOldFontCommand}
  \end{syntax}

\end{function}



\begin{function}[]
  {
    \DeclareOption
  }
  \begin{syntax}
    \cs{DeclareOption}
  \end{syntax}

\end{function}



\begin{function}[]
  {
    \DeclarePreloadSizes
  }
  \begin{syntax}
    \cs{DeclarePreloadSizes}
  \end{syntax}

\end{function}







\begin{function}[]
  {
    \DeclareRobustCommand
  }
  \begin{syntax}
    \cs{DeclareRobustCommand}
  \end{syntax}

\end{function}



\begin{function}[]
  {
    \DeclareSizeFunction
  }
  \begin{syntax}
    \cs{DeclareSizeFunction}
  \end{syntax}

\end{function}



\begin{function}[]
  {
    \DeclareSymbolFont
  }
  \begin{syntax}
    \cs{DeclareSymbolFont}
  \end{syntax}

\end{function}



\begin{function}[]
  {
    \DeclareSymbolFont@m@dropped
  }
  \begin{syntax}
    \cs{DeclareSymbolFont@m@dropped}
  \end{syntax}

\end{function}



\begin{function}[]
  {
    \DeclareSymbolFontAlphabet
  }
  \begin{syntax}
    \cs{DeclareSymbolFontAlphabet}
  \end{syntax}

\end{function}



\begin{function}[]
  {
    \DeclareSymbolFontAlphabet@
  }
  \begin{syntax}
    \cs{DeclareSymbolFontAlphabet@}
  \end{syntax}

\end{function}



\begin{function}[]
  {
    \DeclareTemplateCode
  }
  \begin{syntax}
    \cs{DeclareTemplateCode}
  \end{syntax}

\end{function}



\begin{function}[]
  {
    \DeclareTemplateCopy
  }
  \begin{syntax}
    \cs{DeclareTemplateCopy}
  \end{syntax}

\end{function}



\begin{function}[]
  {
    \DeclareTemplateInterface
  }
  \begin{syntax}
    \cs{DeclareTemplateInterface}
  \end{syntax}

\end{function}



\begin{function}[]
  {
    \DeclareTextAccent
  }
  \begin{syntax}
    \cs{DeclareTextAccent}
  \end{syntax}

\end{function}



\begin{function}[]
  {
    \DeclareTextAccentDefault
  }
  \begin{syntax}
    \cs{DeclareTextAccentDefault}
  \end{syntax}

\end{function}



\begin{function}[]
  {
    \DeclareTextCommand
  }
  \begin{syntax}
    \cs{DeclareTextCommand}
  \end{syntax}

\end{function}



\begin{function}[]
  {
    \DeclareTextCommandDefault
  }
  \begin{syntax}
    \cs{DeclareTextCommandDefault}
  \end{syntax}

\end{function}



\begin{function}[]
  {
    \DeclareTextComposite
  }
  \begin{syntax}
    \cs{DeclareTextComposite}
  \end{syntax}

\end{function}



\begin{function}[]
  {
    \DeclareTextCompositeCommand
  }
  \begin{syntax}
    \cs{DeclareTextCompositeCommand}
  \end{syntax}

\end{function}



\begin{function}[]
  {
    \DeclareTextFontCommand
  }
  \begin{syntax}
    \cs{DeclareTextFontCommand}
  \end{syntax}

\end{function}



\begin{function}[]
  {
    \DeclareTextSymbol
  }
  \begin{syntax}
    \cs{DeclareTextSymbol}
  \end{syntax}

\end{function}



\begin{function}[]
  {
    \DeclareTextSymbolDefault
  }
  \begin{syntax}
    \cs{DeclareTextSymbolDefault}
  \end{syntax}

\end{function}



\begin{function}[]
  {
    \DeclareTitlecaseExclusions
  }
  \begin{syntax}
    \cs{DeclareTitlecaseExclusions}
  \end{syntax}

\end{function}



\begin{function}[]
  {
    \DeclareTitlecaseMapping
  }
  \begin{syntax}
    \cs{DeclareTitlecaseMapping}
  \end{syntax}

\end{function}



\begin{function}[]
  {
    \DeclareUnicodeCharacter
  }
  \begin{syntax}
    \cs{DeclareUnicodeCharacter}
  \end{syntax}

\end{function}



\begin{function}[]
  {
    \DeclareUnknownKeyHandler
  }
  \begin{syntax}
    \cs{DeclareUnknownKeyHandler}
  \end{syntax}

\end{function}



\begin{function}[]
  {
    \DeclareUppercaseExclusions
  }
  \begin{syntax}
    \cs{DeclareUppercaseExclusions}
  \end{syntax}

\end{function}



\begin{function}[]
  {
    \DeclareUppercaseMapping
  }
  \begin{syntax}
    \cs{DeclareUppercaseMapping}
  \end{syntax}

\end{function}



\begin{function}[]
  {
    \Diamond
  }
  \begin{syntax}
    \cs{Diamond}
  \end{syntax}

\end{function}



\begin{function}[]
  {
    \DisableGenericHook
  }
  \begin{syntax}
    \cs{DisableGenericHook}
  \end{syntax}

\end{function}



\begin{function}[]
  {
    \DisableHook
  }
  \begin{syntax}
    \cs{DisableHook}
  \end{syntax}

\end{function}



\begin{function}[]
  {
    \DiscardShipoutBox
  }
  \begin{syntax}
    \cs{DiscardShipoutBox}
  \end{syntax}

\end{function}



\begin{function}[]
  {
    \DocumentMetadata
  }
  \begin{syntax}
    \cs{DocumentMetadata}
  \end{syntax}

\end{function}



\begin{function}[]
  {
    \E
  }
  \begin{syntax}
    \cs{E}
  \end{syntax}

\end{function}



\begin{function}[]
  {
    \ERROR
  }
  \begin{syntax}
    \cs{ERROR}
  \end{syntax}

\end{function}



\begin{function}[]
  {
    \ERRORmissingcells
  }
  \begin{syntax}
    \cs{ERRORmissingcells}
  \end{syntax}

\end{function}



\begin{function}[]
  {
    \ERRORnewtaggingsocket
  }
  \begin{syntax}
    \cs{ERRORnewtaggingsocket}
  \end{syntax}

\end{function}



\begin{function}[]
  {
    \ERRORusetaggingsocket
  }
  \begin{syntax}
    \cs{ERRORusetaggingsocket}
  \end{syntax}

\end{function}



\begin{function}[]
  {
    \EditInstance
  }
  \begin{syntax}
    \cs{EditInstance}
  \end{syntax}

\end{function}



\begin{function}[]
  {
    \EditTemplateDefaults
  }
  \begin{syntax}
    \cs{EditTemplateDefaults}
  \end{syntax}

\end{function}




\begin{function}[]
  {
    \ExecuteOptions
  }
  \begin{syntax}
    \cs{ExecuteOptions}
  \end{syntax}

\end{function}



\begin{function}[]
  {
    \ExpandArgs
  }
  \begin{syntax}
    \cs{ExpandArgs}
  \end{syntax}

\end{function}



\begin{function}[]
  {
    \ExplLoaderFileDate
  }
  \begin{syntax}
    \cs{ExplLoaderFileDate}
  \end{syntax}

\end{function}



\begin{function}[]
  {
    \ExplSyntaxOff
  }
  \begin{syntax}
    \cs{ExplSyntaxOff}
  \end{syntax}

\end{function}



\begin{function}[]
  {
    \ExplSyntaxOn
  }
  \begin{syntax}
    \cs{ExplSyntaxOn}
  \end{syntax}

\end{function}



\begin{function}[]
  {
    \FirstMark
  }
  \begin{syntax}
    \cs{FirstMark}
  \end{syntax}

\end{function}



\begin{function}[]
  {
    \G@refundefinedtrue
  }
  \begin{syntax}
    \cs{G@refundefinedtrue}
  \end{syntax}

\end{function}



\begin{function}[]
  {
    \GenericError
  }
  \begin{syntax}
    \cs{GenericError}
  \end{syntax}

\end{function}



\begin{function}[]
  {
    \GenericInfo
  }
  \begin{syntax}
    \cs{GenericInfo}
  \end{syntax}

\end{function}



\begin{function}[]
  {
    \GenericWarning
  }
  \begin{syntax}
    \cs{GenericWarning}
  \end{syntax}

\end{function}



\begin{function}[]
  {
    \GetDocumentProperty
  }
  \begin{syntax}
    \cs{GetDocumentProperty}
  \end{syntax}

\end{function}



\begin{function}[]
  {
    \H
  }
  \begin{syntax}
    \cs{H}
  \end{syntax}

\end{function}



\begin{function}[]
  {
    \I
  }
  \begin{syntax}
    \cs{I}
  \end{syntax}

\end{function}



\begin{function}[]
  {
    \IJ
  }
  \begin{syntax}
    \cs{IJ}
  \end{syntax}

\end{function}



\begin{function}[]
  {
    \IeC
  }
  \begin{syntax}
    \cs{IeC}
  \end{syntax}

\end{function}





\begin{function}[TF]
  {
    \IfBlank
  }
  \begin{syntax}
    \cs{IfBlankTF} \Arg{data} \Arg{true} \Arg{false}
  \end{syntax}

\end{function}



\begin{function}[TF]
  {
    \IfBoolean
  }
  \begin{syntax}
    \cs{IfBoolean} \meta{boolean} \Arg{true} \Arg{false}
  \end{syntax}

\end{function}




\begin{function}[TF]
  {
    \IfClassAtLeast
  }
  \begin{syntax}
    \cs{IfClassAtLeastTF}
  \end{syntax}

\end{function}



\begin{function}[TF]
  {
    \IfClassLoaded
  }
  \begin{syntax}
    \cs{IfClassLoadedTF}
  \end{syntax}

\end{function}



\begin{function}[TF]
  {
    \IfClassLoadedWithOptions
  }
  \begin{syntax}
    \cs{IfClassLoadedWithOptionsTF}
  \end{syntax}

\end{function}



\begin{function}[TF]
  {
    \IfDocumentMetadata
  }
  \begin{syntax}
    \cs{IfDocumentMetadataTF}
  \end{syntax}

\end{function}


\begin{function}[TF]
  {
    \IfExplAtLeast
  }
  \begin{syntax}
    \cs{IfExplAtLeastTF}
  \end{syntax}

\end{function}




\begin{function}[TF]
  {
    \IfFileAtLeast
  }
  \begin{syntax}
    \cs{IfFileAtLeastTF}
  \end{syntax}

\end{function}


\begin{function}[]
  {
    \IfFileExists
  }
  \begin{syntax}
    \cs{IfFileExists}
  \end{syntax}

\end{function}



\begin{function}[int]
  {
    \IfFileExists@
  }
  \begin{syntax}
    \cs{IfFileExists@}
  \end{syntax}

\end{function}



\begin{function}[int]
  {
    \IfFileExists@@
  }
  \begin{syntax}
    \cs{IfFileExists@@}
  \end{syntax}

\end{function}



\begin{function}[TF]
  {
    \IfFileLoaded
  }
  \begin{syntax}
    \cs{IfFileLoadedTF}
  \end{syntax}

\end{function}



\begin{function}[TF]
  {
    \IfFontSeriesContext
  }
  \begin{syntax}
    \cs{IfFontSeriesContextTF}
  \end{syntax}

\end{function}



\begin{function}[]
  {
    \IfFormatAtLeastF
  }
  \begin{syntax}
    \cs{IfFormatAtLeastF}
  \end{syntax}

\end{function}



\begin{function}[]
  {
    \IfFormatAtLeastT
  }
  \begin{syntax}
    \cs{IfFormatAtLeastT}
  \end{syntax}

\end{function}



\begin{function}[]
  {
    \IfFormatAtLeastTF
  }
  \begin{syntax}
    \cs{IfFormatAtLeastTF}
  \end{syntax}

\end{function}



\begin{function}[]
  {
    \IfHookEmptyF
  }
  \begin{syntax}
    \cs{IfHookEmptyF}
  \end{syntax}

\end{function}



\begin{function}[]
  {
    \IfHookEmptyT
  }
  \begin{syntax}
    \cs{IfHookEmptyT}
  \end{syntax}

\end{function}



\begin{function}[]
  {
    \IfHookEmptyTF
  }
  \begin{syntax}
    \cs{IfHookEmptyTF}
  \end{syntax}

\end{function}



\begin{function}[]
  {
    \IfHookExistsTF
  }
  \begin{syntax}
    \cs{IfHookExistsTF}
  \end{syntax}

\end{function}



\begin{function}[]
  {
    \IfInstanceExistsF
  }
  \begin{syntax}
    \cs{IfInstanceExistsF}
  \end{syntax}

\end{function}



\begin{function}[]
  {
    \IfInstanceExistsT
  }
  \begin{syntax}
    \cs{IfInstanceExistsT}
  \end{syntax}

\end{function}



\begin{function}[]
  {
    \IfInstanceExistsTF
  }
  \begin{syntax}
    \cs{IfInstanceExistsTF}
  \end{syntax}

\end{function}



\begin{function}[]
  {
    \IfLabelExistsF
  }
  \begin{syntax}
    \cs{IfLabelExistsF}
  \end{syntax}

\end{function}



\begin{function}[]
  {
    \IfLabelExistsT
  }
  \begin{syntax}
    \cs{IfLabelExistsT}
  \end{syntax}

\end{function}



\begin{function}[]
  {
    \IfLabelExistsTF
  }
  \begin{syntax}
    \cs{IfLabelExistsTF}
  \end{syntax}

\end{function}



\begin{function}[]
  {
    \IfMarksEqualF
  }
  \begin{syntax}
    \cs{IfMarksEqualF}
  \end{syntax}

\end{function}



\begin{function}[]
  {
    \IfMarksEqualT
  }
  \begin{syntax}
    \cs{IfMarksEqualT}
  \end{syntax}

\end{function}



\begin{function}[]
  {
    \IfMarksEqualTF
  }
  \begin{syntax}
    \cs{IfMarksEqualTF}
  \end{syntax}

\end{function}



\begin{function}[]
  {
    \IfNoValueF
  }
  \begin{syntax}
    \cs{IfNoValueF}
  \end{syntax}

\end{function}



\begin{function}[]
  {
    \IfNoValueT
  }
  \begin{syntax}
    \cs{IfNoValueT}
  \end{syntax}

\end{function}



\begin{function}[]
  {
    \IfNoValueTF
  }
  \begin{syntax}
    \cs{IfNoValueTF}
  \end{syntax}

\end{function}



\begin{function}[]
  {
    \IfPDFManagementActiveF
  }
  \begin{syntax}
    \cs{IfPDFManagementActiveF}
  \end{syntax}

\end{function}



\begin{function}[]
  {
    \IfPDFManagementActiveT
  }
  \begin{syntax}
    \cs{IfPDFManagementActiveT}
  \end{syntax}

\end{function}



\begin{function}[]
  {
    \IfPDFManagementActiveTF
  }
  \begin{syntax}
    \cs{IfPDFManagementActiveTF}
  \end{syntax}

\end{function}



\begin{function}[]
  {
    \IfPackageAtLeastF
  }
  \begin{syntax}
    \cs{IfPackageAtLeastF}
  \end{syntax}

\end{function}



\begin{function}[]
  {
    \IfPackageAtLeastT
  }
  \begin{syntax}
    \cs{IfPackageAtLeastT}
  \end{syntax}

\end{function}



\begin{function}[]
  {
    \IfPackageAtLeastTF
  }
  \begin{syntax}
    \cs{IfPackageAtLeastTF}
  \end{syntax}

\end{function}



\begin{function}[]
  {
    \IfPackageLoadedF
  }
  \begin{syntax}
    \cs{IfPackageLoadedF}
  \end{syntax}

\end{function}



\begin{function}[]
  {
    \IfPackageLoadedT
  }
  \begin{syntax}
    \cs{IfPackageLoadedT}
  \end{syntax}

\end{function}



\begin{function}[]
  {
    \IfPackageLoadedTF
  }
  \begin{syntax}
    \cs{IfPackageLoadedTF}
  \end{syntax}

\end{function}



\begin{function}[]
  {
    \IfPackageLoadedWithOptionsF
  }
  \begin{syntax}
    \cs{IfPackageLoadedWithOptionsF}
  \end{syntax}

\end{function}



\begin{function}[]
  {
    \IfPackageLoadedWithOptionsT
  }
  \begin{syntax}
    \cs{IfPackageLoadedWithOptionsT}
  \end{syntax}

\end{function}



\begin{function}[]
  {
    \IfPackageLoadedWithOptionsTF
  }
  \begin{syntax}
    \cs{IfPackageLoadedWithOptionsTF}
  \end{syntax}

\end{function}



\begin{function}[]
  {
    \IfPropertyExistsF
  }
  \begin{syntax}
    \cs{IfPropertyExistsF}
  \end{syntax}

\end{function}



\begin{function}[]
  {
    \IfPropertyExistsT
  }
  \begin{syntax}
    \cs{IfPropertyExistsT}
  \end{syntax}

\end{function}



\begin{function}[]
  {
    \IfPropertyExistsTF
  }
  \begin{syntax}
    \cs{IfPropertyExistsTF}
  \end{syntax}

\end{function}



\begin{function}[]
  {
    \IfPropertyRecordedF
  }
  \begin{syntax}
    \cs{IfPropertyRecordedF}
  \end{syntax}

\end{function}



\begin{function}[]
  {
    \IfPropertyRecordedT
  }
  \begin{syntax}
    \cs{IfPropertyRecordedT}
  \end{syntax}

\end{function}



\begin{function}[]
  {
    \IfPropertyRecordedTF
  }
  \begin{syntax}
    \cs{IfPropertyRecordedTF}
  \end{syntax}

\end{function}



\begin{function}[]
  {
    \IfSocketExistsF
  }
  \begin{syntax}
    \cs{IfSocketExistsF}
  \end{syntax}

\end{function}



\begin{function}[]
  {
    \IfSocketExistsT
  }
  \begin{syntax}
    \cs{IfSocketExistsT}
  \end{syntax}

\end{function}



\begin{function}[]
  {
    \IfSocketExistsTF
  }
  \begin{syntax}
    \cs{IfSocketExistsTF}
  \end{syntax}

\end{function}



\begin{function}[]
  {
    \IfSocketPlugAssignedF
  }
  \begin{syntax}
    \cs{IfSocketPlugAssignedF}
  \end{syntax}

\end{function}



\begin{function}[]
  {
    \IfSocketPlugAssignedT
  }
  \begin{syntax}
    \cs{IfSocketPlugAssignedT}
  \end{syntax}

\end{function}



\begin{function}[]
  {
    \IfSocketPlugAssignedTF
  }
  \begin{syntax}
    \cs{IfSocketPlugAssignedTF}
  \end{syntax}

\end{function}



\begin{function}[]
  {
    \IfSocketPlugExistsF
  }
  \begin{syntax}
    \cs{IfSocketPlugExistsF}
  \end{syntax}

\end{function}



\begin{function}[]
  {
    \IfSocketPlugExistsT
  }
  \begin{syntax}
    \cs{IfSocketPlugExistsT}
  \end{syntax}

\end{function}



\begin{function}[]
  {
    \IfSocketPlugExistsTF
  }
  \begin{syntax}
    \cs{IfSocketPlugExistsTF}
  \end{syntax}

\end{function}




\begin{function}[]
  {
    \IfValueF
  }
  \begin{syntax}
    \cs{IfValueF}
  \end{syntax}

\end{function}



\begin{function}[]
  {
    \IfValueT
  }
  \begin{syntax}
    \cs{IfValueT}
  \end{syntax}

\end{function}



\begin{function}[]
  {
    \IfValueTF
  }
  \begin{syntax}
    \cs{IfValueTF}
  \end{syntax}

\end{function}



\begin{function}[]
  {
    \IncludeInRelease
  }
  \begin{syntax}
    \cs{IncludeInRelease}
  \end{syntax}

\end{function}



\begin{function}[]
  {
    \IndentBox
  }
  \begin{syntax}
    \cs{IndentBox}
  \end{syntax}

\end{function}



\begin{function}[]
  {
    \InputIfFileExists
  }
  \begin{syntax}
    \cs{InputIfFileExists}
  \end{syntax}

\end{function}



\begin{function}[]
  {
    \InsertMark
  }
  \begin{syntax}
    \cs{InsertMark}
  \end{syntax}

\end{function}



\begin{function}[]
  {
    \InstanceValue
  }
  \begin{syntax}
    \cs{InstanceValue}
  \end{syntax}

\end{function}



\begin{function}[]
  {
    \J
  }
  \begin{syntax}
    \cs{J}
  \end{syntax}

\end{function}



\begin{function}[]
  {
    \Join
  }
  \begin{syntax}
    \cs{Join}
  \end{syntax}

\end{function}



\begin{function}[]
  {
    \KeyValue
  }
  \begin{syntax}
    \cs{KeyValue}
  \end{syntax}

\end{function}



\begin{function}[]
  {
    \L
  }
  \begin{syntax}
    \cs{L}
  \end{syntax}

\end{function}








\begin{function}[]
  {
    \LastDeclaredEncoding
  }
  \begin{syntax}
    \cs{LastDeclaredEncoding}
  \end{syntax}

\end{function}



\begin{function}[]
  {
    \LastMark
  }
  \begin{syntax}
    \cs{LastMark}
  \end{syntax}

\end{function}



\begin{function}[]
  {
    \LinkTargetOff
  }
  \begin{syntax}
    \cs{LinkTargetOff}
  \end{syntax}

\end{function}



\begin{function}[]
  {
    \LinkTargetOn
  }
  \begin{syntax}
    \cs{LinkTargetOn}
  \end{syntax}

\end{function}



\begin{function}[]
  {
    \LoadClass
  }
  \begin{syntax}
    \cs{LoadClass}
  \end{syntax}

\end{function}



\begin{function}[]
  {
    \LoadClassWithOptions
  }
  \begin{syntax}
    \cs{LoadClassWithOptions}
  \end{syntax}

\end{function}



\begin{function}[]
  {
    \LoadFontDefinitionFile
  }
  \begin{syntax}
    \cs{LoadFontDefinitionFile}
  \end{syntax}

\end{function}



\begin{function}[]
  {
    \LogDocumentProperties
  }
  \begin{syntax}
    \cs{LogDocumentProperties}
  \end{syntax}

\end{function}



\begin{function}[]
  {
    \LogHook
  }
  \begin{syntax}
    \cs{LogHook}
  \end{syntax}

\end{function}



\begin{function}[]
  {
    \LogSocket
  }
  \begin{syntax}
    \cs{LogSocket}
  \end{syntax}

\end{function}



\begin{function}[]
  {
    \MakeLinkTarget
  }
  \begin{syntax}
    \cs{MakeLinkTarget}
  \end{syntax}

\end{function}



\begin{function}[]
  {
    \MakeRobust
  }
  \begin{syntax}
    \cs{MakeRobust}
  \end{syntax}

\end{function}



\begin{function}[]
  {
    \MakeUppercase
  }
  \begin{syntax}
    \cs{MakeUppercase}
  \end{syntax}

\end{function}



\begin{function}[]
  {
    \MathCollectFalse
  }
  \begin{syntax}
    \cs{MathCollectFalse}
  \end{syntax}

\end{function}



\begin{function}[]
  {
    \MathCollectTrue
  }
  \begin{syntax}
    \cs{MathCollectTrue}
  \end{syntax}

\end{function}



\begin{function}[]
  {
    \MathMLarg
  }
  \begin{syntax}
    \cs{MathMLarg}
  \end{syntax}

\end{function}



\begin{function}[]
  {
    \MathMLintent
  }
  \begin{syntax}
    \cs{MathMLintent}
  \end{syntax}

\end{function}



\begin{function}[]
  {
    \MessageBreak
  }
  \begin{syntax}
    \cs{MessageBreak}
  \end{syntax}

\end{function}



\begin{function}[]
  {
    \NG
  }
  \begin{syntax}
    \cs{NG}
  \end{syntax}

\end{function}



\begin{function}[]
  {
    \NeedsDocumentMetadata
  }
  \begin{syntax}
    \cs{NeedsDocumentMetadata}
  \end{syntax}

\end{function}



\begin{function}[]
  {
    \NeedsTeXFormat
  }
  \begin{syntax}
    \cs{NeedsTeXFormat}
  \end{syntax}

\end{function}



\begin{function}[]
  {
    \NewCommandCopy
  }
  \begin{syntax}
    \cs{NewCommandCopy}
  \end{syntax}

\end{function}



\begin{function}[]
  {
    \NewDocumentCommand
  }
  \begin{syntax}
    \cs{NewDocumentCommand}
  \end{syntax}

\end{function}



\begin{function}[]
  {
    \NewDocumentEnvironment
  }
  \begin{syntax}
    \cs{NewDocumentEnvironment}
  \end{syntax}

\end{function}



\begin{function}[]
  {
    \NewEnvironmentCopy
  }
  \begin{syntax}
    \cs{NewEnvironmentCopy}
  \end{syntax}

\end{function}



\begin{function}[]
  {
    \NewExpandableDocumentCommand
  }
  \begin{syntax}
    \cs{NewExpandableDocumentCommand}
  \end{syntax}

\end{function}



\begin{function}[]
  {
    \NewHook
  }
  \begin{syntax}
    \cs{NewHook}
  \end{syntax}

\end{function}



\begin{function}[]
  {
    \NewHookWithArguments
  }
  \begin{syntax}
    \cs{NewHookWithArguments}
  \end{syntax}

\end{function}



\begin{function}[]
  {
    \NewMarkClass
  }
  \begin{syntax}
    \cs{NewMarkClass}
  \end{syntax}

\end{function}



\begin{function}[]
  {
    \NewMirroredHookPair
  }
  \begin{syntax}
    \cs{NewMirroredHookPair}
  \end{syntax}

\end{function}



\begin{function}[]
  {
    \NewMirroredHookPairWithArguments
  }
  \begin{syntax}
    \cs{NewMirroredHookPairWithArguments}
  \end{syntax}

\end{function}




\begin{function}[]
  {
    \NewProperty
  }
  \begin{syntax}
    \cs{NewProperty}
  \end{syntax}

\end{function}



\begin{function}[]
  {
    \NewReversedHook
  }
  \begin{syntax}
    \cs{NewReversedHook}
  \end{syntax}

\end{function}



\begin{function}[]
  {
    \NewReversedHookWithArguments
  }
  \begin{syntax}
    \cs{NewReversedHookWithArguments}
  \end{syntax}

\end{function}



\begin{function}[]
  {
    \NewSocket
  }
  \begin{syntax}
    \cs{NewSocket}
  \end{syntax}

\end{function}



\begin{function}[]
  {
    \NewSocketPlug
  }
  \begin{syntax}
    \cs{NewSocketPlug}
  \end{syntax}

\end{function}



\begin{function}[]
  {
    \NewStructureName
  }
  \begin{syntax}
    \cs{NewStructureName}
  \end{syntax}

\end{function}



\begin{function}[]
  {
    \NewTaggingSocket
  }
  \begin{syntax}
    \cs{NewTaggingSocket}
  \end{syntax}

\end{function}



\begin{function}[]
  {
    \NewTaggingSocketPlug
  }
  \begin{syntax}
    \cs{NewTaggingSocketPlug}
  \end{syntax}

\end{function}



\begin{function}[]
  {
    \NewTemplateType
  }
  \begin{syntax}
    \cs{NewTemplateType}
  \end{syntax}

\end{function}



\begin{function}[]
  {
    \NextLinkTarget
  }
  \begin{syntax}
    \cs{NextLinkTarget}
  \end{syntax}

\end{function}



\begin{function}[]
  {
    \NoCaseChange
  }
  \begin{syntax}
    \cs{NoCaseChange}
  \end{syntax}

\end{function}



\begin{function}[]
  {
    \NoValue
  }
  \begin{syntax}
    \cs{NoValue}
  \end{syntax}

\end{function}



\begin{function}[]
  {
    \O
  }
  \begin{syntax}
    \cs{O}
  \end{syntax}

\end{function}



\begin{function}[]
  {
    \OE
  }
  \begin{syntax}
    \cs{OE}
  \end{syntax}

\end{function}



\begin{function}[]
  {
    \OmitIndent
  }
  \begin{syntax}
    \cs{OmitIndent}
  \end{syntax}

\end{function}



\begin{function}[]
  {
    \OptionNotUsed
  }
  \begin{syntax}
    \cs{OptionNotUsed}
  \end{syntax}

\end{function}



\begin{function}[]
  {
    \P
  }
  \begin{syntax}
    \cs{P}
  \end{syntax}

\end{function}



\begin{function}[]
  {
    \PackageError
  }
  \begin{syntax}
    \cs{PackageError}
  \end{syntax}

\end{function}



\begin{function}[]
  {
    \PackageInfo
  }
  \begin{syntax}
    \cs{PackageInfo}
  \end{syntax}

\end{function}



\begin{function}[]
  {
    \PackageNote
  }
  \begin{syntax}
    \cs{PackageNote}
  \end{syntax}

\end{function}



\begin{function}[]
  {
    \PackageNoteNoLine
  }
  \begin{syntax}
    \cs{PackageNoteNoLine}
  \end{syntax}

\end{function}



\begin{function}[]
  {
    \PackageWarning
  }
  \begin{syntax}
    \cs{PackageWarning}
  \end{syntax}

\end{function}



\begin{function}[]
  {
    \PackageWarningNoLine
  }
  \begin{syntax}
    \cs{PackageWarningNoLine}
  \end{syntax}

\end{function}



\begin{function}[]
  {
    \PassOptionsToClass
  }
  \begin{syntax}
    \cs{PassOptionsToClass}
  \end{syntax}

\end{function}



\begin{function}[]
  {
    \PassOptionsToPackage
  }
  \begin{syntax}
    \cs{PassOptionsToPackage}
  \end{syntax}

\end{function}



\begin{function}[]
  {
    \PopDefaultHookLabel
  }
  \begin{syntax}
    \cs{PopDefaultHookLabel}
  \end{syntax}

\end{function}



\begin{function}[]
  {
    \Pr
  }
  \begin{syntax}
    \cs{Pr}
  \end{syntax}

\end{function}



\begin{function}[]
  {
    \PreviousTotalPages
  }
  \begin{syntax}
    \cs{PreviousTotalPages}
  \end{syntax}

\end{function}



\begin{function}[]
  {
    \ProcessKeyOptions
  }
  \begin{syntax}
    \cs{ProcessKeyOptions}
  \end{syntax}

\end{function}



\begin{function}[]
  {
    \ProcessList
  }
  \begin{syntax}
    \cs{ProcessList}
  \end{syntax}

\end{function}



\begin{function}[]
  {
    \ProcessOptions
  }
  \begin{syntax}
    \cs{ProcessOptions}
  \end{syntax}

\end{function}



\begin{function}[]
  {
    \ProcessedArgument
  }
  \begin{syntax}
    \cs{ProcessedArgument}
  \end{syntax}

\end{function}



\begin{function}[]
  {
    \ProvideDocumentCommand
  }
  \begin{syntax}
    \cs{ProvideDocumentCommand}
  \end{syntax}

\end{function}



\begin{function}[]
  {
    \ProvideDocumentEnvironment
  }
  \begin{syntax}
    \cs{ProvideDocumentEnvironment}
  \end{syntax}

\end{function}



\begin{function}[]
  {
    \ProvideExpandableDocumentCommand
  }
  \begin{syntax}
    \cs{ProvideExpandableDocumentCommand}
  \end{syntax}

\end{function}



\begin{function}[]
  {
    \ProvideHook
  }
  \begin{syntax}
    \cs{ProvideHook}
  \end{syntax}

\end{function}



\begin{function}[]
  {
    \ProvideMirroredHookPair
  }
  \begin{syntax}
    \cs{ProvideMirroredHookPair}
  \end{syntax}

\end{function}



\begin{function}[]
  {
    \ProvideReversedHook
  }
  \begin{syntax}
    \cs{ProvideReversedHook}
  \end{syntax}

\end{function}



\begin{function}[]
  {
    \ProvideTextCommand
  }
  \begin{syntax}
    \cs{ProvideTextCommand}
  \end{syntax}

\end{function}



\begin{function}[]
  {
    \ProvideTextCommandDefault
  }
  \begin{syntax}
    \cs{ProvideTextCommandDefault}
  \end{syntax}

\end{function}



\begin{function}[]
  {
    \ProvidesClass
  }
  \begin{syntax}
    \cs{ProvidesClass}
  \end{syntax}

\end{function}



\begin{function}[]
  {
    \ProvidesFile
  }
  \begin{syntax}
    \cs{ProvidesFile}
  \end{syntax}

\end{function}



\begin{function}[]
  {
    \ProvidesPackage
  }
  \begin{syntax}
    \cs{ProvidesPackage}
  \end{syntax}

\end{function}



\begin{function}[]
  {
    \PushDefaultHookLabel
  }
  \begin{syntax}
    \cs{PushDefaultHookLabel}
  \end{syntax}

\end{function}



\begin{function}[]
  {
    \RawIndent
  }
  \begin{syntax}
    \cs{RawIndent}
  \end{syntax}

\end{function}



\begin{function}[]
  {
    \RawNoindent
  }
  \begin{syntax}
    \cs{RawNoindent}
  \end{syntax}

\end{function}



\begin{function}[]
  {
    \RawParEnd
  }
  \begin{syntax}
    \cs{RawParEnd}
  \end{syntax}

\end{function}



\begin{function}[]
  {
    \RawShipout
  }
  \begin{syntax}
    \cs{RawShipout}
  \end{syntax}

\end{function}



\begin{function}[]
  {
    \ReadonlyShipoutCounter
  }
  \begin{syntax}
    \cs{ReadonlyShipoutCounter}
  \end{syntax}

\end{function}



\begin{function}[]
  {
    \RecordProperties
  }
  \begin{syntax}
    \cs{RecordProperties}
  \end{syntax}

\end{function}



\begin{function}[]
  {
    \Ref
  }
  \begin{syntax}
    \cs{Ref}
  \end{syntax}

\end{function}



\begin{function}[]
  {
    \RefProperty
  }
  \begin{syntax}
    \cs{RefProperty}
  \end{syntax}

\end{function}



\begin{function}[]
  {
    \RefUndefinedWarn
  }
  \begin{syntax}
    \cs{RefUndefinedWarn}
  \end{syntax}

\end{function}



\begin{function}[]
  {
    \RemoveFromHook
  }
  \begin{syntax}
    \cs{RemoveFromHook}
  \end{syntax}

\end{function}



\begin{function}[]
  {
    \RenewCommandCopy
  }
  \begin{syntax}
    \cs{RenewCommandCopy}
  \end{syntax}

\end{function}



\begin{function}[]
  {
    \RenewDocumentCommand
  }
  \begin{syntax}
    \cs{RenewDocumentCommand}
  \end{syntax}

\end{function}



\begin{function}[]
  {
    \RenewDocumentEnvironment
  }
  \begin{syntax}
    \cs{RenewDocumentEnvironment}
  \end{syntax}

\end{function}



\begin{function}[]
  {
    \RenewEnvironmentCopy
  }
  \begin{syntax}
    \cs{RenewEnvironmentCopy}
  \end{syntax}

\end{function}



\begin{function}[]
  {
    \RenewExpandableDocumentCommand
  }
  \begin{syntax}
    \cs{RenewExpandableDocumentCommand}
  \end{syntax}

\end{function}



\begin{function}[]
  {
    \RequirePackage
  }
  \begin{syntax}
    \cs{RequirePackage}
  \end{syntax}

\end{function}



\begin{function}[]
  {
    \RequirePackageWithOptions
  }
  \begin{syntax}
    \cs{RequirePackageWithOptions}
  \end{syntax}

\end{function}



\begin{function}[]
  {
    \RestoreLastSkip
  }
  \begin{syntax}
    \cs{RestoreLastSkip}
  \end{syntax}

\end{function}



\begin{function}[]
  {
    \ResumeTagging
  }
  \begin{syntax}
    \cs{ResumeTagging}
  \end{syntax}

\end{function}



\begin{function}[]
  {
    \ReverseBoolean
  }
  \begin{syntax}
    \cs{ReverseBoolean}
  \end{syntax}

\end{function}



\begin{function}[]
  {
    \Roman
  }
  \begin{syntax}
    \cs{Roman}
  \end{syntax}

\end{function}



\begin{function}[]
  {
    \S
  }
  \begin{syntax}
    \cs{S}
  \end{syntax}

\end{function}



\begin{function}[]
  {
    \SS
  }
  \begin{syntax}
    \cs{SS}
  \end{syntax}

\end{function}



\begin{function}[]
  {
    \SaveLastSkip
  }
  \begin{syntax}
    \cs{SaveLastSkip}
  \end{syntax}

\end{function}



\begin{function}[]
  {
    \SetDefaultHookLabel
  }
  \begin{syntax}
    \cs{SetDefaultHookLabel}
  \end{syntax}

\end{function}



\begin{function}[]
  {
    \SetKeys
  }
  \begin{syntax}
    \cs{SetKeys}
  \end{syntax}

\end{function}



\begin{function}[]
  {
    \SetKnownTemplateKeys
  }
  \begin{syntax}
    \cs{SetKnownTemplateKeys}
  \end{syntax}

\end{function}



\begin{function}[]
  {
    \SetMathAlphabet
  }
  \begin{syntax}
    \cs{SetMathAlphabet}
  \end{syntax}

\end{function}



\begin{function}[]
  {
    \SetMathAlphabet@
  }
  \begin{syntax}
    \cs{SetMathAlphabet@}
  \end{syntax}

\end{function}



\begin{function}[]
  {
    \SetProperty
  }
  \begin{syntax}
    \cs{SetProperty}
  \end{syntax}

\end{function}



\begin{function}[]
  {
    \SetSymbolFont
  }
  \begin{syntax}
    \cs{SetSymbolFont}
  \end{syntax}

\end{function}



\begin{function}[]
  {
    \SetSymbolFont@
  }
  \begin{syntax}
    \cs{SetSymbolFont@}
  \end{syntax}

\end{function}



\begin{function}[]
  {
    \SetSymbolFont@m@dropped
  }
  \begin{syntax}
    \cs{SetSymbolFont@m@dropped}
  \end{syntax}

\end{function}



\begin{function}[]
  {
    \SetTemplateKeys
  }
  \begin{syntax}
    \cs{SetTemplateKeys}
  \end{syntax}

\end{function}



\begin{function}[]
  {
    \ShipoutBox
  }
  \begin{syntax}
    \cs{ShipoutBox}
  \end{syntax}

\end{function}



\begin{function}[]
  {
    \ShowCommand
  }
  \begin{syntax}
    \cs{ShowCommand}
  \end{syntax}

\end{function}



\begin{function}[]
  {
    \ShowDocumentProperties
  }
  \begin{syntax}
    \cs{ShowDocumentProperties}
  \end{syntax}

\end{function}



\begin{function}[]
  {
    \ShowEnvironment
  }
  \begin{syntax}
    \cs{ShowEnvironment}
  \end{syntax}

\end{function}



\begin{function}[]
  {
    \ShowFloat
  }
  \begin{syntax}
    \cs{ShowFloat}
  \end{syntax}

\end{function}



\begin{function}[]
  {
    \ShowHook
  }
  \begin{syntax}
    \cs{ShowHook}
  \end{syntax}

\end{function}



\begin{function}[]
  {
    \ShowInstanceValues
  }
  \begin{syntax}
    \cs{ShowInstanceValues}
  \end{syntax}

\end{function}



\begin{function}[]
  {
    \ShowMarksAt
  }
  \begin{syntax}
    \cs{ShowMarksAt}
  \end{syntax}

\end{function}



\begin{function}[]
  {
    \ShowSocket
  }
  \begin{syntax}
    \cs{ShowSocket}
  \end{syntax}

\end{function}



\begin{function}[]
  {
    \ShowTemplateCode
  }
  \begin{syntax}
    \cs{ShowTemplateCode}
  \end{syntax}

\end{function}



\begin{function}[]
  {
    \ShowTemplateDefaults
  }
  \begin{syntax}
    \cs{ShowTemplateDefaults}
  \end{syntax}

\end{function}



\begin{function}[]
  {
    \ShowTemplateInterface
  }
  \begin{syntax}
    \cs{ShowTemplateInterface}
  \end{syntax}

\end{function}



\begin{function}[]
  {
    \ShowTemplateVariables
  }
  \begin{syntax}
    \cs{ShowTemplateVariables}
  \end{syntax}

\end{function}



\begin{function}[]
  {
    \SplitArgument
  }
  \begin{syntax}
    \cs{SplitArgument}
  \end{syntax}

\end{function}



\begin{function}[]
  {
    \SplitList
  }
  \begin{syntax}
    \cs{SplitList}
  \end{syntax}

\end{function}



\begin{function}[]
  {
    \SuspendTagging
  }
  \begin{syntax}
    \cs{SuspendTagging}
  \end{syntax}

\end{function}



\begin{function}[]
  {
    \T
  }
  \begin{syntax}
    \cs{T}
  \end{syntax}

\end{function}



\begin{function}[]
  {
    \TH
  }
  \begin{syntax}
    \cs{TH}
  \end{syntax}

\end{function}





\begin{function}[]
  {
    \TextOrMath
  }
  \begin{syntax}
    \cs{TextOrMath}
  \end{syntax}

\end{function}



\begin{function}[]
  {
    \TextSymbolUnavailable
  }
  \begin{syntax}
    \cs{TextSymbolUnavailable}
  \end{syntax}

\end{function}



\begin{function}[]
  {
    \TopMark
  }
  \begin{syntax}
    \cs{TopMark}
  \end{syntax}

\end{function}



\begin{function}[]
  {
    \TrimSpaces
  }
  \begin{syntax}
    \cs{TrimSpaces}
  \end{syntax}

\end{function}



\begin{function}[]
  {
    \UTFviii@four@octets
  }
  \begin{syntax}
    \cs{UTFviii@four@octets}
  \end{syntax}

\end{function}



\begin{function}[]
  {
    \UTFviii@four@octets@@
  }
  \begin{syntax}
    \cs{UTFviii@four@octets@@}
  \end{syntax}

\end{function}



\begin{function}[]
  {
    \UTFviii@invalid@err
  }
  \begin{syntax}
    \cs{UTFviii@invalid@err}
  \end{syntax}

\end{function}



\begin{function}[]
  {
    \UTFviii@invalid@err@@
  }
  \begin{syntax}
    \cs{UTFviii@invalid@err@@}
  \end{syntax}

\end{function}



\begin{function}[]
  {
    \UTFviii@three@octets
  }
  \begin{syntax}
    \cs{UTFviii@three@octets}
  \end{syntax}

\end{function}



\begin{function}[]
  {
    \UTFviii@three@octets@@
  }
  \begin{syntax}
    \cs{UTFviii@three@octets@@}
  \end{syntax}

\end{function}



\begin{function}[]
  {
    \UTFviii@two@octets
  }
  \begin{syntax}
    \cs{UTFviii@two@octets}
  \end{syntax}

\end{function}



\begin{function}[]
  {
    \UTFviii@two@octets@@
  }
  \begin{syntax}
    \cs{UTFviii@two@octets@@}
  \end{syntax}

\end{function}



\begin{function}[]
  {
    \UTFviii@undefined@err
  }
  \begin{syntax}
    \cs{UTFviii@undefined@err}
  \end{syntax}

\end{function}



\begin{function}[]
  {
    \UTFviii@undefined@err@@
  }
  \begin{syntax}
    \cs{UTFviii@undefined@err@@}
  \end{syntax}

\end{function}



\begin{function}[]
  {
    \Ucharcat
  }
  \begin{syntax}
    \cs{Ucharcat}
  \end{syntax}

\end{function}



\begin{function}[]
  {
    \Umathaxis
  }
  \begin{syntax}
    \cs{Umathaxis}
  \end{syntax}

\end{function}



\begin{function}[]
  {
    \Umathcode
  }
  \begin{syntax}
    \cs{Umathcode}
  \end{syntax}

\end{function}



\begin{function}[]
  {
    \Umathsubshiftdown
  }
  \begin{syntax}
    \cs{Umathsubshiftdown}
  \end{syntax}

\end{function}



\begin{function}[]
  {
    \Umathsubtopmax
  }
  \begin{syntax}
    \cs{Umathsubtopmax}
  \end{syntax}

\end{function}



\begin{function}[]
  {
    \Umathsupbottommin
  }
  \begin{syntax}
    \cs{Umathsupbottommin}
  \end{syntax}

\end{function}



\begin{function}[]
  {
    \Umathsupshiftup
  }
  \begin{syntax}
    \cs{Umathsupshiftup}
  \end{syntax}

\end{function}



\begin{function}[]
  {
    \UndeclareTextCommand
  }
  \begin{syntax}
    \cs{UndeclareTextCommand}
  \end{syntax}

\end{function}



\begin{function}[]
  {
    \UnicodeEncodingName
  }
  \begin{syntax}
    \cs{UnicodeEncodingName}
  \end{syntax}

\end{function}



\begin{function}[]
  {
    \UnusedTemplateKeys
  }
  \begin{syntax}
    \cs{UnusedTemplateKeys}
  \end{syntax}

\end{function}



\begin{function}[]
  {
    \UseExpandableTaggingSocket
  }
  \begin{syntax}
    \cs{UseExpandableTaggingSocket}
  \end{syntax}

\end{function}



\begin{function}[]
  {
    \UseHook
  }
  \begin{syntax}
    \cs{UseHook}
  \end{syntax}

\end{function}



\begin{function}[]
  {
    \UseHookWithArguments
  }
  \begin{syntax}
    \cs{UseHookWithArguments}
  \end{syntax}

\end{function}



\begin{function}[]
  {
    \UseInstance
  }
  \begin{syntax}
    \cs{UseInstance}
  \end{syntax}

\end{function}



\begin{function}[]
  {
    \UseLegacyTextSymbols
  }
  \begin{syntax}
    \cs{UseLegacyTextSymbols}
  \end{syntax}

\end{function}



\begin{function}[]
  {
    \UseName
  }
  \begin{syntax}
    \cs{UseName}
  \end{syntax}

\end{function}



\begin{function}[]
  {
    \UseOneTimeHook
  }
  \begin{syntax}
    \cs{UseOneTimeHook}
  \end{syntax}

\end{function}



\begin{function}[]
  {
    \UseOneTimeHookWithArguments
  }
  \begin{syntax}
    \cs{UseOneTimeHookWithArguments}
  \end{syntax}

\end{function}



\begin{function}[]
  {
    \UseRawInputEncoding
  }
  \begin{syntax}
    \cs{UseRawInputEncoding}
  \end{syntax}

\end{function}



\begin{function}[]
  {
    \UseSocket
  }
  \begin{syntax}
    \cs{UseSocket}
  \end{syntax}

\end{function}



\begin{function}[]
  {
    \UseStructureName
  }
  \begin{syntax}
    \cs{UseStructureName}
  \end{syntax}

\end{function}



\begin{function}[]
  {
    \UseTaggingSocket
  }
  \begin{syntax}
    \cs{UseTaggingSocket}
  \end{syntax}

\end{function}



\begin{function}[]
  {
    \UseTemplate
  }
  \begin{syntax}
    \cs{UseTemplate}
  \end{syntax}

\end{function}



\begin{function}[]
  {
    \UseTextAccent
  }
  \begin{syntax}
    \cs{UseTextAccent}
  \end{syntax}

\end{function}



\begin{function}[]
  {
    \UseTextSymbol
  }
  \begin{syntax}
    \cs{UseTextSymbol}
  \end{syntax}

\end{function}



\begin{function}[]
  {
    \XeTeXcharclass
  }
  \begin{syntax}
    \cs{XeTeXcharclass}
  \end{syntax}

\end{function}



\begin{function}[]
  {
    \XeTeXdashbreakstate
  }
  \begin{syntax}
    \cs{XeTeXdashbreakstate}
  \end{syntax}

\end{function}



\begin{function}[]
  {
    \XeTeXmathcode
  }
  \begin{syntax}
    \cs{XeTeXmathcode}
  \end{syntax}

\end{function}



\begin{function}[]
  {
    \XeTeXrevision
  }
  \begin{syntax}
    \cs{XeTeXrevision}
  \end{syntax}

\end{function}



\begin{function}[]
  {
    \XeTeXuseglyphmetrics
  }
  \begin{syntax}
    \cs{XeTeXuseglyphmetrics}
  \end{syntax}

\end{function}



\begin{function}[]
  {
    \XeTeXversion
  }
  \begin{syntax}
    \cs{XeTeXversion}
  \end{syntax}

\end{function}



\begin{function}[]
  {
    \Z
  }
  \begin{syntax}
    \cs{Z}
  \end{syntax}

\end{function}



\begin{function}[]
  {
    \\
  }
  \begin{syntax}
    \cs{\textbackslash}
  \end{syntax}

\end{function}





\begin{function}[]
  {
    \^
  }
  \begin{syntax}
%    \cs{^}
  \end{syntax}

\end{function}





\begin{function}[]
  {
    \_
  }
  \begin{syntax}
    \cs{\_}
  \end{syntax}

\end{function}



\begin{function}[]
  {
    \`
  }
  \begin{syntax}
    \cs{`}
  \end{syntax}

\end{function}



\begin{function}[]
  {
    \a
  }
  \begin{syntax}
    \cs{a}
  \end{syntax}

\end{function}



\begin{function}[]
  {
    \aa
  }
  \begin{syntax}
    \cs{aa}
  \end{syntax}

\end{function}



\begin{function}[]
  {
    \abovedisplayshortskip
  }
  \begin{syntax}
    \cs{abovedisplayshortskip}
  \end{syntax}

\end{function}



\begin{function}[]
  {
    \abovedisplayskip
  }
  \begin{syntax}
    \cs{abovedisplayskip}
  \end{syntax}

\end{function}



\begin{function}[]
  {
    \accent
  }
  \begin{syntax}
    \cs{accent}
  \end{syntax}

\end{function}



\begin{function}[]
  {
    \accent@spacefactor
  }
  \begin{syntax}
    \cs{accent@spacefactor}
  \end{syntax}

\end{function}





\begin{function}[]
  {
    \active@math@prime
  }
  \begin{syntax}
    \cs{active@math@prime}
  \end{syntax}

\end{function}



\begin{function}[]
  {
    \add@accent
  }
  \begin{syntax}
    \cs{add@accent}
  \end{syntax}

\end{function}



\begin{function}[]
  {
    \add@percent@to@temptokena
  }
  \begin{syntax}
    \cs{add@percent@to@temptokena}
  \end{syntax}

\end{function}



\begin{function}[]
  {
    \addcontentsline
  }
  \begin{syntax}
    \cs{addcontentsline}
  \end{syntax}

\end{function}



\begin{function}[]
  {
    \addpenalty
  }
  \begin{syntax}
    \cs{addpenalty}
  \end{syntax}

\end{function}



\begin{function}[]
  {
    \addto@hook
  }
  \begin{syntax}
    \cs{addto@hook}
  \end{syntax}

\end{function}



\begin{function}[]
  {
    \addtocontents
  }
  \begin{syntax}
    \cs{addtocontents}
  \end{syntax}

\end{function}







\begin{function}[]
  {
    \addvspace
  }
  \begin{syntax}
    \cs{addvspace}
  \end{syntax}

\end{function}



\begin{function}[]
  {
    \adjdemerits
  }
  \begin{syntax}
    \cs{adjdemerits}
  \end{syntax}

\end{function}




\begin{function}[]
  {
    \ae
  }
  \begin{syntax}
    \cs{ae}
  \end{syntax}

\end{function}



\begin{function}[]
  {
    \afterassignment
  }
  \begin{syntax}
    \cs{afterassignment}
  \end{syntax}

\end{function}



\begin{function}[]
  {
    \aftergroup
  }
  \begin{syntax}
    \cs{aftergroup}
  \end{syntax}

\end{function}





\begin{function}[]
  {
    \allowbreak
  }
  \begin{syntax}
    \cs{allowbreak}
  \end{syntax}

\end{function}



\begin{function}[]
  {
    \alph
  }
  \begin{syntax}
    \cs{alph}
  \end{syntax}

\end{function}



\begin{function}[]
  {
    \alpha@elt
  }
  \begin{syntax}
    \cs{alpha@elt}
  \end{syntax}

\end{function}



\begin{function}[]
  {
    \alpha@list
  }
  \begin{syntax}
    \cs{alpha@list}
  \end{syntax}

\end{function}



\begin{function}[]
  {
    \and
  }
  \begin{syntax}
    \cs{and}
  \end{syntax}

\end{function}



\begin{function}[]
  {
    \arabic
  }
  \begin{syntax}
    \cs{arabic}
  \end{syntax}

\end{function}





\begin{function}[]
  {
    \arg
  }
  \begin{syntax}
    \cs{arg}
  \end{syntax}

\end{function}



\begin{function}[]
  {
    \array
  }
  \begin{syntax}
    \cs{array}
  \end{syntax}

\end{function}



\begin{function}[]
  {
    \arraycolsep
  }
  \begin{syntax}
    \cs{arraycolsep}
  \end{syntax}

\end{function}



\begin{function}[]
  {
    \arrayrulewidth
  }
  \begin{syntax}
    \cs{arrayrulewidth}
  \end{syntax}

\end{function}



\begin{function}[]
  {
    \arraystretch
  }
  \begin{syntax}
    \cs{arraystretch}
  \end{syntax}

\end{function}



\begin{function}[]
  {
    \asciispace
  }
  \begin{syntax}
    \cs{asciispace}
  \end{syntax}

\end{function}




\begin{function}[]
  {
    \attribute
  }
  \begin{syntax}
    \cs{attribute}
  \end{syntax}

\end{function}



\begin{function}[]
  {
    \attributedef
  }
  \begin{syntax}
    \cs{attributedef}
  \end{syntax}

\end{function}



\begin{function}[]
  {
    \attributezero
  }
  \begin{syntax}
    \cs{attributezero}
  \end{syntax}

\end{function}



\begin{function}[]
  {
    \author
  }
  \begin{syntax}
    \cs{author}
  \end{syntax}

\end{function}





\begin{function}[]
  {
    \baselineskip
  }
  \begin{syntax}
    \cs{baselineskip}
  \end{syntax}

\end{function}



\begin{function}[]
  {
    \baselinestretch
  }
  \begin{syntax}
    \cs{baselinestretch}
  \end{syntax}

\end{function}



\begin{function}[]
  {
    \batchmode
  }
  \begin{syntax}
    \cs{batchmode}
  \end{syntax}

\end{function}





\begin{function}[]
  {
    \begingroup
  }
  \begin{syntax}
    \cs{begingroup}
  \end{syntax}

\end{function}



\begin{function}[]
  {
    \belowdisplayshortskip
  }
  \begin{syntax}
    \cs{belowdisplayshortskip}
  \end{syntax}

\end{function}



\begin{function}[]
  {
    \belowdisplayskip
  }
  \begin{syntax}
    \cs{belowdisplayskip}
  \end{syntax}

\end{function}



\begin{function}[]
  {
    \best@size
  }
  \begin{syntax}
    \cs{best@size}
  \end{syntax}

\end{function}



\begin{function}[]
  {
    \bezier
  }
  \begin{syntax}
    \cs{bezier}
  \end{syntax}

\end{function}



\begin{function}[]
  {
    \bfdef@ult
  }
  \begin{syntax}
    \cs{bfdef@ult}
  \end{syntax}

\end{function}



\begin{function}[]
  {
    \bfdefault
  }
  \begin{syntax}
    \cs{bfdefault}
  \end{syntax}

\end{function}



\begin{function}[]
  {
    \bfdefault@previous
  }
  \begin{syntax}
    \cs{bfdefault@previous}
  \end{syntax}

\end{function}



\begin{function}[]
  {
    \bfseries
  }
  \begin{syntax}
    \cs{bfseries}
  \end{syntax}

\end{function}



\begin{function}[]
  {
    \bfseries@rm
  }
  \begin{syntax}
    \cs{bfseries@rm}
  \end{syntax}

\end{function}



\begin{function}[]
  {
    \bfseries@rm@kernel
  }
  \begin{syntax}
    \cs{bfseries@rm@kernel}
  \end{syntax}

\end{function}



\begin{function}[]
  {
    \bfseries@sf
  }
  \begin{syntax}
    \cs{bfseries@sf}
  \end{syntax}

\end{function}



\begin{function}[]
  {
    \bfseries@sf@kernel
  }
  \begin{syntax}
    \cs{bfseries@sf@kernel}
  \end{syntax}

\end{function}



\begin{function}[]
  {
    \bfseries@tt
  }
  \begin{syntax}
    \cs{bfseries@tt}
  \end{syntax}

\end{function}



\begin{function}[]
  {
    \bfseries@tt@kernel
  }
  \begin{syntax}
    \cs{bfseries@tt@kernel}
  \end{syntax}

\end{function}



\begin{function}[]
  {
    \bgroup
  }
  \begin{syntax}
    \cs{bgroup}
  \end{syntax}

\end{function}



\begin{function}[]
  {
    \bibcite
  }
  \begin{syntax}
    \cs{bibcite}
  \end{syntax}

\end{function}



\begin{function}[]
  {
    \bibdata
  }
  \begin{syntax}
    \cs{bibdata}
  \end{syntax}

\end{function}



\begin{function}[]
  {
    \bibitem
  }
  \begin{syntax}
    \cs{bibitem}
  \end{syntax}

\end{function}



\begin{function}[]
  {
    \bibliography
  }
  \begin{syntax}
    \cs{bibliography}
  \end{syntax}

\end{function}



\begin{function}[]
  {
    \bibliographystyle
  }
  \begin{syntax}
    \cs{bibliographystyle}
  \end{syntax}

\end{function}



\begin{function}[]
  {
    \bibstyle
  }
  \begin{syntax}
    \cs{bibstyle}
  \end{syntax}

\end{function}




\begin{function}[]
  {
    \bigbreak
  }
  \begin{syntax}
    \cs{bigbreak}
  \end{syntax}

\end{function}








\begin{function}[]
  {
    \binoppenalty
  }
  \begin{syntax}
    \cs{binoppenalty}
  \end{syntax}

\end{function}



\begin{function}[]
  {
    \bm@b
  }
  \begin{syntax}
    \cs{bm@b}
  \end{syntax}

\end{function}



\begin{function}[]
  {
    \bm@c
  }
  \begin{syntax}
    \cs{bm@c}
  \end{syntax}

\end{function}



\begin{function}[]
  {
    \bm@l
  }
  \begin{syntax}
    \cs{bm@l}
  \end{syntax}

\end{function}



\begin{function}[]
  {
    \bm@r
  }
  \begin{syntax}
    \cs{bm@r}
  \end{syntax}

\end{function}



\begin{function}[]
  {
    \bm@s
  }
  \begin{syntax}
    \cs{bm@s}
  \end{syntax}

\end{function}



\begin{function}[]
  {
    \bm@t
  }
  \begin{syntax}
    \cs{bm@t}
  \end{syntax}

\end{function}



\begin{function}[]
  {
    \bmod
  }
  \begin{syntax}
    \cs{bmod}
  \end{syntax}

\end{function}



\begin{function}[]
  {
    \boldmath
  }
  \begin{syntax}
    \cs{boldmath}
  \end{syntax}

\end{function}



\begin{function}[]
  {
    \bordermatrix
  }
  \begin{syntax}
    \cs{bordermatrix}
  \end{syntax}

\end{function}



\begin{function}[]
  {
    \botfigrule
  }
  \begin{syntax}
    \cs{botfigrule}
  \end{syntax}

\end{function}



\begin{function}[]
  {
    \bottomfraction
  }
  \begin{syntax}
    \cs{bottomfraction}
  \end{syntax}

\end{function}



\begin{function}[]
  {
    \box
  }
  \begin{syntax}
    \cs{box}
  \end{syntax}

\end{function}



\begin{function}[]
  {
    \boxmaxdepth
  }
  \begin{syntax}
    \cs{boxmaxdepth}
  \end{syntax}

\end{function}



\begin{function}[]
  {
    \brace
  }
  \begin{syntax}
    \cs{brace}
  \end{syntax}

\end{function}



\begin{function}[]
  {
    \brack
  }
  \begin{syntax}
    \cs{brack}
  \end{syntax}

\end{function}



\begin{function}[]
  {
    \break
  }
  \begin{syntax}
    \cs{break}
  \end{syntax}

\end{function}



\begin{function}[]
  {
    \brokenpenalty
  }
  \begin{syntax}
    \cs{brokenpenalty}
  \end{syntax}

\end{function}



\begin{function}[]
  {
    \buildrel
  }
  \begin{syntax}
    \cs{buildrel}
  \end{syntax}

\end{function}



\begin{function}[]
  {
    \bx@A
  }
  \begin{syntax}
    \cs{bx@A}
  \end{syntax}

\end{function}



\begin{function}[]
  {
    \bx@AA
  }
  \begin{syntax}
    \cs{bx@AA}
  \end{syntax}

\end{function}



\begin{function}[]
  {
    \bx@B
  }
  \begin{syntax}
    \cs{bx@B}
  \end{syntax}

\end{function}



\begin{function}[]
  {
    \bx@BB
  }
  \begin{syntax}
    \cs{bx@BB}
  \end{syntax}

\end{function}



\begin{function}[]
  {
    \bx@C
  }
  \begin{syntax}
    \cs{bx@C}
  \end{syntax}

\end{function}



\begin{function}[]
  {
    \bx@CC
  }
  \begin{syntax}
    \cs{bx@CC}
  \end{syntax}

\end{function}



\begin{function}[]
  {
    \bx@D
  }
  \begin{syntax}
    \cs{bx@D}
  \end{syntax}

\end{function}



\begin{function}[]
  {
    \bx@DD
  }
  \begin{syntax}
    \cs{bx@DD}
  \end{syntax}

\end{function}



\begin{function}[]
  {
    \bx@E
  }
  \begin{syntax}
    \cs{bx@E}
  \end{syntax}

\end{function}



\begin{function}[]
  {
    \bx@EE
  }
  \begin{syntax}
    \cs{bx@EE}
  \end{syntax}

\end{function}



\begin{function}[]
  {
    \bx@F
  }
  \begin{syntax}
    \cs{bx@F}
  \end{syntax}

\end{function}



\begin{function}[]
  {
    \bx@FF
  }
  \begin{syntax}
    \cs{bx@FF}
  \end{syntax}

\end{function}



\begin{function}[]
  {
    \bx@G
  }
  \begin{syntax}
    \cs{bx@G}
  \end{syntax}

\end{function}



\begin{function}[]
  {
    \bx@GG
  }
  \begin{syntax}
    \cs{bx@GG}
  \end{syntax}

\end{function}



\begin{function}[]
  {
    \bx@H
  }
  \begin{syntax}
    \cs{bx@H}
  \end{syntax}

\end{function}



\begin{function}[]
  {
    \bx@HH
  }
  \begin{syntax}
    \cs{bx@HH}
  \end{syntax}

\end{function}



\begin{function}[]
  {
    \bx@I
  }
  \begin{syntax}
    \cs{bx@I}
  \end{syntax}

\end{function}



\begin{function}[]
  {
    \bx@II
  }
  \begin{syntax}
    \cs{bx@II}
  \end{syntax}

\end{function}



\begin{function}[]
  {
    \bx@J
  }
  \begin{syntax}
    \cs{bx@J}
  \end{syntax}

\end{function}



\begin{function}[]
  {
    \bx@JJ
  }
  \begin{syntax}
    \cs{bx@JJ}
  \end{syntax}

\end{function}



\begin{function}[]
  {
    \bx@K
  }
  \begin{syntax}
    \cs{bx@K}
  \end{syntax}

\end{function}



\begin{function}[]
  {
    \bx@KK
  }
  \begin{syntax}
    \cs{bx@KK}
  \end{syntax}

\end{function}



\begin{function}[]
  {
    \bx@L
  }
  \begin{syntax}
    \cs{bx@L}
  \end{syntax}

\end{function}



\begin{function}[]
  {
    \bx@LL
  }
  \begin{syntax}
    \cs{bx@LL}
  \end{syntax}

\end{function}



\begin{function}[]
  {
    \bx@M
  }
  \begin{syntax}
    \cs{bx@M}
  \end{syntax}

\end{function}



\begin{function}[]
  {
    \bx@MM
  }
  \begin{syntax}
    \cs{bx@MM}
  \end{syntax}

\end{function}



\begin{function}[]
  {
    \bx@N
  }
  \begin{syntax}
    \cs{bx@N}
  \end{syntax}

\end{function}



\begin{function}[]
  {
    \bx@NN
  }
  \begin{syntax}
    \cs{bx@NN}
  \end{syntax}

\end{function}



\begin{function}[]
  {
    \bx@O
  }
  \begin{syntax}
    \cs{bx@O}
  \end{syntax}

\end{function}



\begin{function}[]
  {
    \bx@OO
  }
  \begin{syntax}
    \cs{bx@OO}
  \end{syntax}

\end{function}



\begin{function}[]
  {
    \bx@P
  }
  \begin{syntax}
    \cs{bx@P}
  \end{syntax}

\end{function}



\begin{function}[]
  {
    \bx@PP
  }
  \begin{syntax}
    \cs{bx@PP}
  \end{syntax}

\end{function}



\begin{function}[]
  {
    \bx@Q
  }
  \begin{syntax}
    \cs{bx@Q}
  \end{syntax}

\end{function}



\begin{function}[]
  {
    \bx@QQ
  }
  \begin{syntax}
    \cs{bx@QQ}
  \end{syntax}

\end{function}



\begin{function}[]
  {
    \bx@R
  }
  \begin{syntax}
    \cs{bx@R}
  \end{syntax}

\end{function}



\begin{function}[]
  {
    \bx@RR
  }
  \begin{syntax}
    \cs{bx@RR}
  \end{syntax}

\end{function}



\begin{function}[]
  {
    \bx@S
  }
  \begin{syntax}
    \cs{bx@S}
  \end{syntax}

\end{function}



\begin{function}[]
  {
    \bx@SS
  }
  \begin{syntax}
    \cs{bx@SS}
  \end{syntax}

\end{function}



\begin{function}[]
  {
    \bx@T
  }
  \begin{syntax}
    \cs{bx@T}
  \end{syntax}

\end{function}



\begin{function}[]
  {
    \bx@TT
  }
  \begin{syntax}
    \cs{bx@TT}
  \end{syntax}

\end{function}



\begin{function}[]
  {
    \bx@U
  }
  \begin{syntax}
    \cs{bx@U}
  \end{syntax}

\end{function}



\begin{function}[]
  {
    \bx@UU
  }
  \begin{syntax}
    \cs{bx@UU}
  \end{syntax}

\end{function}



\begin{function}[]
  {
    \bx@V
  }
  \begin{syntax}
    \cs{bx@V}
  \end{syntax}

\end{function}



\begin{function}[]
  {
    \bx@VV
  }
  \begin{syntax}
    \cs{bx@VV}
  \end{syntax}

\end{function}



\begin{function}[]
  {
    \bx@W
  }
  \begin{syntax}
    \cs{bx@W}
  \end{syntax}

\end{function}



\begin{function}[]
  {
    \bx@WW
  }
  \begin{syntax}
    \cs{bx@WW}
  \end{syntax}

\end{function}



\begin{function}[]
  {
    \bx@X
  }
  \begin{syntax}
    \cs{bx@X}
  \end{syntax}

\end{function}



\begin{function}[]
  {
    \bx@XX
  }
  \begin{syntax}
    \cs{bx@XX}
  \end{syntax}

\end{function}



\begin{function}[]
  {
    \bx@Y
  }
  \begin{syntax}
    \cs{bx@Y}
  \end{syntax}

\end{function}



\begin{function}[]
  {
    \bx@YY
  }
  \begin{syntax}
    \cs{bx@YY}
  \end{syntax}

\end{function}



\begin{function}[]
  {
    \bx@Z
  }
  \begin{syntax}
    \cs{bx@Z}
  \end{syntax}

\end{function}



\begin{function}[]
  {
    \bx@ZZ
  }
  \begin{syntax}
    \cs{bx@ZZ}
  \end{syntax}

\end{function}







\begin{function}[]
  {
    \c@dbltopnumber
  }
  \begin{syntax}
    \cs{c@dbltopnumber}
  \end{syntax}

\end{function}



\begin{function}[]
  {
    \c@equation
  }
  \begin{syntax}
    \cs{c@equation}
  \end{syntax}

\end{function}



\begin{function}[]
  {
    \c@errorcontextlines
  }
  \begin{syntax}
    \cs{c@errorcontextlines}
  \end{syntax}

\end{function}



\begin{function}[]
  {
    \c@footnote
  }
  \begin{syntax}
    \cs{c@footnote}
  \end{syntax}

\end{function}



\begin{function}[]
  {
    \c@localmathalphabets
  }
  \begin{syntax}
    \cs{c@localmathalphabets}
  \end{syntax}

\end{function}



\begin{function}[]
  {
    \c@mpfootnote
  }
  \begin{syntax}
    \cs{c@mpfootnote}
  \end{syntax}

\end{function}



\begin{function}[]
  {
    \c@page
  }
  \begin{syntax}
    \cs{c@page}
  \end{syntax}

\end{function}



\begin{function}[]
  {
    \c@secnumdepth
  }
  \begin{syntax}
    \cs{c@secnumdepth}
  \end{syntax}

\end{function}



\begin{function}[]
  {
    \c@tocdepth
  }
  \begin{syntax}
    \cs{c@tocdepth}
  \end{syntax}

\end{function}



\begin{function}[]
  {
    \c@topnumber
  }
  \begin{syntax}
    \cs{c@topnumber}
  \end{syntax}

\end{function}



\begin{function}[]
  {
    \c@totalnumber
  }
  \begin{syntax}
    \cs{c@totalnumber}
  \end{syntax}

\end{function}



\begin{function}[]
  {
    \c@totalpages
  }
  \begin{syntax}
    \cs{c@totalpages}
  \end{syntax}

\end{function}



\begin{function}[]
  {
    \calculate@math@sizes
  }
  \begin{syntax}
    \cs{calculate@math@sizes}
  \end{syntax}

\end{function}



\begin{function}[]
  {
    \capitalacute
  }
  \begin{syntax}
    \cs{capitalacute}
  \end{syntax}

\end{function}



\begin{function}[]
  {
    \capitalbreve
  }
  \begin{syntax}
    \cs{capitalbreve}
  \end{syntax}

\end{function}



\begin{function}[]
  {
    \capitalcaron
  }
  \begin{syntax}
    \cs{capitalcaron}
  \end{syntax}

\end{function}



\begin{function}[]
  {
    \capitalcedilla
  }
  \begin{syntax}
    \cs{capitalcedilla}
  \end{syntax}

\end{function}



\begin{function}[]
  {
    \capitalcircumflex
  }
  \begin{syntax}
    \cs{capitalcircumflex}
  \end{syntax}

\end{function}



\begin{function}[]
  {
    \capitaldieresis
  }
  \begin{syntax}
    \cs{capitaldieresis}
  \end{syntax}

\end{function}



\begin{function}[]
  {
    \capitaldotaccent
  }
  \begin{syntax}
    \cs{capitaldotaccent}
  \end{syntax}

\end{function}



\begin{function}[]
  {
    \capitalgrave
  }
  \begin{syntax}
    \cs{capitalgrave}
  \end{syntax}

\end{function}



\begin{function}[]
  {
    \capitalhungarumlaut
  }
  \begin{syntax}
    \cs{capitalhungarumlaut}
  \end{syntax}

\end{function}



\begin{function}[]
  {
    \capitalmacron
  }
  \begin{syntax}
    \cs{capitalmacron}
  \end{syntax}

\end{function}



\begin{function}[]
  {
    \capitalnewtie
  }
  \begin{syntax}
    \cs{capitalnewtie}
  \end{syntax}

\end{function}



\begin{function}[]
  {
    \capitalogonek
  }
  \begin{syntax}
    \cs{capitalogonek}
  \end{syntax}

\end{function}



\begin{function}[]
  {
    \capitalring
  }
  \begin{syntax}
    \cs{capitalring}
  \end{syntax}

\end{function}



\begin{function}[]
  {
    \capitaltie
  }
  \begin{syntax}
    \cs{capitaltie}
  \end{syntax}

\end{function}



\begin{function}[]
  {
    \capitaltilde
  }
  \begin{syntax}
    \cs{capitaltilde}
  \end{syntax}

\end{function}



\begin{function}[]
  {
    \caption
  }
  \begin{syntax}
    \cs{caption}
  \end{syntax}

\end{function}



\begin{function}[]
  {
    \cases
  }
  \begin{syntax}
    \cs{cases}
  \end{syntax}

\end{function}





\begin{function}[]
  {
    \catcodetable
  }
  \begin{syntax}
    \cs{catcodetable}
  \end{syntax}

\end{function}



\begin{function}[]
  {
    \catcodetable@atletter
  }
  \begin{syntax}
    \cs{catcodetable@atletter}
  \end{syntax}

\end{function}



\begin{function}[]
  {
    \catcodetable@initex
  }
  \begin{syntax}
    \cs{catcodetable@initex}
  \end{syntax}

\end{function}



\begin{function}[]
  {
    \catcodetable@latex
  }
  \begin{syntax}
    \cs{catcodetable@latex}
  \end{syntax}

\end{function}



\begin{function}[]
  {
    \catcodetable@string
  }
  \begin{syntax}
    \cs{catcodetable@string}
  \end{syntax}

\end{function}



\begin{function}[]
  {
    \cdp@elt
  }
  \begin{syntax}
    \cs{cdp@elt}
  \end{syntax}

\end{function}



\begin{function}[]
  {
    \cdp@list
  }
  \begin{syntax}
    \cs{cdp@list}
  \end{syntax}

\end{function}



\begin{function}[]
  {
    \center
  }
  \begin{syntax}
    \cs{center}
  \end{syntax}

\end{function}



\begin{function}[]
  {
    \centering
  }
  \begin{syntax}
    \cs{centering}
  \end{syntax}

\end{function}





\begin{function}[]
  {
    \cf@encoding
  }
  \begin{syntax}
    \cs{cf@encoding}
  \end{syntax}

\end{function}



\begin{function}[]
  {
    \ch@ck
  }
  \begin{syntax}
    \cs{ch@ck}
  \end{syntax}

\end{function}



\begin{function}[]
  {
    \char
  }
  \begin{syntax}
    \cs{char}
  \end{syntax}

\end{function}



\begin{function}[]
  {
    \char@if@alph
  }
  \begin{syntax}
    \cs{char@if@alph}
  \end{syntax}

\end{function}



\begin{function}[]
  {
    \char`
  }
  \begin{syntax}
    \cs{char`}
  \end{syntax}

\end{function}





\begin{function}[]
  {
    \chardef@text@cmd
  }
  \begin{syntax}
    \cs{chardef@text@cmd}
  \end{syntax}

\end{function}



\begin{function}[]
  {
    \charsubdef
  }
  \begin{syntax}
    \cs{charsubdef}
  \end{syntax}

\end{function}



\begin{function}[]
  {
    \charzero
  }
  \begin{syntax}
    \cs{charzero}
  \end{syntax}

\end{function}



\begin{function}[]
  {
    \check@command
  }
  \begin{syntax}
    \cs{check@command}
  \end{syntax}

\end{function}



\begin{function}[]
  {
    \check@icl
  }
  \begin{syntax}
    \cs{check@icl}
  \end{syntax}

\end{function}



\begin{function}[]
  {
    \check@icr
  }
  \begin{syntax}
    \cs{check@icr}
  \end{syntax}

\end{function}



\begin{function}[]
  {
    \check@mathfonts
  }
  \begin{syntax}
    \cs{check@mathfonts}
  \end{syntax}

\end{function}



\begin{function}[]
  {
    \check@nocorr@
  }
  \begin{syntax}
    \cs{check@nocorr@}
  \end{syntax}

\end{function}



\begin{function}[]
  {
    \check@range
  }
  \begin{syntax}
    \cs{check@range}
  \end{syntax}

\end{function}



\begin{function}[]
  {
    \check@single
  }
  \begin{syntax}
    \cs{check@single}
  \end{syntax}

\end{function}



\begin{function}[]
  {
    \choose
  }
  \begin{syntax}
    \cs{choose}
  \end{syntax}

\end{function}



\begin{function}[]
  {
    \circle
  }
  \begin{syntax}
    \cs{circle}
  \end{syntax}

\end{function}



\begin{function}[]
  {
    \citation
  }
  \begin{syntax}
    \cs{citation}
  \end{syntax}

\end{function}



\begin{function}[]
  {
    \cite
  }
  \begin{syntax}
    \cs{cite}
  \end{syntax}

\end{function}



\begin{function}[]
  {
    \cl@@ckpt
  }
  \begin{syntax}
    \cs{cl@@ckpt}
  \end{syntax}

\end{function}



\begin{function}[]
  {
    \cl@page
  }
  \begin{syntax}
    \cs{cl@page}
  \end{syntax}

\end{function}





\begin{function}[]
  {
    \cleaders
  }
  \begin{syntax}
    \cs{cleaders}
  \end{syntax}

\end{function}



\begin{function}[]
  {
    \cleardoublepage
  }
  \begin{syntax}
    \cs{cleardoublepage}
  \end{syntax}

\end{function}



\begin{function}[]
  {
    \clearpage
  }
  \begin{syntax}
    \cs{clearpage}
  \end{syntax}

\end{function}



\begin{function}[]
  {
    \cline
  }
  \begin{syntax}
    \cs{cline}
  \end{syntax}

\end{function}



\begin{function}[]
  {
    \closein
  }
  \begin{syntax}
    \cs{closein}
  \end{syntax}

\end{function}



\begin{function}[]
  {
    \closeout
  }
  \begin{syntax}
    \cs{closeout}
  \end{syntax}

\end{function}



\begin{function}[]
  {
    \clubpenalty
  }
  \begin{syntax}
    \cs{clubpenalty}
  \end{syntax}

\end{function}



\begin{function}[]
  {
    \col@number
  }
  \begin{syntax}
    \cs{col@number}
  \end{syntax}

\end{function}



\begin{function}[]
  {
    \color@begingroup
  }
  \begin{syntax}
    \cs{color@begingroup}
  \end{syntax}

\end{function}



\begin{function}[]
  {
    \color@endbox
  }
  \begin{syntax}
    \cs{color@endbox}
  \end{syntax}

\end{function}



\begin{function}[]
  {
    \color@endgroup
  }
  \begin{syntax}
    \cs{color@endgroup}
  \end{syntax}

\end{function}



\begin{function}[]
  {
    \color@hbox
  }
  \begin{syntax}
    \cs{color@hbox}
  \end{syntax}

\end{function}



\begin{function}[]
  {
    \color@setgroup
  }
  \begin{syntax}
    \cs{color@setgroup}
  \end{syntax}

\end{function}



\begin{function}[]
  {
    \color@vbox
  }
  \begin{syntax}
    \cs{color@vbox}
  \end{syntax}

\end{function}



\begin{function}[]
  {
    \columnsep
  }
  \begin{syntax}
    \cs{columnsep}
  \end{syntax}

\end{function}



\begin{function}[]
  {
    \columnseprule
  }
  \begin{syntax}
    \cs{columnseprule}
  \end{syntax}

\end{function}



\begin{function}[]
  {
    \columnwidth
  }
  \begin{syntax}
    \cs{columnwidth}
  \end{syntax}

\end{function}



\begin{function}[]
  {
    \conditionally@traceoff
  }
  \begin{syntax}
    \cs{conditionally@traceoff}
  \end{syntax}

\end{function}



\begin{function}[]
  {
    \conditionally@traceon
  }
  \begin{syntax}
    \cs{conditionally@traceon}
  \end{syntax}

\end{function}



\begin{function}[]
  {
    \contentsline
  }
  \begin{syntax}
    \cs{contentsline}
  \end{syntax}

\end{function}



\begin{function}[]
  {
    \copy
  }
  \begin{syntax}
    \cs{copy}
  \end{syntax}

\end{function}



\begin{function}[]
  {
    \copy@kernel@robust@command
  }
  \begin{syntax}
    \cs{copy@kernel@robust@command}
  \end{syntax}

\end{function}



\begin{function}[]
  {
    \copyright
  }
  \begin{syntax}
    \cs{copyright}
  \end{syntax}

\end{function}





\begin{function}[]
  {
    \count
  }
  \begin{syntax}
    \cs{count}
  \end{syntax}

\end{function}



\begin{function}[]
  {
    \count@
  }
  \begin{syntax}
    \cs{count@}
  \end{syntax}

\end{function}





\begin{function}[]
  {
    \counterwithin
  }
  \begin{syntax}
    \cs{counterwithin}
  \end{syntax}

\end{function}



\begin{function}[]
  {
    \counterwithout
  }
  \begin{syntax}
    \cs{counterwithout}
  \end{syntax}

\end{function}



\begin{function}[]
  {
    \cr
  }
  \begin{syntax}
    \cs{cr}
  \end{syntax}

\end{function}



\begin{function}[]
  {
    \crcr
  }
  \begin{syntax}
    \cs{crcr}
  \end{syntax}

\end{function}






\begin{function}[]
  {
    \csname
  }
  \begin{syntax}
    \cs{csname}
  \end{syntax}

\end{function}



\begin{function}[]
  {
    \csstring
  }
  \begin{syntax}
    \cs{csstring}
  \end{syntax}

\end{function}



\begin{function}[]
  {
    \curr@fontshape
  }
  \begin{syntax}
    \cs{curr@fontshape}
  \end{syntax}

\end{function}



\begin{function}[]
  {
    \curr@math@size
  }
  \begin{syntax}
    \cs{curr@math@size}
  \end{syntax}

\end{function}



\begin{function}[]
  {
    \currentgrouplevel
  }
  \begin{syntax}
    \cs{currentgrouplevel}
  \end{syntax}

\end{function}



\begin{function}[]
  {
    \currentgrouptype
  }
  \begin{syntax}
    \cs{currentgrouptype}
  \end{syntax}

\end{function}



\begin{function}[]
  {
    \d
  }
  \begin{syntax}
    \cs{d}
  \end{syntax}

\end{function}



\begin{function}[]
  {
    \dag
  }
  \begin{syntax}
    \cs{dag}
  \end{syntax}

\end{function}



\begin{function}[]
  {
    \dagger
  }
  \begin{syntax}
    \cs{dagger}
  \end{syntax}

\end{function}



\begin{function}[]
  {
    \dashbox
  }
  \begin{syntax}
    \cs{dashbox}
  \end{syntax}

\end{function}



\begin{function}[]
  {
    \date
  }
  \begin{syntax}
    \cs{date}
  \end{syntax}

\end{function}



\begin{function}[]
  {
    \day
  }
  \begin{syntax}
    \cs{day}
  \end{syntax}

\end{function}



\begin{function}[]
  {
    \dblfigrule
  }
  \begin{syntax}
    \cs{dblfigrule}
  \end{syntax}

\end{function}



\begin{function}[]
  {
    \dblfloatpagefraction
  }
  \begin{syntax}
    \cs{dblfloatpagefraction}
  \end{syntax}

\end{function}



\begin{function}[]
  {
    \dblfloatsep
  }
  \begin{syntax}
    \cs{dblfloatsep}
  \end{syntax}

\end{function}



\begin{function}[]
  {
    \dbltextfloatsep
  }
  \begin{syntax}
    \cs{dbltextfloatsep}
  \end{syntax}

\end{function}



\begin{function}[]
  {
    \dbltopfraction
  }
  \begin{syntax}
    \cs{dbltopfraction}
  \end{syntax}

\end{function}



\begin{function}[]
  {
    \ddag
  }
  \begin{syntax}
    \cs{ddag}
  \end{syntax}

\end{function}



\begin{function}[]
  {
    \ddagger
  }
  \begin{syntax}
    \cs{ddagger}
  \end{syntax}

\end{function}



\begin{function}[]
  {
    \deadcycles
  }
  \begin{syntax}
    \cs{deadcycles}
  \end{syntax}

\end{function}



\begin{function}[]
  {
    \declare@commandcopy
  }
  \begin{syntax}
    \cs{declare@commandcopy}
  \end{syntax}

\end{function}



\begin{function}[]
  {
    \declare@commandcopy@do
  }
  \begin{syntax}
    \cs{declare@commandcopy@do}
  \end{syntax}

\end{function}



\begin{function}[]
  {
    \declare@commandcopy@let
  }
  \begin{syntax}
    \cs{declare@commandcopy@let}
  \end{syntax}

\end{function}



\begin{function}[]
  {
    \declare@environmentcopy
  }
  \begin{syntax}
    \cs{declare@environmentcopy}
  \end{syntax}

\end{function}



\begin{function}[]
  {
    \declare@file@substitution
  }
  \begin{syntax}
    \cs{declare@file@substitution}
  \end{syntax}

\end{function}



\begin{function}[]
  {
    \declare@robustcommand
  }
  \begin{syntax}
    \cs{declare@robustcommand}
  \end{syntax}

\end{function}



\begin{function}[]
  {
    \declare@robustcommand@auxi
  }
  \begin{syntax}
    \cs{declare@robustcommand@auxi}
  \end{syntax}

\end{function}



\begin{function}[]
  {
    \declare@robustcommand@auxii
  }
  \begin{syntax}
    \cs{declare@robustcommand@auxii}
  \end{syntax}

\end{function}



\begin{function}[]
  {
    \declare@robustcommand@auxiii
  }
  \begin{syntax}
    \cs{declare@robustcommand@auxiii}
  \end{syntax}

\end{function}



\begin{function}[]
  {
    \def
  }
  \begin{syntax}
    \cs{def}
  \end{syntax}

\end{function}



\begin{function}[]
  {
    \default@M
  }
  \begin{syntax}
    \cs{default@M}
  \end{syntax}

\end{function}



\begin{function}[]
  {
    \default@T
  }
  \begin{syntax}
    \cs{default@T}
  \end{syntax}

\end{function}



\begin{function}[]
  {
    \default@ds
  }
  \begin{syntax}
    \cs{default@ds}
  \end{syntax}

\end{function}



\begin{function}[]
  {
    \default@family
  }
  \begin{syntax}
    \cs{default@family}
  \end{syntax}

\end{function}



\begin{function}[]
  {
    \default@series
  }
  \begin{syntax}
    \cs{default@series}
  \end{syntax}

\end{function}



\begin{function}[]
  {
    \default@shape
  }
  \begin{syntax}
    \cs{default@shape}
  \end{syntax}

\end{function}



\begin{function}[]
  {
    \defaulthyphenchar
  }
  \begin{syntax}
    \cs{defaulthyphenchar}
  \end{syntax}

\end{function}



\begin{function}[]
  {
    \defaultscriptratio
  }
  \begin{syntax}
    \cs{defaultscriptratio}
  \end{syntax}

\end{function}



\begin{function}[]
  {
    \defaultscriptscriptratio
  }
  \begin{syntax}
    \cs{defaultscriptscriptratio}
  \end{syntax}

\end{function}



\begin{function}[]
  {
    \defaultskewchar
  }
  \begin{syntax}
    \cs{defaultskewchar}
  \end{syntax}

\end{function}



\begin{function}[]
  {
    \define@newfont
  }
  \begin{syntax}
    \cs{define@newfont}
  \end{syntax}

\end{function}



\begin{function}[]
  {
    \deg
  }
  \begin{syntax}
    \cs{deg}
  \end{syntax}

\end{function}



\begin{function}[]
  {
    \delayed@f@adjustment
  }
  \begin{syntax}
    \cs{delayed@f@adjustment}
  \end{syntax}

\end{function}



\begin{function}[]
  {
    \delayed@merge@font@series
  }
  \begin{syntax}
    \cs{delayed@merge@font@series}
  \end{syntax}

\end{function}



\begin{function}[]
  {
    \delayed@merge@font@shape
  }
  \begin{syntax}
    \cs{delayed@merge@font@shape}
  \end{syntax}

\end{function}



\begin{function}[]
  {
    \delcode
  }
  \begin{syntax}
    \cs{delcode}
  \end{syntax}

\end{function}



\begin{function}[]
  {
    \delimiter
  }
  \begin{syntax}
    \cs{delimiter}
  \end{syntax}

\end{function}



\begin{function}[]
  {
    \delimiterfactor
  }
  \begin{syntax}
    \cs{delimiterfactor}
  \end{syntax}

\end{function}



\begin{function}[]
  {
    \delimitershortfall
  }
  \begin{syntax}
    \cs{delimitershortfall}
  \end{syntax}

\end{function}



\begin{function}[]
  {
    \depth
  }
  \begin{syntax}
    \cs{depth}
  \end{syntax}

\end{function}



\begin{function}[]
  {
    \det
  }
  \begin{syntax}
    \cs{det}
  \end{syntax}

\end{function}



\begin{function}[]
  {
    \detokenize
  }
  \begin{syntax}
    \cs{detokenize}
  \end{syntax}

\end{function}





\begin{function}[]
  {
    \dh
  }
  \begin{syntax}
    \cs{dh}
  \end{syntax}

\end{function}



\begin{function}[]
  {
    \dim
  }
  \begin{syntax}
    \cs{dim}
  \end{syntax}

\end{function}



\begin{function}[]
  {
    \dimen
  }
  \begin{syntax}
    \cs{dimen}
  \end{syntax}

\end{function}



\begin{function}[]
  {
    \dimen0
  }
  \begin{syntax}
    \cs{dimen0}
  \end{syntax}

\end{function}



\begin{function}[]
  {
    \dimen@
  }
  \begin{syntax}
    \cs{dimen@}
  \end{syntax}

\end{function}



\begin{function}[]
  {
    \dimen@i
  }
  \begin{syntax}
    \cs{dimen@i}
  \end{syntax}

\end{function}



\begin{function}[]
  {
    \dimen@ii
  }
  \begin{syntax}
    \cs{dimen@ii}
  \end{syntax}

\end{function}





\begin{function}[]
  {
    \dimenzero
  }
  \begin{syntax}
    \cs{dimenzero}
  \end{syntax}

\end{function}



\begin{function}[]
  {
    \dimeval
  }
  \begin{syntax}
    \cs{dimeval}
  \end{syntax}

\end{function}



\begin{function}[]
  {
    \dimexpr
  }
  \begin{syntax}
    \cs{dimexpr}
  \end{syntax}

\end{function}



\begin{function}[]
  {
    \directlua
  }
  \begin{syntax}
    \cs{directlua}
  \end{syntax}

\end{function}



\begin{function}[]
  {
    \disable@package@load
  }
  \begin{syntax}
    \cs{disable@package@load}
  \end{syntax}

\end{function}



\begin{function}[]
  {
    \discretionary
  }
  \begin{syntax}
    \cs{discretionary}
  \end{syntax}

\end{function}



\begin{function}[]
  {
    \displ@y
  }
  \begin{syntax}
    \cs{displ@y}
  \end{syntax}

\end{function}



\begin{function}[]
  {
    \displaylines
  }
  \begin{syntax}
    \cs{displaylines}
  \end{syntax}

\end{function}





\begin{function}[]
  {
    \displaystyle
  }
  \begin{syntax}
    \cs{displaystyle}
  \end{syntax}

\end{function}



\begin{function}[]
  {
    \displaywidowpenalty
  }
  \begin{syntax}
    \cs{displaywidowpenalty}
  \end{syntax}

\end{function}



\begin{function}[]
  {
    \displaywidth
  }
  \begin{syntax}
    \cs{displaywidth}
  \end{syntax}

\end{function}





\begin{function}[]
  {
    \dj
  }
  \begin{syntax}
    \cs{dj}
  \end{syntax}

\end{function}



\begin{function}[]
  {
    \do
  }
  \begin{syntax}
    \cs{do}
  \end{syntax}

\end{function}



\begin{function}[]
  {
    \do@add@percent@to@temptokena
  }
  \begin{syntax}
    \cs{do@add@percent@to@temptokena}
  \end{syntax}

\end{function}



\begin{function}[]
  {
    \do@emfont@update
  }
  \begin{syntax}
    \cs{do@emfont@update}
  \end{syntax}

\end{function}



\begin{function}[]
  {
    \do@noligs
  }
  \begin{syntax}
    \cs{do@noligs}
  \end{syntax}

\end{function}



\begin{function}[]
  {
    \do@subst@correction
  }
  \begin{syntax}
    \cs{do@subst@correction}
  \end{syntax}

\end{function}






\begin{function}[]
  {
    \document@default@language
  }
  \begin{syntax}
    \cs{document@default@language}
  \end{syntax}

\end{function}



\begin{function}[]
  {
    \document@select@group
  }
  \begin{syntax}
    \cs{document@select@group}
  \end{syntax}

\end{function}



\begin{function}[]
  {
    \documentclass
  }
  \begin{syntax}
    \cs{documentclass}
  \end{syntax}

\end{function}



\begin{function}[deprecated]
  {
    \documentstyle
  }
  \begin{syntax}
    \cs{documentstyle}
  \end{syntax}

\end{function}



\begin{function}[]
  {
    \dollardollar@begin
  }
  \begin{syntax}
    \cs{dollardollar@begin}
  \end{syntax}

\end{function}



\begin{function}[]
  {
    \dollardollar@end
  }
  \begin{syntax}
    \cs{dollardollar@end}
  \end{syntax}

\end{function}



\begin{function}[]
  {
    \dont@add@percent@to@temptokena
  }
  \begin{syntax}
    \cs{dont@add@percent@to@temptokena}
  \end{syntax}

\end{function}



\begin{function}[]
  {
    \dorestore@version
  }
  \begin{syntax}
    \cs{dorestore@version}
  \end{syntax}

\end{function}





\begin{function}[]
  {
    \dotfill
  }
  \begin{syntax}
    \cs{dotfill}
  \end{syntax}

\end{function}



\begin{function}[]
  {
    \dots
  }
  \begin{syntax}
    \cs{dots}
  \end{syntax}

\end{function}



\begin{function}[]
  {
    \doublehyphendemerits
  }
  \begin{syntax}
    \cs{doublehyphendemerits}
  \end{syntax}

\end{function}



\begin{function}[]
  {
    \doublerulesep
  }
  \begin{syntax}
    \cs{doublerulesep}
  \end{syntax}

\end{function}



\begin{function}[]
  {
    \dp
  }
  \begin{syntax}
    \cs{dp}
  \end{syntax}

\end{function}



\begin{function}[]
  {
    \ds@
  }
  \begin{syntax}
    \cs{ds@}
  \end{syntax}

\end{function}



\begin{function}[]
  {
    \dt@pfalse
  }
  \begin{syntax}
    \cs{dt@pfalse}
  \end{syntax}

\end{function}



\begin{function}[]
  {
    \dt@ptrue
  }
  \begin{syntax}
    \cs{dt@ptrue}
  \end{syntax}

\end{function}



\begin{function}[]
  {
    \dump
  }
  \begin{syntax}
    \cs{dump}
  \end{syntax}

\end{function}



\begin{function}[]
  {
    \e@alloc@attribute@count
  }
  \begin{syntax}
    \cs{e@alloc@attribute@count}
  \end{syntax}

\end{function}



\begin{function}[]
  {
    \e@alloc@bytecode@count
  }
  \begin{syntax}
    \cs{e@alloc@bytecode@count}
  \end{syntax}

\end{function}



\begin{function}[]
  {
    \e@alloc@ccodetable@count
  }
  \begin{syntax}
    \cs{e@alloc@ccodetable@count}
  \end{syntax}

\end{function}



\begin{function}[]
  {
    \e@alloc@chardef
  }
  \begin{syntax}
    \cs{e@alloc@chardef}
  \end{syntax}

\end{function}



\begin{function}[]
  {
    \e@alloc@intercharclass@top
  }
  \begin{syntax}
    \cs{e@alloc@intercharclass@top}
  \end{syntax}

\end{function}



\begin{function}[]
  {
    \e@alloc@luachunk@count
  }
  \begin{syntax}
    \cs{e@alloc@luachunk@count}
  \end{syntax}

\end{function}



\begin{function}[]
  {
    \e@alloc@luafunction@count
  }
  \begin{syntax}
    \cs{e@alloc@luafunction@count}
  \end{syntax}

\end{function}



\begin{function}[]
  {
    \e@alloc@whatsit@count
  }
  \begin{syntax}
    \cs{e@alloc@whatsit@count}
  \end{syntax}

\end{function}





\begin{function}[]
  {
    \e@mathgroup@top
  }
  \begin{syntax}
    \cs{e@mathgroup@top}
  \end{syntax}

\end{function}



\begin{function}[]
  {
    \eTeXversion
  }
  \begin{syntax}
    \cs{eTeXversion}
  \end{syntax}

\end{function}



\begin{function}[]
  {
    \edef
  }
  \begin{syntax}
    \cs{edef}
  \end{syntax}

\end{function}



\begin{function}[]
  {
    \egroup
  }
  \begin{syntax}
    \cs{egroup}
  \end{syntax}

\end{function}



\begin{function}[]
  {
    \eject
  }
  \begin{syntax}
    \cs{eject}
  \end{syntax}

\end{function}



\begin{function}[]
  {
    \else
  }
  \begin{syntax}
    \cs{else}
  \end{syntax}

\end{function}



\begin{function}[]
  {
    \em
  }
  \begin{syntax}
    \cs{em}
  \end{syntax}

\end{function}



\begin{function}[]
  {
    \em@currfont
  }
  \begin{syntax}
    \cs{em@currfont}
  \end{syntax}

\end{function}



\begin{function}[]
  {
    \em@force
  }
  \begin{syntax}
    \cs{em@force}
  \end{syntax}

\end{function}



\begin{function}[]
  {
    \emergencystretch
  }
  \begin{syntax}
    \cs{emergencystretch}
  \end{syntax}

\end{function}



\begin{function}[]
  {
    \emfontdeclare@clist
  }
  \begin{syntax}
    \cs{emfontdeclare@clist}
  \end{syntax}

\end{function}



\begin{function}[]
  {
    \emforce
  }
  \begin{syntax}
    \cs{emforce}
  \end{syntax}

\end{function}



\begin{function}[]
  {
    \eminnershape
  }
  \begin{syntax}
    \cs{eminnershape}
  \end{syntax}

\end{function}



\begin{function}[]
  {
    \emph
  }
  \begin{syntax}
    \cs{emph}
  \end{syntax}

\end{function}



\begin{function}[]
  {
    \empty
  }
  \begin{syntax}
    \cs{empty}
  \end{syntax}

\end{function}



\begin{function}[]
  {
    \empty@sfcnt
  }
  \begin{syntax}
    \cs{empty@sfcnt}
  \end{syntax}

\end{function}



\begin{function}[]
  {
    \emreset
  }
  \begin{syntax}
    \cs{emreset}
  \end{syntax}

\end{function}



\begin{function}[]
  {
    \enc@update
  }
  \begin{syntax}
    \cs{enc@update}
  \end{syntax}

\end{function}



\begin{function}[]
  {
    \encodingdefault
  }
  \begin{syntax}
    \cs{encodingdefault}
  \end{syntax}

\end{function}




\begin{function}[]
  {
    \end@dblfloat
  }
  \begin{syntax}
    \cs{end@dblfloat}
  \end{syntax}

\end{function}



\begin{function}[]
  {
    \end@float
  }
  \begin{syntax}
    \cs{end@float}
  \end{syntax}

\end{function}



\begin{function}[]
  {
    \endarray
  }
  \begin{syntax}
    \cs{endarray}
  \end{syntax}

\end{function}



\begin{function}[]
  {
    \endcenter
  }
  \begin{syntax}
    \cs{endcenter}
  \end{syntax}

\end{function}



\begin{function}[]
  {
    \endcsname
  }
  \begin{syntax}
    \cs{endcsname}
  \end{syntax}

\end{function}







\begin{function}[]
  {
    \endenumerate
  }
  \begin{syntax}
    \cs{endenumerate}
  \end{syntax}

\end{function}



\begin{function}[]
  {
    \endeqnarray
  }
  \begin{syntax}
    \cs{endeqnarray}
  \end{syntax}

\end{function}





\begin{function}[]
  {
    \endflushleft
  }
  \begin{syntax}
    \cs{endflushleft}
  \end{syntax}

\end{function}



\begin{function}[]
  {
    \endflushright
  }
  \begin{syntax}
    \cs{endflushright}
  \end{syntax}

\end{function}



\begin{function}[]
  {
    \endgraf
  }
  \begin{syntax}
    \cs{endgraf}
  \end{syntax}

\end{function}



\begin{function}[]
  {
    \endgroup
  }
  \begin{syntax}
    \cs{endgroup}
  \end{syntax}

\end{function}



\begin{function}[]
  {
    \endinput
  }
  \begin{syntax}
    \cs{endinput}
  \end{syntax}

\end{function}



\begin{function}[]
  {
    \enditemize
  }
  \begin{syntax}
    \cs{enditemize}
  \end{syntax}

\end{function}



\begin{function}[]
  {
    \endline
  }
  \begin{syntax}
    \cs{endline}
  \end{syntax}

\end{function}



\begin{function}[]
  {
    \endlinechar
  }
  \begin{syntax}
    \cs{endlinechar}
  \end{syntax}

\end{function}




\begin{function}[]
  {
    \endlrbox
  }
  \begin{syntax}
    \cs{endlrbox}
  \end{syntax}

\end{function}





\begin{function}[]
  {
    \endminipage
  }
  \begin{syntax}
    \cs{endminipage}
  \end{syntax}

\end{function}



\begin{function}[]
  {
    \endpicture
  }
  \begin{syntax}
    \cs{endpicture}
  \end{syntax}

\end{function}



\begin{function}[]
  {
    \endsloppypar
  }
  \begin{syntax}
    \cs{endsloppypar}
  \end{syntax}

\end{function}



\begin{function}[]
  {
    \endtabbing
  }
  \begin{syntax}
    \cs{endtabbing}
  \end{syntax}

\end{function}



\begin{function}[]
  {
    \endtabular
  }
  \begin{syntax}
    \cs{endtabular}
  \end{syntax}

\end{function}



\begin{function}[]
  {
    \endtrivlist
  }
  \begin{syntax}
    \cs{endtrivlist}
  \end{syntax}

\end{function}



\begin{function}[]
  {
    \endverbatim
  }
  \begin{syntax}
    \cs{endverbatim}
  \end{syntax}

\end{function}



\begin{function}[]
  {
    \enlargethispage
  }
  \begin{syntax}
    \cs{enlargethispage}
  \end{syntax}

\end{function}



\begin{function}[]
  {
    \enskip
  }
  \begin{syntax}
    \cs{enskip}
  \end{syntax}

\end{function}



\begin{function}[]
  {
    \enspace
  }
  \begin{syntax}
    \cs{enspace}
  \end{syntax}

\end{function}



\begin{function}[]
  {
    \ensuremath
  }
  \begin{syntax}
    \cs{ensuremath}
  \end{syntax}

\end{function}



\begin{function}[]
  {
    \enumerate
  }
  \begin{syntax}
    \cs{enumerate}
  \end{syntax}

\end{function}



\begin{function}[]
  {
    \eqnarray
  }
  \begin{syntax}
    \cs{eqnarray}
  \end{syntax}

\end{function}







\begin{function}[]
  {
    \errhelp
  }
  \begin{syntax}
    \cs{errhelp}
  \end{syntax}

\end{function}



\begin{function}[]
  {
    \errmessage
  }
  \begin{syntax}
    \cs{errmessage}
  \end{syntax}

\end{function}



\begin{function}[]
  {
    \error@fontshape
  }
  \begin{syntax}
    \cs{error@fontshape}
  \end{syntax}

\end{function}



\begin{function}[]
  {
    \errorcontextlines
  }
  \begin{syntax}
    \cs{errorcontextlines}
  \end{syntax}

\end{function}



\begin{function}[]
  {
    \errorstopmode
  }
  \begin{syntax}
    \cs{errorstopmode}
  \end{syntax}

\end{function}



\begin{function}[]
  {
    \escapechar
  }
  \begin{syntax}
    \cs{escapechar}
  \end{syntax}

\end{function}



\begin{function}[]
  {
    \evensidemargin
  }
  \begin{syntax}
    \cs{evensidemargin}
  \end{syntax}

\end{function}



\begin{function}[]
  {
    \every@math@size
  }
  \begin{syntax}
    \cs{every@math@size}
  \end{syntax}

\end{function}



\begin{function}[]
  {
    \everycr
  }
  \begin{syntax}
    \cs{everycr}
  \end{syntax}

\end{function}



\begin{function}[]
  {
    \everydisplay
  }
  \begin{syntax}
    \cs{everydisplay}
  \end{syntax}

\end{function}



\begin{function}[]
  {
    \everyjob
  }
  \begin{syntax}
    \cs{everyjob}
  \end{syntax}

\end{function}



\begin{function}[]
  {
    \everymath
  }
  \begin{syntax}
    \cs{everymath}
  \end{syntax}

\end{function}



\begin{function}[]
  {
    \everypar
  }
  \begin{syntax}
    \cs{everypar}
  \end{syntax}

\end{function}



\begin{function}[]
  {
    \execute@size@function
  }
  \begin{syntax}
    \cs{execute@size@function}
  \end{syntax}

\end{function}



\begin{function}[]
  {
    \exhyphenpenalty
  }
  \begin{syntax}
    \cs{exhyphenpenalty}
  \end{syntax}

\end{function}



\begin{function}[]
  {
    \exp
  }
  \begin{syntax}
    \cs{exp}
  \end{syntax}

\end{function}



\begin{function}[]
  {
    \expand@font@defaults
  }
  \begin{syntax}
    \cs{expand@font@defaults}
  \end{syntax}

\end{function}



\begin{function}[]
  {
    \expandableinput
  }
  \begin{syntax}
    \cs{expandableinput}
  \end{syntax}

\end{function}



\begin{function}[]
  {
    \expandafter
  }
  \begin{syntax}
    \cs{expandafter}
  \end{syntax}

\end{function}



\begin{function}[]
  {
    \expanded
  }
  \begin{syntax}
    \cs{expanded}
  \end{syntax}

\end{function}



\begin{function}[]
  {
    \external@font
  }
  \begin{syntax}
    \cs{external@font}
  \end{syntax}

\end{function}



\begin{function}[]
  {
    \extracolsep
  }
  \begin{syntax}
    \cs{extracolsep}
  \end{syntax}

\end{function}



\begin{function}[]
  {
    \extract@alph@from@version
  }
  \begin{syntax}
    \cs{extract@alph@from@version}
  \end{syntax}

\end{function}



\begin{function}[]
  {
    \extract@font
  }
  \begin{syntax}
    \cs{extract@font}
  \end{syntax}

\end{function}



\begin{function}[]
  {
    \extract@fontinfo
  }
  \begin{syntax}
    \cs{extract@fontinfo}
  \end{syntax}

\end{function}



\begin{function}[]
  {
    \extract@rangefontinfo
  }
  \begin{syntax}
    \cs{extract@rangefontinfo}
  \end{syntax}

\end{function}



\begin{function}[]
  {
    \extract@sizefn
  }
  \begin{syntax}
    \cs{extract@sizefn}
  \end{syntax}

\end{function}





\begin{function}[]
  {
    \f@baselineskip
  }
  \begin{syntax}
    \cs{f@baselineskip}
  \end{syntax}

\end{function}



\begin{function}[]
  {
    \f@depth
  }
  \begin{syntax}
    \cs{f@depth}
  \end{syntax}

\end{function}



\begin{function}[]
  {
    \f@encoding
  }
  \begin{syntax}
    \cs{f@encoding}
  \end{syntax}

\end{function}



\begin{function}[]
  {
    \f@family
  }
  \begin{syntax}
    \cs{f@family}
  \end{syntax}

\end{function}



\begin{function}[]
  {
    \f@linespread
  }
  \begin{syntax}
    \cs{f@linespread}
  \end{syntax}

\end{function}



\begin{function}[]
  {
    \f@series
  }
  \begin{syntax}
    \cs{f@series}
  \end{syntax}

\end{function}



\begin{function}[]
  {
    \f@series@saved
  }
  \begin{syntax}
    \cs{f@series@saved}
  \end{syntax}

\end{function}



\begin{function}[]
  {
    \f@shape
  }
  \begin{syntax}
    \cs{f@shape}
  \end{syntax}

\end{function}



\begin{function}[]
  {
    \f@shape@saved
  }
  \begin{syntax}
    \cs{f@shape@saved}
  \end{syntax}

\end{function}



\begin{function}[]
  {
    \f@size
  }
  \begin{syntax}
    \cs{f@size}
  \end{syntax}

\end{function}



\begin{function}[]
  {
    \f@user@size
  }
  \begin{syntax}
    \cs{f@user@size}
  \end{syntax}

\end{function}



\begin{function}[]
  {
    \fam
  }
  \begin{syntax}
    \cs{fam}
  \end{syntax}

\end{function}



\begin{function}[]
  {
    \familydefault
  }
  \begin{syntax}
    \cs{familydefault}
  \end{syntax}

\end{function}



\begin{function}[]
  {
    \fbox
  }
  \begin{syntax}
    \cs{fbox}
  \end{syntax}

\end{function}



\begin{function}[]
  {
    \fboxrule
  }
  \begin{syntax}
    \cs{fboxrule}
  \end{syntax}

\end{function}



\begin{function}[]
  {
    \fboxsep
  }
  \begin{syntax}
    \cs{fboxsep}
  \end{syntax}

\end{function}



\begin{function}[]
  {
    \fi
  }
  \begin{syntax}
    \cs{fi}
  \end{syntax}

\end{function}



\begin{function}[]
  {
    \filbreak
  }
  \begin{syntax}
    \cs{filbreak}
  \end{syntax}

\end{function}



\begin{function}[]
  {
    \filec@ntents
  }
  \begin{syntax}
    \cs{filec@ntents}
  \end{syntax}

\end{function}



\begin{function}[]
  {
    \filec@ntents@checkdir
  }
  \begin{syntax}
    \cs{filec@ntents@checkdir}
  \end{syntax}

\end{function}



\begin{function}[]
  {
    \filec@ntents@force
  }
  \begin{syntax}
    \cs{filec@ntents@force}
  \end{syntax}

\end{function}



\begin{function}[]
  {
    \filec@ntents@noheader
  }
  \begin{syntax}
    \cs{filec@ntents@noheader}
  \end{syntax}

\end{function}



\begin{function}[]
  {
    \filec@ntents@nosearch
  }
  \begin{syntax}
    \cs{filec@ntents@nosearch}
  \end{syntax}

\end{function}



\begin{function}[]
  {
    \filec@ntents@nowarn
  }
  \begin{syntax}
    \cs{filec@ntents@nowarn}
  \end{syntax}

\end{function}



\begin{function}[]
  {
    \filec@ntents@opt
  }
  \begin{syntax}
    \cs{filec@ntents@opt}
  \end{syntax}

\end{function}



\begin{function}[]
  {
    \filec@ntents@overwrite
  }
  \begin{syntax}
    \cs{filec@ntents@overwrite}
  \end{syntax}

\end{function}



\begin{function}[]
  {
    \filec@ntents@warning
  }
  \begin{syntax}
    \cs{filec@ntents@warning}
  \end{syntax}

\end{function}



\begin{function}[]
  {
    \filec@ntents@where
  }
  \begin{syntax}
    \cs{filec@ntents@where}
  \end{syntax}

\end{function}




\begin{function}[]
  {
    \filename@area
  }
  \begin{syntax}
    \cs{filename@area}
  \end{syntax}

\end{function}



\begin{function}[]
  {
    \filename@base
  }
  \begin{syntax}
    \cs{filename@base}
  \end{syntax}

\end{function}



\begin{function}[]
  {
    \filename@dot
  }
  \begin{syntax}
    \cs{filename@dot}
  \end{syntax}

\end{function}



\begin{function}[]
  {
    \filename@dots
  }
  \begin{syntax}
    \cs{filename@dots}
  \end{syntax}

\end{function}



\begin{function}[]
  {
    \filename@ext
  }
  \begin{syntax}
    \cs{filename@ext}
  \end{syntax}

\end{function}



\begin{function}[]
  {
    \filename@parse
  }
  \begin{syntax}
    \cs{filename@parse}
  \end{syntax}

\end{function}



\begin{function}[]
  {
    \filename@path
  }
  \begin{syntax}
    \cs{filename@path}
  \end{syntax}

\end{function}



\begin{function}[]
  {
    \filename@simple
  }
  \begin{syntax}
    \cs{filename@simple}
  \end{syntax}

\end{function}



\begin{function}[]
  {
    \filesize
  }
  \begin{syntax}
    \cs{filesize}
  \end{syntax}

\end{function}



\begin{function}[]
  {
    \fill
  }
  \begin{syntax}
    \cs{fill}
  \end{syntax}

\end{function}



\begin{function}[]
  {
    \finalhyphendemerits
  }
  \begin{syntax}
    \cs{finalhyphendemerits}
  \end{syntax}

\end{function}





\begin{function}[]
  {
    \finph@nt
  }
  \begin{syntax}
    \cs{finph@nt}
  \end{syntax}

\end{function}



\begin{function}[]
  {
    \finsm@sh
  }
  \begin{syntax}
    \cs{finsm@sh}
  \end{syntax}

\end{function}



\begin{function}[]
  {
    \fix@penalty
  }
  \begin{syntax}
    \cs{fix@penalty}
  \end{syntax}

\end{function}



\begin{function}[]
  {
    \fixed@sfcnt
  }
  \begin{syntax}
    \cs{fixed@sfcnt}
  \end{syntax}

\end{function}



\begin{function}[]
  {
    \fl@ShowFloat
  }
  \begin{syntax}
    \cs{fl@ShowFloat}
  \end{syntax}

\end{function}



\begin{function}[]
  {
    \fl@trace
  }
  \begin{syntax}
    \cs{fl@trace}
  \end{syntax}

\end{function}



\begin{function}[]
  {
    \fl@tracemessage
  }
  \begin{syntax}
    \cs{fl@tracemessage}
  \end{syntax}

\end{function}



\begin{function}[]
  {
    \fl@traceval
  }
  \begin{syntax}
    \cs{fl@traceval}
  \end{syntax}

\end{function}



\begin{function}[]
  {
    \float@count
  }
  \begin{syntax}
    \cs{float@count}
  \end{syntax}

\end{function}



\begin{function}[]
  {
    \floatingpenalty
  }
  \begin{syntax}
    \cs{floatingpenalty}
  \end{syntax}

\end{function}



\begin{function}[]
  {
    \floatpagefraction
  }
  \begin{syntax}
    \cs{floatpagefraction}
  \end{syntax}

\end{function}



\begin{function}[]
  {
    \floatsep
  }
  \begin{syntax}
    \cs{floatsep}
  \end{syntax}

\end{function}



\begin{function}[]
  {
    \flushbottom
  }
  \begin{syntax}
    \cs{flushbottom}
  \end{syntax}

\end{function}



\begin{function}[]
  {
    \flushleft
  }
  \begin{syntax}
    \cs{flushleft}
  \end{syntax}

\end{function}



\begin{function}[]
  {
    \flushright
  }
  \begin{syntax}
    \cs{flushright}
  \end{syntax}

\end{function}






\begin{function}[]
  {
    \fmtversion
  }
  \begin{syntax}
    \cs{fmtversion}
  \end{syntax}

\end{function}



\begin{function}[]
  {
    \fnsymbol
  }
  \begin{syntax}
    \cs{fnsymbol}
  \end{syntax}

\end{function}



\begin{function}[]
  {
    \font
  }
  \begin{syntax}
    \cs{font}
  \end{syntax}

\end{function}



\begin{function}[]
  {
    \font@info
  }
  \begin{syntax}
    \cs{font@info}
  \end{syntax}

\end{function}



\begin{function}[]
  {
    \font@name
  }
  \begin{syntax}
    \cs{font@name}
  \end{syntax}

\end{function}



\begin{function}[]
  {
    \font@submax
  }
  \begin{syntax}
    \cs{font@submax}
  \end{syntax}

\end{function}



\begin{function}[]
  {
    \fontdimen
  }
  \begin{syntax}
    \cs{fontdimen}
  \end{syntax}

\end{function}



\begin{function}[]
  {
    \fontencoding
  }
  \begin{syntax}
    \cs{fontencoding}
  \end{syntax}

\end{function}



\begin{function}[]
  {
    \fontfamily
  }
  \begin{syntax}
    \cs{fontfamily}
  \end{syntax}

\end{function}



\begin{function}[]
  {
    \fontname
  }
  \begin{syntax}
    \cs{fontname}
  \end{syntax}

\end{function}



\begin{function}[]
  {
    \fontseries
  }
  \begin{syntax}
    \cs{fontseries}
  \end{syntax}

\end{function}



\begin{function}[]
  {
    \fontseriesforce
  }
  \begin{syntax}
    \cs{fontseriesforce}
  \end{syntax}

\end{function}



\begin{function}[]
  {
    \fontshape
  }
  \begin{syntax}
    \cs{fontshape}
  \end{syntax}

\end{function}



\begin{function}[]
  {
    \fontshapeforce
  }
  \begin{syntax}
    \cs{fontshapeforce}
  \end{syntax}

\end{function}



\begin{function}[]
  {
    \fontsize
  }
  \begin{syntax}
    \cs{fontsize}
  \end{syntax}

\end{function}



\begin{function}[]
  {
    \fontsubfuzz
  }
  \begin{syntax}
    \cs{fontsubfuzz}
  \end{syntax}

\end{function}



\begin{function}[]
  {
    \footins
  }
  \begin{syntax}
    \cs{footins}
  \end{syntax}

\end{function}



\begin{function}[]
  {
    \footnote
  }
  \begin{syntax}
    \cs{footnote}
  \end{syntax}

\end{function}



\begin{function}[]
  {
    \footnotemark
  }
  \begin{syntax}
    \cs{footnotemark}
  \end{syntax}

\end{function}



\begin{function}[]
  {
    \footnoterule
  }
  \begin{syntax}
    \cs{footnoterule}
  \end{syntax}

\end{function}



\begin{function}[]
  {
    \footnotesep
  }
  \begin{syntax}
    \cs{footnotesep}
  \end{syntax}

\end{function}



\begin{function}[]
  {
    \footnotesize
  }
  \begin{syntax}
    \cs{footnotesize}
  \end{syntax}

\end{function}



\begin{function}[]
  {
    \footnotetext
  }
  \begin{syntax}
    \cs{footnotetext}
  \end{syntax}

\end{function}



\begin{function}[]
  {
    \footref
  }
  \begin{syntax}
    \cs{footref}
  \end{syntax}

\end{function}



\begin{function}[]
  {
    \footskip
  }
  \begin{syntax}
    \cs{footskip}
  \end{syntax}

\end{function}



\begin{function}[]
  {
    \fpeval
  }
  \begin{syntax}
    \cs{fpeval}
  \end{syntax}

\end{function}



\begin{function}[]
  {
    \frac
  }
  \begin{syntax}
    \cs{frac}
  \end{syntax}

\end{function}



\begin{function}[]
  {
    \frame
  }
  \begin{syntax}
    \cs{frame}
  \end{syntax}

\end{function}



\begin{function}[]
  {
    \framebox
  }
  \begin{syntax}
    \cs{framebox}
  \end{syntax}

\end{function}



\begin{function}[]
  {
    \freeze@math@version
  }
  \begin{syntax}
    \cs{freeze@math@version}
  \end{syntax}

\end{function}



\begin{function}[]
  {
    \frenchspacing
  }
  \begin{syntax}
    \cs{frenchspacing}
  \end{syntax}

\end{function}



\begin{function}[]
  {
    \frozen@everydisplay
  }
  \begin{syntax}
    \cs{frozen@everydisplay}
  \end{syntax}

\end{function}



\begin{function}[]
  {
    \frozen@everymath
  }
  \begin{syntax}
    \cs{frozen@everymath}
  \end{syntax}

\end{function}



\begin{function}[]
  {
    \fussy
  }
  \begin{syntax}
    \cs{fussy}
  \end{syntax}

\end{function}



\begin{function}[]
  {
    \futurelet
  }
  \begin{syntax}
    \cs{futurelet}
  \end{syntax}

\end{function}



\begin{function}[]
  {
    \g@addto@macro
  }
  \begin{syntax}
    \cs{g@addto@macro}
  \end{syntax}

\end{function}



\begin{function}[]
  {
    \gcd
  }
  \begin{syntax}
    \cs{gcd}
  \end{syntax}

\end{function}



\begin{function}[]
  {
    \gdef
  }
  \begin{syntax}
    \cs{gdef}
  \end{syntax}

\end{function}



\begin{function}[]
  {
    \gen@sfcnt
  }
  \begin{syntax}
    \cs{gen@sfcnt}
  \end{syntax}

\end{function}



\begin{function}[]
  {
    \genb@sfcnt
  }
  \begin{syntax}
    \cs{genb@sfcnt}
  \end{syntax}

\end{function}



\begin{function}[]
  {
    \genb@x
  }
  \begin{syntax}
    \cs{genb@x}
  \end{syntax}

\end{function}



\begin{function}[]
  {
    \genb@y
  }
  \begin{syntax}
    \cs{genb@y}
  \end{syntax}

\end{function}



\begin{function}[]
  {
    \get@cdp
  }
  \begin{syntax}
    \cs{get@cdp}
  \end{syntax}

\end{function}



\begin{function}[]
  {
    \get@external@font
  }
  \begin{syntax}
    \cs{get@external@font}
  \end{syntax}

\end{function}



\begin{function}[]
  {
    \getanddefine@fonts
  }
  \begin{syntax}
    \cs{getanddefine@fonts}
  \end{syntax}

\end{function}



\begin{function}[]
  {
    \glb@currsize
  }
  \begin{syntax}
    \cs{glb@currsize}
  \end{syntax}

\end{function}



\begin{function}[]
  {
    \glb@settings
  }
  \begin{syntax}
    \cs{glb@settings}
  \end{syntax}

\end{function}



\begin{function}[]
  {
    \global
  }
  \begin{syntax}
    \cs{global}
  \end{syntax}

\end{function}



\begin{function}[]
  {
    \globaldefs
  }
  \begin{syntax}
    \cs{globaldefs}
  \end{syntax}

\end{function}



\begin{function}[]
  {
    \glossary
  }
  \begin{syntax}
    \cs{glossary}
  \end{syntax}

\end{function}



\begin{function}[]
  {
    \glossaryentry
  }
  \begin{syntax}
    \cs{glossaryentry}
  \end{syntax}

\end{function}



\begin{function}[]
  {
    \gluestretchorder
  }
  \begin{syntax}
    \cs{gluestretchorder}
  \end{syntax}

\end{function}





\begin{function}[]
  {
    \goodbreak
  }
  \begin{syntax}
    \cs{goodbreak}
  \end{syntax}

\end{function}



\begin{function}[]
  {
    \group@elt
  }
  \begin{syntax}
    \cs{group@elt}
  \end{syntax}

\end{function}



\begin{function}[]
  {
    \group@list
  }
  \begin{syntax}
    \cs{group@list}
  \end{syntax}

\end{function}



\begin{function}[]
  {
    \h@false
  }
  \begin{syntax}
    \cs{h@false}
  \end{syntax}

\end{function}



\begin{function}[]
  {
    \h@true
  }
  \begin{syntax}
    \cs{h@true}
  \end{syntax}

\end{function}



\begin{function}[]
  {
    \halign
  }
  \begin{syntax}
    \cs{halign}
  \end{syntax}

\end{function}



\begin{function}[]
  {
    \hangindent
  }
  \begin{syntax}
    \cs{hangindent}
  \end{syntax}

\end{function}



\begin{function}[]
  {
    \hb@xt@
  }
  \begin{syntax}
    \cs{hb@xt@}
  \end{syntax}

\end{function}



\begin{function}[]
  {
    \hbadness
  }
  \begin{syntax}
    \cs{hbadness}
  \end{syntax}

\end{function}



\begin{function}[]
  {
    \hbox
  }
  \begin{syntax}
    \cs{hbox}
  \end{syntax}

\end{function}



\begin{function}[]
  {
    \headheight
  }
  \begin{syntax}
    \cs{headheight}
  \end{syntax}

\end{function}



\begin{function}[]
  {
    \headsep
  }
  \begin{syntax}
    \cs{headsep}
  \end{syntax}

\end{function}



\begin{function}[]
  {
    \height
  }
  \begin{syntax}
    \cs{height}
  \end{syntax}

\end{function}



\begin{function}[]
  {
    \hexnumber@
  }
  \begin{syntax}
    \cs{hexnumber@}
  \end{syntax}

\end{function}



\begin{function}[]
  {
    \hfil
  }
  \begin{syntax}
    \cs{hfil}
  \end{syntax}

\end{function}




\begin{function}[]
  {
    \hfill
  }
  \begin{syntax}
    \cs{hfill}
  \end{syntax}

\end{function}



\begin{function}[]
  {
    \hfuzz
  }
  \begin{syntax}
    \cs{hfuzz}
  \end{syntax}

\end{function}



\begin{function}[]
  {
    \hgl@
  }
  \begin{syntax}
    \cs{hgl@}
  \end{syntax}

\end{function}



\begin{function}[]
  {
    \hglue
  }
  \begin{syntax}
    \cs{hglue}
  \end{syntax}

\end{function}



\begin{function}[]
  {
    \hideoutput
  }
  \begin{syntax}
    \cs{hideoutput}
  \end{syntax}

\end{function}



\begin{function}[]
  {
    \hideskip
  }
  \begin{syntax}
    \cs{hideskip}
  \end{syntax}

\end{function}



\begin{function}[]
  {
    \hidewidth
  }
  \begin{syntax}
    \cs{hidewidth}
  \end{syntax}

\end{function}



\begin{function}[]
  {
    \hline
  }
  \begin{syntax}
    \cs{hline}
  \end{syntax}

\end{function}



\begin{function}[]
  {
    \hmode@bgroup
  }
  \begin{syntax}
    \cs{hmode@bgroup}
  \end{syntax}

\end{function}



\begin{function}[]
  {
    \hmode@start@before@group
  }
  \begin{syntax}
    \cs{hmode@start@before@group}
  \end{syntax}

\end{function}



\begin{function}[]
  {
    \hom
  }
  \begin{syntax}
    \cs{hom}
  \end{syntax}

\end{function}



\begin{function}[]
  {
    \hook_disable_generic:n
  }
  \begin{syntax}
    \cs{hook_disable_generic:n}
  \end{syntax}

\end{function}



\begin{function}[]
  {
    \hphantom
  }
  \begin{syntax}
    \cs{hphantom}
  \end{syntax}

\end{function}



\begin{function}[]
  {
    \hrule
  }
  \begin{syntax}
    \cs{hrule}
  \end{syntax}

\end{function}



\begin{function}[]
  {
    \hrulefill
  }
  \begin{syntax}
    \cs{hrulefill}
  \end{syntax}

\end{function}



\begin{function}[]
  {
    \hsize
  }
  \begin{syntax}
    \cs{hsize}
  \end{syntax}

\end{function}



\begin{function}[]
  {
    \hskip
  }
  \begin{syntax}
    \cs{hskip}
  \end{syntax}

\end{function}



\begin{function}[]
  {
    \hspace
  }
  \begin{syntax}
    \cs{hspace}
  \end{syntax}

\end{function}



\begin{function}[]
  {
    \hss
  }
  \begin{syntax}
    \cs{hss}
  \end{syntax}

\end{function}



\begin{function}[]
  {
    \ht
  }
  \begin{syntax}
    \cs{ht}
  \end{syntax}

\end{function}



\begin{function}[]
  {
    \ht0
  }
  \begin{syntax}
    \cs{ht0}
  \end{syntax}

\end{function}



\begin{function}[]
  {
    \hyper@nopatch@longtable
  }
  \begin{syntax}
    \cs{hyper@nopatch@longtable}
  \end{syntax}

\end{function}



\begin{function}[]
  {
    \hyphenchar
  }
  \begin{syntax}
    \cs{hyphenchar}
  \end{syntax}

\end{function}



\begin{function}[]
  {
    \hyphenpenalty
  }
  \begin{syntax}
    \cs{hyphenpenalty}
  \end{syntax}

\end{function}



\begin{function}[]
  {
    \i
  }
  \begin{syntax}
    \cs{i}
  \end{syntax}

\end{function}



\begin{function}[]
  {
    \ialign
  }
  \begin{syntax}
    \cs{ialign}
  \end{syntax}

\end{function}



\begin{function}[]
  {
    \if
  }
  \begin{syntax}
    \cs{if}
  \end{syntax}

\end{function}



\begin{function}[]
  {
    \if@afterindent
  }
  \begin{syntax}
    \cs{if@afterindent}
  \end{syntax}

\end{function}



\begin{function}[]
  {
    \if@compatibility
  }
  \begin{syntax}
    \cs{if@compatibility}
  \end{syntax}

\end{function}



\begin{function}[]
  {
    \if@endpe
  }
  \begin{syntax}
    \cs{if@endpe}
  \end{syntax}

\end{function}



\begin{function}[]
  {
    \if@eqnsw
  }
  \begin{syntax}
    \cs{if@eqnsw}
  \end{syntax}

\end{function}



\begin{function}[]
  {
    \if@fcolmade
  }
  \begin{syntax}
    \cs{if@fcolmade}
  \end{syntax}

\end{function}



\begin{function}[]
  {
    \if@filesw
  }
  \begin{syntax}
    \cs{if@filesw}
  \end{syntax}

\end{function}



\begin{function}[]
  {
    \if@firstamp
  }
  \begin{syntax}
    \cs{if@firstamp}
  \end{syntax}

\end{function}



\begin{function}[]
  {
    \if@firstcolumn
  }
  \begin{syntax}
    \cs{if@firstcolumn}
  \end{syntax}

\end{function}



\begin{function}[]
  {
    \if@font@series@context
  }
  \begin{syntax}
    \cs{if@font@series@context}
  \end{syntax}

\end{function}



\begin{function}[]
  {
    \if@forced@series
  }
  \begin{syntax}
    \cs{if@forced@series}
  \end{syntax}

\end{function}




\begin{function}[]
  {
    \if@in@minipage@env
  }
  \begin{syntax}
    \cs{if@in@minipage@env}
  \end{syntax}

\end{function}



\begin{function}[]
  {
    \if@includeinrelease
  }
  \begin{syntax}
    \cs{if@includeinrelease}
  \end{syntax}

\end{function}



\begin{function}[]
  {
    \if@inlabel
  }
  \begin{syntax}
    \cs{if@inlabel}
  \end{syntax}

\end{function}



\begin{function}[]
  {
    \if@insert
  }
  \begin{syntax}
    \cs{if@insert}
  \end{syntax}

\end{function}



\begin{function}[]
  {
    \if@listfiles@hashes
  }
  \begin{syntax}
    \cs{if@listfiles@hashes}
  \end{syntax}

\end{function}



\begin{function}[]
  {
    \if@listfiles@sizes
  }
  \begin{syntax}
    \cs{if@listfiles@sizes}
  \end{syntax}

\end{function}



\begin{function}[]
  {
    \if@minipage
  }
  \begin{syntax}
    \cs{if@minipage}
  \end{syntax}

\end{function}



\begin{function}[]
  {
    \if@mparswitch
  }
  \begin{syntax}
    \cs{if@mparswitch}
  \end{syntax}

\end{function}



\begin{function}[]
  {
    \if@negarg
  }
  \begin{syntax}
    \cs{if@negarg}
  \end{syntax}

\end{function}



\begin{function}[]
  {
    \if@newlist
  }
  \begin{syntax}
    \cs{if@newlist}
  \end{syntax}

\end{function}



\begin{function}[]
  {
    \if@nmbrlist
  }
  \begin{syntax}
    \cs{if@nmbrlist}
  \end{syntax}

\end{function}



\begin{function}[]
  {
    \if@noitemarg
  }
  \begin{syntax}
    \cs{if@noitemarg}
  \end{syntax}

\end{function}




\begin{function}[]
  {
    \if@noparlist
  }
  \begin{syntax}
    \cs{if@noparlist}
  \end{syntax}

\end{function}



\begin{function}[]
  {
    \if@noskipsec
  }
  \begin{syntax}
    \cs{if@noskipsec}
  \end{syntax}

\end{function}



\begin{function}[]
  {
    \if@ovb
  }
  \begin{syntax}
    \cs{if@ovb}
  \end{syntax}

\end{function}



\begin{function}[]
  {
    \if@ovhline
  }
  \begin{syntax}
    \cs{if@ovhline}
  \end{syntax}

\end{function}



\begin{function}[]
  {
    \if@ovl
  }
  \begin{syntax}
    \cs{if@ovl}
  \end{syntax}

\end{function}



\begin{function}[]
  {
    \if@ovr
  }
  \begin{syntax}
    \cs{if@ovr}
  \end{syntax}

\end{function}



\begin{function}[]
  {
    \if@ovt
  }
  \begin{syntax}
    \cs{if@ovt}
  \end{syntax}

\end{function}



\begin{function}[]
  {
    \if@ovvline
  }
  \begin{syntax}
    \cs{if@ovvline}
  \end{syntax}

\end{function}



\begin{function}[]
  {
    \if@partsw
  }
  \begin{syntax}
    \cs{if@partsw}
  \end{syntax}

\end{function}



\begin{function}[]
  {
    \if@pboxsw
  }
  \begin{syntax}
    \cs{if@pboxsw}
  \end{syntax}

\end{function}



\begin{function}[]
  {
    \if@reversemargin
  }
  \begin{syntax}
    \cs{if@reversemargin}
  \end{syntax}

\end{function}



\begin{function}[]
  {
    \if@rjfield
  }
  \begin{syntax}
    \cs{if@rjfield}
  \end{syntax}

\end{function}





\begin{function}[]
  {
    \if@specialpage
  }
  \begin{syntax}
    \cs{if@specialpage}
  \end{syntax}

\end{function}



\begin{function}[]
  {
    \if@tempswa
  }
  \begin{syntax}
    \cs{if@tempswa}
  \end{syntax}

\end{function}



\begin{function}[]
  {
    \if@test
  }
  \begin{syntax}
    \cs{if@test}
  \end{syntax}

\end{function}



\begin{function}[]
  {
    \if@twocolumn
  }
  \begin{syntax}
    \cs{if@twocolumn}
  \end{syntax}

\end{function}



\begin{function}[]
  {
    \if@twoside
  }
  \begin{syntax}
    \cs{if@twoside}
  \end{syntax}

\end{function}



\begin{function}[]
  {
    \ifcase
  }
  \begin{syntax}
    \cs{ifcase}
  \end{syntax}

\end{function}



\begin{function}[]
  {
    \ifcat
  }
  \begin{syntax}
    \cs{ifcat}
  \end{syntax}

\end{function}



\begin{function}[]
  {
    \ifcsname
  }
  \begin{syntax}
    \cs{ifcsname}
  \end{syntax}

\end{function}



\begin{function}[]
  {
    \ifdefined
  }
  \begin{syntax}
    \cs{ifdefined}
  \end{syntax}

\end{function}



\begin{function}[]
  {
    \ifdim
  }
  \begin{syntax}
    \cs{ifdim}
  \end{syntax}

\end{function}



\begin{function}[]
  {
    \ifdt@p
  }
  \begin{syntax}
    \cs{ifdt@p}
  \end{syntax}

\end{function}



\begin{function}[]
  {
    \ifeof
  }
  \begin{syntax}
    \cs{ifeof}
  \end{syntax}

\end{function}



\begin{function}[]
  {
    \iffalse
  }
  \begin{syntax}
    \cs{iffalse}
  \end{syntax}

\end{function}



\begin{function}[]
  {
    \ifh@
  }
  \begin{syntax}
    \cs{ifh@}
  \end{syntax}

\end{function}



\begin{function}[]
  {
    \ifhmode
  }
  \begin{syntax}
    \cs{ifhmode}
  \end{syntax}

\end{function}



\begin{function}[]
  {
    \ifin@
  }
  \begin{syntax}
    \cs{ifin@}
  \end{syntax}

\end{function}



\begin{function}[]
  {
    \ifincsname
  }
  \begin{syntax}
    \cs{ifincsname}
  \end{syntax}

\end{function}



\begin{function}[]
  {
    \ifinner
  }
  \begin{syntax}
    \cs{ifinner}
  \end{syntax}

\end{function}



\begin{function}[]
  {
    \ifmath@fonts
  }
  \begin{syntax}
    \cs{ifmath@fonts}
  \end{syntax}

\end{function}



\begin{function}[]
  {
    \ifmaybe@ic
  }
  \begin{syntax}
    \cs{ifmaybe@ic}
  \end{syntax}

\end{function}



\begin{function}[]
  {
    \ifmmode
  }
  \begin{syntax}
    \cs{ifmmode}
  \end{syntax}

\end{function}



\begin{function}[]
  {
    \ifnot@nil
  }
  \begin{syntax}
    \cs{ifnot@nil}
  \end{syntax}

\end{function}



\begin{function}[]
  {
    \ifnum
  }
  \begin{syntax}
    \cs{ifnum}
  \end{syntax}

\end{function}



\begin{function}[]
  {
    \ifodd
  }
  \begin{syntax}
    \cs{ifodd}
  \end{syntax}

\end{function}



\begin{function}[]
  {
    \iftrue
  }
  \begin{syntax}
    \cs{iftrue}
  \end{syntax}

\end{function}



\begin{function}[]
  {
    \ifv@
  }
  \begin{syntax}
    \cs{ifv@}
  \end{syntax}

\end{function}



\begin{function}[]
  {
    \ifvbox
  }
  \begin{syntax}
    \cs{ifvbox}
  \end{syntax}

\end{function}



\begin{function}[]
  {
    \ifvmode
  }
  \begin{syntax}
    \cs{ifvmode}
  \end{syntax}

\end{function}



\begin{function}[]
  {
    \ifvoid
  }
  \begin{syntax}
    \cs{ifvoid}
  \end{syntax}

\end{function}



\begin{function}[]
  {
    \ifx
  }
  \begin{syntax}
    \cs{ifx}
  \end{syntax}

\end{function}



\begin{function}[]
  {
    \ignorespaces
  }
  \begin{syntax}
    \cs{ignorespaces}
  \end{syntax}

\end{function}




\begin{function}[]
  {
    \ij
  }
  \begin{syntax}
    \cs{ij}
  \end{syntax}

\end{function}



\begin{function}[]
  {
    \immediate
  }
  \begin{syntax}
    \cs{immediate}
  \end{syntax}

\end{function}



\begin{function}[]
  {
    \in@
  }
  \begin{syntax}
    \cs{in@}
  \end{syntax}

\end{function}



\begin{function}[]
  {
    \in@@
  }
  \begin{syntax}
    \cs{in@@}
  \end{syntax}

\end{function}



\begin{function}[]
  {
    \in@false
  }
  \begin{syntax}
    \cs{in@false}
  \end{syntax}

\end{function}



\begin{function}[]
  {
    \in@true
  }
  \begin{syntax}
    \cs{in@true}
  \end{syntax}

\end{function}



\begin{function}[]
  {
    \include
  }
  \begin{syntax}
    \cs{include}
  \end{syntax}

\end{function}



\begin{function}[]
  {
    \includeonly
  }
  \begin{syntax}
    \cs{includeonly}
  \end{syntax}

\end{function}



\begin{function}[]
  {
    \indent
  }
  \begin{syntax}
    \cs{indent}
  \end{syntax}

\end{function}



\begin{function}[]
  {
    \index
  }
  \begin{syntax}
    \cs{index}
  \end{syntax}

\end{function}



\begin{function}[]
  {
    \indexentry
  }
  \begin{syntax}
    \cs{indexentry}
  \end{syntax}

\end{function}



\begin{function}[]
  {
    \inf
  }
  \begin{syntax}
    \cs{inf}
  \end{syntax}

\end{function}



\begin{function}[]
  {
    \init@restore@glb@settings
  }
  \begin{syntax}
    \cs{init@restore@glb@settings}
  \end{syntax}

\end{function}



\begin{function}[]
  {
    \init@restore@version
  }
  \begin{syntax}
    \cs{init@restore@version}
  \end{syntax}

\end{function}



\begin{function}[]
  {
    \init@series@setup
  }
  \begin{syntax}
    \cs{init@series@setup}
  \end{syntax}

\end{function}



\begin{function}[]
  {
    \initcatcodetable
  }
  \begin{syntax}
    \cs{initcatcodetable}
  \end{syntax}

\end{function}



\begin{function}[]
  {
    \input
  }
  \begin{syntax}
    \cs{input}
  \end{syntax}

\end{function}



\begin{function}[]
  {
    \input@path
  }
  \begin{syntax}
    \cs{input@path}
  \end{syntax}

\end{function}



\begin{function}[]
  {
    \inputencodingname
  }
  \begin{syntax}
    \cs{inputencodingname}
  \end{syntax}

\end{function}



\begin{function}[]
  {
    \inputlineno
  }
  \begin{syntax}
    \cs{inputlineno}
  \end{syntax}

\end{function}





\begin{function}[]
  {
    \insert
  }
  \begin{syntax}
    \cs{insert}
  \end{syntax}

\end{function}



\begin{function}[]
  {
    \install@mathalphabet
  }
  \begin{syntax}
    \cs{install@mathalphabet}
  \end{syntax}

\end{function}



\begin{function}[]
  {
    \interdisplaylinepenalty
  }
  \begin{syntax}
    \cs{interdisplaylinepenalty}
  \end{syntax}

\end{function}



\begin{function}[]
  {
    \interfootnotelinepenalty
  }
  \begin{syntax}
    \cs{interfootnotelinepenalty}
  \end{syntax}

\end{function}



\begin{function}[]
  {
    \interlinepenalty
  }
  \begin{syntax}
    \cs{interlinepenalty}
  \end{syntax}

\end{function}



\begin{function}[]
  {
    \inteval
  }
  \begin{syntax}
    \cs{inteval}
  \end{syntax}

\end{function}



\begin{function}[]
  {
    \intextsep
  }
  \begin{syntax}
    \cs{intextsep}
  \end{syntax}

\end{function}



\begin{function}[]
  {
    \is@range
  }
  \begin{syntax}
    \cs{is@range}
  \end{syntax}

\end{function}



\begin{function}[]
  {
    \itdefault
  }
  \begin{syntax}
    \cs{itdefault}
  \end{syntax}

\end{function}



\begin{function}[]
  {
    \item
  }
  \begin{syntax}
    \cs{item}
  \end{syntax}

\end{function}



\begin{function}[]
  {
    \itemindent
  }
  \begin{syntax}
    \cs{itemindent}
  \end{syntax}

\end{function}



\begin{function}[]
  {
    \itemize
  }
  \begin{syntax}
    \cs{itemize}
  \end{syntax}

\end{function}





\begin{function}[]
  {
    \iterate
  }
  \begin{syntax}
    \cs{iterate}
  \end{syntax}

\end{function}



\begin{function}[]
  {
    \itshape
  }
  \begin{syntax}
    \cs{itshape}
  \end{syntax}

\end{function}



\begin{function}[]
  {
    \j
  }
  \begin{syntax}
    \cs{j}
  \end{syntax}

\end{function}



\begin{function}[]
  {
    \jobname
  }
  \begin{syntax}
    \cs{jobname}
  \end{syntax}

\end{function}



\begin{function}[]
  {
    \jot
  }
  \begin{syntax}
    \cs{jot}
  \end{syntax}

\end{function}



\begin{function}[]
  {
    \k
  }
  \begin{syntax}
    \cs{k}
  \end{syntax}

\end{function}



\begin{function}[]
  {
    \kanjiskip
  }
  \begin{syntax}
    \cs{kanjiskip}
  \end{syntax}

\end{function}



\begin{function}[]
  {
    \ker
  }
  \begin{syntax}
    \cs{ker}
  \end{syntax}

\end{function}



\begin{function}[]
  {
    \kern
  }
  \begin{syntax}
    \cs{kern}
  \end{syntax}

\end{function}



\begin{function}[]
  {
    \kernel@ifnextchar
  }
  \begin{syntax}
    \cs{kernel@ifnextchar}
  \end{syntax}

\end{function}



\begin{function}[]
  {
    \kernel@make@fragile
  }
  \begin{syntax}
    \cs{kernel@make@fragile}
  \end{syntax}

\end{function}



\begin{function}[]
  {
    \kerneltmpDoNotUse
  }
  \begin{syntax}
    \cs{kerneltmpDoNotUse}
  \end{syntax}

\end{function}



\begin{function}[]
  {
    \kill
  }
  \begin{syntax}
    \cs{kill}
  \end{syntax}

\end{function}



\begin{function}[]
  {
    \l
  }
  \begin{syntax}
    \cs{l}
  \end{syntax}

\end{function}



\begin{function}[]
  {
    \l@ngrel@x
  }
  \begin{syntax}
    \cs{l@ngrel@x}
  \end{syntax}

\end{function}



\begin{function}[]
  {
    \l@nohyphenation
  }
  \begin{syntax}
    \cs{l@nohyphenation}
  \end{syntax}

\end{function}



\begin{function}[]
  {
    \label
  }
  \begin{syntax}
    \cs{label}
  \end{syntax}

\end{function}



\begin{function}[]
  {
    \labelformat
  }
  \begin{syntax}
    \cs{labelformat}
  \end{syntax}

\end{function}



\begin{function}[]
  {
    \labelsep
  }
  \begin{syntax}
    \cs{labelsep}
  \end{syntax}

\end{function}



\begin{function}[]
  {
    \labelwidth
  }
  \begin{syntax}
    \cs{labelwidth}
  \end{syntax}

\end{function}



\begin{function}[]
  {
    \language
  }
  \begin{syntax}
    \cs{language}
  \end{syntax}

\end{function}



\begin{function}[]
  {
    \last@fontshape
  }
  \begin{syntax}
    \cs{last@fontshape}
  \end{syntax}

\end{function}



\begin{function}[]
  {
    \lastbox
  }
  \begin{syntax}
    \cs{lastbox}
  \end{syntax}

\end{function}



\begin{function}[]
  {
    \lastnamedcs
  }
  \begin{syntax}
    \cs{lastnamedcs}
  \end{syntax}

\end{function}



\begin{function}[]
  {
    \lastnodetype
  }
  \begin{syntax}
    \cs{lastnodetype}
  \end{syntax}

\end{function}



\begin{function}[]
  {
    \lastpenalty
  }
  \begin{syntax}
    \cs{lastpenalty}
  \end{syntax}

\end{function}



\begin{function}[]
  {
    \lastskip
  }
  \begin{syntax}
    \cs{lastskip}
  \end{syntax}

\end{function}



\begin{function}[]
  {
    \latelua
  }
  \begin{syntax}
    \cs{latelua}
  \end{syntax}

\end{function}



\begin{function}[]
  {
    \lbrace
  }
  \begin{syntax}
    \cs{lbrace}
  \end{syntax}

\end{function}



\begin{function}[]
  {
    \lbrack
  }
  \begin{syntax}
    \cs{lbrack}
  \end{syntax}

\end{function}



\begin{function}[]
  {
    \lccode
  }
  \begin{syntax}
    \cs{lccode}
  \end{syntax}

\end{function}



\begin{function}[]
  {
    \ldots
  }
  \begin{syntax}
    \cs{ldots}
  \end{syntax}

\end{function}



\begin{function}[]
  {
    \leaders
  }
  \begin{syntax}
    \cs{leaders}
  \end{syntax}

\end{function}



\begin{function}[]
  {
    \leadsto
  }
  \begin{syntax}
    \cs{leadsto}
  \end{syntax}

\end{function}



\begin{function}[]
  {
    \leavevmode
  }
  \begin{syntax}
    \cs{leavevmode}
  \end{syntax}

\end{function}



\begin{function}[]
  {
    \leavevmode@ifvmode
  }
  \begin{syntax}
    \cs{leavevmode@ifvmode}
  \end{syntax}

\end{function}



\begin{function}[]
  {
    \left
  }
  \begin{syntax}
    \cs{left}
  \end{syntax}

\end{function}



\begin{function}[]
  {
    \lefteqn
  }
  \begin{syntax}
    \cs{lefteqn}
  \end{syntax}

\end{function}











\begin{function}[]
  {
    \leftmark
  }
  \begin{syntax}
    \cs{leftmark}
  \end{syntax}

\end{function}



\begin{function}[]
  {
    \leftskip
  }
  \begin{syntax}
    \cs{leftskip}
  \end{syntax}

\end{function}



\begin{function}[]
  {
    \legacyoldstylenums
  }
  \begin{syntax}
    \cs{legacyoldstylenums}
  \end{syntax}

\end{function}





\begin{function}[]
  {
    \let
  }
  \begin{syntax}
    \cs{let}
  \end{syntax}

\end{function}





\begin{function}[]
  {
    \lhd
  }
  \begin{syntax}
    \cs{lhd}
  \end{syntax}

\end{function}










\begin{function}[]
  {
    \line
  }
  \begin{syntax}
    \cs{line}
  \end{syntax}

\end{function}




\begin{function}[]
  {
    \linepenalty
  }
  \begin{syntax}
    \cs{linepenalty}
  \end{syntax}

\end{function}



\begin{function}[]
  {
    \lineskip
  }
  \begin{syntax}
    \cs{lineskip}
  \end{syntax}

\end{function}



\begin{function}[]
  {
    \lineskiplimit
  }
  \begin{syntax}
    \cs{lineskiplimit}
  \end{syntax}

\end{function}



\begin{function}[]
  {
    \linespread
  }
  \begin{syntax}
    \cs{linespread}
  \end{syntax}

\end{function}



\begin{function}[]
  {
    \linethickness
  }
  \begin{syntax}
    \cs{linethickness}
  \end{syntax}

\end{function}



\begin{function}[]
  {
    \linewidth
  }
  \begin{syntax}
    \cs{linewidth}
  \end{syntax}

\end{function}





\begin{function}[]
  {
    \listfiles
  }
  \begin{syntax}
    \cs{listfiles}
  \end{syntax}

\end{function}








\begin{function}[]
  {
    \load@onefile@withoptions
  }
  \begin{syntax}
    \cs{load@onefile@withoptions}
  \end{syntax}

\end{function}



\begin{function}[]
  {
    \load@onefilewithoptions
  }
  \begin{syntax}
    \cs{load@onefilewithoptions}
  \end{syntax}

\end{function}





\begin{function}[]
  {
    \loggingall
  }
  \begin{syntax}
    \cs{loggingall}
  \end{syntax}

\end{function}



\begin{function}[]
  {
    \loggingoutput
  }
  \begin{syntax}
    \cs{loggingoutput}
  \end{syntax}

\end{function}



\begin{function}[]
  {
    \long
  }
  \begin{syntax}
    \cs{long}
  \end{syntax}

\end{function}



\begin{function}[]
  {
    \loop
  }
  \begin{syntax}
    \cs{loop}
  \end{syntax}

\end{function}



\begin{function}[]
  {
    \lower
  }
  \begin{syntax}
    \cs{lower}
  \end{syntax}

\end{function}



\begin{function}[]
  {
    \lower@bound
  }
  \begin{syntax}
    \cs{lower@bound}
  \end{syntax}

\end{function}



\begin{function}[]
  {
    \lower@bound0
  }
  \begin{syntax}
    \cs{lower@bound0}
  \end{syntax}

\end{function}



\begin{function}[]
  {
    \lowercase
  }
  \begin{syntax}
    \cs{lowercase}
  \end{syntax}

\end{function}



\begin{function}[]
  {
    \lq
  }
  \begin{syntax}
    \cs{lq}
  \end{syntax}

\end{function}



\begin{function}[]
  {
    \lrbox
  }
  \begin{syntax}
    \cs{lrbox}
  \end{syntax}

\end{function}



\begin{function}[]
  {
    \ltcmddate
  }
  \begin{syntax}
    \cs{ltcmddate}
  \end{syntax}

\end{function}



\begin{function}[]
  {
    \ltcmdversion
  }
  \begin{syntax}
    \cs{ltcmdversion}
  \end{syntax}

\end{function}



\begin{function}[int]
  {
    \ltx@sh@ft
  }
  \begin{syntax}
    \cs{ltx@sh@ft}
  \end{syntax}

\end{function}



\begin{function}[int]
  {
    \ltx@star@counter
  }
  \begin{syntax}
    \cs{ltx@star@counter}
  \end{syntax}

\end{function}



\begin{function}[]
  {
    \luabytecode
  }
  \begin{syntax}
    \cs{luabytecode}
  \end{syntax}

\end{function}



\begin{function}[]
  {
    \luachunk
  }
  \begin{syntax}
    \cs{luachunk}
  \end{syntax}

\end{function}



\begin{function}[]
  {
    \luadef
  }
  \begin{syntax}
    \cs{luadef}
  \end{syntax}

\end{function}



\begin{function}[]
  {
    \luafunction
  }
  \begin{syntax}
    \cs{luafunction}
  \end{syntax}

\end{function}



\begin{function}[]
  {
    \luatexversion
  }
  \begin{syntax}
    \cs{luatexversion}
  \end{syntax}

\end{function}



\begin{function}[]
  {
    \m@ne
  }
  \begin{syntax}
    \cs{m@ne}
  \end{syntax}
  Constant (actually a counter) representing the integer \texttt{-1}
\end{function}



\begin{function}[]
  {
    \m@th
  }
  \begin{syntax}
    \cs{m@th}
  \end{syntax}

\end{function}



\begin{function}[]
  {
    \magstep
  }
  \begin{syntax}
    \cs{magstep}
  \end{syntax}

\end{function}



\begin{function}[]
  {
    \magstephalf
  }
  \begin{syntax}
    \cs{magstephalf}
  \end{syntax}

\end{function}



\begin{function}[]
  {
    \makeatletter
  }
  \begin{syntax}
    \cs{makeatletter}
  \end{syntax}

\end{function}



\begin{function}[]
  {
    \makeatother
  }
  \begin{syntax}
    \cs{makeatother}
  \end{syntax}

\end{function}



\begin{function}[]
  {
    \makebox
  }
  \begin{syntax}
    \cs{makebox}
  \end{syntax}

\end{function}



\begin{function}[]
  {
    \makeglossary
  }
  \begin{syntax}
    \cs{makeglossary}
  \end{syntax}

\end{function}



\begin{function}[]
  {
    \makeindex
  }
  \begin{syntax}
    \cs{makeindex}
  \end{syntax}

\end{function}



\begin{function}[]
  {
    \makelabel
  }
  \begin{syntax}
    \cs{makelabel}
  \end{syntax}

\end{function}



\begin{function}[]
  {
    \makeph@nt
  }
  \begin{syntax}
    \cs{makeph@nt}
  \end{syntax}

\end{function}



\begin{function}[]
  {
    \makesm@sh
  }
  \begin{syntax}
    \cs{makesm@sh}
  \end{syntax}

\end{function}



\begin{function}[]
  {
    \mandatory@arg
  }
  \begin{syntax}
    \cs{mandatory@arg}
  \end{syntax}

\end{function}



\begin{function}[]
  {
    \marginpar
  }
  \begin{syntax}
    \cs{marginpar}
  \end{syntax}

\end{function}



\begin{function}[]
  {
    \marginparpush
  }
  \begin{syntax}
    \cs{marginparpush}
  \end{syntax}

\end{function}



\begin{function}[]
  {
    \marginparsep
  }
  \begin{syntax}
    \cs{marginparsep}
  \end{syntax}

\end{function}



\begin{function}[]
  {
    \marginparwidth
  }
  \begin{syntax}
    \cs{marginparwidth}
  \end{syntax}

\end{function}



\begin{function}[]
  {
    \markboth
  }
  \begin{syntax}
    \cs{markboth}
  \end{syntax}

\end{function}



\begin{function}[]
  {
    \markright
  }
  \begin{syntax}
    \cs{markright}
  \end{syntax}

\end{function}



\begin{function}[]
  {
    \marks
  }
  \begin{syntax}
    \cs{marks}
  \end{syntax}

\end{function}





\begin{function}[]
  {
    \math@bgroup
  }
  \begin{syntax}
    \cs{math@bgroup}
  \end{syntax}

\end{function}



\begin{function}[]
  {
    \math@egroup
  }
  \begin{syntax}
    \cs{math@egroup}
  \end{syntax}

\end{function}



\begin{function}[]
  {
    \math@fonts
  }
  \begin{syntax}
    \cs{math@fonts}
  \end{syntax}

\end{function}



\begin{function}[]
  {
    \math@fontsfalse
  }
  \begin{syntax}
    \cs{math@fontsfalse}
  \end{syntax}

\end{function}



\begin{function}[]
  {
    \math@fontstrue
  }
  \begin{syntax}
    \cs{math@fontstrue}
  \end{syntax}

\end{function}



\begin{function}[]
  {
    \math@version
  }
  \begin{syntax}
    \cs{math@version}
  \end{syntax}

\end{function}



\begin{function}[]
  {
    \mathaccent
  }
  \begin{syntax}
    \cs{mathaccent}
  \end{syntax}

\end{function}



\begin{function}[]
  {
    \mathalpha
  }
  \begin{syntax}
    \cs{mathalpha}
  \end{syntax}

\end{function}



\begin{function}[]
  {
    \mathbf
  }
  \begin{syntax}
    \cs{mathbf}
  \end{syntax}

\end{function}



\begin{function}[]
  {
    \mathbin
  }
  \begin{syntax}
    \cs{mathbin}
  \end{syntax}

\end{function}



\begin{function}[]
  {
    \mathchar
  }
  \begin{syntax}
    \cs{mathchar}
  \end{syntax}

\end{function}



\begin{function}[]
  {
    \mathchar@type
  }
  \begin{syntax}
    \cs{mathchar@type}
  \end{syntax}

\end{function}



\begin{function}[]
  {
    \mathchardef
  }
  \begin{syntax}
    \cs{mathchardef}
  \end{syntax}

\end{function}



\begin{function}[]
  {
    \mathcharzero
  }
  \begin{syntax}
    \cs{mathcharzero}
  \end{syntax}

\end{function}



\begin{function}[]
  {
    \mathchoice
  }
  \begin{syntax}
    \cs{mathchoice}
  \end{syntax}

\end{function}



\begin{function}[]
  {
    \mathclose
  }
  \begin{syntax}
    \cs{mathclose}
  \end{syntax}

\end{function}



\begin{function}[]
  {
    \mathdefaultsmode
  }
  \begin{syntax}
    \cs{mathdefaultsmode}
  \end{syntax}

\end{function}



\begin{function}[]
  {
    \mathdollar
  }
  \begin{syntax}
    \cs{mathdollar}
  \end{syntax}

\end{function}



\begin{function}[]
  {
    \mathellipsis
  }
  \begin{syntax}
    \cs{mathellipsis}
  \end{syntax}

\end{function}



\begin{function}[]
  {
    \mathgroup
  }
  \begin{syntax}
    \cs{mathgroup}
  \end{syntax}

\end{function}



\begin{function}[]
  {
    \mathhexbox
  }
  \begin{syntax}
    \cs{mathhexbox}
  \end{syntax}

\end{function}



\begin{function}[]
  {
    \mathit
  }
  \begin{syntax}
    \cs{mathit}
  \end{syntax}

\end{function}



\begin{function}[]
  {
    \mathop
  }
  \begin{syntax}
    \cs{mathop}
  \end{syntax}

\end{function}



\begin{function}[]
  {
    \mathopen
  }
  \begin{syntax}
    \cs{mathopen}
  \end{syntax}

\end{function}



\begin{function}[]
  {
    \mathord
  }
  \begin{syntax}
    \cs{mathord}
  \end{syntax}

\end{function}



\begin{function}[]
  {
    \mathpalette
  }
  \begin{syntax}
    \cs{mathpalette}
  \end{syntax}

\end{function}



\begin{function}[]
  {
    \mathparagraph
  }
  \begin{syntax}
    \cs{mathparagraph}
  \end{syntax}

\end{function}



\begin{function}[]
  {
    \mathph@nt
  }
  \begin{syntax}
    \cs{mathph@nt}
  \end{syntax}

\end{function}



\begin{function}[]
  {
    \mathpunct
  }
  \begin{syntax}
    \cs{mathpunct}
  \end{syntax}

\end{function}



\begin{function}[]
  {
    \mathrel
  }
  \begin{syntax}
    \cs{mathrel}
  \end{syntax}

\end{function}



\begin{function}[]
  {
    \mathrm
  }
  \begin{syntax}
    \cs{mathrm}
  \end{syntax}

\end{function}



\begin{function}[]
  {
    \mathsection
  }
  \begin{syntax}
    \cs{mathsection}
  \end{syntax}

\end{function}



\begin{function}[]
  {
    \mathsf
  }
  \begin{syntax}
    \cs{mathsf}
  \end{syntax}

\end{function}



\begin{function}[]
  {
    \mathsm@sh
  }
  \begin{syntax}
    \cs{mathsm@sh}
  \end{syntax}

\end{function}



\begin{function}[]
  {
    \mathsterling
  }
  \begin{syntax}
    \cs{mathsterling}
  \end{syntax}

\end{function}



\begin{function}[]
  {
    \mathstrut
  }
  \begin{syntax}
    \cs{mathstrut}
  \end{syntax}

\end{function}



\begin{function}[]
  {
    \mathsurround
  }
  \begin{syntax}
    \cs{mathsurround}
  \end{syntax}

\end{function}



\begin{function}[]
  {
    \mathtt
  }
  \begin{syntax}
    \cs{mathtt}
  \end{syntax}

\end{function}



\begin{function}[]
  {
    \mathversion
  }
  \begin{syntax}
    \cs{mathversion}
  \end{syntax}

\end{function}



\begin{function}[]
  {
    \matrix
  }
  \begin{syntax}
    \cs{matrix}
  \end{syntax}

\end{function}





\begin{function}[]
  {
    \maxdeadcycles
  }
  \begin{syntax}
    \cs{maxdeadcycles}
  \end{syntax}

\end{function}



\begin{function}[]
  {
    \maxdepth
  }
  \begin{syntax}
    \cs{maxdepth}
  \end{syntax}

\end{function}



\begin{function}[]
  {
    \maxdimen
  }
  \begin{syntax}
    \cs{maxdimen}
  \end{syntax}

\end{function}



\begin{function}[]
  {
    \maybe@ic
  }
  \begin{syntax}
    \cs{maybe@ic}
  \end{syntax}

\end{function}



\begin{function}[]
  {
    \maybe@ic@
  }
  \begin{syntax}
    \cs{maybe@ic@}
  \end{syntax}

\end{function}



\begin{function}[]
  {
    \maybe@icfalse
  }
  \begin{syntax}
    \cs{maybe@icfalse}
  \end{syntax}

\end{function}



\begin{function}[]
  {
    \maybe@ictrue
  }
  \begin{syntax}
    \cs{maybe@ictrue}
  \end{syntax}

\end{function}



\begin{function}[]
  {
    \maybe@load@fontshape
  }
  \begin{syntax}
    \cs{maybe@load@fontshape}
  \end{syntax}

\end{function}



\begin{function}[]
  {
    \maybe@update@bfseries@defaults
  }
  \begin{syntax}
    \cs{maybe@update@bfseries@defaults}
  \end{syntax}

\end{function}



\begin{function}[]
  {
    \maybe@update@mdseries@defaults
  }
  \begin{syntax}
    \cs{maybe@update@mdseries@defaults}
  \end{syntax}

\end{function}



\begin{function}[]
  {
    \mb@b
  }
  \begin{syntax}
    \cs{mb@b}
  \end{syntax}

\end{function}



\begin{function}[]
  {
    \mb@l
  }
  \begin{syntax}
    \cs{mb@l}
  \end{syntax}

\end{function}



\begin{function}[]
  {
    \mb@r
  }
  \begin{syntax}
    \cs{mb@r}
  \end{syntax}

\end{function}



\begin{function}[]
  {
    \mb@t
  }
  \begin{syntax}
    \cs{mb@t}
  \end{syntax}

\end{function}



\begin{function}[]
  {
    \mbox
  }
  \begin{syntax}
    \cs{mbox}
  \end{syntax}

\end{function}



\begin{function}[]
  {
    \mddef@ult
  }
  \begin{syntax}
    \cs{mddef@ult}
  \end{syntax}

\end{function}



\begin{function}[]
  {
    \mddefault
  }
  \begin{syntax}
    \cs{mddefault}
  \end{syntax}

\end{function}



\begin{function}[]
  {
    \mddefault@previous
  }
  \begin{syntax}
    \cs{mddefault@previous}
  \end{syntax}

\end{function}



\begin{function}[]
  {
    \mdseries
  }
  \begin{syntax}
    \cs{mdseries}
  \end{syntax}

\end{function}



\begin{function}[]
  {
    \mdseries@rm
  }
  \begin{syntax}
    \cs{mdseries@rm}
  \end{syntax}

\end{function}



\begin{function}[]
  {
    \mdseries@sf
  }
  \begin{syntax}
    \cs{mdseries@sf}
  \end{syntax}

\end{function}



\begin{function}[]
  {
    \mdseries@tt
  }
  \begin{syntax}
    \cs{mdseries@tt}
  \end{syntax}

\end{function}



\begin{function}[]
  {
    \meaning
  }
  \begin{syntax}
    \cs{meaning}
  \end{syntax}

\end{function}



\begin{function}[]
  {
    \medbreak
  }
  \begin{syntax}
    \cs{medbreak}
  \end{syntax}

\end{function}



\begin{function}[]
  {
    \medmuskip
  }
  \begin{syntax}
    \cs{medmuskip}
  \end{syntax}

\end{function}








\begin{function}[]
  {
    \merge@font@series
  }
  \begin{syntax}
    \cs{merge@font@series}
  \end{syntax}

\end{function}



\begin{function}[]
  {
    \merge@font@series@
  }
  \begin{syntax}
    \cs{merge@font@series@}
  \end{syntax}

\end{function}



\begin{function}[]
  {
    \merge@font@series@without@substitution
  }
  \begin{syntax}
    \cs{merge@font@series@without@substitution}
  \end{syntax}

\end{function}



\begin{function}[]
  {
    \merge@font@series@without@substitution@
  }
  \begin{syntax}
    \cs{merge@font@series@without@substitution@}
  \end{syntax}

\end{function}



\begin{function}[]
  {
    \merge@font@shape
  }
  \begin{syntax}
    \cs{merge@font@shape}
  \end{syntax}

\end{function}



\begin{function}[]
  {
    \merge@font@shape@
  }
  \begin{syntax}
    \cs{merge@font@shape@}
  \end{syntax}

\end{function}



\begin{function}[]
  {
    \merge@font@shape@without@substitution
  }
  \begin{syntax}
    \cs{merge@font@shape@without@substitution}
  \end{syntax}

\end{function}



\begin{function}[]
  {
    \merge@font@shape@without@substitution@
  }
  \begin{syntax}
    \cs{merge@font@shape@without@substitution@}
  \end{syntax}

\end{function}





\begin{function}[]
  {
    \mho
  }
  \begin{syntax}
    \cs{mho}
  \end{syntax}

\end{function}





\begin{function}[]
  {
    \minipage
  }
  \begin{syntax}
    \cs{minipage}
  \end{syntax}

\end{function}



\begin{function}[]
  {
    \mkern
  }
  \begin{syntax}
    \cs{mkern}
  \end{syntax}

\end{function}



\begin{function}[]
  {
    \month
  }
  \begin{syntax}
    \cs{month}
  \end{syntax}

\end{function}



\begin{function}[]
  {
    \month/
  }
  \begin{syntax}
    \cs{month/}
  \end{syntax}

\end{function}



\begin{function}[]
  {
    \mskip
  }
  \begin{syntax}
    \cs{mskip}
  \end{syntax}

\end{function}



\begin{function}[]
  {
    \mubyte
  }
  \begin{syntax}
    \cs{mubyte}
  \end{syntax}

\end{function}



\begin{function}[]
  {
    \multicolumn
  }
  \begin{syntax}
    \cs{multicolumn}
  \end{syntax}

\end{function}



\begin{function}[]
  {
    \multiply
  }
  \begin{syntax}
    \cs{multiply}
  \end{syntax}

\end{function}



\begin{function}[]
  {
    \multiput
  }
  \begin{syntax}
    \cs{multiput}
  \end{syntax}

\end{function}



\begin{function}[]
  {
    \multispan
  }
  \begin{syntax}
    \cs{multispan}
  \end{syntax}

\end{function}



\begin{function}[]
  {
    \muskip
  }
  \begin{syntax}
    \cs{muskip}
  \end{syntax}

\end{function}



\begin{function}[]
  {
    \muskip_eval:n
  }
  \begin{syntax}
    \cs{muskip_eval:n}
  \end{syntax}

\end{function}



\begin{function}[]
  {
    \muskip_new:N
  }
  \begin{syntax}
    \cs{muskip_new:N}
  \end{syntax}

\end{function}



\begin{function}[]
  {
    \muskipdef
  }
  \begin{syntax}
    \cs{muskipdef}
  \end{syntax}

\end{function}



\begin{function}[]
  {
    \muskipzero
  }
  \begin{syntax}
    \cs{muskipzero}
  \end{syntax}

\end{function}



\begin{function}[]
  {
    \n
  }
  \begin{syntax}
    \cs{n}
  \end{syntax}

\end{function}



\begin{function}[]
  {
    \narrower
  }
  \begin{syntax}
    \cs{narrower}
  \end{syntax}

\end{function}



\begin{function}[]
  {
    \negmedspace
  }
  \begin{syntax}
    \cs{negmedspace}
  \end{syntax}

\end{function}



\begin{function}[]
  {
    \negthickspace
  }
  \begin{syntax}
    \cs{negthickspace}
  \end{syntax}

\end{function}



\begin{function}[]
  {
    \negthinspace
  }
  \begin{syntax}
    \cs{negthinspace}
  \end{syntax}

\end{function}



\begin{function}[]
  {
    \new@command
  }
  \begin{syntax}
    \cs{new@command}
  \end{syntax}

\end{function}



\begin{function}[]
  {
    \new@environment
  }
  \begin{syntax}
    \cs{new@environment}
  \end{syntax}

\end{function}



\begin{function}[]
  {
    \new@label@record
  }
  \begin{syntax}
    \cs{new@label@record}
  \end{syntax}

\end{function}



\begin{function}[]
  {
    \new@mathalphabet
  }
  \begin{syntax}
    \cs{new@mathalphabet}
  \end{syntax}

\end{function}



\begin{function}[]
  {
    \new@mathgroup
  }
  \begin{syntax}
    \cs{new@mathgroup}
  \end{syntax}

\end{function}



\begin{function}[]
  {
    \new@mathversion
  }
  \begin{syntax}
    \cs{new@mathversion}
  \end{syntax}

\end{function}





\begin{function}[]
  {
    \new@symbolfont
  }
  \begin{syntax}
    \cs{new@symbolfont}
  \end{syntax}

\end{function}



\begin{function}[]
  {
    \newXeTeXintercharclass
  }
  \begin{syntax}
    \cs{newXeTeXintercharclass}
  \end{syntax}

\end{function}



\begin{function}[]
  {
    \newattribute
  }
  \begin{syntax}
    \cs{newattribute}
  \end{syntax}

\end{function}





\begin{function}[]
  {
    \newcatcodetable
  }
  \begin{syntax}
    \cs{newcatcodetable}
  \end{syntax}

\end{function}



\begin{function}[]
  {
    \newcommand
  }
  \begin{syntax}
    \cs{newcommand}
  \end{syntax}

\end{function}







\begin{function}[]
  {
    \newcounteralias
  }
  \begin{syntax}
    \cs{newcounteralias}
  \end{syntax}

\end{function}






\begin{function}[]
  {
    \newenvironment
  }
  \begin{syntax}
    \cs{newenvironment}
  \end{syntax}

\end{function}



\begin{function}[]
  {
    \newfam
  }
  \begin{syntax}
    \cs{newfam}
  \end{syntax}

\end{function}



\begin{function}[]
  {
    \newfont
  }
  \begin{syntax}
    \cs{newfont}
  \end{syntax}

\end{function}



\begin{function}[]
  {
    \newhelp
  }
  \begin{syntax}
    \cs{newhelp}
  \end{syntax}

\end{function}



\begin{function}[]
  {
    \newif
  }
  \begin{syntax}
    \cs{newif}
  \end{syntax}

\end{function}





\begin{function}[]
  {
    \newlabel
  }
  \begin{syntax}
    \cs{newlabel}
  \end{syntax}

\end{function}







\begin{function}[]
  {
    \newline
  }
  \begin{syntax}
    \cs{newline}
  \end{syntax}

\end{function}



\begin{function}[]
  {
    \newlinechar
  }
  \begin{syntax}
    \cs{newlinechar}
  \end{syntax}

\end{function}



\begin{function}[]
  {
    \newluabytecode
  }
  \begin{syntax}
    \cs{newluabytecode}
  \end{syntax}

\end{function}



\begin{function}[]
  {
    \newluachunkname
  }
  \begin{syntax}
    \cs{newluachunkname}
  \end{syntax}

\end{function}



\begin{function}[]
  {
    \newluacmd
  }
  \begin{syntax}
    \cs{newluacmd}
  \end{syntax}

\end{function}



\begin{function}[]
  {
    \newluafunction
  }
  \begin{syntax}
    \cs{newluafunction}
  \end{syntax}

\end{function}



\begin{function}[]
  {
    \newmarks
  }
  \begin{syntax}
    \cs{newmarks}
  \end{syntax}

\end{function}





\begin{function}[]
  {
    \newpage
  }
  \begin{syntax}
    \cs{newpage}
  \end{syntax}

\end{function}



\begin{function}[]
  {
    \newprotectedluacmd
  }
  \begin{syntax}
    \cs{newprotectedluacmd}
  \end{syntax}

\end{function}





\begin{function}[]
  {
    \newsavebox
  }
  \begin{syntax}
    \cs{newsavebox}
  \end{syntax}

\end{function}





\begin{function}[]
  {
    \newtheorem
  }
  \begin{syntax}
    \cs{newtheorem}
  \end{syntax}

\end{function}



\begin{function}[]
  {
    \newtie
  }
  \begin{syntax}
    \cs{newtie}
  \end{syntax}

\end{function}





\begin{function}[]
  {
    \newwhatsit
  }
  \begin{syntax}
    \cs{newwhatsit}
  \end{syntax}

\end{function}





\begin{function}[]
  {
    \nfss@catcodes
  }
  \begin{syntax}
    \cs{nfss@catcodes}
  \end{syntax}

\end{function}



\begin{function}[]
  {
    \nfss@text
  }
  \begin{syntax}
    \cs{nfss@text}
  \end{syntax}

\end{function}



\begin{function}[]
  {
    \ng
  }
  \begin{syntax}
    \cs{ng}
  \end{syntax}

\end{function}



\begin{function}[]
  {
    \no@alphabet@error
  }
  \begin{syntax}
    \cs{no@alphabet@error}
  \end{syntax}

\end{function}



\begin{function}[]
  {
    \noaccents@
  }
  \begin{syntax}
    \cs{noaccents@}
  \end{syntax}

\end{function}



\begin{function}[]
  {
    \noalign
  }
  \begin{syntax}
    \cs{noalign}
  \end{syntax}

\end{function}



\begin{function}[]
  {
    \nobreak
  }
  \begin{syntax}
    \cs{nobreak}
  \end{syntax}

\end{function}



\begin{function}[]
  {
    \nobreakdashes
  }
  \begin{syntax}
    \cs{nobreakdashes}
  \end{syntax}

\end{function}



\begin{function}[]
  {
    \nobreakspace
  }
  \begin{syntax}
    \cs{nobreakspace}
  \end{syntax}

\end{function}



\begin{function}[]
  {
    \nocite
  }
  \begin{syntax}
    \cs{nocite}
  \end{syntax}

\end{function}



\begin{function}[]
  {
    \nocorr
  }
  \begin{syntax}
    \cs{nocorr}
  \end{syntax}

\end{function}



\begin{function}[]
  {
    \nocorrlist
  }
  \begin{syntax}
    \cs{nocorrlist}
  \end{syntax}

\end{function}



\begin{function}[]
  {
    \noexpand
  }
  \begin{syntax}
    \cs{noexpand}
  \end{syntax}

\end{function}



\begin{function}[]
  {
    \nofiles
  }
  \begin{syntax}
    \cs{nofiles}
  \end{syntax}

\end{function}



\begin{function}[]
  {
    \noindent
  }
  \begin{syntax}
    \cs{noindent}
  \end{syntax}

\end{function}



\begin{function}[]
  {
    \nointerlineskip
  }
  \begin{syntax}
    \cs{nointerlineskip}
  \end{syntax}

\end{function}



\begin{function}[]
  {
    \nolimits
  }
  \begin{syntax}
    \cs{nolimits}
  \end{syntax}

\end{function}



\begin{function}[]
  {
    \non@alpherr
  }
  \begin{syntax}
    \cs{non@alpherr}
  \end{syntax}

\end{function}



\begin{function}[]
  {
    \nonfrenchspacing
  }
  \begin{syntax}
    \cs{nonfrenchspacing}
  \end{syntax}

\end{function}



\begin{function}[]
  {
    \nonscript
  }
  \begin{syntax}
    \cs{nonscript}
  \end{syntax}

\end{function}



\begin{function}[]
  {
    \nonumber
  }
  \begin{syntax}
    \cs{nonumber}
  \end{syntax}

\end{function}




\begin{function}[]
  {
    \noprotrusion
  }
  \begin{syntax}
    \cs{noprotrusion}
  \end{syntax}

\end{function}



\begin{function}[]
  {
    \normalbaselines
  }
  \begin{syntax}
    \cs{normalbaselines}
  \end{syntax}

\end{function}



\begin{function}[]
  {
    \normalbaselineskip
  }
  \begin{syntax}
    \cs{normalbaselineskip}
  \end{syntax}

\end{function}



\begin{function}[]
  {
    \normalcolor
  }
  \begin{syntax}
    \cs{normalcolor}
  \end{syntax}

\end{function}



\begin{function}[]
  {
    \normalfont
  }
  \begin{syntax}
    \cs{normalfont}
  \end{syntax}

\end{function}



\begin{function}[]
  {
    \normallineskip
  }
  \begin{syntax}
    \cs{normallineskip}
  \end{syntax}

\end{function}



\begin{function}[]
  {
    \normallineskiplimit
  }
  \begin{syntax}
    \cs{normallineskiplimit}
  \end{syntax}

\end{function}



\begin{function}[]
  {
    \normalmarginpar
  }
  \begin{syntax}
    \cs{normalmarginpar}
  \end{syntax}

\end{function}



\begin{function}[]
  {
    \normalsfcodes
  }
  \begin{syntax}
    \cs{normalsfcodes}
  \end{syntax}

\end{function}



\begin{function}[]
  {
    \normalshape
  }
  \begin{syntax}
    \cs{normalshape}
  \end{syntax}

\end{function}



\begin{function}[]
  {
    \normalsize
  }
  \begin{syntax}
    \cs{normalsize}
  \end{syntax}

\end{function}



\begin{function}[]
  {
    \not@base
  }
  \begin{syntax}
    \cs{not@base}
  \end{syntax}

\end{function}



\begin{function}[]
  {
    \not@math@alphabet
  }
  \begin{syntax}
    \cs{not@math@alphabet}
  \end{syntax}

\end{function}



\begin{function}[]
  {
    \now@and@everyjob
  }
  \begin{syntax}
    \cs{now@and@everyjob}
  \end{syntax}

\end{function}



\begin{function}[]
  {
    \null
  }
  \begin{syntax}
    \cs{null}
  \end{syntax}

\end{function}



\begin{function}[]
  {
    \nulldelimiterspace
  }
  \begin{syntax}
    \cs{nulldelimiterspace}
  \end{syntax}

\end{function}



\begin{function}[]
  {
    \nullfont
  }
  \begin{syntax}
    \cs{nullfont}
  \end{syntax}

\end{function}



\begin{function}[]
  {
    \number
  }
  \begin{syntax}
    \cs{number}
  \end{syntax}

\end{function}



\begin{function}[]
  {
    \numberline
  }
  \begin{syntax}
    \cs{numberline}
  \end{syntax}

\end{function}



\begin{function}[]
  {
    \numexpr
  }
  \begin{syntax}
    \cs{numexpr}
  \end{syntax}

\end{function}



\begin{function}[]
  {
    \o
  }
  \begin{syntax}
    \cs{o}
  \end{syntax}

\end{function}



\begin{function}[]
  {
    \o@lign
  }
  \begin{syntax}
    \cs{o@lign}
  \end{syntax}

\end{function}



\begin{function}[]
  {
    \oalign
  }
  \begin{syntax}
    \cs{oalign}
  \end{syntax}

\end{function}



\begin{function}[]
  {
    \obeycr
  }
  \begin{syntax}
    \cs{obeycr}
  \end{syntax}

\end{function}



\begin{function}[]
  {
    \obeyedline
  }
  \begin{syntax}
    \cs{obeyedline}
  \end{syntax}

\end{function}



\begin{function}[]
  {
    \obeyedspace
  }
  \begin{syntax}
    \cs{obeyedspace}
  \end{syntax}

\end{function}



\begin{function}[]
  {
    \obeylines
  }
  \begin{syntax}
    \cs{obeylines}
  \end{syntax}

\end{function}



\begin{function}[]
  {
    \obeyspaces
  }
  \begin{syntax}
    \cs{obeyspaces}
  \end{syntax}

\end{function}



\begin{function}[]
  {
    \oddsidemargin
  }
  \begin{syntax}
    \cs{oddsidemargin}
  \end{syntax}

\end{function}



\begin{function}[]
  {
    \oe
  }
  \begin{syntax}
    \cs{oe}
  \end{syntax}

\end{function}



\begin{function}[]
  {
    \of
  }
  \begin{syntax}
    \cs{of}
  \end{syntax}

\end{function}



\begin{function}[]
  {
    \offinterlineskip
  }
  \begin{syntax}
    \cs{offinterlineskip}
  \end{syntax}

\end{function}



\begin{function}[]
  {
    \oldstylenums
  }
  \begin{syntax}
    \cs{oldstylenums}
  \end{syntax}

\end{function}



\begin{function}[]
  {
    \omit
  }
  \begin{syntax}
    \cs{omit}
  \end{syntax}

\end{function}



\begin{function}[]
  {
    \on@line
  }
  \begin{syntax}
    \cs{on@line}
  \end{syntax}

\end{function}



\begin{function}[]
  {
    \onecolumn
  }
  \begin{syntax}
    \cs{onecolumn}
  \end{syntax}

\end{function}



\begin{function}[]
  {
    \ooalign
  }
  \begin{syntax}
    \cs{ooalign}
  \end{syntax}

\end{function}





\begin{function}[]
  {
    \openout15
  }
  \begin{syntax}
    \cs{openout15}
  \end{syntax}

\end{function}



\begin{function}[]
  {
    \openup
  }
  \begin{syntax}
    \cs{openup}
  \end{syntax}

\end{function}



\begin{function}[]
  {
    \operator@font
  }
  \begin{syntax}
    \cs{operator@font}
  \end{syntax}

\end{function}



\begin{function}[]
  {
    \optional@arg
  }
  \begin{syntax}
    \cs{optional@arg}
  \end{syntax}

\end{function}



\begin{function}[]
  {
    \or
  }
  \begin{syntax}
    \cs{or}
  \end{syntax}

\end{function}



\begin{function}[]
  {
    \outer@nobreak
  }
  \begin{syntax}
    \cs{outer@nobreak}
  \end{syntax}

\end{function}



\begin{function}[]
  {
    \output
  }
  \begin{syntax}
    \cs{output}
  \end{syntax}

\end{function}



\begin{function}[]
  {
    \outputmode
  }
  \begin{syntax}
    \cs{outputmode}
  \end{syntax}

\end{function}



\begin{function}[]
  {
    \outputpenalty
  }
  \begin{syntax}
    \cs{outputpenalty}
  \end{syntax}

\end{function}



\begin{function}[]
  {
    \oval
  }
  \begin{syntax}
    \cs{oval}
  \end{syntax}

\end{function}



\begin{function}[]
  {
    \over
  }
  \begin{syntax}
    \cs{over}
  \end{syntax}

\end{function}



\begin{function}[]
  {
    \overfullrule
  }
  \begin{syntax}
    \cs{overfullrule}
  \end{syntax}

\end{function}



\begin{function}[]
  {
    \p@
  }
  \begin{syntax}
    \cs{p@}
  \end{syntax}

\end{function}



\begin{function}[]
  {
    \p@equation
  }
  \begin{syntax}
    \cs{p@equation}
  \end{syntax}

\end{function}





\begin{function}[]
  {
    \pagenumbering
  }
  \begin{syntax}
    \cs{pagenumbering}
  \end{syntax}

\end{function}



\begin{function}[]
  {
    \pageref
  }
  \begin{syntax}
    \cs{pageref}
  \end{syntax}

\end{function}



\begin{function}[]
  {
    \pageshrink
  }
  \begin{syntax}
    \cs{pageshrink}
  \end{syntax}

\end{function}



\begin{function}[]
  {
    \pagestyle
  }
  \begin{syntax}
    \cs{pagestyle}
  \end{syntax}

\end{function}



\begin{function}[]
  {
    \pagetotal
  }
  \begin{syntax}
    \cs{pagetotal}
  \end{syntax}

\end{function}



\begin{function}[]
  {
    \paperheight
  }
  \begin{syntax}
    \cs{paperheight}
  \end{syntax}

\end{function}



\begin{function}[]
  {
    \paperwidth
  }
  \begin{syntax}
    \cs{paperwidth}
  \end{syntax}

\end{function}



\begin{function}[]
  {
    \par
  }
%  \begin{syntax}
%    \cs{par}
%  \end{syntax}

\end{function}



\begin{function}[]
  {
    \par@deathcycles
  }
  \begin{syntax}
    \cs{par@deathcycles}
  \end{syntax}

\end{function}



\begin{function}[]
  {
    \para_end:
  }
  \begin{syntax}
    \cs{para_end:}
  \end{syntax}

\end{function}



\begin{function}[]
  {
    \para_raw_indent:
  }
  \begin{syntax}
    \cs{para_raw_indent:}
  \end{syntax}

\end{function}



\begin{function}[]
  {
    \para_raw_noindent:
  }
  \begin{syntax}
    \cs{para_raw_noindent:}
  \end{syntax}

\end{function}



\begin{function}[]
  {
    \paragraphmark
  }
  \begin{syntax}
    \cs{paragraphmark}
  \end{syntax}

\end{function}



\begin{function}[]
  {
    \parbox
  }
  \begin{syntax}
    \cs{parbox}
  \end{syntax}

\end{function}



\begin{function}[]
  {
    \parfillskip
  }
  \begin{syntax}
    \cs{parfillskip}
  \end{syntax}

\end{function}



\begin{function}[]
  {
    \parindent
  }
  \begin{syntax}
    \cs{parindent}
  \end{syntax}

\end{function}






\begin{function}[]
  {
    \parseunicodedataI
  }
  \begin{syntax}
    \cs{parseunicodedataI}
  \end{syntax}

\end{function}



\begin{function}[]
  {
    \parseunicodedataII
  }
  \begin{syntax}
    \cs{parseunicodedataII}
  \end{syntax}

\end{function}



\begin{function}[]
  {
    \parseunicodedataIII
  }
  \begin{syntax}
    \cs{parseunicodedataIII}
  \end{syntax}

\end{function}



\begin{function}[]
  {
    \parseunicodedataIV
  }
  \begin{syntax}
    \cs{parseunicodedataIV}
  \end{syntax}

\end{function}



\begin{function}[]
  {
    \parseunicodedataV
  }
  \begin{syntax}
    \cs{parseunicodedataV}
  \end{syntax}

\end{function}



\begin{function}[]
  {
    \parshape
  }
  \begin{syntax}
    \cs{parshape}
  \end{syntax}

\end{function}



\begin{function}[]
  {
    \parskip
  }
  \begin{syntax}
    \cs{parskip}
  \end{syntax}

\end{function}



\begin{function}[]
  {
    \partokencontext
  }
  \begin{syntax}
    \cs{partokencontext}
  \end{syntax}

\end{function}





\begin{function}[]
  {
    \patch@level
  }
  \begin{syntax}
    \cs{patch@level}
  \end{syntax}

\end{function}



\begin{function}[]
  {
    \pdfannot@link@off@@
  }
  \begin{syntax}
    \cs{pdfannot@link@off@@}
  \end{syntax}

\end{function}



\begin{function}[]
  {
    \pdfannot@link@on@@
  }
  \begin{syntax}
    \cs{pdfannot@link@on@@}
  \end{syntax}

\end{function}



\begin{function}[]
  {
    \pdfextension
  }
  \begin{syntax}
    \cs{pdfextension}
  \end{syntax}

\end{function}



\begin{function}[]
  {
    \pdffilesize
  }
  \begin{syntax}
    \cs{pdffilesize}
  \end{syntax}

\end{function}



\begin{function}[]
  {
    \pdfgentounicode
  }
  \begin{syntax}
    \cs{pdfgentounicode}
  \end{syntax}

\end{function}



\begin{function}[]
  {
    \pdfglyphtounicode
  }
  \begin{syntax}
    \cs{pdfglyphtounicode}
  \end{syntax}

\end{function}



\begin{function}[]
  {
    \pdfhorigin
  }
  \begin{syntax}
    \cs{pdfhorigin}
  \end{syntax}

\end{function}



\begin{function}[]
  {
    \pdftexrevision
  }
  \begin{syntax}
    \cs{pdftexrevision}
  \end{syntax}

\end{function}



\begin{function}[]
  {
    \pdftexversion
  }
  \begin{syntax}
    \cs{pdftexversion}
  \end{syntax}

\end{function}



\begin{function}[]
  {
    \pdfvariable
  }
  \begin{syntax}
    \cs{pdfvariable}
  \end{syntax}

\end{function}



\begin{function}[]
  {
    \pdfvorigin
  }
  \begin{syntax}
    \cs{pdfvorigin}
  \end{syntax}

\end{function}



\begin{function}[]
  {
    \penalty
  }
  \begin{syntax}
    \cs{penalty}
  \end{syntax}

\end{function}



\begin{function}[]
  {
    \ph@nt
  }
  \begin{syntax}
    \cs{ph@nt}
  \end{syntax}

\end{function}



\begin{function}[]
  {
    \phantom
  }
  \begin{syntax}
    \cs{phantom}
  \end{syntax}

\end{function}



\begin{function}[]
  {
    \pickup@font
  }
  \begin{syntax}
    \cs{pickup@font}
  \end{syntax}

\end{function}



\begin{function}[]
  {
    \pictur@
  }
  \begin{syntax}
    \cs{pictur@}
  \end{syntax}

\end{function}



\begin{function}[]
  {
    \picture
  }
  \begin{syntax}
    \cs{picture}
  \end{syntax}

\end{function}




\begin{function}[]
  {
    \poptabs
  }
  \begin{syntax}
    \cs{poptabs}
  \end{syntax}

\end{function}



\begin{function}[]
  {
    \postdisplaypenalty
  }
  \begin{syntax}
    \cs{postdisplaypenalty}
  \end{syntax}

\end{function}



\begin{function}[]
  {
    \pounds
  }
  \begin{syntax}
    \cs{pounds}
  \end{syntax}

\end{function}



\begin{function}[]
  {
    \pr@@@s
  }
  \begin{syntax}
    \cs{pr@@@s}
  \end{syntax}

\end{function}



\begin{function}[]
  {
    \pr@@@t
  }
  \begin{syntax}
    \cs{pr@@@t}
  \end{syntax}

\end{function}



\begin{function}[]
  {
    \pr@m@s
  }
  \begin{syntax}
    \cs{pr@m@s}
  \end{syntax}

\end{function}



\begin{function}[]
  {
    \pr@tectedrel@x
  }
  \begin{syntax}
    \cs{pr@tectedrel@x}
  \end{syntax}

\end{function}



\begin{function}[]
  {
    \predisplaypenalty
  }
  \begin{syntax}
    \cs{predisplaypenalty}
  \end{syntax}

\end{function}



\begin{function}[]
  {
    \prepare@family@series@update
  }
  \begin{syntax}
    \cs{prepare@family@series@update}
  \end{syntax}

\end{function}



\begin{function}[]
  {
    \pretolerance
  }
  \begin{syntax}
    \cs{pretolerance}
  \end{syntax}

\end{function}



\begin{function}[]
  {
    \prevdepth
  }
  \begin{syntax}
    \cs{prevdepth}
  \end{syntax}

\end{function}



\begin{function}[]
  {
    \prim@s
  }
  \begin{syntax}
    \cs{prim@s}
  \end{syntax}

\end{function}



\begin{function}[]
  {
    \prime
  }
  \begin{syntax}
    \cs{prime}
  \end{syntax}

\end{function}



\begin{function}[]
  {
    \process@table
  }
  \begin{syntax}
    \cs{process@table}
  \end{syntax}

\end{function}



\begin{function}[]
  {
    \protect
  }
  \begin{syntax}
    \cs{protect}
  \end{syntax}

\end{function}



\begin{function}[]
  {
    \protected
  }
  \begin{syntax}
    \cs{protected}
  \end{syntax}

\end{function}



\begin{function}[]
  {
    \protected@edef
  }
  \begin{syntax}
    \cs{protected@edef}
  \end{syntax}

\end{function}



\begin{function}[]
  {
    \protected@file@percent
  }
  \begin{syntax}
    \cs{protected@file@percent}
  \end{syntax}

\end{function}





\begin{function}[]
  {
    \protected@write
  }
  \begin{syntax}
    \cs{protected@write}
  \end{syntax}

\end{function}



\begin{function}[]
  {
    \protected@xdef
  }
  \begin{syntax}
    \cs{protected@xdef}
  \end{syntax}

\end{function}



\begin{function}[]
  {
    \provide@command
  }
  \begin{syntax}
    \cs{provide@command}
  \end{syntax}

\end{function}



\begin{function}[]
  {
    \providecommand
  }
  \begin{syntax}
    \cs{providecommand}
  \end{syntax}

\end{function}



\begin{function}[]
  {
    \ps@empty
  }
  \begin{syntax}
    \cs{ps@empty}
  \end{syntax}

\end{function}



\begin{function}[]
  {
    \ps@plain
  }
  \begin{syntax}
    \cs{ps@plain}
  \end{syntax}

\end{function}



\begin{function}[]
  {
    \pushtabs
  }
  \begin{syntax}
    \cs{pushtabs}
  \end{syntax}

\end{function}



\begin{function}[]
  {
    \put
  }
  \begin{syntax}
    \cs{put}
  \end{syntax}

\end{function}



\begin{function}[]
  {
    \q@curr@file
  }
  \begin{syntax}
    \cs{q@curr@file}
  \end{syntax}

\end{function}



\begin{function}[]
  {
    \qbezier
  }
  \begin{syntax}
    \cs{qbezier}
  \end{syntax}

\end{function}



\begin{function}[]
  {
    \qbeziermax
  }
  \begin{syntax}
    \cs{qbeziermax}
  \end{syntax}

\end{function}



\begin{function}[]
  {
    \qquad
  }
  \begin{syntax}
    \cs{qquad}
  \end{syntax}

\end{function}



\begin{function}[]
  {
    \quad
  }
  \begin{syntax}
    \cs{quad}
  \end{syntax}

\end{function}



\begin{function}[]
  {
    \quote@@name
  }
  \begin{syntax}
    \cs{quote@@name}
  \end{syntax}

\end{function}



\begin{function}[]
  {
    \quote@name
  }
  \begin{syntax}
    \cs{quote@name}
  \end{syntax}

\end{function}



\begin{function}[]
  {
    \r
  }
  \begin{syntax}
    \cs{r}
  \end{syntax}

\end{function}



\begin{function}[]
  {
    \r@@t
  }
  \begin{syntax}
    \cs{r@@t}
  \end{syntax}

\end{function}



\begin{function}[]
  {
    \radical
  }
  \begin{syntax}
    \cs{radical}
  \end{syntax}

\end{function}



\begin{function}[]
  {
    \raggedbottom
  }
  \begin{syntax}
    \cs{raggedbottom}
  \end{syntax}

\end{function}



\begin{function}[]
  {
    \raggedleft
  }
  \begin{syntax}
    \cs{raggedleft}
  \end{syntax}

\end{function}



\begin{function}[]
  {
    \raggedright
  }
  \begin{syntax}
    \cs{raggedright}
  \end{syntax}

\end{function}



\begin{function}[]
  {
    \raise
  }
  \begin{syntax}
    \cs{raise}
  \end{syntax}

\end{function}



\begin{function}[]
  {
    \raisebox
  }
  \begin{syntax}
    \cs{raisebox}
  \end{syntax}

\end{function}



\begin{function}[]
  {
    \rbrace
  }
  \begin{syntax}
    \cs{rbrace}
  \end{syntax}

\end{function}



\begin{function}[]
  {
    \rbrack
  }
  \begin{syntax}
    \cs{rbrack}
  \end{syntax}

\end{function}





\begin{function}[]
  {
    \reenable@package@load
  }
  \begin{syntax}
    \cs{reenable@package@load}
  \end{syntax}

\end{function}



\begin{function}[]
  {
    \ref
  }
  \begin{syntax}
    \cs{ref}
  \end{syntax}

\end{function}




\begin{function}[]
  {
    \reinstall@nfss@defs
  }
  \begin{syntax}
    \cs{reinstall@nfss@defs}
  \end{syntax}

\end{function}



\begin{function}[]
  {
    \relax
  }
  \begin{syntax}
    \cs{relax}
  \end{syntax}

\end{function}



\begin{function}[]
  {
    \relpenalty
  }
  \begin{syntax}
    \cs{relpenalty}
  \end{syntax}

\end{function}



\begin{function}[]
  {
    \rem@pt
  }
  \begin{syntax}
    \cs{rem@pt}
  \end{syntax}

\end{function}



\begin{function}[]
  {
    \remove@angles
  }
  \begin{syntax}
    \cs{remove@angles}
  \end{syntax}

\end{function}



\begin{function}[]
  {
    \remove@star
  }
  \begin{syntax}
    \cs{remove@star}
  \end{syntax}

\end{function}



\begin{function}[]
  {
    \remove@to@nnil
  }
  \begin{syntax}
    \cs{remove@to@nnil}
  \end{syntax}

\end{function}



\begin{function}[]
  {
    \removelastskip
  }
  \begin{syntax}
    \cs{removelastskip}
  \end{syntax}

\end{function}



\begin{function}[]
  {
    \renew@command
  }
  \begin{syntax}
    \cs{renew@command}
  \end{syntax}

\end{function}



\begin{function}[]
  {
    \renew@environment
  }
  \begin{syntax}
    \cs{renew@environment}
  \end{syntax}

\end{function}



\begin{function}[]
  {
    \renewcommand
  }
  \begin{syntax}
    \cs{renewcommand}
  \end{syntax}

\end{function}



\begin{function}[]
  {
    \renewenvironment
  }
  \begin{syntax}
    \cs{renewenvironment}
  \end{syntax}

\end{function}



\begin{function}[]
  {
    \repeat
  }
  \begin{syntax}
    \cs{repeat}
  \end{syntax}

\end{function}



\begin{function}[]
  {
    \requested@test@context
  }
  \begin{syntax}
    \cs{requested@test@context}
  \end{syntax}

\end{function}





\begin{function}[]
  {
    \requestedpatchdate
  }
  \begin{syntax}
    \cs{requestedpatchdate}
  \end{syntax}

\end{function}



\begin{function}[]
  {
    \reserved@a
  }
  \begin{syntax}
    \cs{reserved@a}
  \end{syntax}

\end{function}



\begin{function}[]
  {
    \reserved@b
  }
  \begin{syntax}
    \cs{reserved@b}
  \end{syntax}

\end{function}



\begin{function}[]
  {
    \reserved@c
  }
  \begin{syntax}
    \cs{reserved@c}
  \end{syntax}

\end{function}



\begin{function}[]
  {
    \reserved@d
  }
  \begin{syntax}
    \cs{reserved@d}
  \end{syntax}

\end{function}



\begin{function}[]
  {
    \reserved@e
  }
  \begin{syntax}
    \cs{reserved@e}
  \end{syntax}

\end{function}



\begin{function}[]
  {
    \reserved@f
  }
  \begin{syntax}
    \cs{reserved@f}
  \end{syntax}

\end{function}



\begin{function}[]
  {
    \reset@font
  }
  \begin{syntax}
    \cs{reset@font}
  \end{syntax}

\end{function}



\begin{function}[]
  {
    \restglb@settings
  }
  \begin{syntax}
    \cs{restglb@settings}
  \end{syntax}

\end{function}



\begin{function}[]
  {
    \restore@mathversion
  }
  \begin{syntax}
    \cs{restore@mathversion}
  \end{syntax}

\end{function}



\begin{function}[]
  {
    \restore@protect
  }
  \begin{syntax}
    \cs{restore@protect}
  \end{syntax}

\end{function}



\begin{function}[]
  {
    \restorecr
  }
  \begin{syntax}
    \cs{restorecr}
  \end{syntax}

\end{function}



\begin{function}[]
  {
    \reversemarginpar
  }
  \begin{syntax}
    \cs{reversemarginpar}
  \end{syntax}

\end{function}



\begin{function}[]
  {
    \rhd
  }
  \begin{syntax}
    \cs{rhd}
  \end{syntax}

\end{function}



\begin{function}[]
  {
    \right
  }
  \begin{syntax}
    \cs{right}
  \end{syntax}

\end{function}





\begin{function}[]
  {
    \rightmark
  }
  \begin{syntax}
    \cs{rightmark}
  \end{syntax}

\end{function}



\begin{function}[]
  {
    \rightskip
  }
  \begin{syntax}
    \cs{rightskip}
  \end{syntax}

\end{function}





\begin{function}[]
  {
    \rmdef@ult
  }
  \begin{syntax}
    \cs{rmdef@ult}
  \end{syntax}

\end{function}



\begin{function}[]
  {
    \rmdefault
  }
  \begin{syntax}
    \cs{rmdefault}
  \end{syntax}

\end{function}



\begin{function}[]
  {
    \rmfamily
  }
  \begin{syntax}
    \cs{rmfamily}
  \end{syntax}

\end{function}



\begin{function}[]
  {
    \rmsubstdefault
  }
  \begin{syntax}
    \cs{rmsubstdefault}
  \end{syntax}

\end{function}



\begin{function}[]
  {
    \robust@command@act
  }
  \begin{syntax}
    \cs{robust@command@act}
  \end{syntax}

\end{function}



\begin{function}[]
  {
    \robust@command@act@chk@args
  }
  \begin{syntax}
    \cs{robust@command@act@chk@args}
  \end{syntax}

\end{function}



\begin{function}[]
  {
    \robust@command@act@do
  }
  \begin{syntax}
    \cs{robust@command@act@do}
  \end{syntax}

\end{function}



\begin{function}[]
  {
    \robust@command@act@end
  }
  \begin{syntax}
    \cs{robust@command@act@end}
  \end{syntax}

\end{function}



\begin{function}[]
  {
    \robust@command@act@loop
  }
  \begin{syntax}
    \cs{robust@command@act@loop}
  \end{syntax}

\end{function}



\begin{function}[]
  {
    \robust@command@act@loop@aux
  }
  \begin{syntax}
    \cs{robust@command@act@loop@aux}
  \end{syntax}

\end{function}



\begin{function}[]
  {
    \robust@command@chk@safe
  }
  \begin{syntax}
    \cs{robust@command@chk@safe}
  \end{syntax}

\end{function}



\begin{function}[]
  {
    \roman
  }
  \begin{syntax}
    \cs{roman}
  \end{syntax}

\end{function}



\begin{function}[]
  {
    \romannumeral
  }
  \begin{syntax}
    \cs{romannumeral}
  \end{syntax}

\end{function}



\begin{function}[]
  {
    \root
  }
  \begin{syntax}
    \cs{root}
  \end{syntax}

\end{function}



\begin{function}[]
  {
    \rootbox
  }
  \begin{syntax}
    \cs{rootbox}
  \end{syntax}

\end{function}



\begin{function}[]
  {
    \rq
  }
  \begin{syntax}
    \cs{rq}
  \end{syntax}

\end{function}



\begin{function}[]
  {
    \rule
  }
  \begin{syntax}
    \cs{rule}
  \end{syntax}

\end{function}




\begin{function}[]
  {
    \savebox
  }
  \begin{syntax}
    \cs{savebox}
  \end{syntax}

\end{function}



\begin{function}[]
  {
    \savecatcodetable
  }
  \begin{syntax}
    \cs{savecatcodetable}
  \end{syntax}

\end{function}



\begin{function}[]
  {
    \saved@reqcolroom
  }
  \begin{syntax}
    \cs{saved@reqcolroom}
  \end{syntax}

\end{function}



\begin{function}[]
  {
    \saved@space@catcode
  }
  \begin{syntax}
    \cs{saved@space@catcode}
  \end{syntax}

\end{function}



\begin{function}[]
  {
    \sb
  }
  \begin{syntax}
    \cs{sb}
  \end{syntax}

\end{function}



\begin{function}[]
  {
    \sbox
  }
  \begin{syntax}
    \cs{sbox}
  \end{syntax}

\end{function}



\begin{function}[]
  {
    \scan_new:N
  }
  \begin{syntax}
    \cs{scan_new:N}
  \end{syntax}

\end{function}



\begin{function}[]
  {
    \scan_stop:
  }
  \begin{syntax}
    \cs{scan_stop:}
  \end{syntax}

\end{function}



\begin{function}[]
  {
    \scdefault
  }
  \begin{syntax}
    \cs{scdefault}
  \end{syntax}

\end{function}



\begin{function}[]
  {
    \scriptfont
  }
  \begin{syntax}
    \cs{scriptfont}
  \end{syntax}

\end{function}



\begin{function}[]
  {
    \scriptfont@name
  }
  \begin{syntax}
    \cs{scriptfont@name}
  \end{syntax}

\end{function}



\begin{function}[]
  {
    \scriptscriptfont
  }
  \begin{syntax}
    \cs{scriptscriptfont}
  \end{syntax}

\end{function}



\begin{function}[]
  {
    \scriptscriptstyle
  }
  \begin{syntax}
    \cs{scriptscriptstyle}
  \end{syntax}

\end{function}



\begin{function}[]
  {
    \scriptspace
  }
  \begin{syntax}
    \cs{scriptspace}
  \end{syntax}

\end{function}



\begin{function}[]
  {
    \scriptstyle
  }
  \begin{syntax}
    \cs{scriptstyle}
  \end{syntax}

\end{function}



\begin{function}[]
  {
    \scshape
  }
  \begin{syntax}
    \cs{scshape}
  \end{syntax}

\end{function}





\begin{function}[]
  {
    \secdef
  }
  \begin{syntax}
    \cs{secdef}
  \end{syntax}

\end{function}



\begin{function}[]
  {
    \sectionmark
  }
  \begin{syntax}
    \cs{sectionmark}
  \end{syntax}

\end{function}



\begin{function}[]
  {
    \select@group
  }
  \begin{syntax}
    \cs{select@group}
  \end{syntax}

\end{function}



\begin{function}[]
  {
    \selectfont
  }
  \begin{syntax}
    \cs{selectfont}
  \end{syntax}

\end{function}



\begin{function}[]
  {
    \selectfont`
  }
  \begin{syntax}
    \cs{selectfont`}
  \end{syntax}

\end{function}



\begin{function}[]
  {
    \selectfonte
  }
  \begin{syntax}
    \cs{selectfonte}
  \end{syntax}

\end{function}



\begin{function}[]
  {
    \series@check@toks
  }
  \begin{syntax}
    \cs{series@check@toks}
  \end{syntax}

\end{function}



\begin{function}[]
  {
    \series@drop@one@m
  }
  \begin{syntax}
    \cs{series@drop@one@m}
  \end{syntax}

\end{function}



\begin{function}[]
  {
    \series@maybe@drop@one@m
  }
  \begin{syntax}
    \cs{series@maybe@drop@one@m}
  \end{syntax}

\end{function}



\begin{function}[]
  {
    \series@maybe@drop@one@m@x
  }
  \begin{syntax}
    \cs{series@maybe@drop@one@m@x}
  \end{syntax}

\end{function}



\begin{function}[]
  {
    \seriesdefault
  }
  \begin{syntax}
    \cs{seriesdefault}
  \end{syntax}

\end{function}



\begin{function}[]
  {
    \seriesdefault@kernel
  }
  \begin{syntax}
    \cs{seriesdefault@kernel}
  \end{syntax}

\end{function}



\begin{function}[]
  {
    \set@@mathdelimiter
  }
  \begin{syntax}
    \cs{set@@mathdelimiter}
  \end{syntax}

\end{function}



\begin{function}[]
  {
    \set@color
  }
  \begin{syntax}
    \cs{set@color}
  \end{syntax}

\end{function}



\begin{function}[]
  {
    \set@curr@file
  }
  \begin{syntax}
    \cs{set@curr@file}
  \end{syntax}

\end{function}



\begin{function}[]
  {
    \set@curr@file@aux
  }
  \begin{syntax}
    \cs{set@curr@file@aux}
  \end{syntax}

\end{function}



\begin{function}[]
  {
    \set@curr@file@nosearch
  }
  \begin{syntax}
    \cs{set@curr@file@nosearch}
  \end{syntax}

\end{function}



\begin{function}[]
  {
    \set@current@meta@family
  }
  \begin{syntax}
    \cs{set@current@meta@family}
  \end{syntax}

\end{function}



\begin{function}[]
  {
    \set@display@protect
  }
  \begin{syntax}
    \cs{set@display@protect}
  \end{syntax}

\end{function}



\begin{function}[]
  {
    \set@fontsize
  }
  \begin{syntax}
    \cs{set@fontsize}
  \end{syntax}

\end{function}



\begin{function}[]
  {
    \set@mathaccent
  }
  \begin{syntax}
    \cs{set@mathaccent}
  \end{syntax}

\end{function}



\begin{function}[]
  {
    \set@mathchar
  }
  \begin{syntax}
    \cs{set@mathchar}
  \end{syntax}

\end{function}



\begin{function}[]
  {
    \set@mathdelimiter
  }
  \begin{syntax}
    \cs{set@mathdelimiter}
  \end{syntax}

\end{function}



\begin{function}[]
  {
    \set@mathsymbol
  }
  \begin{syntax}
    \cs{set@mathsymbol}
  \end{syntax}

\end{function}



\begin{function}[]
  {
    \set@simple@size@args
  }
  \begin{syntax}
    \cs{set@simple@size@args}
  \end{syntax}

\end{function}



\begin{function}[]
  {
    \set@size@funct@args
  }
  \begin{syntax}
    \cs{set@size@funct@args}
  \end{syntax}

\end{function}



\begin{function}[]
  {
    \set@size@funct@args@
  }
  \begin{syntax}
    \cs{set@size@funct@args@}
  \end{syntax}

\end{function}



\begin{function}[]
  {
    \set@target@series
  }
  \begin{syntax}
    \cs{set@target@series}
  \end{syntax}

\end{function}



\begin{function}[]
  {
    \set@typeset@protect
  }
  \begin{syntax}
    \cs{set@typeset@protect}
  \end{syntax}

\end{function}



\begin{function}[]
  {
    \setattribute
  }
  \begin{syntax}
    \cs{setattribute}
  \end{syntax}

\end{function}



\begin{function}[]
  {
    \setbox
  }
  \begin{syntax}
    \cs{setbox}
  \end{syntax}

\end{function}






\begin{function}[]
  {
    \setrangecatcode
  }
  \begin{syntax}
    \cs{setrangecatcode}
  \end{syntax}

\end{function}



\begin{function}[]
  {
    \settodepth
  }
  \begin{syntax}
    \cs{settodepth}
  \end{syntax}

\end{function}





\begin{function}[]
  {
    \sf@size
  }
  \begin{syntax}
    \cs{sf@size}
  \end{syntax}

\end{function}



\begin{function}[]
  {
    \sfcode
  }
  \begin{syntax}
    \cs{sfcode}
  \end{syntax}

\end{function}




\begin{function}[]
  {
    \sfdef@ult
  }
  \begin{syntax}
    \cs{sfdef@ult}
  \end{syntax}

\end{function}



\begin{function}[]
  {
    \sfdefault
  }
  \begin{syntax}
    \cs{sfdefault}
  \end{syntax}

\end{function}



\begin{function}[]
  {
    \sffamily
  }
  \begin{syntax}
    \cs{sffamily}
  \end{syntax}

\end{function}



\begin{function}[]
  {
    \sfsubstdefault
  }
  \begin{syntax}
    \cs{sfsubstdefault}
  \end{syntax}

\end{function}



\begin{function}[]
  {
    \sh@ft
  }
  \begin{syntax}
    \cs{sh@ft}
  \end{syntax}

\end{function}



\begin{function}[]
  {
    \shapedefault
  }
  \begin{syntax}
    \cs{shapedefault}
  \end{syntax}

\end{function}



\begin{function}[]
  {
    \shipout
  }
  \begin{syntax}
    \cs{shipout}
  \end{syntax}

\end{function}



\begin{function}[]
  {
    \shortstack
  }
  \begin{syntax}
    \cs{shortstack}
  \end{syntax}

\end{function}



\begin{function}[]
  {
    \show
  }
  \begin{syntax}
    \cs{show}
  \end{syntax}

\end{function}



\begin{function}[]
  {
    \show@kernel@robust@command
  }
  \begin{syntax}
    \cs{show@kernel@robust@command}
  \end{syntax}

\end{function}



\begin{function}[]
  {
    \show@release@info
  }
  \begin{syntax}
    \cs{show@release@info}
  \end{syntax}

\end{function}



\begin{function}[]
  {
    \showbox
  }
  \begin{syntax}
    \cs{showbox}
  \end{syntax}

\end{function}



\begin{function}[]
  {
    \showboxbreadth
  }
  \begin{syntax}
    \cs{showboxbreadth}
  \end{syntax}

\end{function}



\begin{function}[]
  {
    \showboxdepth
  }
  \begin{syntax}
    \cs{showboxdepth}
  \end{syntax}

\end{function}



\begin{function}[]
  {
    \showhyphens
  }
  \begin{syntax}
    \cs{showhyphens}
  \end{syntax}

\end{function}



\begin{function}[]
  {
    \showoutput
  }
  \begin{syntax}
    \cs{showoutput}
  \end{syntax}

\end{function}



\begin{function}[]
  {
    \showoverfull
  }
  \begin{syntax}
    \cs{showoverfull}
  \end{syntax}

\end{function}



\begin{function}[]
  {
    \showtokens
  }
  \begin{syntax}
    \cs{showtokens}
  \end{syntax}

\end{function}





\begin{function}[]
  {
    \sixt@@n
  }
  \begin{syntax}
    \cs{sixt@@n}
  \end{syntax}

\end{function}



\begin{function}[]
  {
    \size@update
  }
  \begin{syntax}
    \cs{size@update}
  \end{syntax}

\end{function}



\begin{function}[]
  {
    \sizefn@info
  }
  \begin{syntax}
    \cs{sizefn@info}
  \end{syntax}

\end{function}



\begin{function}[]
  {
    \skip
  }
  \begin{syntax}
    \cs{skip}
  \end{syntax}

\end{function}



\begin{function}[]
  {
    \skip@
  }
  \begin{syntax}
    \cs{skip@}
  \end{syntax}

\end{function}



\begin{function}[]
  {
    \skipdef
  }
  \begin{syntax}
    \cs{skipdef}
  \end{syntax}

\end{function}



\begin{function}[]
  {
    \skipeval
  }
  \begin{syntax}
    \cs{skipeval}
  \end{syntax}

\end{function}



\begin{function}[]
  {
    \skipzero
  }
  \begin{syntax}
    \cs{skipzero}
  \end{syntax}

\end{function}



\begin{function}[]
  {
    \slash
  }
  \begin{syntax}
    \cs{slash}
  \end{syntax}

\end{function}



\begin{function}[]
  {
    \sldefault
  }
  \begin{syntax}
    \cs{sldefault}
  \end{syntax}

\end{function}



\begin{function}[]
  {
    \sloppy
  }
  \begin{syntax}
    \cs{sloppy}
  \end{syntax}

\end{function}



\begin{function}[]
  {
    \sloppypar
  }
  \begin{syntax}
    \cs{sloppypar}
  \end{syntax}

\end{function}



\begin{function}[]
  {
    \slshape
  }
  \begin{syntax}
    \cs{slshape}
  \end{syntax}

\end{function}



\begin{function}[]
  {
    \smallbreak
  }
  \begin{syntax}
    \cs{smallbreak}
  \end{syntax}

\end{function}







\begin{function}[]
  {
    \smash
  }
  \begin{syntax}
    \cs{smash}
  \end{syntax}

\end{function}



\begin{function}[]
  {
    \sourceLaTeXdate
  }
  \begin{syntax}
    \cs{sourceLaTeXdate}
  \end{syntax}

\end{function}



\begin{function}[]
  {
    \sp
  }
  \begin{syntax}
    \cs{sp}
  \end{syntax}

\end{function}



\begin{function}[]
  {
    \sp@ce@skip
  }
  \begin{syntax}
    \cs{sp@ce@skip}
  \end{syntax}

\end{function}



\begin{function}[]
  {
    \sp@n
  }
  \begin{syntax}
    \cs{sp@n}
  \end{syntax}

\end{function}




\begin{function}[]
  {
    \spacefactor
  }
  \begin{syntax}
    \cs{spacefactor}
  \end{syntax}

\end{function}



\begin{function}[]
  {
    \spaceskip
  }
  \begin{syntax}
    \cs{spaceskip}
  \end{syntax}

\end{function}



\begin{function}[int]
  {
    \span
  }
  \begin{syntax}
    \cs{span}
  \end{syntax}
  %
  Unless you are programming low-level table macros you don't want to
  know and even then you may not want to know. If you do, see the
  \TeX{}book.
\end{function}



\begin{function}[]
  {
    \special
  }
  \begin{syntax}
    \cs{special}
  \end{syntax}

\end{function}



\begin{function}[]
  {
    \split@name
  }
  \begin{syntax}
    \cs{split@name}
  \end{syntax}

\end{function}



\begin{function}[]
  {
    \splitmaxdepth
  }
  \begin{syntax}
    \cs{splitmaxdepth}
  \end{syntax}

\end{function}



\begin{function}[]
  {
    \splittopskip
  }
  \begin{syntax}
    \cs{splittopskip}
  \end{syntax}

\end{function}





\begin{function}[]
  {
    \sqsubset
  }
  \begin{syntax}
    \cs{sqsubset}
  \end{syntax}

\end{function}



\begin{function}[]
  {
    \sqsupset
  }
  \begin{syntax}
    \cs{sqsupset}
  \end{syntax}

\end{function}



\begin{function}[]
  {
    \ss
  }
  \begin{syntax}
    \cs{ss}
  \end{syntax}

\end{function}



\begin{function}[]
  {
    \sscdefault
  }
  \begin{syntax}
    \cs{sscdefault}
  \end{syntax}

\end{function}



\begin{function}[]
  {
    \sscshape
  }
  \begin{syntax}
    \cs{sscshape}
  \end{syntax}

\end{function}



\begin{function}[]
  {
    \ssf@size
  }
  \begin{syntax}
    \cs{ssf@size}
  \end{syntax}

\end{function}



\begin{function}[]
  {
    \stackrel
  }
  \begin{syntax}
    \cs{stackrel}
  \end{syntax}

\end{function}



\begin{function}[]
  {
    \stockheight
  }
  \begin{syntax}
    \cs{stockheight}
  \end{syntax}

\end{function}



\begin{function}[]
  {
    \stockwidth
  }
  \begin{syntax}
    \cs{stockwidth}
  \end{syntax}

\end{function}



\begin{function}[]
  {
    \stop
  }
  \begin{syntax}
    \cs{stop}
  \end{syntax}

\end{function}



\begin{function}[]
  {
    \storedpar
  }
  \begin{syntax}
    \cs{storedpar}
  \end{syntax}

\end{function}



\begin{function}[]
  {
    \stretch
  }
  \begin{syntax}
    \cs{stretch}
  \end{syntax}

\end{function}



\begin{function}[]
  {
    \string
  }
  \begin{syntax}
    \cs{string}
  \end{syntax}

\end{function}



\begin{function}[]
  {
    \string@makeletter
  }
  \begin{syntax}
    \cs{string@makeletter}
  \end{syntax}

\end{function}



\begin{function}[]
  {
    \strip@prefix
  }
  \begin{syntax}
    \cs{strip@prefix}
  \end{syntax}

\end{function}



\begin{function}[]
  {
    \strip@pt
  }
  \begin{syntax}
    \cs{strip@pt}
  \end{syntax}

\end{function}



\begin{function}[]
  {
    \strut
  }
  \begin{syntax}
    \cs{strut}
  \end{syntax}

\end{function}



\begin{function}[]
  {
    \strutbox
  }
  \begin{syntax}
    \cs{strutbox}
  \end{syntax}

\end{function}



\begin{function}[]
  {
    \sub@sfcnt
  }
  \begin{syntax}
    \cs{sub@sfcnt}
  \end{syntax}

\end{function}



\begin{function}[]
  {
    \subf@sfcnt
  }
  \begin{syntax}
    \cs{subf@sfcnt}
  \end{syntax}

\end{function}



\begin{function}[]
  {
    \subparagraphmark
  }
  \begin{syntax}
    \cs{subparagraphmark}
  \end{syntax}

\end{function}



\begin{function}[]
  {
    \subsectionmark
  }
  \begin{syntax}
    \cs{subsectionmark}
  \end{syntax}

\end{function}



\begin{function}[]
  {
    \subst@correction
  }
  \begin{syntax}
    \cs{subst@correction}
  \end{syntax}

\end{function}



\begin{function}[]
  {
    \subsubsectionmark
  }
  \begin{syntax}
    \cs{subsubsectionmark}
  \end{syntax}

\end{function}



\begin{function}[]
  {
    \sup
  }
  \begin{syntax}
    \cs{sup}
  \end{syntax}

\end{function}



\begin{function}[]
  {
    \suppressfloats
  }
  \begin{syntax}
    \cs{suppressfloats}
  \end{syntax}

\end{function}



\begin{function}[]
  {
    \sw@slant
  }
  \begin{syntax}
    \cs{sw@slant}
  \end{syntax}

\end{function}



\begin{function}[]
  {
    \swdefault
  }
  \begin{syntax}
    \cs{swdefault}
  \end{syntax}

\end{function}



\begin{function}[]
  {
    \swshape
  }
  \begin{syntax}
    \cs{swshape}
  \end{syntax}

\end{function}



\begin{function}[]
  {
    \symbol
  }
  \begin{syntax}
    \cs{symbol}
  \end{syntax}

\end{function}



\begin{function}[]
  {
    \symletters
  }
  \begin{syntax}
    \cs{symletters}
  \end{syntax}

\end{function}



\begin{function}[]
  {
    \t
  }
  \begin{syntax}
    \cs{t}
  \end{syntax}

\end{function}



\begin{function}[]
  {
    \t@st@ic
  }
  \begin{syntax}
    \cs{t@st@ic}
  \end{syntax}

\end{function}



\begin{function}[]
  {
    \tabbing
  }
  \begin{syntax}
    \cs{tabbing}
  \end{syntax}

\end{function}



\begin{function}[]
  {
    \tabbingsep
  }
  \begin{syntax}
    \cs{tabbingsep}
  \end{syntax}

\end{function}



\begin{function}[]
  {
    \tabcolsep
  }
  \begin{syntax}
    \cs{tabcolsep}
  \end{syntax}

\end{function}



\begin{function}[]
  {
    \tabskip
  }
  \begin{syntax}
    \cs{tabskip}
  \end{syntax}

\end{function}



\begin{function}[]
  {
    \tabular
  }
  \begin{syntax}
    \cs{tabular}
  \end{syntax}

\end{function}



\begin{function}[]
  {
    \tabularnewline
  }
  \begin{syntax}
    \cs{tabularnewline}
  \end{syntax}

\end{function}






\begin{function}[]
  {
    \target@series@value
  }
  \begin{syntax}
    \cs{target@series@value}
  \end{syntax}

\end{function}



\begin{function}[]
  {
    \tc@check@accent
  }
  \begin{syntax}
    \cs{tc@check@accent}
  \end{syntax}

\end{function}



\begin{function}[]
  {
    \tc@check@symbol
  }
  \begin{syntax}
    \cs{tc@check@symbol}
  \end{syntax}

\end{function}



\begin{function}[]
  {
    \tc@errorwarn
  }
  \begin{syntax}
    \cs{tc@errorwarn}
  \end{syntax}

\end{function}



\begin{function}[]
  {
    \tc@fake@euro
  }
  \begin{syntax}
    \cs{tc@fake@euro}
  \end{syntax}

\end{function}



\begin{function}[]
  {
    \tc@oldstylesubst
  }
  \begin{syntax}
    \cs{tc@oldstylesubst}
  \end{syntax}

\end{function}



\begin{function}[]
  {
    \tc@subst
  }
  \begin{syntax}
    \cs{tc@subst}
  \end{syntax}

\end{function}



\begin{function}[]
  {
    \tc@swap@accent
  }
  \begin{syntax}
    \cs{tc@swap@accent}
  \end{syntax}

\end{function}



\begin{function}[]
  {
    \tencirc
  }
  \begin{syntax}
    \cs{tencirc}
  \end{syntax}

\end{function}



\begin{function}[]
  {
    \tencircw
  }
  \begin{syntax}
    \cs{tencircw}
  \end{syntax}

\end{function}



\begin{function}[]
  {
    \tenln
  }
  \begin{syntax}
    \cs{tenln}
  \end{syntax}

\end{function}



\begin{function}[]
  {
    \tenlnw
  }
  \begin{syntax}
    \cs{tenlnw}
  \end{syntax}

\end{function}



\begin{function}[]
  {
    \test@font@series@context
  }
  \begin{syntax}
    \cs{test@font@series@context}
  \end{syntax}

\end{function}



\begin{function}[]
  {
    \text@command
  }
  \begin{syntax}
    \cs{text@command}
  \end{syntax}

\end{function}



\begin{function}[]
  {
    \textacutedbl
  }
  \begin{syntax}
    \cs{textacutedbl}
  \end{syntax}

\end{function}



\begin{function}[]
  {
    \textascendercompwordmark
  }
  \begin{syntax}
    \cs{textascendercompwordmark}
  \end{syntax}

\end{function}



\begin{function}[]
  {
    \textasciiacute
  }
  \begin{syntax}
    \cs{textasciiacute}
  \end{syntax}

\end{function}



\begin{function}[]
  {
    \textasciibreve
  }
  \begin{syntax}
    \cs{textasciibreve}
  \end{syntax}

\end{function}



\begin{function}[]
  {
    \textasciicaron
  }
  \begin{syntax}
    \cs{textasciicaron}
  \end{syntax}

\end{function}



\begin{function}[]
  {
    \textasciicircum
  }
  \begin{syntax}
    \cs{textasciicircum}
  \end{syntax}

\end{function}



\begin{function}[]
  {
    \textasciidieresis
  }
  \begin{syntax}
    \cs{textasciidieresis}
  \end{syntax}

\end{function}



\begin{function}[]
  {
    \textasciigrave
  }
  \begin{syntax}
    \cs{textasciigrave}
  \end{syntax}

\end{function}



\begin{function}[]
  {
    \textasciimacron
  }
  \begin{syntax}
    \cs{textasciimacron}
  \end{syntax}

\end{function}





\begin{function}[]
  {
    \textasteriskcentered
  }
  \begin{syntax}
    \cs{textasteriskcentered}
  \end{syntax}

\end{function}



\begin{function}[]
  {
    \textbackslash
  }
  \begin{syntax}
    \cs{textbackslash}
  \end{syntax}

\end{function}



\begin{function}[]
  {
    \textbaht
  }
  \begin{syntax}
    \cs{textbaht}
  \end{syntax}

\end{function}



\begin{function}[]
  {
    \textbar
  }
  \begin{syntax}
    \cs{textbar}
  \end{syntax}

\end{function}



\begin{function}[]
  {
    \textbardbl
  }
  \begin{syntax}
    \cs{textbardbl}
  \end{syntax}

\end{function}



\begin{function}[]
  {
    \textbf
  }
  \begin{syntax}
    \cs{textbf}
  \end{syntax}

\end{function}



\begin{function}[]
  {
    \textbigcircle
  }
  \begin{syntax}
    \cs{textbigcircle}
  \end{syntax}

\end{function}



\begin{function}[]
  {
    \textblank
  }
  \begin{syntax}
    \cs{textblank}
  \end{syntax}

\end{function}



\begin{function}[]
  {
    \textborn
  }
  \begin{syntax}
    \cs{textborn}
  \end{syntax}

\end{function}



\begin{function}[]
  {
    \textbraceleft
  }
  \begin{syntax}
    \cs{textbraceleft}
  \end{syntax}

\end{function}



\begin{function}[]
  {
    \textbraceright
  }
  \begin{syntax}
    \cs{textbraceright}
  \end{syntax}

\end{function}



\begin{function}[]
  {
    \textbrokenbar
  }
  \begin{syntax}
    \cs{textbrokenbar}
  \end{syntax}

\end{function}



\begin{function}[]
  {
    \textbullet
  }
  \begin{syntax}
    \cs{textbullet}
  \end{syntax}

\end{function}



\begin{function}[]
  {
    \textcapitalcompwordmark
  }
  \begin{syntax}
    \cs{textcapitalcompwordmark}
  \end{syntax}

\end{function}



\begin{function}[]
  {
    \textcelsius
  }
  \begin{syntax}
    \cs{textcelsius}
  \end{syntax}

\end{function}



\begin{function}[]
  {
    \textcent
  }
  \begin{syntax}
    \cs{textcent}
  \end{syntax}

\end{function}



\begin{function}[]
  {
    \textcentoldstyle
  }
  \begin{syntax}
    \cs{textcentoldstyle}
  \end{syntax}

\end{function}



\begin{function}[]
  {
    \textcircled
  }
  \begin{syntax}
    \cs{textcircled}
  \end{syntax}

\end{function}



\begin{function}[]
  {
    \textcircledP
  }
  \begin{syntax}
    \cs{textcircledP}
  \end{syntax}

\end{function}



\begin{function}[]
  {
    \textcolonmonetary
  }
  \begin{syntax}
    \cs{textcolonmonetary}
  \end{syntax}

\end{function}



\begin{function}[]
  {
    \textcommaabove
  }
  \begin{syntax}
    \cs{textcommaabove}
  \end{syntax}

\end{function}



\begin{function}[]
  {
    \textcommabelow
  }
  \begin{syntax}
    \cs{textcommabelow}
  \end{syntax}

\end{function}



\begin{function}[]
  {
    \textcompsubstdefault
  }
  \begin{syntax}
    \cs{textcompsubstdefault}
  \end{syntax}

\end{function}



\begin{function}[]
  {
    \textcompwordmark
  }
  \begin{syntax}
    \cs{textcompwordmark}
  \end{syntax}

\end{function}



\begin{function}[]
  {
    \textcopyleft
  }
  \begin{syntax}
    \cs{textcopyleft}
  \end{syntax}

\end{function}



\begin{function}[]
  {
    \textcopyright
  }
  \begin{syntax}
    \cs{textcopyright}
  \end{syntax}

\end{function}



\begin{function}[]
  {
    \textcurrency
  }
  \begin{syntax}
    \cs{textcurrency}
  \end{syntax}

\end{function}



\begin{function}[]
  {
    \textdagger
  }
  \begin{syntax}
    \cs{textdagger}
  \end{syntax}

\end{function}



\begin{function}[]
  {
    \textdaggerdbl
  }
  \begin{syntax}
    \cs{textdaggerdbl}
  \end{syntax}

\end{function}



\begin{function}[]
  {
    \textdblhyphen
  }
  \begin{syntax}
    \cs{textdblhyphen}
  \end{syntax}

\end{function}



\begin{function}[]
  {
    \textdblhyphenchar
  }
  \begin{syntax}
    \cs{textdblhyphenchar}
  \end{syntax}

\end{function}



\begin{function}[]
  {
    \textdegree
  }
  \begin{syntax}
    \cs{textdegree}
  \end{syntax}

\end{function}



\begin{function}[]
  {
    \textdied
  }
  \begin{syntax}
    \cs{textdied}
  \end{syntax}

\end{function}



\begin{function}[]
  {
    \textdiscount
  }
  \begin{syntax}
    \cs{textdiscount}
  \end{syntax}

\end{function}



\begin{function}[]
  {
    \textdiv
  }
  \begin{syntax}
    \cs{textdiv}
  \end{syntax}

\end{function}



\begin{function}[]
  {
    \textdivorced
  }
  \begin{syntax}
    \cs{textdivorced}
  \end{syntax}

\end{function}



\begin{function}[]
  {
    \textdollar
  }
  \begin{syntax}
    \cs{textdollar}
  \end{syntax}

\end{function}



\begin{function}[]
  {
    \textdollaroldstyle
  }
  \begin{syntax}
    \cs{textdollaroldstyle}
  \end{syntax}

\end{function}



\begin{function}[]
  {
    \textdong
  }
  \begin{syntax}
    \cs{textdong}
  \end{syntax}

\end{function}



\begin{function}[]
  {
    \textdownarrow
  }
  \begin{syntax}
    \cs{textdownarrow}
  \end{syntax}

\end{function}



\begin{function}[]
  {
    \texteightoldstyle
  }
  \begin{syntax}
    \cs{texteightoldstyle}
  \end{syntax}

\end{function}



\begin{function}[]
  {
    \textellipsis
  }
  \begin{syntax}
    \cs{textellipsis}
  \end{syntax}

\end{function}



\begin{function}[]
  {
    \textemdash
  }
  \begin{syntax}
    \cs{textemdash}
  \end{syntax}

\end{function}



\begin{function}[]
  {
    \textendash
  }
  \begin{syntax}
    \cs{textendash}
  \end{syntax}

\end{function}



\begin{function}[]
  {
    \textestimated
  }
  \begin{syntax}
    \cs{textestimated}
  \end{syntax}

\end{function}



\begin{function}[]
  {
    \texteuro
  }
  \begin{syntax}
    \cs{texteuro}
  \end{syntax}

\end{function}



\begin{function}[]
  {
    \textexclamdown
  }
  \begin{syntax}
    \cs{textexclamdown}
  \end{syntax}

\end{function}



\begin{function}[]
  {
    \textfiveoldstyle
  }
  \begin{syntax}
    \cs{textfiveoldstyle}
  \end{syntax}

\end{function}



\begin{function}[]
  {
    \textfloatsep
  }
  \begin{syntax}
    \cs{textfloatsep}
  \end{syntax}

\end{function}



\begin{function}[]
  {
    \textflorin
  }
  \begin{syntax}
    \cs{textflorin}
  \end{syntax}

\end{function}



\begin{function}[]
  {
    \textfont
  }
  \begin{syntax}
    \cs{textfont}
  \end{syntax}

\end{function}



\begin{function}[]
  {
    \textfont@name
  }
  \begin{syntax}
    \cs{textfont@name}
  \end{syntax}

\end{function}



\begin{function}[]
  {
    \textfouroldstyle
  }
  \begin{syntax}
    \cs{textfouroldstyle}
  \end{syntax}

\end{function}



\begin{function}[]
  {
    \textfraction
  }
  \begin{syntax}
    \cs{textfraction}
  \end{syntax}

\end{function}



\begin{function}[]
  {
    \textfractionsolidus
  }
  \begin{syntax}
    \cs{textfractionsolidus}
  \end{syntax}

\end{function}



\begin{function}[]
  {
    \textgravedbl
  }
  \begin{syntax}
    \cs{textgravedbl}
  \end{syntax}

\end{function}



\begin{function}[]
  {
    \textgreater
  }
  \begin{syntax}
    \cs{textgreater}
  \end{syntax}

\end{function}



\begin{function}[]
  {
    \textguarani
  }
  \begin{syntax}
    \cs{textguarani}
  \end{syntax}

\end{function}



\begin{function}[]
  {
    \textheight
  }
  \begin{syntax}
    \cs{textheight}
  \end{syntax}

\end{function}



\begin{function}[]
  {
    \textinterrobang
  }
  \begin{syntax}
    \cs{textinterrobang}
  \end{syntax}

\end{function}



\begin{function}[]
  {
    \textinterrobangdown
  }
  \begin{syntax}
    \cs{textinterrobangdown}
  \end{syntax}

\end{function}



\begin{function}[]
  {
    \textit
  }
  \begin{syntax}
    \cs{textit}
  \end{syntax}

\end{function}



\begin{function}[]
  {
    \textlangle
  }
  \begin{syntax}
    \cs{textlangle}
  \end{syntax}

\end{function}



\begin{function}[]
  {
    \textlbrackdbl
  }
  \begin{syntax}
    \cs{textlbrackdbl}
  \end{syntax}

\end{function}



\begin{function}[]
  {
    \textleaf
  }
  \begin{syntax}
    \cs{textleaf}
  \end{syntax}

\end{function}



\begin{function}[]
  {
    \textleftarrow
  }
  \begin{syntax}
    \cs{textleftarrow}
  \end{syntax}

\end{function}



\begin{function}[]
  {
    \textlegacyasteriskcentered
  }
  \begin{syntax}
    \cs{textlegacyasteriskcentered}
  \end{syntax}

\end{function}



\begin{function}[]
  {
    \textlegacybardbl
  }
  \begin{syntax}
    \cs{textlegacybardbl}
  \end{syntax}

\end{function}



\begin{function}[]
  {
    \textlegacybullet
  }
  \begin{syntax}
    \cs{textlegacybullet}
  \end{syntax}

\end{function}



\begin{function}[]
  {
    \textlegacydagger
  }
  \begin{syntax}
    \cs{textlegacydagger}
  \end{syntax}

\end{function}



\begin{function}[]
  {
    \textlegacydaggerdbl
  }
  \begin{syntax}
    \cs{textlegacydaggerdbl}
  \end{syntax}

\end{function}



\begin{function}[]
  {
    \textlegacyparagraph
  }
  \begin{syntax}
    \cs{textlegacyparagraph}
  \end{syntax}

\end{function}



\begin{function}[]
  {
    \textlegacyperiodcentered
  }
  \begin{syntax}
    \cs{textlegacyperiodcentered}
  \end{syntax}

\end{function}



\begin{function}[]
  {
    \textlegacysection
  }
  \begin{syntax}
    \cs{textlegacysection}
  \end{syntax}

\end{function}



\begin{function}[]
  {
    \textless
  }
  \begin{syntax}
    \cs{textless}
  \end{syntax}

\end{function}



\begin{function}[]
  {
    \textlira
  }
  \begin{syntax}
    \cs{textlira}
  \end{syntax}

\end{function}



\begin{function}[]
  {
    \textlnot
  }
  \begin{syntax}
    \cs{textlnot}
  \end{syntax}

\end{function}



\begin{function}[]
  {
    \textlquill
  }
  \begin{syntax}
    \cs{textlquill}
  \end{syntax}

\end{function}



\begin{function}[]
  {
    \textmarried
  }
  \begin{syntax}
    \cs{textmarried}
  \end{syntax}

\end{function}



\begin{function}[]
  {
    \textmd
  }
  \begin{syntax}
    \cs{textmd}
  \end{syntax}

\end{function}



\begin{function}[]
  {
    \textmho
  }
  \begin{syntax}
    \cs{textmho}
  \end{syntax}

\end{function}



\begin{function}[]
  {
    \textminus
  }
  \begin{syntax}
    \cs{textminus}
  \end{syntax}

\end{function}



\begin{function}[]
  {
    \textmu
  }
  \begin{syntax}
    \cs{textmu}
  \end{syntax}

\end{function}



\begin{function}[]
  {
    \textmusicalnote
  }
  \begin{syntax}
    \cs{textmusicalnote}
  \end{syntax}

\end{function}



\begin{function}[]
  {
    \textnaira
  }
  \begin{syntax}
    \cs{textnaira}
  \end{syntax}

\end{function}



\begin{function}[]
  {
    \textnineoldstyle
  }
  \begin{syntax}
    \cs{textnineoldstyle}
  \end{syntax}

\end{function}



\begin{function}[]
  {
    \textnormal
  }
  \begin{syntax}
    \cs{textnormal}
  \end{syntax}

\end{function}



\begin{function}[]
  {
    \textnumero
  }
  \begin{syntax}
    \cs{textnumero}
  \end{syntax}

\end{function}



\begin{function}[]
  {
    \textohm
  }
  \begin{syntax}
    \cs{textohm}
  \end{syntax}

\end{function}



\begin{function}[]
  {
    \textonehalf
  }
  \begin{syntax}
    \cs{textonehalf}
  \end{syntax}

\end{function}



\begin{function}[]
  {
    \textoneoldstyle
  }
  \begin{syntax}
    \cs{textoneoldstyle}
  \end{syntax}

\end{function}



\begin{function}[]
  {
    \textonequarter
  }
  \begin{syntax}
    \cs{textonequarter}
  \end{syntax}

\end{function}



\begin{function}[]
  {
    \textonesuperior
  }
  \begin{syntax}
    \cs{textonesuperior}
  \end{syntax}

\end{function}



\begin{function}[]
  {
    \textopenbullet
  }
  \begin{syntax}
    \cs{textopenbullet}
  \end{syntax}

\end{function}



\begin{function}[]
  {
    \textordfeminine
  }
  \begin{syntax}
    \cs{textordfeminine}
  \end{syntax}

\end{function}



\begin{function}[]
  {
    \textordmasculine
  }
  \begin{syntax}
    \cs{textordmasculine}
  \end{syntax}

\end{function}



\begin{function}[]
  {
    \textparagraph
  }
  \begin{syntax}
    \cs{textparagraph}
  \end{syntax}

\end{function}



\begin{function}[]
  {
    \textperiodcentered
  }
  \begin{syntax}
    \cs{textperiodcentered}
  \end{syntax}

\end{function}



\begin{function}[]
  {
    \textpertenthousand
  }
  \begin{syntax}
    \cs{textpertenthousand}
  \end{syntax}

\end{function}



\begin{function}[]
  {
    \textperthousand
  }
  \begin{syntax}
    \cs{textperthousand}
  \end{syntax}

\end{function}



\begin{function}[]
  {
    \textpeso
  }
  \begin{syntax}
    \cs{textpeso}
  \end{syntax}

\end{function}



\begin{function}[]
  {
    \textpilcrow
  }
  \begin{syntax}
    \cs{textpilcrow}
  \end{syntax}

\end{function}



\begin{function}[]
  {
    \textpm
  }
  \begin{syntax}
    \cs{textpm}
  \end{syntax}

\end{function}



\begin{function}[]
  {
    \textquestiondown
  }
  \begin{syntax}
    \cs{textquestiondown}
  \end{syntax}

\end{function}



\begin{function}[]
  {
    \textquotedblleft
  }
  \begin{syntax}
    \cs{textquotedblleft}
  \end{syntax}

\end{function}



\begin{function}[]
  {
    \textquotedblright
  }
  \begin{syntax}
    \cs{textquotedblright}
  \end{syntax}

\end{function}



\begin{function}[]
  {
    \textquoteleft
  }
  \begin{syntax}
    \cs{textquoteleft}
  \end{syntax}

\end{function}



\begin{function}[]
  {
    \textquoteright
  }
  \begin{syntax}
    \cs{textquoteright}
  \end{syntax}

\end{function}



\begin{function}[]
  {
    \textquotesingle
  }
  \begin{syntax}
    \cs{textquotesingle}
  \end{syntax}

\end{function}



\begin{function}[]
  {
    \textquotestraightbase
  }
  \begin{syntax}
    \cs{textquotestraightbase}
  \end{syntax}

\end{function}



\begin{function}[]
  {
    \textquotestraightdblbase
  }
  \begin{syntax}
    \cs{textquotestraightdblbase}
  \end{syntax}

\end{function}



\begin{function}[]
  {
    \textrangle
  }
  \begin{syntax}
    \cs{textrangle}
  \end{syntax}

\end{function}



\begin{function}[]
  {
    \textrbrackdbl
  }
  \begin{syntax}
    \cs{textrbrackdbl}
  \end{syntax}

\end{function}



\begin{function}[]
  {
    \textrecipe
  }
  \begin{syntax}
    \cs{textrecipe}
  \end{syntax}

\end{function}



\begin{function}[]
  {
    \textreferencemark
  }
  \begin{syntax}
    \cs{textreferencemark}
  \end{syntax}

\end{function}



\begin{function}[]
  {
    \textregistered
  }
  \begin{syntax}
    \cs{textregistered}
  \end{syntax}

\end{function}



\begin{function}[]
  {
    \textrightarrow
  }
  \begin{syntax}
    \cs{textrightarrow}
  \end{syntax}

\end{function}



\begin{function}[]
  {
    \textrm
  }
  \begin{syntax}
    \cs{textrm}
  \end{syntax}

\end{function}



\begin{function}[]
  {
    \textrquill
  }
  \begin{syntax}
    \cs{textrquill}
  \end{syntax}

\end{function}



\begin{function}[]
  {
    \textsc
  }
  \begin{syntax}
    \cs{textsc}
  \end{syntax}

\end{function}



\begin{function}[]
  {
    \textsection
  }
  \begin{syntax}
    \cs{textsection}
  \end{syntax}

\end{function}



\begin{function}[]
  {
    \textservicemark
  }
  \begin{syntax}
    \cs{textservicemark}
  \end{syntax}

\end{function}



\begin{function}[]
  {
    \textsevenoldstyle
  }
  \begin{syntax}
    \cs{textsevenoldstyle}
  \end{syntax}

\end{function}



\begin{function}[]
  {
    \textsf
  }
  \begin{syntax}
    \cs{textsf}
  \end{syntax}

\end{function}



\begin{function}[]
  {
    \textsixoldstyle
  }
  \begin{syntax}
    \cs{textsixoldstyle}
  \end{syntax}

\end{function}



\begin{function}[]
  {
    \textsl
  }
  \begin{syntax}
    \cs{textsl}
  \end{syntax}

\end{function}



\begin{function}[]
  {
    \textssc
  }
  \begin{syntax}
    \cs{textssc}
  \end{syntax}

\end{function}



\begin{function}[]
  {
    \textsterling
  }
  \begin{syntax}
    \cs{textsterling}
  \end{syntax}

\end{function}



\begin{function}[]
  {
    \textstyle
  }
  \begin{syntax}
    \cs{textstyle}
  \end{syntax}

\end{function}



\begin{function}[]
  {
    \textsubscript
  }
  \begin{syntax}
    \cs{textsubscript}
  \end{syntax}

\end{function}



\begin{function}[]
  {
    \textsubscript@offset
  }
  \begin{syntax}
    \cs{textsubscript@offset}
  \end{syntax}

\end{function}



\begin{function}[]
  {
    \textsubscript@space
  }
  \begin{syntax}
    \cs{textsubscript@space}
  \end{syntax}

\end{function}



\begin{function}[]
  {
    \textsuperscript
  }
  \begin{syntax}
    \cs{textsuperscript}
  \end{syntax}

\end{function}



\begin{function}[]
  {
    \textsuperscript@offset
  }
  \begin{syntax}
    \cs{textsuperscript@offset}
  \end{syntax}

\end{function}



\begin{function}[]
  {
    \textsuperscript@space
  }
  \begin{syntax}
    \cs{textsuperscript@space}
  \end{syntax}

\end{function}



\begin{function}[]
  {
    \textsurd
  }
  \begin{syntax}
    \cs{textsurd}
  \end{syntax}

\end{function}



\begin{function}[]
  {
    \textsw
  }
  \begin{syntax}
    \cs{textsw}
  \end{syntax}

\end{function}



\begin{function}[]
  {
    \textthreeoldstyle
  }
  \begin{syntax}
    \cs{textthreeoldstyle}
  \end{syntax}

\end{function}



\begin{function}[]
  {
    \textthreequarters
  }
  \begin{syntax}
    \cs{textthreequarters}
  \end{syntax}

\end{function}



\begin{function}[]
  {
    \textthreequartersemdash
  }
  \begin{syntax}
    \cs{textthreequartersemdash}
  \end{syntax}

\end{function}



\begin{function}[]
  {
    \textthreesuperior
  }
  \begin{syntax}
    \cs{textthreesuperior}
  \end{syntax}

\end{function}



\begin{function}[]
  {
    \texttildelow
  }
  \begin{syntax}
    \cs{texttildelow}
  \end{syntax}

\end{function}



\begin{function}[]
  {
    \texttimes
  }
  \begin{syntax}
    \cs{texttimes}
  \end{syntax}

\end{function}



\begin{function}[]
  {
    \texttrademark
  }
  \begin{syntax}
    \cs{texttrademark}
  \end{syntax}

\end{function}



\begin{function}[]
  {
    \texttt
  }
  \begin{syntax}
    \cs{texttt}
  \end{syntax}

\end{function}



\begin{function}[]
  {
    \texttwelveudash
  }
  \begin{syntax}
    \cs{texttwelveudash}
  \end{syntax}

\end{function}



\begin{function}[]
  {
    \texttwooldstyle
  }
  \begin{syntax}
    \cs{texttwooldstyle}
  \end{syntax}

\end{function}



\begin{function}[]
  {
    \texttwosuperior
  }
  \begin{syntax}
    \cs{texttwosuperior}
  \end{syntax}

\end{function}



\begin{function}[]
  {
    \textulc
  }
  \begin{syntax}
    \cs{textulc}
  \end{syntax}

\end{function}



\begin{function}[]
  {
    \textunderscore
  }
  \begin{syntax}
    \cs{textunderscore}
  \end{syntax}

\end{function}



\begin{function}[]
  {
    \textup
  }
  \begin{syntax}
    \cs{textup}
  \end{syntax}

\end{function}



\begin{function}[]
  {
    \textuparrow
  }
  \begin{syntax}
    \cs{textuparrow}
  \end{syntax}

\end{function}



\begin{function}[]
  {
    \textvisiblespace
  }
  \begin{syntax}
    \cs{textvisiblespace}
  \end{syntax}

\end{function}



\begin{function}[]
  {
    \textwidth
  }
  \begin{syntax}
    \cs{textwidth}
  \end{syntax}

\end{function}



\begin{function}[]
  {
    \textwon
  }
  \begin{syntax}
    \cs{textwon}
  \end{syntax}

\end{function}



\begin{function}[]
  {
    \textyen
  }
  \begin{syntax}
    \cs{textyen}
  \end{syntax}

\end{function}



\begin{function}[]
  {
    \textzerooldstyle
  }
  \begin{syntax}
    \cs{textzerooldstyle}
  \end{syntax}

\end{function}



\begin{function}[]
  {
    \tf@size
  }
  \begin{syntax}
    \cs{tf@size}
  \end{syntax}

\end{function}



\begin{function}[]
  {
    \th
  }
  \begin{syntax}
    \cs{th}
  \end{syntax}

\end{function}



\begin{function}[]
  {
    \thanks
  }
  \begin{syntax}
    \cs{thanks}
  \end{syntax}

\end{function}



\begin{function}[]
  {
    \the
  }
  \begin{syntax}
    \cs{the}
  \end{syntax}

\end{function}



\begin{function}[]
  {
    \theequation
  }
  \begin{syntax}
    \cs{theequation}
  \end{syntax}

\end{function}



\begin{function}[]
  {
    \thefootnote
  }
  \begin{syntax}
    \cs{thefootnote}
  \end{syntax}

\end{function}



\begin{function}[]
  {
    \thempfn
  }
  \begin{syntax}
    \cs{thempfn}
  \end{syntax}

\end{function}



\begin{function}[]
  {
    \thempfootnote
  }
  \begin{syntax}
    \cs{thempfootnote}
  \end{syntax}

\end{function}



\begin{function}[]
  {
    \thepage
  }
  \begin{syntax}
    \cs{thepage}
  \end{syntax}

\end{function}



\begin{function}[]
  {
    \thetotalpages
  }
  \begin{syntax}
    \cs{thetotalpages}
  \end{syntax}

\end{function}



\begin{function}[]
  {
    \thicklines
  }
  \begin{syntax}
    \cs{thicklines}
  \end{syntax}

\end{function}



\begin{function}[]
  {
    \thickmuskip
  }
  \begin{syntax}
    \cs{thickmuskip}
  \end{syntax}

\end{function}



\begin{function}[]
  {
    \thickspace
  }
  \begin{syntax}
    \cs{thickspace}
  \end{syntax}

\end{function}



\begin{function}[]
  {
    \thinlines
  }
  \begin{syntax}
    \cs{thinlines}
  \end{syntax}

\end{function}



\begin{function}[]
  {
    \thinmuskip
  }
  \begin{syntax}
    \cs{thinmuskip}
  \end{syntax}

\end{function}



\begin{function}[]
  {
    \thinspace
  }
  \begin{syntax}
    \cs{thinspace}
  \end{syntax}

\end{function}



\begin{function}[]
  {
    \thispagestyle
  }
  \begin{syntax}
    \cs{thispagestyle}
  \end{syntax}

\end{function}



\begin{function}[]
  {
    \thr@@
  }
  \begin{syntax}
    \cs{thr@@}
  \end{syntax}

\end{function}



\begin{function}[]
  {
    \time
  }
  \begin{syntax}
    \cs{time}
  \end{syntax}

\end{function}



\begin{function}[]
  {
    \title
  }
  \begin{syntax}
    \cs{title}
  \end{syntax}

\end{function}



\begin{function}[]
  {
    \tmspace
  }
  \begin{syntax}
    \cs{tmspace}
  \end{syntax}

\end{function}



\begin{function}[]
  {
    \today
  }
  \begin{syntax}
    \cs{today}
  \end{syntax}

\end{function}



\begin{function}[]
  {
    \toks
  }
  \begin{syntax}
    \cs{toks}
  \end{syntax}

\end{function}



\begin{function}[]
  {
    \toks@
  }
  \begin{syntax}
    \cs{toks@}
  \end{syntax}

\end{function}



\begin{function}[]
  {
    \toksdef
  }
  \begin{syntax}
    \cs{toksdef}
  \end{syntax}

\end{function}



\begin{function}[]
  {
    \tokszero
  }
  \begin{syntax}
    \cs{tokszero}
  \end{syntax}

\end{function}



\begin{function}[]
  {
    \tolerance
  }
  \begin{syntax}
    \cs{tolerance}
  \end{syntax}

\end{function}



\begin{function}[]
  {
    \topfigrule
  }
  \begin{syntax}
    \cs{topfigrule}
  \end{syntax}

\end{function}



\begin{function}[]
  {
    \topfraction
  }
  \begin{syntax}
    \cs{topfraction}
  \end{syntax}

\end{function}



\begin{function}[]
  {
    \topmargin
  }
  \begin{syntax}
    \cs{topmargin}
  \end{syntax}

\end{function}




\begin{function}[]
  {
    \topskip
  }
  \begin{syntax}
    \cs{topskip}
  \end{syntax}

\end{function}



\begin{function}[]
  {
    \totalheight
  }
  \begin{syntax}
    \cs{totalheight}
  \end{syntax}

\end{function}



\begin{function}[]
  {
    \tracingall
  }
  \begin{syntax}
    \cs{tracingall}
  \end{syntax}

\end{function}



\begin{function}[]
  {
    \tracingassigns
  }
  \begin{syntax}
    \cs{tracingassigns}
  \end{syntax}

\end{function}



\begin{function}[]
  {
    \tracingcommands
  }
  \begin{syntax}
    \cs{tracingcommands}
  \end{syntax}

\end{function}



\begin{function}[]
  {
    \tracingfonts
  }
  \begin{syntax}
    \cs{tracingfonts}
  \end{syntax}

\end{function}



\begin{function}[]
  {
    \tracinggroups
  }
  \begin{syntax}
    \cs{tracinggroups}
  \end{syntax}

\end{function}



\begin{function}[]
  {
    \tracingifs
  }
  \begin{syntax}
    \cs{tracingifs}
  \end{syntax}

\end{function}



\begin{function}[]
  {
    \tracinglostchars
  }
  \begin{syntax}
    \cs{tracinglostchars}
  \end{syntax}

\end{function}



\begin{function}[]
  {
    \tracingmacros
  }
  \begin{syntax}
    \cs{tracingmacros}
  \end{syntax}

\end{function}



\begin{function}[]
  {
    \tracingnesting
  }
  \begin{syntax}
    \cs{tracingnesting}
  \end{syntax}

\end{function}



\begin{function}[]
  {
    \tracingnone
  }
  \begin{syntax}
    \cs{tracingnone}
  \end{syntax}

\end{function}



\begin{function}[]
  {
    \tracingonline
  }
  \begin{syntax}
    \cs{tracingonline}
  \end{syntax}

\end{function}



\begin{function}[]
  {
    \tracingoutput
  }
  \begin{syntax}
    \cs{tracingoutput}
  \end{syntax}

\end{function}



\begin{function}[]
  {
    \tracingpages
  }
  \begin{syntax}
    \cs{tracingpages}
  \end{syntax}

\end{function}



\begin{function}[]
  {
    \tracingparagraphs
  }
  \begin{syntax}
    \cs{tracingparagraphs}
  \end{syntax}

\end{function}



\begin{function}[]
  {
    \tracingrestores
  }
  \begin{syntax}
    \cs{tracingrestores}
  \end{syntax}

\end{function}



\begin{function}[]
  {
    \tracingscantokens
  }
  \begin{syntax}
    \cs{tracingscantokens}
  \end{syntax}

\end{function}



\begin{function}[]
  {
    \tracingstacklevels
  }
  \begin{syntax}
    \cs{tracingstacklevels}
  \end{syntax}

\end{function}



\begin{function}[]
  {
    \tracingstats
  }
  \begin{syntax}
    \cs{tracingstats}
  \end{syntax}

\end{function}



\begin{function}[]
  {
    \transform@scriptfont
  }
  \begin{syntax}
    \cs{transform@scriptfont}
  \end{syntax}

\end{function}



\begin{function}[]
  {
    \trivlist
  }
  \begin{syntax}
    \cs{trivlist}
  \end{syntax}

\end{function}



\begin{function}[]
  {
    \try@load@fontshape
  }
  \begin{syntax}
    \cs{try@load@fontshape}
  \end{syntax}

\end{function}



\begin{function}[]
  {
    \try@simple@size
  }
  \begin{syntax}
    \cs{try@simple@size}
  \end{syntax}

\end{function}



\begin{function}[]
  {
    \try@simples
  }
  \begin{syntax}
    \cs{try@simples}
  \end{syntax}

\end{function}



\begin{function}[]
  {
    \try@size@range
  }
  \begin{syntax}
    \cs{try@size@range}
  \end{syntax}

\end{function}



\begin{function}[]
  {
    \try@size@substitution
  }
  \begin{syntax}
    \cs{try@size@substitution}
  \end{syntax}

\end{function}



\begin{function}[]
  {
    \tryif@simple
  }
  \begin{syntax}
    \cs{tryif@simple}
  \end{syntax}

\end{function}



\begin{function}[]
  {
    \ttdef@ult
  }
  \begin{syntax}
    \cs{ttdef@ult}
  \end{syntax}

\end{function}



\begin{function}[]
  {
    \ttdefault
  }
  \begin{syntax}
    \cs{ttdefault}
  \end{syntax}

\end{function}



\begin{function}[]
  {
    \ttfamily
  }
  \begin{syntax}
    \cs{ttfamily}
  \end{syntax}

\end{function}



\begin{function}[]
  {
    \ttsubstdefault
  }
  \begin{syntax}
    \cs{ttsubstdefault}
  \end{syntax}

\end{function}



\begin{function}[]
  {
    \tw@
  }
  \begin{syntax}
    \cs{tw@}
  \end{syntax}

\end{function}



\begin{function}[]
  {
    \two@digits
  }
  \begin{syntax}
    \cs{two@digits}
  \end{syntax}

\end{function}



\begin{function}[]
  {
    \twocolumn
  }
  \begin{syntax}
    \cs{twocolumn}
  \end{syntax}

\end{function}



\begin{function}[]
  {
    \typein
  }
  \begin{syntax}
    \cs{typein}
  \end{syntax}

\end{function}



\begin{function}[]
  {
    \typeout
  }
  \begin{syntax}
    \cs{typeout}
  \end{syntax}

\end{function}



\begin{function}[]
  {
    \u
  }
  \begin{syntax}
    \cs{u}
  \end{syntax}

\end{function}



\begin{function}[]
  {
    \uccode
  }
  \begin{syntax}
    \cs{uccode}
  \end{syntax}

\end{function}



\begin{function}[]
  {
    \uchyph
  }
  \begin{syntax}
    \cs{uchyph}
  \end{syntax}

\end{function}



\begin{function}[]
  {
    \ulcdefault
  }
  \begin{syntax}
    \cs{ulcdefault}
  \end{syntax}

\end{function}



\begin{function}[]
  {
    \ulcshape
  }
  \begin{syntax}
    \cs{ulcshape}
  \end{syntax}

\end{function}



\begin{function}[]
  {
    \unboldmath
  }
  \begin{syntax}
    \cs{unboldmath}
  \end{syntax}

\end{function}



\begin{function}[]
  {
    \unconditionally@reset@math@version
  }
  \begin{syntax}
    \cs{unconditionally@reset@math@version}
  \end{syntax}

\end{function}



\begin{function}[]
  {
    \undeclare@file@substitution
  }
  \begin{syntax}
    \cs{undeclare@file@substitution}
  \end{syntax}

\end{function}



\begin{function}[]
  {
    \undefined
  }
  \begin{syntax}
    \cs{undefined}
  \end{syntax}

\end{function}



\begin{function}[]
  {
    \undefinedpagestyle
  }
  \begin{syntax}
    \cs{undefinedpagestyle}
  \end{syntax}

\end{function}



\begin{function}[]
  {
    \underbar
  }
  \begin{syntax}
    \cs{underbar}
  \end{syntax}

\end{function}



\begin{function}[]
  {
    \underline
  }
  \begin{syntax}
    \cs{underline}
  \end{syntax}

\end{function}



\begin{function}[]
  {
    \unexpanded
  }
  \begin{syntax}
    \cs{unexpanded}
  \end{syntax}

\end{function}



\begin{function}[]
  {
    \unhbox
  }
  \begin{syntax}
    \cs{unhbox}
  \end{syntax}

\end{function}



\begin{function}[]
  {
    \unhbox0
  }
  \begin{syntax}
    \cs{unhbox0}
  \end{syntax}

\end{function}



\begin{function}[]
  {
    \unhcopy
  }
  \begin{syntax}
    \cs{unhcopy}
  \end{syntax}

\end{function}



\begin{function}[]
  {
    \unicodedataline
  }
  \begin{syntax}
    \cs{unicodedataline}
  \end{syntax}

\end{function}



\begin{function}[]
  {
    \unicoderead
  }
  \begin{syntax}
    \cs{unicoderead}
  \end{syntax}

\end{function}



\begin{function}[]
  {
    \unitlength
  }
  \begin{syntax}
    \cs{unitlength}
  \end{syntax}

\end{function}



\begin{function}[]
  {
    \unkern
  }
  \begin{syntax}
    \cs{unkern}
  \end{syntax}

\end{function}



\begin{function}[]
  {
    \unless
  }
  \begin{syntax}
    \cs{unless}
  \end{syntax}

\end{function}



\begin{function}[]
  {
    \unlhd
  }
  \begin{syntax}
    \cs{unlhd}
  \end{syntax}

\end{function}



\begin{function}[]
  {
    \unpenalty
  }
  \begin{syntax}
    \cs{unpenalty}
  \end{syntax}

\end{function}



\begin{function}[]
  {
    \unqu@tefilef@und
  }
  \begin{syntax}
    \cs{unqu@tefilef@und}
  \end{syntax}

\end{function}



\begin{function}[]
  {
    \unquote@name
  }
  \begin{syntax}
    \cs{unquote@name}
  \end{syntax}

\end{function}



\begin{function}[]
  {
    \unrestored@protected@xdef
  }
  \begin{syntax}
    \cs{unrestored@protected@xdef}
  \end{syntax}

\end{function}



\begin{function}[]
  {
    \unrhd
  }
  \begin{syntax}
    \cs{unrhd}
  \end{syntax}

\end{function}



\begin{function}[]
  {
    \unsetattribute
  }
  \begin{syntax}
    \cs{unsetattribute}
  \end{syntax}

\end{function}



\begin{function}[]
  {
    \unskip
  }
  \begin{syntax}
    \cs{unskip}
  \end{syntax}

\end{function}



\begin{function}[]
  {
    \unvbox
  }
  \begin{syntax}
    \cs{unvbox}
  \end{syntax}

\end{function}



\begin{function}[]
  {
    \unvcopy
  }
  \begin{syntax}
    \cs{unvcopy}
  \end{syntax}

\end{function}



\begin{function}[]
  {
    \update@series@target@value
  }
  \begin{syntax}
    \cs{update@series@target@value}
  \end{syntax}

\end{function}



\begin{function}[]
  {
    \updefault
  }
  \begin{syntax}
    \cs{updefault}
  \end{syntax}

\end{function}



\begin{function}[]
  {
    \upper@bound
  }
  \begin{syntax}
    \cs{upper@bound}
  \end{syntax}

\end{function}



\begin{function}[]
  {
    \upshape
  }
  \begin{syntax}
    \cs{upshape}
  \end{syntax}

\end{function}



\begin{function}[]
  {
    \use@mathgroup
  }
  \begin{syntax}
    \cs{use@mathgroup}
  \end{syntax}

\end{function}



\begin{function}[]
  {
    \usebox
  }
  \begin{syntax}
    \cs{usebox}
  \end{syntax}

\end{function}



\begin{function}[]
  {
    \usecounter
  }
  \begin{syntax}
    \cs{usecounter}
  \end{syntax}

\end{function}



\begin{function}[]
  {
    \usefont
  }
  \begin{syntax}
    \cs{usefont}
  \end{syntax}

\end{function}



\begin{function}[]
  {
    \usepackage
  }
  \begin{syntax}
    \cs{usepackage}
  \end{syntax}

\end{function}



\begin{function}[]
  {
    \v
  }
  \begin{syntax}
    \cs{v}
  \end{syntax}

\end{function}



\begin{function}[]
  {
    \v@false
  }
  \begin{syntax}
    \cs{v@false}
  \end{syntax}

\end{function}



\begin{function}[]
  {
    \v@true
  }
  \begin{syntax}
    \cs{v@true}
  \end{syntax}

\end{function}



\begin{function}[]
  {
    \vadjust
  }
  \begin{syntax}
    \cs{vadjust}
  \end{syntax}

\end{function}



\begin{function}[]
  {
    \valign
  }
  \begin{syntax}
    \cs{valign}
  \end{syntax}

\end{function}



\begin{function}[]
  {
    \value
  }
  \begin{syntax}
    \cs{value}
  \end{syntax}

\end{function}



\begin{function}[]
  {
    \varepsilon
  }
  \begin{syntax}
    \cs{varepsilon}
  \end{syntax}

\end{function}



\begin{function}[]
  {
    \vbadness
  }
  \begin{syntax}
    \cs{vbadness}
  \end{syntax}

\end{function}



\begin{function}[]
  {
    \vbox
  }
  \begin{syntax}
    \cs{vbox}
  \end{syntax}

\end{function}



\begin{function}[]
  {
    \vcenter
  }
  \begin{syntax}
    \cs{vcenter}
  \end{syntax}

\end{function}



\begin{function}[]
  {
    \vcenter@text
  }
  \begin{syntax}
    \cs{vcenter@text}
  \end{syntax}

\end{function}



\begin{function}[]
  {
    \vcenter@text@auxi
  }
  \begin{syntax}
    \cs{vcenter@text@auxi}
  \end{syntax}

\end{function}



\begin{function}[]
  {
    \vcenter@text@auxii
  }
  \begin{syntax}
    \cs{vcenter@text@auxii}
  \end{syntax}

\end{function}



\begin{function}[]
  {
    \vcenter@text@axis
  }
  \begin{syntax}
    \cs{vcenter@text@axis}
  \end{syntax}

\end{function}



\begin{function}[]
  {
    \vector
  }
  \begin{syntax}
    \cs{vector}
  \end{syntax}

\end{function}



\begin{function}[]
  {
    \verb
  }
  \begin{syntax}
    \cs{verb}
  \end{syntax}

\end{function}



\begin{function}[]
  {
    \verb@balance@group
  }
  \begin{syntax}
    \cs{verb@balance@group}
  \end{syntax}

\end{function}



\begin{function}[]
  {
    \verb@egroup
  }
  \begin{syntax}
    \cs{verb@egroup}
  \end{syntax}

\end{function}



\begin{function}[]
  {
    \verb@eol@error
  }
  \begin{syntax}
    \cs{verb@eol@error}
  \end{syntax}

\end{function}



\begin{function}[]
  {
    \verbatim
  }
  \begin{syntax}
    \cs{verbatim}
  \end{syntax}

\end{function}



\begin{function}[]
  {
    \verbatim@font
  }
  \begin{syntax}
    \cs{verbatim@font}
  \end{syntax}

\end{function}



\begin{function}[]
  {
    \verbatim@nolig@list
  }
  \begin{syntax}
    \cs{verbatim@nolig@list}
  \end{syntax}

\end{function}



\begin{function}[]
  {
    \verbvisiblespace
  }
  \begin{syntax}
    \cs{verbvisiblespace}
  \end{syntax}

\end{function}



\begin{function}[]
  {
    \version@elt
  }
  \begin{syntax}
    \cs{version@elt}
  \end{syntax}

\end{function}



\begin{function}[]
  {
    \version@list
  }
  \begin{syntax}
    \cs{version@list}
  \end{syntax}

\end{function}



\begin{function}[]
  {
    \vfil
  }
  \begin{syntax}
    \cs{vfil}
  \end{syntax}

\end{function}



\begin{function}[]
  {
    \vfilneg
  }
  \begin{syntax}
    \cs{vfilneg}
  \end{syntax}

\end{function}



\begin{function}[]
  {
    \vfuzz
  }
  \begin{syntax}
    \cs{vfuzz}
  \end{syntax}

\end{function}



\begin{function}[]
  {
    \vgl@
  }
  \begin{syntax}
    \cs{vgl@}
  \end{syntax}

\end{function}



\begin{function}[]
  {
    \vglue
  }
  \begin{syntax}
    \cs{vglue}
  \end{syntax}

\end{function}



\begin{function}[]
  {
    \vline
  }
  \begin{syntax}
    \cs{vline}
  \end{syntax}

\end{function}



\begin{function}[]
  {
    \voidb@x
  }
  \begin{syntax}
    \cs{voidb@x}
  \end{syntax}

\end{function}



\begin{function}[]
  {
    \vphantom
  }
  \begin{syntax}
    \cs{vphantom}
  \end{syntax}

\end{function}



\begin{function}[]
  {
    \vrule
  }
  \begin{syntax}
    \cs{vrule}
  \end{syntax}

\end{function}



\begin{function}[]
  {
    \vsize
  }
  \begin{syntax}
    \cs{vsize}
  \end{syntax}

\end{function}



\begin{function}[]
  {
    \vskip
  }
  \begin{syntax}
    \cs{vskip}
  \end{syntax}

\end{function}



\begin{function}[]
  {
    \vspace
  }
  \begin{syntax}
    \cs{vspace}
  \end{syntax}

\end{function}



\begin{function}[]
  {
    \vsplit
  }
  \begin{syntax}
    \cs{vsplit}
  \end{syntax}

\end{function}



\begin{function}[]
  {
    \vss
  }
  \begin{syntax}
    \cs{vss}
  \end{syntax}

\end{function}



\begin{function}[]
  {
    \vtop
  }
  \begin{syntax}
    \cs{vtop}
  \end{syntax}

\end{function}



\begin{function}[]
  {
    \wd
  }
  \begin{syntax}
    \cs{wd}
  \end{syntax}

\end{function}



\begin{function}[]
  {
    \whatsit
  }
  \begin{syntax}
    \cs{whatsit}
  \end{syntax}

\end{function}



\begin{function}[]
  {
    \widowpenalty
  }
  \begin{syntax}
    \cs{widowpenalty}
  \end{syntax}

\end{function}



\begin{function}[]
  {
    \width
  }
  \begin{syntax}
    \cs{width}
  \end{syntax}

\end{function}





\begin{function}[]
  {
    \write
  }
  \begin{syntax}
    \cs{write}
  \end{syntax}

\end{function}



\begin{function}[]
  {
    \wrong@fontshape
  }
  \begin{syntax}
    \cs{wrong@fontshape}
  \end{syntax}

\end{function}



\begin{function}[]
  {
    \x
  }
  \begin{syntax}
    \cs{x}
  \end{syntax}

\end{function}



\begin{function}[]
  {
    \x@protect
  }
  \begin{syntax}
    \cs{x@protect}
  \end{syntax}

\end{function}



\begin{function}[]
  {
    \xdef
  }
  \begin{syntax}
    \cs{xdef}
  \end{syntax}

\end{function}



\begin{function}[]
  {
    \xe@alloc@intercharclass
  }
  \begin{syntax}
    \cs{xe@alloc@intercharclass}
  \end{syntax}

\end{function}



\begin{function}[]
  {
    \y
  }
  \begin{syntax}
    \cs{y}
  \end{syntax}

\end{function}



\begin{function}[]
  {
    \year
  }
  \begin{syntax}
    \cs{year}
  \end{syntax}

\end{function}



\begin{function}[]
  {
    \z
  }
  \begin{syntax}
    \cs{z}
  \end{syntax}

\end{function}



\begin{function}[]
  {
    \z@
  }
  \begin{syntax}
    \cs{z@}
  \end{syntax}

\end{function}



\begin{function}[]
  {
    \z@skip
  }
  \begin{syntax}
    \cs{z@skip}
  \end{syntax}

\end{function}



\begin{function}[]
  {
    \zap@space
  }
  \begin{syntax}
    \cs{zap@space}
  \end{syntax}

\end{function}



\begin{function}[]
  {
    \|
  }
  \begin{syntax}
    \cs{|}
  \end{syntax}

\end{function}



%-----------------------------------------------

\begingroup
  \def\endash{--}
  \catcode`\-\active
  \def-{\futurelet\temp\indexdash}
  \def\indexdash{\ifx\temp-\endash\fi}
  \PrintIndex
\endgroup


\end{document}
